% ============ 《调和分析入门》中文翻译 LaTeX 重排版 ============
% 原书: 澤野嘉宏《調和解析への招待 関数の性質を深く理解するために》
%       (SGC ライブラリ-174), サイエンス社, 2022
% 模板: AJbook.cls (李文威《代数学方法》书系模板)
%       ★ 定理编号规则已按原书修改(见 AJbook.cls): 定理/系/引理/命题/定义/例/注意/练习问题
%         各自独立计数器, 按章编号; 字体配置未改动
% 结构: 序言 + 第 1-6 章(含结束语) + 练习题解答 + 参考文献 + 索引
%
% 编号体系 (原书约定):
%   - 章: 第一章, 第二章...(中文数字, 见 titles-setup.tex)
%   - 节: 1.1, 1.2...(章.节); 小节: 不编号 (原书小节无编号)
%   - 定理 1.1 / 定义 2.1 / 例 1.13 / 注意 3.1 / 练习问题 1.21: 每类独立按章计数
%   - 公式: 按章连续编号 (1.1)-(1.16), 显式 \tag{...}
%     (勘误表删去 (1.11)(1.12)(1.13) 编号, 故不用自动计数器)
\documentclass[geometry=a4, fontsetup=font-setup-open.tex, titlesetup=titles-setup, coverpage=cover]{AJbook}

% ============ 宏包 ============
\usepackage{mycommand}      % 常用数学命令(\dd \R \Z \Q \CC \e \ii 等)
\usepackage{myarrows}       % 可伸缩箭头
\usepackage{amssymb}        % \lesssim \gtrsim 等
\usepackage{imakeidx}
\usepackage{hyperref}
\usetikzlibrary{calc}       % 封面坐标计算

% ============ 超链接 ============
\hypersetup{
    pdfauthor={泽野嘉宏},
    pdftitle={调和分析入门——为深入理解函数的性质},
    pdfsubject={调和分析}
}

% ============ 自定义命令 ============
\newcommand{\Fourier}{\mathcal{F}}        % Fourier 变换
\newcommand{\Hilbert}{\mathcal{H}}        % Hilbert 空间/变换
\newcommand{\esssup}{\operatorname{ess\,sup}}
\newcommand{\diam}{\operatorname{diam}}
\newcommand{\dyadic}{\mathrm{dyadic}}     % 下标: M_{\dyadic}

% ============ 索引 ============
\makeindex[columns=2, intoc=true, title={索引}]

% 全局断行弹性: 消除边界性 Overfull (与 Gelbart/古典调和分析项目同款)
\emergencystretch=3em

% ============ 文档开始 ============
\begin{document}
\frontmatter        % 封面(cover.tex) + 目录(模板自动处理, 罗马页码)

% 序言(原书序文)
% ============ 序言(原书序文,原书 pp.1-2 = PDF pp.3-4) ============
\chapter*{序言}\addcontentsline{toc}{chapter}{序言}

本书的目的, 是向对分析学感兴趣的读者传达调和分析学是什么. 调和分析学是以 Lebesgue 积分论为基础的理论, 不理解 Lebesgue 积分论, 就无法达到深刻的理解. 然而, 即使不使用 Lebesgue 积分论也能理解其本质的部分也不在少数, 在第1章中, 我们试图完全不使用 Lebesgue 积分来阐述调和分析. 本书所需要的 Lebesgue 积分知识, 已在第2章前半部分作了整理. 未修过 Lebesgue 积分的读者, 由于第2章前半部分正是关于 Lebesgue 积分的, 可以将其作为事实接受并继续读下去. 即便作为事实接受并继续, 我想也能在一定程度上读下去.

调和分析这个领域实际上非常广阔, 是难以把握的世界. 分析学的领域中有偏微分方程论、算子代数论、复分析、Hilbert 空间论等各种领域, 而调和分析是横跨所有这些领域的学科, 究竟以什么来称之为调和分析, 似乎并不明确. 也许可以称之为一门巧妙利用正负相消, 将由此得到的东西巧妙地累加来进行调节的学问吧?

话虽如此, 为了说明调和分析是什么, 这里原则上将"研究函数性质的数学"称为调和分析. 舍弃函数性质进行抽象化的概念之一有 Banach 空间, 这里也将其作为抽象化的调和分析纳入这一范畴. 原本创立 Banach 空间的 Banach 把一个个函数看作点, 这在当时(恐怕现在也是)是划时代的发现, 所以我认为将 Banach 空间视为调和分析是有意义的.

调和分析所研究的对象是函数的性质, 大致分为函数的微分(质)与函数的积分(量). 正如高中所学的数学那样, 微分与积分是互逆的运算, 但在分析学中这一点并不太被强调. 实际上, 微分和积分各自都有很大的延展, 将它们作为不同的东西来处理才是妥当的. 下面谈谈本书的构成. Lebesgue 积分论是研究调和分析所必需的理论, 但却是困难的理论. 因此, 本书为了不用 Lebesgue 积分也能理解, 决定使用数列空间来概观调和分析. 这就是第1章. 从第2章开始, 转入使用 Lebesgue 积分论的调和分析. 第2章将理解调和分析所必需的 Lebesgue 积分论和泛函分析学作了紧凑的整理. 如果讲解 Lebesgue 积分, 书籍就会变得庞大, 因此这里不作证明, 只记载必要事项. 关于泛函分析学, 以数列空间中触及的内容为基础, 给出了必要事项的证明. 第3章整理了调和分析的重要主题兼技巧的极大算子. 第1章考虑以全体自然数为基调的数列空间, 而第3章考虑 Euclid 空间上的 Lebesgue 空间, 因而叙述变得复杂. 也可以认为正是把那复杂的部分特意抽了出来. 第4章也可以说是考虑第1章所处理的内容在 Lebesgue 空间中会变成怎样. 数学分析中出现的积分算子以 Lebesgue 空间论中所处理的积分算子为基本, 但有时仅使用 Lebesgue 空间论中所处理的不等式, 无法得到令人满意的结论. 利用 Euclid 空间特有的性质可以进行进一步的分析, 将其应用于得到种种结果, 是很重要的事情. 其一个例子就是第1章中处理的 Calderón--Zygmund 分解. 第4章特别考察对流体力学等也很重要的 Hilbert 变换的性质. 第5章考察 Fourier 变换. Fourier 变换是数学分析中重要的积分算子, 从它的定义起就是调和分析学的重要主题. 从 Fourier 变换的定义式来看, 似乎函数不可积就无法定义, 但通过种种方法, 对不可积的函数也能定义. 此时重要的是巧妙地运用拓扑来操作极限等, 本书特别注意到这一侧面, 说明了 Fourier 变换的定义和处理方式等. 由于要处理所谓函数以外对象的 Fourier 变换, 处理会变得困难, 但本书仔细说明了应当注意哪些点来处理. 最后的第6章, 作为 Fourier 变换的应用, 说明了作为 $L^{p}$ 范数改写工具的 Littlewood--Paley 理论. 正如本书也要提到的, $L^{p}$ 范数在 $p = 1, 2$ 以外的情况下难以处理, 但对于其他的 $p$, 通过变形范数有时能看清各种各样的性质. 其一个例子就是 Littlewood--Paley 理论, 本书也兼作与第1章数列空间的比较, 说明了 Littlewood--Paley 理论.

调和分析学中有各种各样的理论, 是非常丰富的世界. 插值理论、函数空间论等作为调和分析学理论的一部分发展至今, 但考虑到入门书的性质, 除插值理论的特殊情形外, 本书不包含这些理论. 请参考卷末列举的书籍. 不过, 本书没有网罗的理论, 其背景也有与本书所说明的类似内容, 需要那种理论的动机, 可以从本书所写的内容中推测出来吧.

本书基于笔者在学習院大学、中央大学、東京都立大学的讲义写成. 执笔之际, 得到了中央大学的波多野修也先生、阿部雄斗先生, 東京都立大学的野ケ山徹先生、服部久美子先生的建议. 在此谨表谢意.

2022年1月

泽野嘉宏


\mainmatter         % 正文(阿拉伯页码从 1 开始)

% 第1章 从数列空间角度概观调和分析 (原书 pp.1-35, PDF pp.7-41)
\chapter{从数列空间角度概观调和分析}
% ============ 第1章 §1.1 ℓ^p 空间 (原书 pp.1-6) ============
\section{\texorpdfstring{$\ell^{p}$ 空间}{ellp 空间}}\label{sec:1.1}
\index{0symbollinf@$\ell^{\infty}$}

\index{0symbollp@$\ell^{p}$}

\index{C@Cauchy--Schwarz 不等式}


设 $N$ 为自然数, 考虑由 $N$ 个实数组成的数列 $a_{1}, a_{2}, \ldots, a_{N}$. 将其视为 $N$ 维向量, 记 $a = (a_{1}, a_{2}, \ldots, a_{N})$. 这样的数列全体亦即向量全体构成线性空间 $\R^{N}$. 加法按分量定义, 数乘通过将该数乘到各分量来定义. 由数列构成的线性空间及其子空间称为数列空间.

\begin{definition}\label{def:1.1}
对 $1 \leq p < \infty$, 定义
\[
\| a \|_{\ell^{p}} = \left(\sum_{n=1}^{N} |a_{n}|^{p}\right)^{\frac{1}{p}},
\]
称之为 $a$ 的 $\ell^{p}$-范数. 另外, $\sum_{j=1}^{N} |a_{j}|^{p}$ 称为 $p$ 乘和.

当把线性空间 $\R^{N}$ 与 $\ell^{p}$-范数成对考虑时, 称 $\R^{N}$ 为有限集合 $\{1, 2, \ldots, N\}$ 上的 $\ell^{p}$ 空间.
\end{definition}

\begin{example}\label{exa:1.1}
当 $\{a_{n}\}_{n=1}^{N}$ 由 $a_{n} = n$ 给定时,
\[
\| \{a_{n}\}_{n=1}^{N} \|_{\ell^{1}} = 1 + 2 + \cdots + N = \frac{N(N+1)}{2},
\]
\[
\| \{a_{n}\}_{n=1}^{N} \|_{\ell^{2}} = \sqrt{1^{2} + 2^{2} + \cdots + N^{2}} = \sqrt{\frac{N(N+1)(2N+1)}{6}}.
\]
\end{example}

下面的不等式将在全书范围内说明 $\ell^{p}$-范数为何重要.

\begin{theorem}\label{thm:1.1}
对两个数列 $a = (a_{1}, a_{2}, \ldots, a_{N})$, $b = (b_{1}, b_{2}, \ldots, b_{N})$, 称为三角不等式的下述不等式成立:
\[
\|a + b\|_{\ell^{p}} \leq \|a\|_{\ell^{p}} + \|b\|_{\ell^{p}}.
\]
\end{theorem}

\begin{proof}
暂时假定 $a_{n}, b_{n} \geq 0$ 来考虑. 如果 $a$, $b$ 中某一方为 $0$, 则不等号变为等号. 基于这一点, 考虑所有分量都不为 $0$ 的两个数列 $a = (a_{1}, a_{2}, \ldots, a_{N})$, $b = (b_{1}, b_{2}, \ldots, b_{N})$. 这里使用三角不等式. 由函数 $\varphi(t) = t^{p}$ 的凸性, 即对 $t_{1}, t_{2} \geq 0$, $0 \leq \theta \leq 1$ 的不等式
\[
\varphi((1-\theta)t_{1} + \theta t_{2}) \leq (1-\theta)\varphi(t_{1}) + \theta\varphi(t_{2}),
\]
对 $\theta = \frac{\|b\|_{\ell^{p}}}{\|a\|_{\ell^{p}} + \|b\|_{\ell^{p}}}$, 有
\begin{align*}
\left(\frac{|a_{n} + b_{n}|}{\|a\|_{\ell^{p}} + \|b\|_{\ell^{p}}}\right)^{p} & \leq \left(\frac{|a_{n}| + |b_{n}|}{\|a\|_{\ell^{p}} + \|b\|_{\ell^{p}}}\right)^{p} \\
& \leq (1-\theta)\left(\frac{|a_{n}|}{\|a\|_{\ell^{p}}}\right)^{p} + \theta\left(\frac{|b_{n}|}{\|b\|_{\ell^{p}}}\right)^{p},
\end{align*}
所以, 将其对 $n = 1, 2, \ldots, N$ 求和并整理, 所期望的不等式成立.
\end{proof}

于是对任意数列 $a$, $b$ 有 $\|a + b\|_{\ell^{p}} \leq \|a\|_{\ell^{p}} + \|b\|_{\ell^{p}}$ 成立. 这个不等式称为 Minkowski 不等式. $p = 2$ 的情形就是通常的三角不等式. 接下来, 为了证明 Hölder 不等式, 先证明 Young 不等式.

\begin{lemma}[Young 不等式]\label{lem:1.1}
设 $p, q > 1$ 满足 $\frac{1}{p} + \frac{1}{q} = 1$. 此时, 对 $\alpha, A, B > 0$,
\[
\frac{\alpha^{p}}{p}A + \frac{\alpha^{-q}}{q}B \geq A^{\frac{1}{p}}B^{\frac{1}{q}}
\]
成立. 当 $\alpha = \left(\frac{B}{A}\right)^{\frac{1}{p+q}}$ 时等号成立. 特别地, 对 $a, b > 0$,
\[
ab \leq \frac{1}{p}a^{p} + \frac{1}{q}b^{q}
\]
成立.
\end{lemma}

\begin{proof}
令 $f(\alpha) = \frac{\alpha^{p}}{p}A + \frac{\alpha^{-q}}{q}B$, 则 $f'(\alpha) = \alpha^{p-1}A - \alpha^{-q-1}B$. 因此, 当 $\alpha = \left(\frac{B}{A}\right)^{\frac{1}{p+q}}$ 时 $f$ 取极小值, 所以
\begin{align*}
f(\alpha) & \geq f\left(\left(\frac{B}{A}\right)^{\frac{1}{p+q}}\right) \\
& = \frac{1}{p}\left(\frac{B}{A}\right)^{\frac{p}{p+q}}A + \frac{1}{q}\left(\frac{B}{A}\right)^{-\frac{q}{p+q}}B \\
& = \frac{1}{p}A^{\frac{q}{p+q}}B^{\frac{p}{p+q}} + \frac{1}{q}A^{\frac{q}{p+q}}B^{\frac{p}{p+q}} \\
& = A^{\frac{1}{p}}B^{\frac{1}{q}}.
\end{align*}
得证.
\end{proof}

应用 Young 不等式可知如下事实. 设 $A, B, C > 0$. 又设 $p, q > 1$ 满足 $\frac{1}{p} + \frac{1}{q} = 1$. 此时, 若对 $\alpha > 0$,
\[
\frac{\alpha^{p}}{p}A + \frac{\alpha^{-q}}{q}B \geq C
\]
成立, 则 $C \leq A^{\frac{1}{p}}B^{\frac{1}{q}}$ 成立. 这个想法在本节剩余部分将频繁使用. 在引理~\ref{lem:1.1} 中, 由关系式 $\frac{1}{p} + \frac{1}{q} = 1$ 联系起来的 $1 < p, q < \infty$ 称为处于调和共轭关系. 进一步, 对 $1 < p < \infty$, 称 $q = \frac{p}{p-1}$ 为 $p$ 的调和共轭, 也常记为 $p'$. 另外, 约定 $1' = \infty$, $\infty' = 1$.

\begin{example}\label{exa:1.2}
$1$ 的调和共轭是 $\infty$, $2$ 的调和共轭是 $2$, $3$ 的调和共轭是 $\frac{3}{2}$.
\end{example}

对于处于调和共轭关系的 $p$, $q$, 联系 $\ell^{1}$, $\ell^{p}$, $\ell^{q}$ 的 Hölder 不等式是重要的.

\begin{theorem}[Hölder 不等式]\label{thm:1.2}
设 $a_{1}, a_{2}, \ldots, a_{N}, b_{1}, b_{2}, \ldots, b_{N}$ 为实数. 又对处于调和共轭关系的 $1 < p, q < \infty$, 不等式
\[
\left| \sum_{n=1}^{N} a_{n}b_{n} \right| \leq \left(\sum_{n=1}^{N} |a_{n}|^{p}\right)^{\frac{1}{p}}\left(\sum_{n=1}^{N} |b_{n}|^{q}\right)^{\frac{1}{q}}
\]
成立.
\end{theorem}

$p = q = 2$ 的情形称为 Cauchy--Schwarz 不等式.

\begin{proof}
必要时将 $a_{j}$ 替换为 $|a_{j}|$, 将 $b_{j}$ 替换为 $|b_{j}|$ 来考虑, 可设 $a_{j}, b_{j} \geq 0$. 实际上, 由三角不等式,
\[
\left| \sum_{n=1}^{N} a_{n}b_{n} \right| \leq \sum_{n=1}^{N} |a_{n}b_{n}|
\]
就是这个缘故. 以 $\alpha > 0$ 为常数, 由引理~\ref{lem:1.1} 的不等式, 令 $a = \alpha a_{n}$, $b = \alpha^{-1}b_{n}$, 得
\begin{equation}\tag{1.1}\label{eq:1.1}
a_{n}b_{n} \leq \frac{\alpha^{p}}{p}a_{n}^{p} + \frac{\alpha^{-q}}{q}b_{n}^{q}
\end{equation}
将其对 $n = 1, 2, \ldots, N$ 求和.
\begin{equation}\tag{1.2}\label{eq:1.2}
\sum_{n=1}^{N} a_{n}b_{n} \leq \frac{\alpha^{p}}{p}\sum_{n=1}^{N} a_{n}^{p} + \frac{\alpha^{-q}}{q}\sum_{n=1}^{N} b_{n}^{q}.
\end{equation}
这里, 将右边的式子看作 $\alpha$ 的函数, 令
\[
g(\alpha) = \frac{\alpha^{p}}{p}\sum_{n=1}^{N} a_{n}^{p} + \frac{\alpha^{-q}}{q}\sum_{n=1}^{N} b_{n}^{q}.
\]
由引理~\ref{lem:1.1}, $g(\alpha) \geq \sum_{n=1}^{N} a_{n}b_{n}$, 所以利用 $g(\alpha)$ 的最小值,
\[
\sum_{n=1}^{N} a_{n}b_{n} \leq \left(\sum_{n=1}^{N} a_{n}^{p}\right)^{\frac{1}{p}}\left(\sum_{n=1}^{N} b_{n}^{q}\right)^{\frac{1}{q}}
\]
得证.
\end{proof}

来考察不等式的等号成立条件吧. 虽说并非所有不等式的等号成立条件都能考察, 但在能够考察等号成立条件的时候, 考察一下是很重要的.

考虑最简单的最初的定理的情形. 由于式~\eqref{eq:1.1} 中等号必须成立, 所以参考 $f(a) = \frac{1}{p}a^{p} + \frac{1}{q}b^{q} - ab$ 的极小值在哪里取到,
\begin{equation}\tag{1.3}\label{eq:1.3}
\alpha a_{n} = \left(\alpha^{-1}b_{n}\right)^{\frac{1}{p-1}}, \quad n = 1, 2, \ldots, N
\end{equation}
必须成立. 这里, 式~\eqref{eq:1.2} 的右边成为 Hölder 不等式的右边, 所以 $\alpha$ 取使式~\eqref{eq:1.2} 的右边最小的值.

反过来, 若条件 ``存在某个 $\beta$ 使得 $a_{n} = \beta(b_{n})^{\frac{1}{p-1}}, n = 1, 2, \ldots, N$ 成立'' 成立, 则令 $\beta = \alpha^{-\frac{p}{p-1}}$, 关系式~\eqref{eq:1.3} 就出现了. 利用它可知 Hölder 不等式的等号成立. 即使 $a_{n}$, $b_{n}$ 是一般的实数, 只要仔细考虑符号, 情况也是相同的. 也就是说, $a_{n}b_{n} \geq 0$ 对所有 $n = 1, 2, \ldots, N$ 成立, 或 $a_{n}b_{n} \leq 0$ 对所有 $n = 1, 2, \ldots, N$ 成立, 并且存在某个 $\beta$ 使得 $|a_{n}| = \beta|b_{n}|^{\frac{1}{p-1}}, n = 1, 2, \ldots, N$ 成立, 这种情形就是等号成立条件.

\begin{definition}\label{def:1.2}
对数列 $a = \{a_{n}\}_{n=1}^{N}$, 定义 $\|a\|_{\ell^{\infty}} = \sup_{n = 1, 2, \ldots, N} |a_{n}|$.
\end{definition}

有时也给 $\R^{N}$ 配备 $\ell^{\infty}$-范数.

可以看出, 到目前为止的命题对 $p = 1, \infty$ 也成立.

数列的 Hölder 不等式还可以变形为积分形式. 设 $f, g: [a, b] \to \R$ 为连续有界函数. 由引理~\ref{lem:1.1}, 对任意 $A > 0$,
\[
|f(x)g(x)| \leq \frac{A^{p}}{p}|f(x)|^{p} + \frac{A^{-q}}{q}|g(x)|^{q}
\]
成立. 将其在 $[a, b]$ 上积分,
\[
\int_{a}^{b} |f(x)g(x)| \dd x \leq \frac{A^{p}}{p}\int_{a}^{b} |f(x)|^{p} \dd x + \frac{A^{-q}}{q}\int_{a}^{b} |g(x)|^{q} \dd x
\]
成立. 由引理~\ref{lem:1.1}, 选取 $A$ 使右边最小, 得
\begin{equation}\tag{1.4}\label{eq:1.4}
\int_{a}^{b} |f(x)g(x)| \dd x \leq \left(\int_{a}^{b} |f(x)|^{p} \dd x\right)^{\frac{1}{p}}\left(\int_{a}^{b} |g(x)|^{q} \dd x\right)^{\frac{1}{q}}.
\end{equation}
通过替换幂次, 对满足 $\frac{1}{r} = \frac{1}{p} + \frac{1}{q}$ 的 $p, q, r \geq 1$, 得
\begin{equation}\tag{1.5}\label{eq:1.5}
\left(\int_{a}^{b} |f(x)g(x)|^{r} \dd x\right)^{\frac{1}{r}} \leq \left(\int_{a}^{b} |f(x)|^{p} \dd x\right)^{\frac{1}{p}}\left(\int_{a}^{b} |g(x)|^{q} \dd x\right)^{\frac{1}{q}}.
\end{equation}
式~\eqref{eq:1.4} 与式~\eqref{eq:1.5} 称为 Hölder 不等式的积分形式. 即使在积分的情形, $p = q = 2$ 时也称为 Cauchy--Schwarz 不等式. 例如, 在 Hölder 不等式~\eqref{eq:1.5} 中, 令 $p = 3$, $q = 6$, $r = 2$, $f = |\varphi|^{\frac{1}{3}}$, $g = |\varphi|^{\frac{2}{3}}$, 则对函数 $\varphi: [0, 1] \to \R$ 得积分不等式
\begin{equation}\tag{1.6}\label{eq:1.6}
\left(\int_{0}^{1} |\varphi(x)|^{2} \dd x\right)^{\frac{1}{2}} \leq \left(\int_{0}^{1} |\varphi(x)| \dd x\right)^{\frac{1}{3}}\left(\int_{0}^{1} |\varphi(x)|^{4} \dd x\right)^{\frac{1}{6}}.
\end{equation}
Hölder 不等式的积分形式不限于连续函数, 只要积分能用 Riemann 和的极限表示就足够了, 因此例如对只在有限个点不连续的有界函数等也成立. 进一步, 例如对这种函数与 $1 \leq p \leq p_{0} < \infty$,
\begin{equation}\tag{1.7}\label{eq:1.7}
\left(\int_{0}^{1} |f(x)|^{p} \dd x\right)^{\frac{1}{p}} \leq \left(\int_{0}^{1} |f(x)|^{p_{0}} \dd x\right)^{\frac{1}{p_{0}}}
\end{equation}
成立. 实际上, 在式~\eqref{eq:1.4} 中取 $[a, b] = [0, 1]$, $g = 1$, 则
\[
\int_{0}^{1} |f(x)| \dd x \leq \left(\int_{0}^{1} |f(x)|^{p} \dd x\right)^{\frac{1}{p}},
\]
然后适当替换函数和指数即可.

\begin{exercise}\label{exer:1.1}
设数列 $a = \{a_k\}_{k=1}^{N}$ 由 $a_k = \sqrt{k} \ (k = 1, 2, \ldots, N)$ 给出. 此时, 计算 $\|a\|_{\ell^2}$, $\|a\|_{\ell^4}$.
\end{exercise}

\begin{exercise}\label{exer:1.2}
设 $\vec{a} = (2, 3, 1)$, $\vec{b} = (3, 4, 2)$. 计算 $\vec{a} + \vec{b}$, 并利用其结果证明 $\sqrt{83} \leq \sqrt{14} + \sqrt{29}$.
\end{exercise}

\begin{exercise}\label{exer:1.3}
对 4 个正实数 $a, b, c, d$, 证明不等式 $(ac^{2} + bd^{2})^{3} \leq (a^{3} + b^{3})(c^{3} + d^{3})^{2}$.
\end{exercise}

\begin{exercise}\label{exer:1.4}
对 $x, y \in \R$ 与 $\theta \in (0, 1)$, 证明 $e^{(1-\theta)x + \theta y} \leq (1-\theta)e^{x} + \theta e^{y}$.
\end{exercise}

\begin{exercise}\label{exer:1.5}
对 $1 \leq p < \infty$, $k \in \R$ 与不为零数列的实数列 $a = \{a_k\}_{k=1}^{N}$, 计算 $\|ka\|_{\ell^{p}} \div \|a\|_{\ell^{p}}$.
\end{exercise}

\begin{exercise}\label{exer:1.6}
设 $a, b \in \R$, 证明
\[
\max(a, b) = \frac{a + b + |a - b|}{2}, \quad \min(a, b) = \frac{a + b - |a - b|}{2}.
\]
\end{exercise}

% ============ 第1章 §1.2 分布等式 (原书 pp.6-9) ============
\section{分布等式}\label{sec:1.2}
\index{0symbolchi@$\chi$}

\index{T@特征函数}






这里以 $N$ 为自然数, 考虑由从第 1 项开始的 $N$ 个数组成的数列 $a = \{a_n\}_{n=1}^{N}$. 另外, 一般地, 集合 $E$ 的特征函数是指在 $E$ 上取 $1$, 在 $E$ 外取 $0$ 的函数, 本书中用 $\chi_E$ 表示. 当 $|a_n| \in (t, \infty)$, 即 $|a_n| > t$ 时, $\chi_{(t,\infty)}(|a_n|) = 1$; 当 $|a_n| \notin (t, \infty)$, 即 $|a_n| \leq t$ 时, $\chi_{(t,\infty)}(|a_n|) = 0$. 因此, 各 $a_n$ 的绝对值可看作
\[
|a_n| = \int_{0}^{|a_n|} \dd t = \int_{0}^{\infty} \chi_{(t,\infty)}(|a_n|) \dd t.
\]
将其对 $n = 1, 2, \ldots, N$ 求和,
\[
\sum_{n=1}^{N} |a_n| = \int_{0}^{\infty} \sum_{n=1}^{N} \chi_{(t,\infty)}(|a_n|) \dd t.
\]
一般地, 对集合 $A$, 用 $\#A$ 表示 $A$ 的元素个数. 使用这个记号, 对每个 $t > 0$,
\[
\sum_{n=1}^{N} \chi_{(t,\infty)}(|a_n|) = \# \{1 \leq n \leq N : |a_n| > t\}.
\]
因此, 下面的定理成立.

\begin{theorem}\label{thm:1.3}
对 $N$ 个实数 $a_{1}, a_{2}, \ldots, a_{N}$,
\[
\sum_{n=1}^{N} |a_n| = \int_{0}^{\infty} \# \{1 \leq n \leq N : |a_n| > t\} \dd t
\]
成立.
\end{theorem}

这里, 当 $0 < p < \infty$ 时,
\begin{align*}
p \int_{0}^{\infty} t^{p-1} \# \{1 \leq n \leq N : |a_n| > t\} \dd t & = p \int_{0}^{\infty} t^{p-1} \# \{1 \leq n \leq N : |a_n|^{p} > t^{p}\} \dd t,
\end{align*}
通过变量替换 $t \mapsto s = t^{p}$, 关于 $p$ 乘和得到下面的结果.

\begin{theorem}\label{thm:1.4}
设 $0 < p < \infty$. 对 $N$ 个实数 $a_{1}, a_{2}, \ldots, a_{N}$,
\[
\sum_{n=1}^{N} |a_n|^{p} = p \int_{0}^{\infty} t^{p-1} \# \{1 \leq n \leq N : |a_n| > t\} \dd t
\]
成立.
\end{theorem}

\begin{exercise}\label{exer:1.7}
计算下列广义积分.
\begin{align*}
(1) & \int_{1}^{\infty} \frac{\dd t}{t} & (2) & \int_{1}^{\infty} \frac{\dd t}{t^{2}} & (3) & \int_{1}^{\infty} \frac{\dd t}{\sqrt{t}} \\
(4) & \int_{0}^{1} \frac{\dd t}{t} & (5) & \int_{0}^{1} \frac{\dd t}{t^{2}} & (6) & \int_{0}^{1} \frac{\dd t}{\sqrt{t}}
\end{align*}
\end{exercise}

\begin{exercise}\label{exer:1.8}
对 $f(t) = 3\chi_{[0,4)}(t) - 5\chi_{[5,8)}(t)$, 画出 $y = f(t)$ 的图形, 并计算 $\int_{0}^{\infty} f(t) \dd t$.
\end{exercise}

\begin{exercise}\label{exer:1.9}
回答下列问题.

(1) 对 $N$ 个自然数 $a_{1}, a_{2}, \ldots, a_{N}$, 证明
\[
\sum_{n=1}^{N} |a_n| = \sum_{l=0}^{\infty} \# \{1 \leq n \leq N : |a_n| > l\}
\]
成立.

(2) 设 $n$ 为自然数. 证明 $n!$ 用十进制书写时, 末尾出现的 $0$ 的个数为
\[
\sum_{k=1}^{\infty} \left[\frac{n}{5^{k}}\right].
\]
\end{exercise}

% ============ 第1章 §1.3 算子的范数 (原书 pp.9-14) ============
\section{算子的范数}\label{sec:1.3}
\index{J@基本向量}

\index{S@算子范数}






作为线性代数的重要要素, 可以举出线性空间与线性映射. 作为线性空间, 这里考虑了有限集合 $\{1, 2, \ldots, N\}$ 上的 $\ell^{p}$ 空间, $1 \leq p \leq \infty$. 由 $\ell^{p}$-范数, 向量的大小被确定了. 与之对应, 对未必相等的自然数 $M$, $N$ 与满足 $1 \leq p, q \leq \infty$ 的未必相等的实数 $p$, $q$, 想确定从有限集合 $\{1, 2, \ldots, M\}$ 上的 $\ell^{p}$ 空间到有限集合 $\{1, 2, \ldots, N\}$ 上的 $\ell^{q}$ 空间的线性映射 $T$ 的大小. 考虑到 $T$ 的线性, 具体按如下方式确定.
\[
\| T \|_{\ell^{p} \to \ell^{q}} = \sup_{a \in \ell^{p} \setminus \{0\}} \frac{\| T a \|_{\ell^{q}}}{\| a \|_{\ell^{p}}} = \sup_{a \in \ell^{p}, \| a \|_{\ell^{p}} = 1} \| T a \|_{\ell^{q}}.
\]
将 $\|T\|_{\ell^{p} \to \ell^{q}}$ 称为 $T$ 的从 $(\R^{N}, \ell^{p})$ 到 $(\R^{M}, \ell^{q})$ 的算子范数, 但除去根据上下文显然的要素, 例如也称为从 $\ell^{p}$ 到 $\ell^{q}$ 的算子范数.

利用三角函数, 可以计算配备 $\ell^{2}$ 范数的坐标平面 $\R^{2}$ 之间的线性算子的算子范数.

\begin{example}\label{exa:1.3}
设 $\R^{2}$ 配备范数 $\|(x, y)\| = \sqrt{x^{2} + y^{2}}$. 线性映射 $f: \R^{2} \to \R^{2}$ 由 $f(x, y) = (6x - 2y, 8x - 4y)$ 给出. 来计算 $f$ 的算子范数. 该值作为 $\sqrt{x^{2} + y^{2}} = 1$ 时 $\sqrt{(6x-2y)^{2} + (8x-4y)^{2}}$ 的上限给出. 由于 $(6x-2y)^{2} + (8x-4y)^{2} = 100x^{2} - 88xy + 20y^{2}$, 将约束条件 $x^{2} + y^{2} = 1$ 表示为 $x = \cos\theta$, $y = \sin\theta$ 时, 该式的值可用适当的 $\alpha$ 合成为
\begin{align*}
100\cos^{2}\theta - 88\sin\theta\cos\theta + 20\sin^{2}\theta & = 60 + 40\cos 2\theta - 44\sin 2\theta \\
& = 60 + 4\sqrt{221}\sin(2\theta + \alpha).
\end{align*}
因此, $f$ 的算子范数为 $2\sqrt{15 + \sqrt{221}} = \sqrt{26} + \sqrt{34}$.
\end{example}

来确认: 即使是相同的线性算子, 若配备的范数不同, 算子范数也不同.

\begin{example}\label{exa:1.4}
线性映射 $f: \R^{2} \to \R^{2}$ 由 $f(x, y) = (4x - 2y, x - 4y)$ 给出.

(1) 在条件 $x^{2} + y^{2} = 1$ 下求 $(4x-2y)^{2} + (x-4y)^{2}$ 的最大值, 为 $\frac{37 + 3\sqrt{65}}{2}$, 所以若定义域与值域都配备 $\ell^{2}$-范数, 则 $f$ 的算子范数为 $\sqrt{\frac{37 + 3\sqrt{65}}{2}}$.

(2) 在条件 $|x| + |y| = 1$ 下求 $(4x-2y)^{2} + (x-4y)^{2}$ 的最大值, 为 $20$, 所以若定义域配备 $\ell^{1}$-范数, 值域配备 $\ell^{2}$-范数, 则 $f$ 的算子范数为 $2\sqrt{5}$.

(3) 在条件 $x^{2} + y^{2} = 1$ 下求 $|4x - 2y| + |x - 4y|$ 的最大值. 令 $k = |4x - 2y| + |x - 4y|$. 设 $k > 0$. 此时, $xy$ 平面上的方程 $k = |4x - 2y| + |x - 4y|$ 是以 $\pm\left(\frac{2k}{7}, \frac{k}{14}\right)$, $\pm\left(\frac{k}{7}, \frac{2k}{7}\right)$ 为顶点的平行四边形. 连接 $\left(\frac{2k}{7}, \frac{k}{14}\right)$ 与 $\left(\frac{k}{7}, \frac{2k}{7}\right)$ 的直线 $k = 3x + 2y$ 到原点的距离为 $\frac{k}{\sqrt{13}}$, 连接 $\left(\frac{2k}{7}, \frac{k}{14}\right)$ 与 $-\left(\frac{k}{7}, \frac{2k}{7}\right)$ 的直线 $k = 5x - 6y$ 到原点的距离为 $\frac{k}{\sqrt{61}}$. 而且, 给出后者的点位于 $\left(\frac{2k}{7}, \frac{k}{14}\right)$ 与 $-\left(\frac{k}{7}, \frac{2k}{7}\right)$ 这两点之间. 因此, 使 $k = |4x - 2y| + |x - 4y|$ 与 $x^{2} + y^{2} = 1$ 相交的 $k$ 的最大值为 $\sqrt{61}$. 所以, 若定义域配备 $\ell^{2}$-范数, 值域配备 $\ell^{1}$-范数, 则 $f$ 的算子范数为 $\sqrt{61}$.

(4) 在条件 $x^{2} + y^{2} = 1$ 下求 $\max(|4x-2y|, |x-4y|)$ 的最大值, 为 $2\sqrt{5}$, 所以若定义域配备 $\ell^{2}$-范数, 值域配备 $\ell^{\infty}$-范数, 则 $f$ 的算子范数为 $2\sqrt{5}$.
\end{example}

\begin{example}\label{exa:1.5}
定义域与值域都配备 $\ell^{2}$-范数. 求由
\[
(a_{1}, a_{2}, a_{3}, a_{4}) \mapsto \left(\frac{a_{1} + a_{2}}{2}, \frac{a_{1} + a_{2} + a_{3} + a_{4}}{4}, 0, 0\right)
\]
给出的从 $\R^{4}$ 到 $\R^{4}$ 的线性映射的范数.

在条件
\[
a_{1}^{2} + a_{2}^{2} + a_{3}^{2} + a_{4}^{2} = 1
\]
下, 所求范数是
\[
\left(\frac{a_{1} + a_{2}}{2}\right)^{2} + \left(\frac{a_{1} + a_{2} + a_{3} + a_{4}}{4}\right)^{2}
\]
的最大值的平方根. 令 $a_{1} = r\cos\theta$, $a_{2} = r\sin\theta$, $a_{3} = R\cos\phi$, $a_{4} = R\sin\phi$, 则条件可改写为 $r^{2} + R^{2} = 1$, 所求范数的平方是
\[
\left(\frac{r\cos\theta + r\sin\theta}{2}\right)^{2} + \left(\frac{r\cos\theta + r\sin\theta + R\cos\phi + R\sin\phi}{4}\right)^{2}
\]
关于 $R$, $r$, $\theta$, $\phi$ 的最大值. 关于这个量, 由三角函数的合成可知,
\[
\frac{r^{2}}{2} + \frac{(r + R)^{2}}{8}
\]
是上界. 于是, 令 $r = \cos\tau$, $R = \sin\tau$, 则
\[
\frac{r^{2}}{2} + \frac{(r + R)^{2}}{8} = \frac{5}{8}\cos^{2}\tau + \frac{1}{4}\sin\tau\cos\tau + \frac{1}{8}\sin^{2}\tau.
\]
由半角及倍角公式,
\[
\frac{r^{2}}{2} + \frac{(r + R)^{2}}{8} = \frac{3}{8} + \frac{1}{8}\sin 2\tau + \frac{1}{4}\cos 2\tau,
\]
所以所求最大值为
\[
\sqrt{\frac{3 + \sqrt{5}}{8}} = \frac{1 + \sqrt{5}}{4}.
\]
\end{example}

用 $\mathrm{diag}(\alpha_{1}, \alpha_{2}, \ldots, \alpha_{N})$ 表示对角成分为 $\alpha_{1}, \alpha_{2}, \ldots, \alpha_{N}$ 的对角矩阵. 来计算对角矩阵生成的线性算子的算子范数.

\begin{example}\label{exa:1.6}
设 $1 \leq p \leq \infty$, $b = (b_{1}, b_{2}, \ldots, b_{N}) \in \R^{N}$ 为常数. 考虑把 $a = (a_{1}, a_{2}, \ldots, a_{N}) \in \R^{N}$ 移到 $m_{b}(a) = (b_{1}a_{1}, b_{2}a_{2}, \ldots, b_{N}a_{N}) \in \R^{N}$ 的映射 $m_{b}$. 用适当的 $1 \leq K \leq N$, 记 $\max_{1 \leq j \leq N} |b_{j}| = |b_{K}|$. 只有第 $K$ 个成分为 $1$ 而其他成分为 $0$ 的数列 (第 $K$ 基本向量) $\mathbf{e}_{K}$ 被映射到 $b_{K}\mathbf{e}_{K}$, 所以 $\|m_{b}\|_{\ell^{p} \to \ell^{p}} \geq |b_{K}|$. 另一方面, 将 $|b_{1}|, |b_{2}|, \ldots, |b_{N}|$ 替换为 $|b_{K}|$ 来考虑, 得 $\|m_{b}(a)\|_{\ell^{p}} \leq |b_{K}| \|a\|_{\ell^{p}}$. 因此, $\|m_{b}\|_{\ell^{p} \to \ell^{p}} = |b_{K}| = \max_{1 \leq j \leq N} |b_{j}| = \|b\|_{\ell^{\infty}}$.
\end{example}

在考虑下一个例子时, 将 $N$ 阶单位矩阵记为 $E$, 来复习一下矩阵的知识. 正交矩阵是指由 $PP^{t} = P^{t}P = E$ 给出的实矩阵. 对所有数列 (向量) $a$, 正交矩阵 $P$ 满足 $\|Pa\|_{\ell^{2}} = \|a\|_{\ell^{2}}$. 对称矩阵是指满足 $A = A^{t}$ 的实矩阵.

\begin{example}\label{exa:1.7}
设 $A$ 为 $N$ 阶实对称矩阵, 则用正交矩阵 $P$ 可将其对角化为 $P^{-1}AP = \mathrm{diag}(\alpha_{1}, \alpha_{2}, \ldots, \alpha_{N})$. 其中设 $|\alpha_{1}| \geq |\alpha_{2}| \geq \cdots \geq |\alpha_{N}|$. 设 $A$, $P^{-1}$, $P^{-1}AP$ 给出的 $\R^{N}$ 上的线性映射分别为 $f$, $f_{P^{-1}}$, $f_{P^{-1}AP}$, 则
\[
\frac{\| f(a) \|_{\ell^{2}}}{\| a \|_{\ell^{2}}} = \frac{\| f_{P^{-1}AP}(f_{P^{-1}}(a)) \|_{\ell^{2}}}{\| f_{P^{-1}}(a) \|_{\ell^{2}}},
\]
所以由例~\ref{exa:1.6} 求得 $\|f\|_{\ell^{2} \to \ell^{2}} = |\alpha_{1}|$.
\end{example}

\begin{exercise}\label{exer:1.10}
在 $\R^{2}$ 上配备范数 $\|(x, y)\| = \sqrt{x^{2} + y^{2}}$. 计算下列算子的范数.

(1) $(x, y) \in \R^{2} \mapsto (y, x) \in \R^{2}.$

(2) $(x, y) \in \R^{2} \mapsto (2x, 2y) \in \R^{2}.$

(3) $(x, y) \in \R^{2} \mapsto (x + y, \sqrt{2}x) \in \R^{2}.$
\end{exercise}

\begin{exercise}\label{exer:1.11}
设 $\R^{2}$ 为配备 $\|(x, y)\| = \sqrt{x^{2} + y^{2}}$ 的 Euclid 空间. 线性映射 $f: \R^{2} \to \R^{2}$ 由 $f(x, y) = (4x - 2y, 8x - 4y)$ 给出.

(1) 计算 $f$ 的算子范数 $\|f\|$.

(2) 计算 $f^{2}$ 的算子范数 $\|f^{2}\|$.

(3) 考察 $\|f^{2}\|$ 与 $\|f\|^{2}$ 的大小关系.
\end{exercise}

\begin{exercise}\label{exer:1.12}
线性映射 $f: \R^{2} \to \R^{2}$ 由 $f(x, y) = (4x - 3y, x - 4y)$ 给出. 求下列情形下 $f$ 的算子范数.

(1) 定义域与值域都配备 $\ell^{2}$-范数的情形.

(2) 定义域与值域都配备 $\ell^{1}$-范数的情形.

(3) 定义域配备 $\ell^{2}$-范数, 值域配备 $\ell^{\infty}$-范数的情形.

(4) 定义域配备 $\ell^{\infty}$-范数, 值域配备 $\ell^{2}$-范数的情形.
\end{exercise}

\begin{exercise}\label{exer:1.13}
求配备 $\ell^{2}$-范数的 $\R^{2}$ 之间的映射 $f(x, y) = (2x - 3y, 3x + 4y)$ 的算子范数. 若定义域配备 $\ell^{\infty}$-范数, 值域配备 $\ell^{2}$-范数, 则 $f$ 的算子范数是多少?
\end{exercise}

\begin{exercise}\label{exer:1.14}
求从配备 $\ell^{2}$-范数的 $\R^{2}$ 到配备 $\ell^{2}$-范数的 $\R^{3}$ 的映射 $f(x, y) = (2x + y, x - 2y, x + y)$ 的算子范数.
\end{exercise}

\begin{exercise}\label{exer:1.15}
求从配备 $\ell^{3}$-范数的 $\R^{2}$ 到 $\R$ 的映射 $f(x, y) = 2x + y$ 的算子范数.
\end{exercise}

\begin{exercise}\label{exer:1.16}
求函数 $f(x, y) = \sqrt[3]{|x+1|^{3} + |y-4|^{3}} - \sqrt[3]{|x+2|^{3} + |y-3|^{3}}$, $x, y \in \R$ 的最大值.
\end{exercise}

% ============ 第1章 §1.4 插值理论 (原书 pp.14-19) ============
\section{插值理论}\label{sec:1.4}
\index{M@Marcinkiewicz 插值定理}




解析学为何复杂, 理由之一是调和级数发散. 也就是说,
\[
\frac{1}{1} + \frac{1}{2} + \cdots + \frac{1}{n} + \cdots = \infty,
\]
所以各种问题就变得困难. 关于这个级数为何复杂, 一个说明是: 若 $1 < p < \infty$,
\[
\frac{1}{1^{p}} + \frac{1}{2^{p}} + \cdots + \frac{1}{n^{p}} + \cdots < \infty.
\]
也就是说, $p = 1$ 的情形是特别的. 更进一步说, 可以认为 $p = 1$ 的情形是病态的. 然而, 在解析学中, 与
\[
\frac{1}{1} + \frac{1}{2} + \cdots + \frac{1}{n} + \cdots = \infty
\]
类似的现象, 例如
\[
\int_{1}^{\infty} \frac{\dd t}{t} = [\log t]_{1}^{\infty} = \infty
\]
会出现, 所以调和级数发散这一事实是重要的.

来稍微定量地把握调和级数发散这一问题. 具体地, 用数学归纳法可以证明
\[
\frac{1}{1} + \frac{1}{2} + \cdots + \frac{1}{n} > \log(n+1) \quad (n \in \N).
\]
这里重要的是, 调和级数虽然发散, 但其发散程度非常弱. 也就是说, 第 $n$ 部分和大于 $\log n$, 但并不是大于 $n$ 的任意幂 $cn^{a}$, $a, c > 0$, 这是应该注意的事实. 也就是说, 虽然
\[
\frac{1}{1} + \frac{1}{2} + \cdots + \frac{1}{n} + \cdots = \infty,
\]
但这个无穷大在某种意义上是很弱的发散, 或许应该称之为``准有限''. 这里, 我们说明: 当有两个``准有限''时, 能够产生有限的插值理论. 下面给出的定理是被称为 Marcinkiewicz 插值定理的 1930 年代的定理.

设 $S$ 为从 $\R^{N}$ 到 $\R^{N}$ 的实线性映射. 将 $a \in \R^{N}$ 经 $S$ 的像记为 $Sa = \{(Sa)_{n}\}_{n=1}^{N}$. 对 $t > 0$, 假定
\[
\# \{1 \leq n \leq N : |(Sa)_{n}| > t\} \leq \min(t^{-1}\|a\|_{\ell^{1}}, t^{-q}\|a\|_{\ell^{q}}^{q})
\]
成立. 来考察
\[
\|Sa\|_{\ell^{p}}^{p} = p\int_{0}^{\infty} t^{p-1} \# \{1 \leq n \leq N : |(Sa)_{n}| > t\} \dd t.
\]
对 $a \in \R^{N}$ 与 $t > 0$, 将 $c^{t} = \{c_{k}^{t}\}_{k=1}^{N}$ 定义为: 在 $a$ 的分量中, 绝对值大于 $t$ 的保持原样, 绝对值不超过 $t$ 的替换为 $0$. 进一步, 令 $b^{t} = a - c^{t}$. 此时, 由于 $(Sa)_{n} = (Sb^{t})_{n} + (Sc^{t})_{n}$, 由三角不等式,
\begin{align*}
\{1 \leq n \leq N : |(Sb^{t})_{n}| \leq 2^{-1}t\} \cap \{1 \leq n \leq N : |(Sc^{t})_{n}| \leq 2^{-1}t\} \subset \{1 \leq n \leq N : |(Sa)_{n}| \leq t\}.
\end{align*}
因此, 设 $1 < p < q$, 由假设得
\begin{align*}
& \# \{1 \leq n \leq N : |(Sa)_{n}| > t\} \\
& \leq \# \{1 \leq n \leq N : |(Sb^{t})_{n}| > 2^{-1}t\} + \# \{1 \leq n \leq N : |(Sc^{t})_{n}| > 2^{-1}t\} \\
& \leq 2^{q}t^{-q}\|b^{t}\|_{\ell^{q}}^{q} + 2t^{-1}\|c^{t}\|_{\ell^{1}}.
\end{align*}
这里,
\begin{align*}
\int_{0}^{\infty} t^{p-2}\|c^{t}\|_{\ell^{1}} \dd t & = \int_{0}^{\infty} t^{p-2}\sum_{j=1}^{N} |a_{j}|\chi_{(t,\infty)}(|a_{j}|) \dd t \\
& = \sum_{j=1}^{N} |a_{j}|\int_{0}^{\infty} t^{p-2}\chi_{(t,\infty)}(|a_{j}|) \dd t \\
& = \frac{1}{p-1}\sum_{j=1}^{N} |a_{j}|^{p}
\end{align*}
以及
\begin{align*}
\int_{0}^{\infty} t^{p-q-1}\|b^{t}\|_{\ell^{q}}^{q} \dd t & = \int_{0}^{\infty} t^{p-q-1}\sum_{j=1}^{N} |a_{j}|^{q}\chi_{[0,t]}(|a_{j}|) \dd t \\
& = \sum_{j=1}^{N} |a_{j}|^{q}\int_{0}^{\infty} t^{p-q-1}\chi_{[0,t]}(|a_{j}|) \dd t \\
& = \sum_{j=1}^{N} |a_{j}|^{q}\int_{|a_{j}|}^{\infty} t^{p-q-1} \dd t \\
& = \frac{1}{q-p}\sum_{j=1}^{N} |a_{j}|^{p}
\end{align*}
所以,
\begin{equation}\tag{1.8}\label{eq:1.8}
\|Sa\|_{\ell^{p}} \leq \left(\frac{2p}{p-1} + \frac{2^{q}p}{q-p}\right)^{\frac{1}{p}}\|a\|_{\ell^{p}}
\end{equation}
成立. 由于 $a$ 是任意的,
\[
\|S\|_{\ell^{p} \to \ell^{p}} \leq \left(\frac{2p}{p-1} + \frac{2^{q}p}{q-p}\right)^{\frac{1}{p}}.
\]
这样, 从关于 $\ell^{1}$, $\ell^{q}$ 已知的不等式得到 $\ell^{p}$, $1 < p < q$ 的不等式的理论, 称为插值理论.

当 $q = \infty$ 时, 将假设改写为 $\|Sa\|_{\ell^{\infty}} \leq 2^{-1}\|a\|_{\ell^{\infty}}$ 以及
\[
\# \{1 \leq n \leq N : |(Sa)_{n}| > t\} \leq t^{-1}\|a\|_{\ell^{1}}.
\]
此时要使用
\begin{align*}
\{1 \leq n \leq N : |(Sb^{t})_{n}| > 2^{-1}t\} & = \{1 \leq n \leq N : 2^{-1}t \geq |(Sb^{t})_{n}| > 2^{-1}t\} \\
& = \emptyset.
\end{align*}
无论 $q = \infty$ 还是 $q \neq \infty$, 都可以在
\[
\# \{1 \leq n \leq N : |(Sa)_{n}| > t\} \leq \min(t^{-1}\|a\|_{\ell^{1}}, t^{-q}\|a\|_{\ell^{q}}^{q})
\]
或 $\|Sa\|_{\ell^{\infty}} \leq \|a\|_{\ell^{\infty}}$ 以及
\[
\# \{1 \leq n \leq N : |(Sa)_{n}| > t\} \leq t^{-1}\|a\|_{\ell^{1}}
\]
的位置, 改用对适当的常数 $K > 0$ 的不等式
\[
\# \{1 \leq n \leq N : |(Sa)_{n}| > t\} \leq K\min(t^{-1}\|a\|_{\ell^{1}}, t^{-q}\|a\|_{\ell^{q}}^{q})
\]
或 $\|Sa\|_{\ell^{\infty}} \leq K\|a\|_{\ell^{\infty}}$ 以及
\[
\# \{1 \leq n \leq N : |(Sa)_{n}| > t\} \leq Kt^{-1}\|a\|_{\ell^{1}}.
\]
此时, 根据 $q$ 是有限还是无限, 使用 $K^{-1}S$ 或 $(2K)^{-1}S$ 来使用插值定理.

\begin{theorem}[Marcinkiewicz 插值定理]\label{thm:1.5}
(1) 设 $1 < p < q < \infty$. 对任意 $t > 0$ 与数列 $\{a_{n}\}_{n=1}^{N}$, 不等式
\[
\# \{1 \leq n \leq N : |(Sa)_{n}| > t\} \leq \min(Kt^{-1}\|a\|_{\ell^{1}}, K^{-q}t^{-q}\|a\|_{\ell^{q}}^{q})
\]
成立的线性算子 $S: \R^{N} \to \R^{N}$, 存在仅依赖于 $p$, $q$ 的适当常数 $c_{p,q} > 0$, 使得 $\|S\|_{\ell^{p} \to \ell^{p}} \leq c_{p,q}K$ 成立.

(2) 设 $1 < p < \infty$. 对满足不等式 $\|S\|_{\ell^{\infty} \to \ell^{\infty}} \leq K$ 且对任意数列 $\{a_{n}\}_{n=1}^{N}$,
\[
\# \{1 \leq n \leq N : |(Sa)_{n}| > t\} \leq Kt^{-1}\|a\|_{\ell^{1}} \quad (t > 0)
\]
成立的线性算子 $S: \R^{N} \to \R^{N}$, 存在仅依赖于 $p$ 的适当常数 $c_{p} > 0$, 使得 $\|S\|_{\ell^{p} \to \ell^{p}} \leq c_{p}K$ 成立.
\end{theorem}

来举出满足定理假设的算子的具体例子.

\begin{example}\label{exa:1.8}
(1) 设数列 $a = \{a_{n}\}_{n=1}^{N}$ 已给出. 对满足 $|a_{n}| > t$ 的 $n = 1, 2, \ldots, N$, 有 $t^{-p}|a_{n}|^{p} > 1$. 因此,
\[
\# \{1 \leq n \leq N : |a_{n}| > t\} = \sum_{n=1}^{N} \chi_{(t,\infty)}(|a_{n}|) \leq t^{-p}\|a\|_{\ell^{p}}^{p}
\]
成立.

(2) 设 $1 \leq p \leq \infty$. 定义域与值域都考虑 $\ell^{p}$-范数, 当线性算子 $S: \R^{N} \to \R^{N}$ 的范数不超过 $1$ 时, 对所有数列 $a$, $\|Sa\|_{\ell^{p}} \leq \|a\|_{\ell^{p}}$ 成立, 从而 $\# \{1 \leq n \leq N : |(Sa)_{n}| > t\} \leq t^{-p}\|a\|_{\ell^{p}}^{p}$ 成立. 特别地, 若对 $p = 1$, $q$ 双方这个条件都满足, 则 Marcinkiewicz 插值定理的假设成立.
\end{example}

\begin{example}\label{exa:1.9}
对实数列 $\{a_{n}\}_{n=1}^{2N}$, 令
\[
b_{n} = \frac{a_{1} + a_{2} + \cdots + a_{2n}}{2n}, \quad (n = 1, 2, \ldots, N), \quad b_{N+1} = \cdots = b_{2N} = 0.
\]
考虑对应 $T_{N}: \{a_{n}\}_{n=1}^{2N} \mapsto \{b_{n}\}_{n=1}^{2N}$.

(1) 对 $n = 1, 2, \ldots, 2N$, $|b_{n}| \leq \frac{1}{2n}\|a\|_{\ell^{1}}$ 成立, 所以
\[
\# \left\{k \in \{1, 2, \ldots, 2N\} : |b_{k}| > \frac{1}{2n}\|a\|_{\ell^{1}}\right\} \leq \# \{1, 2, \ldots, n-1\} = n-1
\]
成立.

(2) 若 $\lambda \geq \frac{1}{2}$, 则出现在左边的集合是空集, 所以
\[
\# \left\{k \in \{1, 2, \ldots, 2N\} : |b_{k}| > \lambda\|a\|_{\ell^{1}}\right\} = \#\emptyset = 0.
\]

(3) 另外, 由于出现在左边的集合是 $\{1, 2, \ldots, 2N\}$ 的子集, 若 $\lambda \leq \frac{1}{2N}$,
\[
\# \left\{k \in \{1, 2, \ldots, 2N\} : |b_{k}| > \lambda\|a\|_{\ell^{1}}\right\} \leq 2N \leq \frac{1}{\lambda}.
\]

(4) 若 $\frac{1}{2N} < \lambda < \frac{1}{2}$, 则存在 $n = 1, 2, \ldots, N-1$ 使得 $\frac{1}{2(n+1)} \leq \lambda < \frac{1}{2n}$, 所以
\begin{align*}
\# \left\{k \in \{1, 2, \ldots, 2N\} : |b_{k}| > \lambda\|a\|_{\ell^{1}}\right\} & \leq \# \left\{k \in \{1, 2, \ldots, 2N\} : |b_{k}| > \frac{1}{2n+2}\|a\|_{\ell^{1}}\right\} \\
& = \# \{1, 2, \ldots, n\} = n \leq 2n \leq \frac{1}{\lambda}.
\end{align*}
综上所述, 对任意 $\lambda > 0$,
\[
\# \left\{k \in \{1, 2, \ldots, 2N\} : |b_{k}| > \lambda\|a\|_{\ell^{1}}\right\} \leq \frac{1}{\lambda}
\]
成立. 也就是说, $T_{N}$ 满足定理~\ref{thm:1.5} 的假设.
\end{example}

来实际说明用定理能得到什么.

\begin{example}\label{exa:1.10}
设 $1 \leq p \leq \infty$. 在定义域、值域 $\R^{2N}$ 上配备 $\ell^{p}$-范数, 考虑例~\ref{exa:1.9} 中考虑的算子 $T_{N}$.

(1) 设 $p = 1$. 对实数列 $a = \{a_{n}\}_{n=1}^{2N}$,
\begin{align*}
\|T_{N}a\|_{\ell^{1}} & = \sum_{n=1}^{N} \frac{|a_{1} + a_{2} + \cdots + a_{2n}|}{2n} \\
& \leq \sum_{n=1}^{N} \frac{|a_{1}| + |a_{2}| + \cdots + |a_{2n}|}{2n} \\
& \leq \sum_{n=1}^{N} \frac{|a_{1}| + |a_{2}| + \cdots + |a_{2N}|}{2n} \\
& = \sum_{n=1}^{N} \frac{\|a\|_{\ell^{1}}}{2n}
\end{align*}
由于存在 $a = e_{1}$ 使上式等号成立, 所以 $\|T_{N}\|_{\ell^{1} \to \ell^{1}} = \sum_{n=1}^{N} \frac{1}{2n}$.

(2) 设 $p = \infty$. 平均不会增大 $\ell^{\infty}$-范数, 所以 $\|T_{N}\|_{\ell^{\infty} \to \ell^{\infty}} = 1$ 成立.

(3) 设 $1 < p < \infty$. 由例~\ref{exa:1.9} 与 Marcinkiewicz 插值定理, 存在不依赖于 $N$ 的常数 $c_{p}$, 使得 $\|T_{N}\|_{\ell^{p} \to \ell^{p}} \leq c_{p}$ 成立.

(4) 实际计算可知, $\|T_{1}\|_{\ell^{2} \to \ell^{2}} = \frac{1}{\sqrt{2}}$, $\|T_{2}\|_{\ell^{2} \to \ell^{2}} = \frac{1 + \sqrt{5}}{4}$ (参照例~\ref{exa:1.5}), $\|T_{3}\|_{\ell^{2} \to \ell^{2}} = 0.86\cdots$. 这样, 对 $j = 1, 2, \ldots$, $\|T_{j}\|_{\ell^{2} \to \ell^{2}}$ 递增, 但可以看出不会超过一定值.
\end{example}

\begin{exercise}\label{exer:1.17}
对实数列 $\{a_{n}\}_{n=1}^{N}$, 令
\[
b_{n} = \frac{a_{1} + a_{2} + \cdots + a_{n}}{n}, \quad (n = 1, 2, \ldots, N),
\]
考虑数列 $\{b_{n}\}_{n=1}^{N}$. 设 $1 \leq p \leq \infty$. 在定义域、值域 $\R^{N}$ 上配备 $\ell^{p}$-范数, 考虑这个对应 $U_{N}: \{a_{n}\}_{n=1}^{N} \mapsto \{b_{n}\}_{n=1}^{N}$.

(1) 证明 $\|U_{N}\|_{\ell^{1} \to \ell^{1}} = \sum_{n=1}^{N} \frac{1}{n}$.

(2) 证明 $\|U_{N}\|_{\ell^{\infty} \to \ell^{\infty}} = 1$.

(3) 设 $\lambda > 0$. 证明 $\# \{k \in \{1, 2, \ldots, N\} : |b_k| > \lambda\|a\|_{\ell^{1}}\} \leq \frac{1}{\lambda}$.

(4) 设 $1 < p < \infty$. 证明存在不依赖于 $N$ 的常数 $c_{p} > 0$, 使得 $\|U_{N}\|_{\ell^{p} \to \ell^{p}} \leq c_{p}$ 成立.

(5) 实际计算并比较 $\|U_{1}\|_{\ell^{2} \to \ell^{2}}$, $\|U_{2}\|_{\ell^{2} \to \ell^{2}}$, $\|U_{3}\|_{\ell^{2} \to \ell^{2}}$. 其中, 关于 $\|U_{3}\|_{\ell^{2} \to \ell^{2}}$, 只要导出能看出这三者大小关系程度的信息即可.
\end{exercise}

% ============ 第1章 §1.5 Rademacher 序列 (原书 pp.19-24) ============
\section{Rademacher 序列}\label{sec:1.5}
\index{0symbolrj@$r_{j}$}

\index{K@Khintchine 不等式}

\index{R@Rademacher 序列}








这里想做的事是 $\ell^{2}$-范数的变形. $a = (a_{1}, a_{2}, \ldots, a_{N})$ 的 $\ell^{2}$-范数由
\[
\| a \|_{\ell^{2}} = \sqrt{\sum_{j=1}^{N} |a_{j}|^{2}}
\]
给出, 但因为有平方根, 在某种意义上是难处理的. 若使用 $\ell^{2}$ 内积, 平方根在表面上就消失了. 实际上, 以 $e_{j} = \{\delta_{j,k}\}_{k=1}^{N}$ 为基本向量,
\[
\| a \|_{\ell^{2}} = \left\| \sum_{j=1}^{N} a_{j}\mathbf{e}_{j} \right\| = \| (a_{1}, a_{2}, \ldots, a_{N}) \|,
\]
就是这个缘故. 这个 $\sum_{j=1}^{N} a_{j}e_{j}$ 的表示关于 $a$ 是线性的, 可以说是便于处理的形式. 这个想法不能再发展一些吗?

一般地, 关于数列的范数和算子范数, 从经验上可以说出下面的事实. 首先, 范数和算子范数有时能求出精确值, 但即使能求出, 其方法也多为艰深的. 其次, 即使范数的值不能精确求出, 也有时能知道一定程度的信息. 实际上, 我们已经看到, 用 Marcinkiewicz 插值定理等, 能求出算子范数的上界. 这里, 基于这个经验法则, 来探索使
\[
C^{-1}\| a \|_{\ell^{2}} \leq ??? \leq C\| a \|_{\ell^{2}}
\]
成立的 $???$ 是什么.

在概率论中出现的作为基本工具的 Rademacher 序列, 在调和分析中也非常强大. 这里, 我们来考察 Rademacher 序列的性质, 并将其应用于函数空间的性质, 给出前面问题的解答.

\begin{definition}[Rademacher 序列]\label{def:1.3}
函数列 $r_{j}$, $j = 1, 2, \ldots$ 由
\[
r_{j}(t) \equiv \sum_{k=1}^{2^{j}} (-1)^{k-1}\chi_{[(k-1)2^{-j}, k2^{-j})}(t) \quad (t \in \R)
\]
定义.
\end{definition}

\begin{example}\label{exa:1.11}
设 $t \in \R$.

(1) $r_{1}(t) = \chi_{[0,1/2)}(t) - \chi_{[1/2,1)}(t)$.

(2) $r_{2}(t) = \chi_{[0,1/4)}(t) - \chi_{[1/4,1/2)}(t) + \chi_{[1/2,3/4)}(t) - \chi_{[3/4,1)}(t)$.
\end{example}

积分在 $[0, 1)$ 上考虑, 注意下面的性质.

\begin{proposition}\label{prop:1.1}
设 $N \in \N$. 对函数列 $f_{j}: \{-1, 1\} \to \R$, $j = 1, 2, \ldots, N$,
\[
\int_{0}^{1} \prod_{j=1}^{N} f_{j}(r_{j}(t)) \dd t = \prod_{j=1}^{N} \int_{0}^{1} f_{j}(r_{j}(t)) \dd t
\]
成立.
\end{proposition}

\begin{proof}
由积分的可加性,
\[
\int_{0}^{1} \prod_{j=1}^{N} f_{j}(r_{j}(t)) \dd t = \sum_{k=1}^{2^{N}} \int_{2^{-N}(k-1)}^{2^{-N}k} \prod_{j=1}^{N} f_{j}(r_{j}(t)) \dd t,
\]
但在 $[2^{-N}(k-1), 2^{-N}k)$ 上函数 $\prod_{j=1}^{N} f_{j}(r_{j}(t))$ 是常数, 所以特别地
\[
\prod_{j=1}^{N} f_{j}(r_{j}(t)) = \prod_{j=1}^{N} f_{j}(r_{j}(2^{-N}(k-1))).
\]
因此,
\begin{align*}
\int_{0}^{1} \prod_{j=1}^{N} f_{j}(r_{j}(t)) \dd t & = \sum_{k=1}^{2^{N}} \prod_{j=1}^{N} f_{j}(r_{j}(2^{-N}(k-1)))\int_{2^{-N}(k-1)}^{2^{-N}k} \dd t \\
& = 2^{-N}\sum_{k=1}^{2^{N}} \prod_{j=1}^{N} f_{j}(r_{j}(2^{-N}(k-1))).
\end{align*}
由于
\[
\{1, -1\}^{N} = \left\{\{r_{j}(2^{-N}(k-1))\}_{j=1}^{N} : k = 1, 2, \ldots, 2^{N}\right\},
\]
用 $m = (m_{1}, m_{2}, \ldots, m_{N})$ 表示, 实际计算得
\begin{align*}
\int_{0}^{1} \prod_{j=1}^{N} f_{j}(r_{j}(t)) \dd t & = \sum_{m \in \{-1, 1\}^{N}} 2^{-N}\prod_{j=1}^{N} f_{j}(m_{j}) \\
& = \prod_{j=1}^{N} \frac{f_{j}(1) + f_{j}(-1)}{2} \\
& = \prod_{j=1}^{N} \int_{0}^{1} f_{j}(r_{j}(t)) \dd t.
\end{align*}
\end{proof}

$\{r_{j}\}_{j=1}^{N}$ 在下面的意义下是正交的. 这也是调和分析学中正负相消的一个例子.

\begin{lemma}\label{lem:1.2}
对任意复数列 $a = \{a_{j}\}_{j=1}^{N}$,
\[
\| a \|_{\ell^{2}} = \left(\int_{0}^{1} \left| \sum_{j=1}^{N} a_{j}r_{j}(t) \right|^{2} \dd t\right)^{\frac{1}{2}}
\]
成立.
\end{lemma}

\begin{proof}
用对 $j, j' = 1, 2, \ldots, N$ 的 Kronecker 记号 $\delta_{j,j'}$ 来证明. 由于 $\{r_{j}\}_{j=1}^{N}$ 满足
\[
\int_{0}^{1} r_{j}(t)r_{j'}(t) \dd t = \delta_{j,j'},
\]
\begin{align*}
\int_{0}^{1} \left| \sum_{j=1}^{N} a_{j}r_{j}(t) \right|^{2} \dd t & = \int_{0}^{1} \sum_{j=1}^{N}\sum_{k=1}^{N} a_{j}r_{j}(t)\overline{a_{k}}r_{k}(t) \dd t \\
& = \sum_{j=1}^{N}\sum_{k=1}^{N} a_{j}\overline{a_{k}}\int_{0}^{1} r_{j}(t)r_{k}(t) \dd t \\
& = \sum_{j=1}^{N} |a_{j}|^{2} \\
& = \| a \|_{\ell^{2}}^{2}
\end{align*}
得证.
\end{proof}

因为 $\cosh t$ 满足 $\cosh t = \frac{e^{t} + e^{-t}}{2}$, 所以正如通过 $e^{t}$, $e^{-t}$ 与 $\cosh t$ 的泰勒展开所能确认的那样,
\begin{equation}\tag{1.9}\label{eq:1.9}
|t|^{p} \leq \max(|t|^{[p]}, |t|^{[p+1]}) \leq (2[p+1])! \cosh t \leq (2[p+1])! e^{t^{2}}
\end{equation}
成立 (参照练习题~\ref{exer:1.19}). 利用它来证明作为变形 $\ell^{p}$-范数的工具的 Khintchine 不等式.

\begin{theorem}[Khintchine 不等式]\label{thm:1.6}
对 $N \in \N$, $1 \leq p < \infty$, 存在不依赖于 $N$ 的常数 $c_{p} > 1$, 使得对任意复数列 $a = \{a_{j}\}_{j=1}^{N}$,
\begin{equation}\tag{1.10}\label{eq:1.10}
c_{p}^{-1}\| a \|_{\ell^{2}} \leq \left(\int_{0}^{1} \left| \sum_{j=1}^{N} a_{j}r_{j}(t) \right|^{p} \dd t\right)^{\frac{1}{p}} \leq c_{p}\| a \|_{\ell^{2}}
\end{equation}
成立.
\end{theorem}

\begin{proof}
来证明式~\eqref{eq:1.10} 中右侧的不等式. 归一化为 $\|a\|_{\ell^{2}} = 1$. 由式~\eqref{eq:1.9} 左侧的不等式与引理~\ref{lem:1.1}, 结合 $r_{j}$ 对称地分布在 $\pm 1$ 上这一事实,
\begin{align*}
\int_{0}^{1} \left| \sum_{j=1}^{N} a_{j}r_{j}(t) \right|^{p} \dd t & \leq (2[p+1])! \int_{0}^{1} \cosh\left(\sum_{j=1}^{N} a_{j}r_{j}(t)\right) \dd t \\
& = (2[p+1])! \int_{0}^{1} \exp\left(\sum_{j=1}^{N} a_{j}r_{j}(t)\right) \dd t \\
& = (2[p+1])! \prod_{j=1}^{N} \int_{0}^{1} \exp(a_{j}r_{j}(t)) \dd t.
\end{align*}
这里, 与刚才一样, 由 $r_{j}$ 分布的对称性,
\begin{align*}
\int_{0}^{1} \exp(a_{j}r_{j}(t)) \dd t & = \int_{0}^{1} \exp(a_{j}r_{1}(t)) \dd t \\
& = \int_{0}^{\frac{1}{2}} \exp(a_{j}) \dd t + \int_{\frac{1}{2}}^{1} \exp(-a_{j}) \dd t \\
& = \cosh(a_{j})
\end{align*}
需要注意. 进一步, 由式~\eqref{eq:1.9} 右侧的不等式与 Rademacher 列的对称性,
\begin{align*}
\left(\int_{0}^{1} \left| \sum_{j=1}^{N} a_{j}r_{j}(t) \right|^{p} \dd t\right)^{\frac{1}{p}} & \leq \sqrt[p]{(2[p+1])!}\left(\prod_{j=1}^{N} \cosh a_{j}\right)^{\frac{1}{p}} \\
& \leq \sqrt[p]{(2[p+1])!}\prod_{j=1}^{N} e^{a_{j}^{2}/p} \\
& = \sqrt[p]{(2[p+1])! e}.
\end{align*}
这样, 式~\eqref{eq:1.10} 右侧的证明就完成了.

在证明式~\eqref{eq:1.10} 左侧的不等式时, 由式~\eqref{eq:1.7} 可假定 $p = 1$. 由引理~\ref{lem:1.2} 与式~\eqref{eq:1.6},
\begin{align*}
\| a \|_{\ell^{2}} & \leq \left(\int_{0}^{1} \left| \sum_{j=1}^{N} a_{j}r_{j}(t) \right| \dd t\right)^{\frac{1}{3}}\left(\int_{0}^{1} \left| \sum_{j=1}^{N} a_{j}r_{j}(t) \right|^{4} \dd t\right)^{\frac{1}{6}} \\
& \leq c_{4}^{\frac{2}{3}}\left(\int_{0}^{1} \left| \sum_{j=1}^{N} a_{j}r_{j}(t) \right| \dd t\right)^{\frac{1}{3}}\| a \|_{\ell^{2}}^{\frac{2}{3}}
\end{align*}
得. 将其整理, 就得到 $p = 1$ 时 (因此, 也是 $p \geq 1$ 时的) 式~\eqref{eq:1.10} 左侧的不等式.
\end{proof}

\begin{exercise}\label{exer:1.18}
求积分 $\int_{0}^{\frac{5}{8}} r_{1}(t)r_{2}(t)r_{3}(t) \dd t$.
\end{exercise}

\begin{exercise}\label{exer:1.19}
对实数 $t$, 通过比较两边的泰勒展开, 证明 $2e^{t^{2}} \geq e^{t} + e^{-t}$.
\end{exercise}

\begin{exercise}\label{exer:1.20}
对实数列 $\{a_{j}\}_{j=1}^{N}$, 求 $\max_{t \in \R}\left|\sum_{j=1}^{N} a_{j}r_{j}(t)\right|$. 将 Khintchine 不等式的各个不等式作比较, 也是一件趣事吧. 【提示】取 $N = 3$, 分 $a_{1}, a_{2}, a_{3}$ 的符号来考虑.
\end{exercise}

% ============ 第1章 §1.6 Calderón-Zygmund 分解 (原书 pp.24-29) ============
\section{Calderón--Zygmund 分解}\label{sec:1.6}
\index{0symbolbad@bad part}

\index{0symbolgood@good part}

\index{E@二进制区间}








这里来解说 Calderón--Zygmund 分解. 正如其名所示, 这个分解是由 Calderón 与 Zygmund 构思的分解. 给定实数列 $a = \{a_{m}\}_{m=1}^{N}$ 时, 考虑截断操作. 也就是说, 在 $a_{m}$, $m = 1, 2, \ldots, N$ 中, 若绝对值大于给定的数 $K > 0$, 则将其替换为 $0$, 这样得到的数列记为 $c = \{c_{m}\}_{m=1}^{N}$. 对 $m = 1, 2, \ldots, N$, 令 $b_{m} = a_{m} - c_{m}$. 这样, 数列 $a$ 就借助 $b = \{b_{m}\}_{m=1}^{N}$ 与 $c$ 分解为 $a = b + c$. $c$ 在一定程度上继承了数列 $a$ 的性质, 但对于 $b$ 则什么也不知道. 若只要求 $c$ 的 $\ell^{\infty}$-范数不超过 $K$, 那么在 $a_{m}$, $m = 1, 2, \ldots, N$ 中, 当绝对值大于给定的数 $K > 0$ 时, 保持符号, 将绝对值替换为 $K$, 其余分量不作特别变更, 这样得到的数列 $d$ 也应该能达到目的. 不过, 若令 $a = d + e$, 则 $e$ 的结构无法再进一步研究, 这一方面与前面的 $b$ 相同. 于是, Calderón 与 Zygmund 提出了研究这个 $b$ 的结构的手段. 这就是 Calderón--Zygmund 分解.

这里以 $N$ 为自然数, 考虑下标从 $0$ 到 $2^{N}-1$ 的长度为 $2^{N}$ 的数列. 给定 $j = 0, 1, \ldots, N$ 与 $k = 0, 1, \ldots, 2^{N}-1$ 时, 定义
\begin{equation}\tag{1.14}\label{eq:1.14}
I(j, k) = \Z \cap [2^{j}k, 2^{j}(k+1)).
\end{equation}
称 $I(j,k)$ 为二进制区间. 另外, 当 $j = 1, 2, \ldots, N$ 时, 令
\[
I^{-}(j, k) = \Z \cap [2^{j}k, 2^{j}(k + 2^{-1})) = I(j-1, 2k),
\]
\[
I^{+}(j, k) = \Z \cap [2^{j}(k + 2^{-1}), 2^{j}(k+1)) = I(j-1, 2k+1).
\]
$I^{\pm}(j,k)$ 也是二进制区间.

\begin{example}\label{exa:1.12}
设 $N = 3$. 具体写出 $I(j,k)$, $j, k = 0, 1, 2$.

(1) $I(0,0) = \{0\}$, $I(0,1) = \{1\}$, $I(0,2) = \{2\}$,

(2) $I(1,0) = \{0,1\}$, $I(1,1) = \{2,3\}$, $I(1,2) = \{4,5\}$,

(3) $I(2,0) = \{0,1,2,3\}$, $I(2,1) = \{4,5,6,7\}$, $I(2,2) = \{8,9,10,11\}$.

进一步, 考察 $I^{\pm}(j,k)$, $j = 1, 2$, $k = 0, 1, 2$ 的情形.

(1) $I^{-}(2,0) = I(1,0)$, $I^{-}(2,1) = I(1,2)$, $I^{-}(2,2) = I(1,4)$,

(2) $I^{+}(2,0) = I(1,1)$, $I^{+}(2,1) = I(1,3)$, $I^{+}(2,2) = I(1,5)$.
\end{example}

下面的定理通过在数直线上画出二进制区间即可明白.

\begin{theorem}\label{thm:1.7}
当 $j, j' = 0, 1, \ldots, N$, $k, k' = 0, 1, \ldots, 2^{N}-1$ 时, $I(j,k)$ 与 $I(j',k')$ 互不相交, 或者一方包含另一方.
\end{theorem}

在下面的定理中, 要使用这个二进制区间的几何结构.

\begin{theorem}[Calderón--Zygmund 区间]\label{thm:1.8}
设 $\lambda > 0$ 与 $a = \{a_{j}\}_{j=0}^{2^{N}-1}$ 已给出. 考虑形如 $I(j,k)$ 的集合, 其中绝对值的平均, 即 $\# I(j,k)$ 个数 $\{|a_{l}|\}_{l \in I(j,k)}$ 的平均超过 $\lambda$. 更详细地说, 设下面出现的 $E$, $\mathcal{U}_{0}$ 的定义中的 $j$, $k$ 满足 $j = 0, 1, \ldots, N$, $k = 0, 1, \ldots$, 考虑
\[
\mathcal{U}_{0} = \left\{I(j,k) : \frac{1}{\# I(j,k)}\sum_{l \in I(j,k)} |a_{l}| > \lambda\right\}.
\]
对这个 $\mathcal{U}_{0}$, 令 $E = \bigcup I(j,k)$.

(1) 若 $j \in I(N,0) \setminus E$, 则 $|a_{j}| \leq \lambda$. 因此, $E \neq \emptyset$ 的充分必要条件是存在满足 $|a_{j}| > \lambda$ 的 $j$.

(2) 对 $\lambda > 0$, $\# E \leq \lambda^{-1}\|a\|_{\ell^{1}}$ 成立.
\end{theorem}

\begin{proof}
(1) 证其逆否命题. 若 $|a_{j}| > \lambda$, 则 $I_{0,j} = \Z \cap [2^{0}j, 2^{0}(j+1)) = \{j\} \in \mathcal{U}_{0}$. 因此, $j \in I_{0,j} \subset E$. 关于充分必要条件, 由于取绝对值不超过 $\lambda$ 的数列的各项的绝对值后再取平均仍不超过 $\lambda$, 所以是显然的.

(2) 利用定理~\ref{thm:1.7}, 从构成 $E$ 的属于 $\mathcal{U}_{0}$ 的区间中取一部分, 可将 $E$ 表示为互不相交的二进制区间族 $\mathcal{U} = \{I(j_{l}, k_{l})\}_{l=1}^{L}$ 的并. 对 $l = 1, 2, \ldots, L$,
\[
\#(I(j_{l}, k_{l})) \leq \lambda^{-1}\sum_{m \in I(j_{l}, k_{l})} |a_{m}|,
\]
所以把它们加起来即得结论.
\end{proof}

利用定理~\ref{thm:1.8} 的证明中出现的 $\mathcal{U}$, 来分解 $a$.

\begin{definition}[Calderón--Zygmund 分解]\label{def:1.4}
将定理~\ref{thm:1.8} 的证明中出现的 $E$ 表示为互不相交的二进制区间族 $\mathcal{U} = \{I(j_{l}, k_{l})\}_{l=1}^{L}$ 的并. 进一步, 构成 $E$ 的区间 $I$ 一定包含在 $I(j_{l}, k_{l})$, $l = 1, 2, \ldots, L$ 中的某一个中, 且可以认为真包含 $I(j_{l}, k_{l})$ 的二进制区间 $J$ 不属于 $\mathcal{U}_{0}$. 对属于 $\mathcal{U}$ 中区间 $I$ 的整数 $k$, 将 $a_{k}$ 替换为属于 $I$ 的数列 $\{a_{l}\}_{l \in I}$ 的平均 $\frac{1}{\# I}\sum_{l \in I} a_{l}$, 这样得到的数列记为 $g = \{g_{j}\}_{j=0}^{2^{N}-1}$. 对不属于 $E$ 的 $k \in I(N, 0)$, 令 $g_{k} = a_{k}$. 其中, 当 $\mathcal{U}$ 为空集时, 令 $g = a$. 另外, 令 $b = \{b_{j}\}_{j=0}^{2^{N}-1} = a - g$, $b_{I} = \{\chi_{I}(j)b_{j}\}_{j=0}^{2^{N}-1}$.

称 $g$ 为好的部分 (good part), $b$ 为坏的部分 (bad part).
\end{definition}

\begin{lemma}\label{lem:1.3}
若 $\lambda > 2^{-N}\|a\|_{\ell^{1}}$, 则 $\|g\|_{\ell^{\infty}} \leq 2\lambda$ 成立.
\end{lemma}

\begin{proof}
由定理~\ref{thm:1.8}(1), 对 $j \in E$ 证明 $|g_{j}| \leq 2\lambda$ 就足够了. 由假设,
$\frac{1}{2^{N}}\left|\sum_{j=0}^{2^{N}-1} a_{j}\right| \leq \frac{1}{2^{N}}\|a\|_{\ell^{1}} \leq \lambda$,
所以 $I(N,0) \notin \mathcal{U}$. 因此, 对每个 $I \in \mathcal{U}$, 可以考虑包含 $I$ 且长度为 $2$ 倍的二进制区间 $J$. 由 $J \notin \mathcal{U}_{0}$, $\frac{1}{2\# I}\sum_{j \in I} |a_{j}| \leq \frac{1}{\# J}\sum_{j \in J} |a_{j}| \leq \lambda$, 所以得到结论.
\end{proof}

\begin{lemma}\label{lem:1.4}
若 $b_{j} \neq 0$, 则存在二进制区间 $I$ 使得 $j \in I \in \mathcal{U}$.
\end{lemma}

\begin{proof}
这是因为 $g$ 与 $a$ 的差异出现在 $\bigcup_{I \in \mathcal{U}} I$ 上.
\end{proof}

\begin{example}\label{exa:1.13}
下面的表是对 $a_{0}$ 到 $a_{15}$ 为
\[
0, \ 4, \ 0, \ 4, \ 2, \ 0, \ 6, \ 0, \ 12, \ 0, \ 2, \ 0, \ 6, \ 6, \ 6, \ 0
\]
给出的数列 $a = \{a_{j}\}_{j=0}^{15}$ 进行如下计算得到的表.

\begin{itemize}
\item 在上段写下 16 项的数列.
\item 从左边数起, 对奇数号的数字, 取其与右邻数字的平均, 在该数字下面写下求得的平均.
\item 从左边数起, 对偶数号的数字, 该数字下面留空.
\item 当数字只剩下一个时结束.
\end{itemize}

\begin{center}
\begin{tabular}{cccccccccccccccc}
\hline
$0$ & $4$ & $0$ & $4$ & $2$ & $0$ & $6$ & $0$ & $12$ & $0$ & $2$ & $0$ & $6$ & $6$ & $6$ & $0$ \\
\hline
$2$ &  & $2$ &  & $1$ &  & $3$ &  & $6$ &  & $1$ &  & $6$ &  & $3$ &  \\
\hline
$2$ &  &  &  & $2$ &  &  &  & $3.5$ &  &  &  & $4.5$ &  &  &  \\
\hline
$2$ &  &  &  &  &  &  &  & $4$ &  &  &  &  &  &  &  \\
\hline
$3$ &  &  &  &  &  &  &  &  &  &  &  &  &  &  &  \\
\hline
\end{tabular}
\end{center}

(1) 设 $\lambda = \pi$. 基于定义~\ref{def:1.4} 求 $\mathcal{U}_{0}$, $\mathcal{U}$, $E$, $g$, $b$, $b_{I}$, $I \in \mathcal{U}$, 如下所示.
\begin{align*}
\mathcal{U}_{0} = \{ & \{1\}, \{3\}, \{6\}, \{8\}, \{12\}, \{13\}, \{14\}, \{8,9\}, \{12,13\}, \\
& \{8,9,10,11\}, \{12,13,14,15\}, \{8,9,10,11,12,13,14,15\}\},
\end{align*}
\[
\mathcal{U} = \{\{1\}, \{3\}, \{6\}, \{8,9,10,11,12,13,14,15\}\},
\]
\[
E = \{1, 3, 6, 8, 9, 10, 11, 12, 13, 14, 15\},
\]
\[
g: 0, 4, 0, 4, 2, 0, 6, 0, 4 = \frac{1}{8}\sum_{j=8}^{15} a_{j}, 4, 4, 4, 4, 4, 4, 4
\]
\[
b: 0, 0, 0, 0, 0, 0, 0, 0, 8, -4, -2, -4, 2, 2, 2, -4
\]
\[
b_{\{1\}} = b_{\{3\}} = b_{\{6\}} = 0, \quad b_{\{8,9,10,11,12,13,14,15\}} = b
\]

(2) 设 $\lambda = 4$. 基于定义~\ref{def:1.4} 求 $\mathcal{U}_{0}$, $\mathcal{U}$, $E$, $g$, $b$, $b_{I}$, $I \in \mathcal{U}$, 如下所示.
\[
\mathcal{U}_{0} = \{\{6\}, \{8\}, \{12\}, \{13\}, \{14\}, \{8,9\}, \{12,13\}, \{12,13,14,15\}\},
\]
\[
\mathcal{U} = \{\{6\}, \{8,9\}, \{12,13,14,15\}\}, \quad E = \{6, 8, 9, 12, 13, 14, 15\},
\]
\[
g: 0, 4, 0, 4, 2, 0, 6, 0, 6, 6, 2, 0, 4.5 = \frac{1}{4}\sum_{j=12}^{15} a_{j}, 4.5, 4.5, 4.5
\]
\[
b: 0, 0, 0, 0, 0, 0, 0, 0, 6, -6, 0, 0, 1.5, 1.5, 1.5, -4.5 = a_{15} - g_{15}
\]
\[
b_{\{6\}} = 0,
\]
\[
b_{\{8,9\}}: 0, 0, 0, 0, 0, 0, 0, 0, 6, -6, 0, 0, 0, 0, 0, 0
\]
\[
b_{\{12,13,14,15\}}: 0, 0, 0, 0, 0, 0, 0, 0, 0, 0, 0, 0, 1.5, 1.5, 1.5, -4.5
\]
\end{example}

\begin{exercise}\label{exer:1.21}
设 $\lambda = 2.5$. 对例~\ref{exa:1.13} 的数列, 基于定义~\ref{def:1.4}, 求 $\mathcal{U}_{0}$, $\mathcal{U}$, $E$, $g = \{g_{j}\}_{j=0}^{15}$, $b$, $b_{I}$, $I \in \mathcal{U}$.
\end{exercise}

% ============ 第1章 §1.7 平均算子与差分算子 (原书 pp.23-26) ============

\section{平均算子与差分算子}\label{sec:1.7}
\index{0symbolDk@$D_{k}$}

\index{0symbolEk@$E_{k}$}

\index{C@差分算子}

\index{P@平均算子}










这里同样取 N 为自然数, 考虑长度 $2^{N}$ 的数列 $a = \{a_{j}\}_{j=0}^{2^{N}-1}$.

用 $[a]$ 表示实数 a 的整数部分. 当整数 j 满足 $0 \leq j < 2^{N}$ 时, 对 $k = 0, 1, \ldots, N$, $[2^{-k}j]$ 是整数, 由这个数的定义可知它满足 $2^{k}[2^{-k}j] \leq j < 2^{k}[2^{-k}j] + 2^{k}$. 定义 $E_{k}(a) = \{E_{k}(a)_{j}\}_{j=0}^{2^{N}-1}$ 为
\[
E_{k}(a)_{j} = \frac{1}{2^{k}}\sum_{l = 2^{k}[2^{-k}j]}^{2^{k}[2^{-k}j] + 2^{k} - 1}a_{l}.
\]
也就是说, 将数列 $a = \{a_{j}\}_{j=0}^{2^{N}-1}$ 等分割成长度 $2^{k}$ 的组, 用 $a_{j}$ 所属组的数列的平均替换各 $a_{j}$, 所得的数列记为 $E_{k}(a) = \{E_{k}(a)_{j}\}_{j=0}^{2^{N}-1}$. 另外, 对 $k = 0, 1, \ldots, N-1$, 令 $D_{k}(a) = E_{k}(a) - E_{k+1}(a)$. 对 $k = 0, 1, \ldots, N-1$, 将 $E_{k}, D_{k}, E_{N}$ 视为线性映射. 称 $E_{k}, k = 0, 1, \ldots, N$ 为平均算子, 称 $D_{k}, k = 0, 1, \ldots, N-1$ 为差分算子. 由定义, $E_{0}$ 是恒等算子.

\begin{example}\label{exa:1.14}
对例~\ref{exa:1.13} 中考虑的数列, 求 $E_{1}(a)$, 得
\[
2, 2, 2, 2, 1, 1, 3, 3, 6, 6, 1, 1, 6, 6, 3, 3.
\]
\end{example}

为了能直观理解这里定义的算子的性质, 让我们考虑简单的例子. 对数列 $a = \{a_{j}\}_{j=0}^{2^{N}-1}$, $b = \{b_{j}\}_{j=0}^{2^{N}-1}$, 令 $\langle a, b \rangle = \langle a, b \rangle_{\ell^{2}} = \sum_{j=0}^{2^{N}-1}a_{j}b_{j}$.

\begin{example}\label{exa:1.15}
对由 8 个元组成的数列 $a = \{a_{0}, a_{1}, \ldots, a_{7}\}$ 施行 $E_{1}$, 得
\begin{align*}
E_{1}(a) &= \left\{\frac{a_{0}+a_{1}}{2}, \frac{a_{0}+a_{1}}{2}, \frac{a_{2}+a_{3}}{2}, \frac{a_{2}+a_{3}}{2},\right. \\
&\qquad \left.\frac{a_{4}+a_{5}}{2}, \frac{a_{4}+a_{5}}{2}, \frac{a_{6}+a_{7}}{2}, \frac{a_{6}+a_{7}}{2}\right\}.
\end{align*}
另外, 对原来的数列 a 施行 $E_{1}$ 所得数列 $E_{1}(a)$ 再施行 $E_{2}$, 所得数列 $E_{2} \circ E_{1}(a)$ 为
\begin{align*}
E_{2}(a) &= \left\{\frac{a_{0}+a_{1}+a_{2}+a_{3}}{4}, \frac{a_{0}+a_{1}+a_{2}+a_{3}}{4},\right. \\
&\qquad \left.\frac{a_{0}+a_{1}+a_{2}+a_{3}}{4}, \frac{a_{0}+a_{1}+a_{2}+a_{3}}{4},\right. \\
&\qquad \left.\frac{a_{4}+a_{5}+a_{6}+a_{7}}{4}, \frac{a_{4}+a_{5}+a_{6}+a_{7}}{4},\right. \\
&\qquad \left.\frac{a_{4}+a_{5}+a_{6}+a_{7}}{4}, \frac{a_{4}+a_{5}+a_{6}+a_{7}}{4}\right\}.
\end{align*}
这意味着 $E_{2} \circ E_{1} = E_{2}$. 又, 数列 $b = \{b_{0}, b_{1}, \ldots, b_{7}\}$ 与 $E_{2}(a)$ 的内积 $\langle E_{2}(a), b \rangle$ 为
\begin{align*}
&\frac{(a_{0}+a_{1}+a_{2}+a_{3})(b_{0}+b_{1}+b_{2}+b_{3})}{4} \\
&\qquad + \frac{(a_{4}+a_{5}+a_{6}+a_{7})(b_{4}+b_{5}+b_{6}+b_{7})}{4},
\end{align*}
由于它对 a, b 是对称的, 它也是 a 与 $E_{2}(b)$ 的内积 $\langle a, E_{2}(b) \rangle$. 也就是说, 有 $\langle E_{2}(a), b \rangle = \langle a, E_{2}(b) \rangle$.
\end{example}

\begin{theorem}\label{thm:1.9}
对实数列 $a = \{a_{j}\}_{j=0}^{2^{N}-1}$, $b = \{b_{j}\}_{j=0}^{2^{N}-1}$ 及 $k, l = 0, 1, \ldots, N$, 以下成立.

(1) $E_{k} \circ E_{l} = E_{\max(k,l)}$.

(2) $\langle E_{k}(a), b \rangle = \langle a, E_{k}(b) \rangle$.

(3) $\|a\|_{\ell^{2}}^{2} = \|E_{N}(a)\|_{\ell^{2}}^{2} + \sum_{k=0}^{N-1}\|D_{k}(a)\|_{\ell^{2}}^{2}$.
\end{theorem}

\begin{proof}
关于 (1), 虽然直观上显然, 但为了追求完美, 给出证明.

(1) 给出 $E_{k}$ 的矩阵, 是把 $(i,j)$ 分量为定值 $2^{-k}$ 的 $2^{k}$ 阶方阵 $A_{k}$ 取 $2^{N-k}$ 个直和得到的. 由于 $(A_{k} \oplus A_{k})A_{k+1} = A_{k+1}(A_{k} \oplus A_{k}) = A_{k+1}$, 对 $k = 0, 1, \ldots, N-1$, 有 $E_{k} \circ E_{k+1} = E_{k+1} \circ E_{k} = E_{k+1}$. 因而, 对 $k = 0, 1, \ldots, N-2$, 由于 $E_{k} \circ (E_{k+1} \circ E_{k+2}) = (E_{k} \circ E_{k+1}) \circ E_{k+2}$, $(E_{k+2} \circ E_{k+1}) \circ E_{k} = E_{k+2} \circ (E_{k+1} \circ E_{k})$, 得
\[
E_{k} \circ E_{k+2} = E_{k} \circ (E_{k+1} \circ E_{k+2}) = E_{k+1} \circ E_{k+2} = E_{k+2}
\]
以及
\[
E_{k+2} \circ E_{k} = (E_{k+2} \circ E_{k+1}) \circ E_{k} = E_{k+2} \circ E_{k+1} = E_{k+2}.
\]
反复进行这样的计算, 对一般的 k, l 即得 (1).

(2) 由 (1) 中所考虑的 $A_{k}$ 是对称矩阵可知显然.

(3) 由 $D_{k}\,(k = 0, 1, \ldots, N-1)$ 与 $E_{N}$ 的定义, 分解为
\begin{align*}
\|E_{N}(a)\|_{\ell^{2}}^{2} + \sum_{k=0}^{N-1}\|D_{k}(a)\|_{\ell^{2}}^{2}
&= \|E_{N}(a)\|_{\ell^{2}}^{2} + \sum_{k=0}^{N-1}\|E_{k+1}(a)\|_{\ell^{2}}^{2} \\
&\quad - 2\sum_{k=0}^{N-1}\langle E_{k}(a), E_{k+1}(a)\rangle_{\ell^{2}} + \sum_{k=0}^{N-1}\|E_{k}(a)\|_{\ell^{2}}^{2}.
\end{align*}
由 (1), (2),
\begin{align*}
\langle E_{k}(a), E_{k+1}(a)\rangle_{\ell^{2}} &= \langle E_{k}(a), E_{k+1} \circ E_{k+1}(a)\rangle_{\ell^{2}} \\
&= \langle E_{k+1} \circ E_{k}(a), E_{k+1}(a)\rangle_{\ell^{2}} \\
&= \langle E_{k+1}(a), E_{k+1}(a)\rangle_{\ell^{2}} = \|E_{k+1}(a)\|_{\ell^{2}}^{2},
\end{align*}
所以
\begin{align*}
\|E_{N}(a)\|_{\ell^{2}}^{2} + \sum_{k=0}^{N-1}\|D_{k}(a)\|_{\ell^{2}}^{2}
&= \|E_{N}(a)\|_{\ell^{2}}^{2} + \sum_{k=0}^{N-1}\|E_{k}(a)\|_{\ell^{2}}^{2} - \sum_{k=0}^{N-1}\|E_{k+1}(a)\|_{\ell^{2}}^{2} \\
&= \|E_{0}(a)\|_{\ell^{2}}^{2}.
\end{align*}
由于 $E_{0}(a) = a$, 得到 (3).
\end{proof}

由定理~\ref{thm:1.9} 可知, $\{E_{N}\} \cup \{D_{k}\}_{k=0}^{N-1}$ 是分解 $\ell^{2}$ 范数的算子族. 在本节收尾之际, 将定理~\ref{thm:1.9} 变形为用内积表示的形式. 一般地, 对实向量 $a, b \in \mathbb{R}^{2^{N}}$, 有
\[
\langle a, b \rangle_{\ell^{2}} = \frac{1}{4}\|a + b\|_{\ell^{2}}^{2} - \frac{1}{4}\|a - b\|_{\ell^{2}}^{2}.
\]
由此与定理~\ref{thm:1.9}(3), 得
\begin{equation}\tag{1.15}
\langle a, b \rangle_{\ell^{2}} = \langle E_{N}(a), E_{N}(b)\rangle_{\ell^{2}} + \sum_{k=0}^{N-1}\langle D_{k}(a), D_{k}(b)\rangle_{\ell^{2}}.
\label{eq:1.15}
\end{equation}

\begin{exercise}\label{exer:1.22}
对例~\ref{exa:1.13} 中考虑的数列, 求 $E_{0}(a)$, $E_{2}(a)$, $E_{3}(a)$, $E_{4}(a)$, $D_{0}(a)$, $D_{1}(a)$, $D_{2}(a)$, $D_{3}(a)$.
\end{exercise}

\begin{exercise}\label{exer:1.23}
设 $a = \{a_{j}\}_{j=0}^{7}$ 为实数列. 令
\begin{align*}
A &= \frac{a_{0}+a_{1}+\cdots+a_{7}}{8}, \\
B &= \frac{a_{0}-a_{1}}{2}, \quad C = \frac{a_{2}-a_{3}}{2}, \quad D = \frac{a_{4}-a_{5}}{2}, \quad E = \frac{a_{6}-a_{7}}{2}, \\
F &= \frac{a_{0}+a_{1}-a_{2}-a_{3}}{4}, \quad G = \frac{a_{4}+a_{5}-a_{6}-a_{7}}{4}, \\
H &= \frac{a_{0}+a_{1}+a_{2}+a_{3}-a_{4}-a_{5}-a_{6}-a_{7}}{8}.
\end{align*}
计算 $8A^{2} + 2B^{2} + 2C^{2} + 2D^{2} + 2E^{2} + 4F^{2} + 4G^{2} + 8H^{2}$.
\end{exercise}

% ============ 第1章 §1.8 鞅变换 (原书 pp.26-29) ============

\section{鞅变换}\label{sec:1.8}
\index{Y@鞅变换}




这里的目的, 是通过组合至今处理过的各种公式, 来考察 $\ell^{p}$, $1 < p < \infty$ 中算子的性质. 即使在 $\ell^{1}$ 中难以分析的算子, 利用 Marcinkiewicz 插值定理也能弄清楚很多东西. Marcinkiewicz 插值定理例如是组合 $\ell^{1}$ 中的``准''有界性与对 1 以外的 p 的 $\ell^{p}$ 中的有界性来使用的. 至于实际上如何使用 Marcinkiewicz 插值定理, 在调和分析的世界中 p = 2 的情形很多. 一般认为, 这是因为在数学的世界中, 与平方有关的各种公式很多, 例如
\begin{align*}
z\overline{z} &= |z|^{2} \quad (z \in \mathbb{C}), \\
\left(\frac{a+b}{2}\right)^{2} + \left(\frac{a-b}{2}\right)^{2} &= \frac{1}{2}(a^{2}+b^{2}) \quad (a, b \in \mathbb{R}).
\end{align*}
§\ref{sec:1.7} 中的 $D_{k}$ 是构成鞅变换的算子, 而鞅变换是金融领域出现的随机积分的原型. 下面考虑的鞅变换, 从调和分析的立场看不会出现积分记号, 但被定位在奇异积分算子的范畴中. 对由 $2^{N}$ 项组成的 N 个数列 $m^{0}, m^{1}, \ldots, m^{N-1}$, 定义鞅变换 $M_{m}$ 如下: 它对由这些数列生成的长度为 $2^{N}$ 的数列 a, 给出同样长度为 $2^{N}$ 的数列 $M_{m}(a)$:
\[
M_{m}(a) = \{M_{m}(a)_{j}\}_{j=0}^{2^{N}-1} = \sum_{k=0}^{N-1}E_{k+1}(m^{k})D_{k}(a),
\]
其中 $E_{k+1}(m^{k})D_{k}(a) = \{E_{k+1}(m^{k})_{j}D_{k}(a)_{j}\}_{j=0}^{2^{N}-1}$.

\begin{theorem}\label{thm:1.10}
对长度 $2^{N}$ 的数列 $a, m^{0}, m^{1}, \ldots, m^{N-1}$, 成立
\[
\|M_{m}(a)\|_{\ell^{2}} \leq \sup_{k=0,1,\ldots,N-1}\|E_{k+1}(m^{k})\|_{\ell^{\infty}}\|a\|_{\ell^{2}}.
\]
\end{theorem}

\begin{proof}
由于 $E_{N} \circ M_{m}(a) = 0$, $D_{k} \circ M_{m}(a) = -E_{k+1}(m^{k})D_{k}(a)$, 由定理~\ref{thm:1.9} 得
\begin{align*}
\|M_{m}(a)\|_{\ell^{2}}^{2} &= \|E_{N} \circ M_{m}(a)\|_{\ell^{2}}^{2} + \sum_{k=0}^{N-1}\|D_{k} \circ M_{m}(a)\|_{\ell^{2}}^{2} \\
&= \sum_{k=0}^{N-1}\|E_{k+1}(m^{k})D_{k}(a)\|_{\ell^{2}}^{2} \\
&\leq \sup_{k=0,1,\ldots,N-1}(\|E_{k+1}(m^{k})\|_{\ell^{\infty}}\|a\|_{\ell^{2}})^{2}.
\end{align*}
\end{proof}

在下面的定理中, Calderón--Zygmund 分解扮演重要角色.

\begin{theorem}\label{thm:1.11}
对 $\lambda > 0$ 与数列 $a = \{a_{j}\}_{j=0}^{2^{N}-1}$, 令 $E_{\lambda}(a) = \{j = 0, 1, \ldots, 2^{N}-1 : |M_{m}(a)_{j}| > \lambda\}$, 则
\[
\sharp(E_{\lambda}(a)) \leq \frac{3}{\lambda}\sup_{k=0,1,\ldots,N-1}\|E_{k+1}(m^{k})\|_{\ell^{\infty}}\|a\|_{\ell^{1}}.
\]
\end{theorem}

\begin{proof}
必要时通过对各 $m^{k}$ 乘以适当的常数倍, 可设
\[
\sup_{k=0,1,\ldots,N-1}\|E_{k+1}(m^{k})\|_{\ell^{\infty}} = 1.
\]
又, 当 $\lambda \leq 2^{-N}\|a\|_{\ell^{1}}$ 时,
\[
\sharp(E_{\lambda}(a)) \leq 2^{N} \leq \frac{2}{\lambda}\|a\|_{\ell^{1}} = \frac{2}{\lambda}\sup_{k=0,1,\ldots,N-1}\|E_{k+1}(m^{k})\|_{\ell^{\infty}}\|a\|_{\ell^{1}},
\]
所以结论显然. 因而可设 $\lambda > 2^{-N}\|a\|_{\ell^{1}}$. 这里使用 Calderón--Zygmund 分解. 按照定理~\ref{thm:1.8}, 将 a 分解为 $a = g + b$. 由于 $\|g\|_{\ell^{2}} \leq \sqrt{\|g\|_{\ell^{1}}\|g\|_{\ell^{\infty}}}$ (参照练习题~\ref{exer:1.25}), $\|g\|_{\ell^{1}} \leq \|a\|_{\ell^{1}}$, $\|g\|_{\ell^{\infty}} \leq 2\lambda$, 由定理~\ref{thm:1.10} 得
\[
\|M_{m}(g)\|_{\ell^{2}}^{2} \leq \|g\|_{\ell^{2}}^{2} \leq 2\lambda\|g\|_{\ell^{1}} \leq 2\lambda\|a\|_{\ell^{1}}.
\]
因此得
\begin{equation}\tag{1.16}
\sharp(E_{\lambda}(g)) = \sum_{j=0}^{2^{N}-1}\chi_{(\lambda,\infty)}(|M_{m}(g)_{j}|) \leq \frac{\|M_{m}(g)\|_{\ell^{2}}^{2}}{\lambda^{2}} \leq \frac{2}{\lambda}\|a\|_{\ell^{1}}.
\label{eq:1.16}
\end{equation}
另一方面, 将 g 表为 $g = \{g_{j}\}_{j=0}^{2^{N}-1}$ 时, 对 $I \in \mathcal{U}$, 数列
\[
b_{I} = \{\chi_{I}(j)(a_{j} - g_{j})\}_{j=0}^{2^{N}-1}
\]
在 I 之外满足 $M_{m}(b_{I}) = 0$. 实际上, 设 $\#I = 2^{j_{0}}$, 若 $k \geq j_{0}$, 则长度为 $2^{k}$ 的二进制区间一旦与 I 相交就会包含 I, 所以 $D_{k}(b_{I}) = 2^{-k}\sum_{j\in I}b_{I}(j) = 0$. 若 $k < j_{0}$, 则包含 I 之外整数的长度为 $2^{k}$ 的二进制区间不与 I 相交. 因此 $D_{k}(b_{I}) = 0$. 由于 $M_{m}(b_{I})$ 可写为 $D_{k}(b_{I})$ 的和, 所以 $M_{m}(b_{I}) = 0$. 因而在
\[
\{0, 1, \ldots, 2^{N}-1\}\setminus E = \{0, 1, \ldots, 2^{N}-1\}\setminus\bigcup_{I\in\mathcal{U}}I
\]
上 $M_{m}(b) = 0$. 所以在 E 之外 $M_{m}(a) = M_{m}(g)$. 因而利用 $E = \bigcup_{I\in\mathcal{U}}I$, $E_{\lambda}(g) = \{j = 0, 1, \ldots, 2^{N}-1 : |M_{m}(g)_{j}| > \lambda\}$, 由式~\eqref{eq:1.16} 得
\[
\sharp(E_{\lambda}(a)) \leq \sharp E + \sharp(E_{\lambda}(a)\setminus E) = \sharp E + \sharp(E_{\lambda}(g)\setminus E) \leq \frac{2}{\lambda}\|a\|_{\ell^{1}}.
\]
\end{proof}

通过对所得不等式进行插值, 得到了下面重要的结论.

\begin{theorem}\label{thm:1.12}
设 $1 < p < \infty$. 此时, 存在仅依赖于 p 的常数 $c_{p}$, 使得对所有长度 $2^{N}$ 的数列 $a, m^{0}, m^{1}, \ldots, m^{N-1}$ 成立
\[
\|M_{m}(a)\|_{\ell^{p}} \leq c_{p}\sup_{k=0,1,\ldots,N-1}\|E_{k+1}(m^{k})\|_{\ell^{\infty}}\|a\|_{\ell^{p}}.
\]
\end{theorem}

\begin{proof}
分三种情形.

(A) 当 $1 < p < 2$ 时, 通过组合式~\eqref{eq:1.8} 与定理~\ref{thm:1.10}, \ref{thm:1.11} 可得结论. 也就是说, 插值理论扮演重要角色.

(B) 当 p = 2 时, 就是定理~\ref{thm:1.10} 本身.

(C) 当 $2 < p < \infty$ 时, 对 $b = \{b_{j}\}_{j=0}^{2^{N}-1}$, 利用
\begin{align*}
\langle M_{m}(a), b \rangle &= \sum_{k=0}^{N-1}\langle E_{k+1}(m^{k})D_{k}(a), b \rangle \\
&= \sum_{k=0}^{N-1}\langle a, E_{k+1}(m^{k})D_{k}(b) \rangle \\
&= \langle a, M_{m}(b) \rangle.
\end{align*}
其中 b 与 Calderón--Zygmund 分解时出现的数列是两回事. 由 Hölder 不等式的等号成立条件 (定理~\ref{thm:1.2} 下方的注意), 得
\begin{align*}
\|M_{m}(a)\|_{\ell^{p}} &= \sup\{|\langle M_{m}(a), b \rangle| : \|b\|_{\ell^{p'}} = 1\} \\
&= \sup\{|\langle a, M_{m}(b) \rangle| : \|b\|_{\ell^{p'}} = 1\}.
\end{align*}
这里, 由 Hölder 不等式,
\[
\|M_{m}(a)\|_{\ell^{p}} \leq \sup\{\|a\|_{\ell^{p}}\|M_{m}(b)\|_{\ell^{p'}} : \|b\|_{\ell^{p'}} = 1\}.
\]
若 $2 < p < \infty$, 则 $1 < p' < 2$, 所以可利用本证明 (A) 的部分, 得
\begin{align*}
\|M_{m}(a)\|_{\ell^{p}} &\leq c_{p'}\sup_{k=0,1,\ldots,N-1}\|E_{k+1}(m^{k})\|_{\ell^{\infty}}\sup\{\|a\|_{\ell^{p}}\|b\|_{\ell^{p'}} : \|b\|_{\ell^{p'}} = 1\} \\
&= c_{p'}\sup_{k=0,1,\ldots,N-1}\|E_{k+1}(m^{k})\|_{\ell^{\infty}}\|a\|_{\ell^{p}}.
\end{align*}
\end{proof}

\begin{exercise}\label{exer:1.24}
设 $1 < p < \infty$. 证明: 此时存在仅依赖于 p 的常数 $c_{p}$, 使得对所有长度 $2^{N}$ 的数列 a 成立
\[
\left\|\sum_{k=0}^{N-1}(-1)^{k}D_{k}(a)\right\|_{\ell^{p}} \leq c_{p}\|a\|_{\ell^{p}}.
\]
\end{exercise}

\begin{exercise}\label{exer:1.25}
对长度 N 的数列 $g = \{g_{j}\}_{j=1}^{N}$, 证明 $\|g\|_{\ell^{2}} \leq \sqrt{\|g\|_{\ell^{1}}\|g\|_{\ell^{\infty}}}$.
\end{exercise}

% ============ 第1章 §1.9 Littlewood--Paley 理论 (原书 pp.30-31) ============

\section{Littlewood--Paley 理论}\label{sec:1.9}

长度 $2^{N}$ 的数列 $a = \{a_{j}\}_{j=0}^{2^{N}-1}$ 的 $\ell^{p}$ 范数为
\[
\|a\|_{\ell^{p}} = \left(\sum_{j=0}^{2^{N}-1}|a_{j}|^{p}\right)^{1/p},
\]
但只要 p 不是 1, 2, 一般地 p 次幂和很难处理. 实际上, 例如若 p = 1, 则不必开根号; 若 p = 2, 则利用 Rademacher 列 $\{r_{j}\}_{j=0}^{2^{N}-1}$, 可表示为
\[
\|a\|_{\ell^{2}} = \left(\int_{0}^{1}\left|\sum_{j=0}^{2^{N}-1}r_{j}(t)a_{j}\right|^{2}\dd t\right)^{1/2}.
\]
但是, 若 $p \neq 1, 2$, 开根号的操作会出现在定义中, 也没有便利的等式. Littlewood--Paley 理论可以说是当 $1 < p < \infty$ 时为变形 $\ell^{p}$ 范数而设的理论. 定义数列 $g(a) = \{g(a)_{j}\}_{j=0}^{2^{N}-1}$ 为
\[
g(a)_{j} = \left(|E_{N}(a)_{j}|^{2} + \sum_{k=0}^{N-1}|D_{k}(a)_{j}|^{2}\right)^{1/2}.
\]

\begin{theorem}[Littlewood--Paley 理论]\label{thm:1.13}
对数列 $a = \{a_{j}\}_{j=0}^{2^{N}-1}$, 当 $1 < p < \infty$ 时, 适当取不依赖于数列长度 $2^{N}$ 及 a 的分量的常数 $C_{p} > 1$, 则成立
\[
C_{p}^{-1}\|a\|_{\ell^{p}} \leq \|g(a)\|_{\ell^{p}} \leq C_{p}\|a\|_{\ell^{p}}.
\]
\end{theorem}

\begin{proof}
证明存在适当常数 $K_{p} > 0$ 使 $\|g(a)\|_{\ell^{p}} \leq K_{p}\|a\|_{\ell^{p}}$. 由 Khintchine 不等式, 对 $j = 0, 1, \ldots, 2^{N}-1$ 成立
\[
g(a)_{j} \leq K_{p}\left\{|E_{N}(a)_{j}| + \left(\int_{0}^{1}\left|\sum_{k=0}^{N-1}r_{k}(t)D_{k}(a)_{j}\right|^{p}\dd t\right)^{1/p}\right\}.
\]
这里, 由于 $r_{k}(t) = \pm 1$, 由定理~\ref{thm:1.12}, 若将 $K_{p}$ 重新取大, 则存在不依赖于 $t \in [0,1]$ 的常数 $K_{p} > 0$, 使
\[
\left\|\sum_{k=0}^{N-1}r_{k}(t)D_{k}(a)\right\|_{\ell^{p}} \leq K_{p}\|a\|_{\ell^{p}}
\]
也成立. 将这个不等式对 t 在 [0,1] 上积分, 代入刚才的不等式所得的常数重新记为 $K_{p}$, 得结论 $\|g(a)\|_{\ell^{p}} \leq K_{p}\|a\|_{\ell^{p}}$.

将这里得到的常数记为 $K_{p}$, 这次证明 $\|a\|_{\ell^{p}} \leq K_{p'}\|g(a)\|_{\ell^{p}}$. 将 b 写为 $b = \{b_{j}\}_{j=0}^{2^{N}-1}$, 由 Hölder 不等式的等号成立条件,
\[
\|a\|_{\ell^{p}} = \sup\left\{\left|\sum_{j=0}^{2^{N}-1}a_{j}b_{j}\right| : \|b\|_{\ell^{p'}} \leq 1\right\}.
\]
由式~\eqref{eq:1.15} 与 Cauchy--Schwarz 不等式,
\begin{align*}
\|a\|_{\ell^{p}} &= \sup\left\{\left|\langle E_{N}(a), E_{N}(b)\rangle + \sum_{k=0}^{N-1}\langle D_{k}(a), D_{k}(b)\rangle\right| : \|b\|_{\ell^{p'}} \leq 1\right\} \\
&\leq \sup\left\{\sum_{j=0}^{2^{N}-1}g(a)_{j}g(b)_{j} : \|b\|_{\ell^{p'}} \leq 1\right\}.
\end{align*}
如前所述, 若 $\|b\|_{\ell^{p'}} \leq 1$, 则 $\|g(b)\|_{\ell^{p'}} \leq K_{p'}\|b\|_{\ell^{p'}} \leq K_{p'}$, 所以由 Hölder 不等式,
\[
\|a\|_{\ell^{p}} \leq \sup\{\|g(a)\|_{\ell^{p}}\|g(b)\|_{\ell^{p'}} : \|b\|_{\ell^{p'}} \leq 1\} \leq K_{p'}\|g(a)\|_{\ell^{p}}.
\]
因此可知, 作为定理结论应取的 $C_{p}$ 可取为 $C_{p} = \max(K_{p}, K_{p'})$.
\end{proof}

\begin{exercise}\label{exer:1.26}
设 $e = (1, 0, 0, \ldots, 0)$ 为初项为 1, 其余 $2^{N}-1$ 项为 0 的数列. 当 $2^{k-1} \leq j < 2^{k}$ 时, 证明 $D_{k}(e)_{j} = 2^{-k-1}$, 并利用它证明 $\|g(e)\|_{\ell^{1}} \geq \frac{N-1}{2}$.
\end{exercise}

% ============ 第1章 §1.10 高维情形的推广 (原书 pp.31-32) ============

\section{高维情形的推广}\label{sec:1.10}
\index{0symbolvecX@$\vec{X}_{k}$}

\index{0symbolDkm@$D_{k}^{m}$}

\index{0symbolE@$E$}

\index{0symbolEk2@$E_{k}$ (高维)}

\index{0symbolEkm@$E_{k}^{m}$}












用整数分量向量作指标编号的数列称为多重数列. 这里取 N, n 为自然数, 考虑由 $2^{nN}$ 个实数组成的多重数列 $a = \{a_{\vec{j}}\}_{\vec{j}\in I(N,0)^{n}}$. 由于这样的数列由 n 维向量组成, 称为 n 重数列或 n 维数列. 区间 $I(N,0)$ 是在式~\eqref{eq:1.14} 中令 j = N, k = 0 定义的. 暂且固定 $k = 0, 1, \ldots, N$. 迄今以 n = 1, 除 $D_{N}$ 外考虑了 $E_{k}, D_{k}$ 等算子, 将把它们作用于第 $m\,(=1, 2, \ldots, n)$ 分量所得的算子记为 $E_{k}^{m}, D_{k}^{m}$. 也就是说, 例如, 对 $\vec{j} = (j_{1}, j_{2}, \ldots, j_{n})$, 定义 $E_{k}^{m}(a) = \{E_{k}^{m}(a)_{\vec{j}}\}_{\vec{j}\in I(N,0)^{n}}$ 为
\[
E_{k}^{m}(a)_{\vec{j}} = \frac{1}{2^{k}}\sum_{l = 2^{k}[2^{-k}j_{m}]}^{2^{k}[2^{-k}j_{m}] + 2^{k} - 1}a_{(j_{1}, j_{2}, \ldots, j_{m-1}, l, j_{m+1}, \ldots, j_{n})}.
\]

继续设 $k = 0, 1, \ldots, N$. 设给定由 D, E 中某一字母组成的 n 个字母的序列 $\vec{X} = (X^{1}, X^{2}, \ldots, X^{n})$. 在这样的序列中, 仅 E 连续出现 n 次的序列记为 E. 一般地, 当 $X^{j} = D$ 时令 $Y^{j} = D_{k}^{j}$, 当 $X^{j} = E$ 时令 $Y^{j} = E_{k+1}^{j}$, 令由 $\vec{X}_{k} = Y^{1} \circ Y^{2} \circ \cdots \circ Y^{n}$ 得到的算子为 $\vec{X}_{k}$. 例如, 当 n = 3 时,
\[
(D, D, E)_{k} = D_{k}^{1} \circ D_{k}^{2} \circ E_{k+1}^{3}, \qquad (E, D, E)_{k} = E_{k+1}^{1} \circ D_{k}^{2} \circ E_{k+1}^{3}.
\]
又, $\mathbf{E}_{k} = (E, E, \ldots, E)_{k}$. 一般地, 对多重数列 a, 由于 $\mathbf{E}_{0}(a) = a$, 分解为
\[
a = \mathbf{E}_{N}(a) + \sum_{k=0}^{N-1}(\mathbf{E}_{k}(a) - \mathbf{E}_{k+1}(a)).
\]
进而, 由于 $D_{k} = E_{k} - E_{k+1}$, 可分解为
\[
\mathbf{E}_{k} - \mathbf{E}_{k+1} = \sum_{\vec{X}\in\{(D,E)\}^{n}\setminus\{\mathbf{E}\}}\vec{X}_{k}.
\]
基于这一点, 对 $\vec{j} \in I(N,0)^{n}$, 定义实数 $g(a)_{\vec{j}} \geq 0$ 为
\[
g(a)_{\vec{j}} = \left(|\mathbf{E}_{N}(a)_{\vec{j}}|^{2} + \sum_{k=1}^{N}\sum_{\vec{X}\in\{(D,E)\}^{n}\setminus\{\mathbf{E}\}}|\vec{X}_{k-1}(a)_{\vec{j}}|^{2}\right)^{1/2}.
\]
记为 $g(a) = \{g(a)_{\vec{j}}\}_{\vec{j}\in I(N,0)^{n}}$. 用与 n = 1 情形相同的证明方法, 可以证明以下结论.

\begin{theorem}\label{thm:1.14}
设 $n, N \in \mathbb{N}$. 存在仅依赖于 p, n 的常数 $C_{p,n}$, 当 $1 < p < \infty$ 时, 对由 $2^{nN}$ 个分量组成的多重数列 $a = \{a_{\vec{j}}\}_{\vec{j}\in I(N,0)^{n}}$, 成立
\[
C_{p,n}^{-1}\|a\|_{\ell^{p}} \leq \|g(a)\|_{\ell^{p}} \leq C_{p,n}\|a\|_{\ell^{p}}.
\]
\end{theorem}

\begin{exercise}\label{exer:1.27}
设 $k, k' = 0, 1, \ldots, N$, n = 2, 证明 $E_{k}^{1} \circ E_{k'}^{2} = E_{k'}^{2} \circ E_{k}^{1}$.
\end{exercise}

% ============ 第1章 §1.11 无穷数列情形的推广 (原书 pp.33-35) ============

\section{无穷数列情形的推广}\label{sec:1.11}
\index{J@基本向量}




设 $N_{0} = \{0, 1, \ldots\}$. n 再次设为自然数. 这次考虑数列项数不是 $2^{nN}$ 个而是无限多个的情形. 即, 考虑表示为 $\{a_{\vec{j}}\}_{\vec{j}\in\mathbb{N}_{0}^{n}}$ 形式的多重数列. 对 $j \in \mathbb{N}$, $\mathbf{e}_{j} = \{\chi_{\{j\}}(k)\}_{k\in\mathbb{N}_{0}}$ 表示基本向量. 将 $E_{k}$ 视为取 $I(k,j)$ ($j = 0, 1, \ldots$) 上的平均, 与此对应, $\vec{X}_{k}, k \in \mathbb{N}$ 也自然地被推广. 对 $1 < p < \infty$, 由于迄今定理中出现的常数 $C_{p}$ 不依赖于 N, 令
\[
g(a)_{\vec{j}} = \left(\sum_{k=1}^{\infty}\sum_{\vec{X}\in\{(D,E)\}^{n}\setminus\{\mathbf{E}\}}|\vec{X}_{k-1}(a)_{\vec{j}}|^{2}\right)^{1/2} \quad (\vec{j} \in \mathbb{N}_{0}^{n}).
\]
对 $1 \leq p \leq \infty$, 多重数列空间 $\ell^{p}(\mathbb{N}_{0}^{n})$ 指具有 $\{a_{\vec{j}}\}_{\vec{j}\in\mathbb{N}_{0}^{n}}$ 形式且满足
\[
\|\{a_{\vec{j}}\}_{\vec{j}\in\mathbb{N}_{0}^{n}}\|_{\ell^{p}} = \left(\sum_{\vec{j}\in\mathbb{N}_{0}^{n}}|a_{\vec{j}}|^{p}\right)^{1/p} < \infty
\]
的数列全体. 当 $p = \infty$ 时, 右边不取 p 次幂和而取上限. 当 n = 1 时, 记为 $\ell^{p}(\mathbb{N}_{0})$. 将指标集合 $N_{0}$ 换成 $\mathbb{N}$ 所得的数列空间 $\ell^{p}(\mathbb{N})$ 也同样定义.

\begin{lemma}\label{lem:1.5}
对 $a = \{a_{\vec{j}}\}_{\vec{j}\in\mathbb{N}_{0}^{n}} \in \ell^{2}(\mathbb{N}_{0}^{n})$, 成立
\[
\|a\|_{\ell^{2}}^{2} = \sum_{k=1}^{\infty}\sum_{\vec{X}\in\{(D,E)\}^{n}\setminus\{\mathbf{E}\}}\|\vec{X}_{k-1}(a)\|_{\ell^{2}}^{2}.
\]
\end{lemma}

\begin{proof}
固定自然数 $N_{0}$. 由 $(A+B)^{2} \leq 2A^{2} + 2B^{2}$, $A, B \in \mathbb{R}$, 与 Cauchy--Schwarz 不等式, 若自然数 N 满足 $N > N_{0}$, 则利用式~\eqref{eq:1.14} 给出的区间 $I(N,0), I(N_{0},0)$, 有
\begin{align*}
\frac{1}{2^{nN}}\left|\sum_{\vec{j}\in I(N,0)^{n}}a_{\vec{j}}\right|^{2}
&\leq \frac{2}{2^{nN}}\left|\sum_{\vec{j}\in I(N_{0},0)^{n}}a_{\vec{j}}\right|^{2} + \frac{2}{2^{nN}}\left|\sum_{\vec{j}\in I(N,0)^{n}\setminus I(N_{0},0)^{n}}a_{\vec{j}}\right|^{2} \\
&\leq \frac{2}{2^{nN}}\left|\sum_{\vec{j}\in I(N_{0},0)^{n}}a_{\vec{j}}\right|^{2} + 2\sum_{\vec{j}\in I(N,0)^{n}\setminus I(N_{0},0)^{n}}|a_{\vec{j}}|^{2}.
\end{align*}
对 N 取上极限, 得
\[
\limsup_{N\to\infty}\frac{1}{2^{nN}}\left|\sum_{\vec{j}\in I(N,0)^{n}}a_{\vec{j}}\right|^{2} \leq 2\sum_{\vec{j}\in\mathbb{N}_{0}^{n}\setminus I(N_{0},0)^{n}}|a_{\vec{j}}|^{2} = 2\sum_{\vec{j}\in\mathbb{N}_{0}^{n}}|a_{\vec{j}}|^{2} - 2\sum_{\vec{j}\in I(N_{0},0)^{n}}|a_{\vec{j}}|^{2}.
\]
再对 $N_{0}$ 取极限, 得
\[
\limsup_{N\to\infty}\frac{1}{2^{nN}}\left|\sum_{\vec{j}\in I(N,0)^{n}}a_{\vec{j}}\right|^{2} = 0 \quad \text{即} \quad \lim_{N\to\infty}\frac{1}{2^{nN}}\left|\sum_{\vec{j}\in I(N,0)^{n}}a_{\vec{j}}\right|^{2} = 0.
\]
这里, 令 $a^{N} = \{\chi_{I(N,0)^{n}}(\vec{j})a_{\vec{j}}\}_{\vec{j}\in\mathbb{N}_{0}^{n}}$, 则由定理~\ref{thm:1.9} 的 n 维版本,
\[
\|a^{N}\|_{\ell^{2}}^{2} = \frac{1}{2^{nN}}\left|\sum_{\vec{j}\in I(N,0)^{n}}a_{\vec{j}}\right|^{2} + \sum_{k=1}^{N}\sum_{\vec{X}\in\{(D,E)\}^{n}\setminus\{\mathbf{E}\}}\|\vec{X}_{k-1}(a)\|_{\ell^{2}}^{2},
\]
令 $N \to \infty$ 即得结论.
\end{proof}

对 $p \in (1, \infty)\setminus\{2\}$ 也仿照刚才的讨论, 在定理~\ref{thm:1.14} 中令 $N \to \infty$, 得下面的定理.

\begin{theorem}\label{thm:1.15}
存在仅依赖于自然数 n 与 $p \in (1, \infty)$ 的常数 $C > 0$, 使得对所有 $a \in \ell^{p}(\mathbb{N}_{0}^{n})$ 成立
\[
C^{-1}\|a\|_{\ell^{p}} \leq \|g(a)\|_{\ell^{p}} \leq C\|a\|_{\ell^{p}}.
\]
\end{theorem}

把``存在常数 $C > 0$, 使得对所有 $a \in \ell^{p}(\mathbb{N}_{0}^{n})$ 成立 $C^{-1}\|a\|_{\ell^{p}} \leq \|g(a)\|_{\ell^{p}} \leq C\|a\|_{\ell^{p}}$''这一状况记为 $\|g(a)\|_{\ell^{p}} \sim \|a\|_{\ell^{p}}$ 是很方便的. 此后, 在类似的情形下也常用 $\sim$ 表示. 又, 当只考虑 $\|g(a)\|_{\ell^{p}} \leq C\|a\|_{\ell^{p}}$ 这种单侧情形时, 记为 $\|g(a)\|_{\ell^{p}} \lesssim \|a\|_{\ell^{p}}$ 或 $\|a\|_{\ell^{p}} \gtrsim \|g(a)\|_{\ell^{p}}$.

\begin{exercise}\label{exer:1.28}
设数列 $\{a_{j}\}_{j=1}^{\infty} \in \ell^{1}(\mathbb{N})$ 对所有 $j \in \mathbb{N}$ 满足 $0 \leq a_{j+1} \leq a_{j}$. 此时, 证明 $\lim_{j\to\infty}ja_{j} = 0$.
\end{exercise}

\begin{exercise}\label{exer:1.29}
设 $1 \leq p \leq \infty$.

(1) 当 $0 \leq j < k < \infty$ 时, 计算 $\|\mathbf{e}_{j} - \mathbf{e}_{k}\|_{\ell^{p}}$.

(2) 证明 $\ell^{p}(\mathbb{N})$ 的单位球 $B = \{a \in \ell^{p}(\mathbb{N}) : \|a\|_{\ell^{p}} \leq 1\}$ 不是紧的.
\end{exercise}

\begin{exercise}\label{exer:1.30}
设 $1 \leq p \leq \infty$. 用 $\operatorname{Span}(\{e_{j}\}_{j=0}^{\infty})$ 表示包含各 $e_{j}$ 的最小线性空间. 此时, 对配备 $\ell^{\infty}$ 范数的 $\operatorname{Span}(\{e_{j}\}_{j=0}^{\infty})$, 证明以下.

(1) 若 $x_{k} = \sum_{l=1}^{k}2^{-l}\mathbf{e}_{l}$, 则当 $1 \leq k < k' < \infty$ 时, $\|x_{k} - x_{k'}\|_{\ell^{\infty}} = 2^{-k-1}$.

(2) 在 $\operatorname{Span}(\{\mathbf{e}_{j}\}_{j=0}^{\infty})$ 中, $\{x_{k}\}_{k=1}^{\infty}$ 不收敛.
\end{exercise}


% 第2章 Lebesgue 积分与泛函分析的基础知识 (原书 pp.36-60, PDF pp.42-66)
\chapter{Lebesgue 积分与泛函分析的基础知识}
% ============ 第2章 §2.1 Lebesgue 测度 (原书 pp.36-43) ============

第1章通过数列概观了调和分析, 但一般地, 调和分析中大量使用积分. 又, 由于必须基于完备性进行讨论, 几乎都使用 Banach 空间. 本章叙述 Lebesgue 积分的基础事项以及函数解析的基础事项. 关于 Lebesgue 积分的基础事项, 不作证明地总结. 详细证明请参考适当的文献. 又, 这一章几乎没有练习题.

\section{Lebesgue 测度}\label{sec:2.1}
\index{0symbolcap@$\bigcap_{\lambda \in \Lambda} U_{\lambda}$}

\index{0symbolcup@$\bigcup_{\lambda \in \Lambda} U_{\lambda}$}

\index{0symbolBxr@$B(x, r)$}

\index{0symbolsupp@$\operatorname{supp}(f)$}

\index{Z@支集}

\index{J@简单函数}

\index{C@长方体}

\index{N@内正则性}

\index{L@Lebesgue 外测度}

\index{L@Lebesgue 可测函数}

\index{L@Lebesgue 可测集合}
























这里, 除证明外, 简洁地总结 Lebesgue 积分的基础事项.

\subsection*{集合的并, 交集}

调和分析亦可说是研究函数性质的学问, 在调和分析中, 集合的操作是必需事项. 本书使用以下记号.

\begin{definition}\label{def:2.1}
(1) 集合的子族, 亦即子集族 $\{U_{\lambda}\}_{\lambda\in\Lambda}\subset\mathbb{R}^{n}$, 是指

(A) $\Lambda$ 是某个集合, 且

(B) 对 $\lambda \in \Lambda$, 给定 $\mathbb{R}^{n}$ 的子集 $U_{\lambda}$.

此时, 称 $\{U_{\lambda}\}_{\lambda\in\Lambda}$ 为 $\mathbb{R}^{n}$ 的子集族.

(2) 对 $\mathbb{R}^{n}$ 的子集族 $\{U_{\lambda}\}_{\lambda\in\Lambda}\subset\mathbb{R}^{n}$, 其并集定义为
\[
\bigcup_{\lambda \in \Lambda}U_{\lambda} = \{x \in \mathbb{R}^{n} : \text{至少存在一个 } \lambda \in \Lambda \text{ 使得 } x \in U_{\lambda}\}.
\]

(3) 对 $\mathbb{R}^{n}$ 的子集族 $\{U_{\lambda}\}_{\lambda\in\Lambda}\subset\mathbb{R}^{n}$, 其交集定义为
\[
\bigcap_{\lambda \in \Lambda}U_{\lambda} = \{x \in \mathbb{R}^{n} : \text{对所有 } \lambda \in \Lambda \text{ 有 } x \in U_{\lambda}\}.
\]
\end{definition}

\subsection*{Lebesgue 可测集合与 Lebesgue 可测函数}

为了研究函数的性质, 时常会用到完备性的讨论, 但如此一来在 Riemann 积分的框架下就不方便了. 这是因为在调和分析中, 几乎所有的情形都是通过由积分定义的范数进行分析的. 因此, 来复习 Lebesgue 可测集合与 Lebesgue 可测函数.

\begin{definition}[直方体]\label{def:2.2}
用满足 $-\infty < a_{j} \leq b_{j} < \infty$ $(j = 1, 2, \ldots, n)$ 的 $2n$ 个实数表为
\[
\prod_{j=1}^{n}[a_{j}, b_{j}) = [a_{1}, b_{1}) \times [a_{2}, b_{2}) \times \cdots \times [a_{n}, b_{n})
\]
的集合称为直方体. 直方体的体积定义为
\[
\left|\prod_{j=1}^{n}[a_{j}, b_{j})\right| = \prod_{j=1}^{n}(b_{j} - a_{j}).
\]
将全体直方体构成的集合记为 $\mathcal{R}$.
\end{definition}

\begin{definition}[Lebesgue 外测度]\label{def:2.3}
设 $A \subset \mathbb{R}^{n}$. $A$ 的 Lebesgue 外测度 $m^{*}(A)$ 是由
\[
m^{*}(A) = \Gamma(A) \equiv \inf\left\{\sum_{j=1}^{\infty}|R_{j}| : \text{各 } R_{j} \text{ 是直方体且 } A \subset \bigcup_{j=1}^{\infty}R_{j}\right\}
\]
给出的属于 $[0, \infty]$ 的量.
\end{definition}

\begin{definition}[Lebesgue 可测集合]\label{def:2.4}
当对 $\mathbb{R}^{n}$ 的所有子集 $F$ 都有 $\Gamma(F) = \Gamma(F \cap E) + \Gamma(F \cap E^{c})$ 成立时, 称 $E$ 为 Lebesgue 可测集合, 或 Lebesgue 集合. 此时规定 $|E| = \Gamma(E)$.
\end{definition}

在 Lebesgue 测度的性质中, 称为内正则性的性质很方便, 复习一下这个性质.

\begin{theorem}[Lebesgue 测度的内正则性]\label{thm:2.1}
Lebesgue 测度的内正则性是指, 对任意可测集合 $A$,
\begin{equation}\tag{2.1}
|A| = \sup\{|K| : K \text{ 是包含于 } A \text{ 的紧集}\}
\label{eq:2.1}
\end{equation}
成立.
\end{theorem}

\begin{proof}
这里不作证明. 请参考适当的文献.
\end{proof}

\begin{definition}\label{def:2.5}
设 $E$ 为 Lebesgue 可测集合. $f: E \to \overline{\mathbb{R}} = \mathbb{R} \cup \{\pm\infty\}$ 是 Lebesgue 可测函数, 是指对任意 $\lambda \in \mathbb{R}$,
\[
f^{-1}(\lambda, \infty] = \{x \in E : f(x) > \lambda\}
\]
是 Lebesgue 可测集合.
\end{definition}

\subsection*{非负值 Lebesgue 可测函数积分的定义}

比区间上的积分更一般地, 考虑集合 $E$ 上的积分. 简单函数是指只取有限个值的可测函数, 且在某个直方体之外为 $0$ 的函数.

\begin{definition}[简单函数的积分]\label{def:2.6}
设 $E$ 为 Lebesgue 可测集合, $f: E \to [0, \infty)$ 为非负值简单函数. 用 $a_{1}, a_{2}, \ldots, a_{N} \in \mathbb{R}$ 与 (不允许有公共部分的) Lebesgue 可测集合 $E_{1}, E_{2}, \ldots, E_{N}$ 表为 $f = \sum_{j=1}^{N}a_{j}\chi_{E_{j}}$ 时,
\[
\int_{E}f(x)\dd x = \sum_{j=1}^{N}a_{j}|E \cap E_{j}| = \sum_{j=1}^{N}a_{j}|E_{j}|
\]
成立.
\end{definition}

非负值函数是指取值于 $[0, \infty]$ 的函数.

\begin{definition}[非负值函数的积分]\label{def:2.7}
设 $E$ 为 Lebesgue 可测集合. 当 $f: E \to [0, \infty]$ 为非负值 Lebesgue 可测函数时,
\[
\int_{E}f(x)\dd x \equiv \sup\left\{\int_{E}g(x)\dd x : g \text{ 是简单函数且对所有 } x \in E \text{ 有 } 0 \leq g(x) \leq f(x)\right\} \quad (\geq 0)
\]
成立. $\int_{E}f(x)\dd x = \infty$ 的可能性也存在.
\end{definition}

\subsection*{单调收敛定理, Fatou 引理}

可测函数这个说法是笼统的, 用容易理解的函数来逼近是调和分析中也常出现的手法之一. 下面的定理经常使用.

\begin{theorem}[单调收敛定理]\label{thm:2.2}
设 $\{f_{j}\}_{j=1}^{\infty}$ 是在 Lebesgue 可测集合 $E$ 上定义的非负值 Lebesgue 可测函数的增列. 用 $\lim_{j\to\infty}f_{j}(x) = f(x)$ $(x \in E)$ 定义函数 $f$ 时,
\begin{equation}\tag{2.2}
\lim_{j \to \infty}\int_{E}f_{j}(x)\dd x = \int_{E}f(x)\dd x
\label{eq:2.2}
\end{equation}
成立.
\end{theorem}

\begin{proof}
这里不作证明. 请参考适当的文献.
\end{proof}

\begin{theorem}[Fatou 引理]\label{thm:2.3}
对在 Lebesgue 可测集合 $E$ 上定义的非负值 Lebesgue 可测函数列 $\{f_{j}\}_{j=1}^{\infty}$,
\begin{equation}\tag{2.3}
\int_{E}\liminf_{j \to \infty}f_{j}(x)\dd x \leq \liminf_{j \to \infty}\int_{E}f_{j}(x)\dd x
\label{eq:2.3}
\end{equation}
成立.
\end{theorem}

\begin{proof}
这里不作证明. 请参考适当的文献.
\end{proof}

给出 Fatou 引理的一个变形. 对函数 $f: (a, b) \to \mathbb{R}$ 与 $t_{0} \in [a, b]$,
\[
\liminf_{u \to t_{0}, u \in (a, b)}f(u) = \sup_{\varepsilon > 0}\left(\inf_{u \in (a, b) \cap (t_{0} - \varepsilon, t_{0} + \varepsilon)}f(u)\right)
\]
成立.

\begin{theorem}\label{thm:2.4}
设 $E$ 为 Lebesgue 可测集合. $F: E \times (a, b) \to [0, \infty]$ 是对每个 $t \in (a, b)$, $F(\cdot, t)$ 与 $\liminf_{u \to t_{0}, u \in (a, b)}F(\cdot, u)$ 成为 Lebesgue 可测的函数. 此时, 对任意 $a \leq t_{0} \leq b$,
\[
\liminf_{t \to t_{0}, t \in (a, b)}\int_{E}F(x, t)\dd x \geq \int_{E}\liminf_{t \to t_{0}}F(x, t)\dd x
\]
成立. (这里, $t_{0} = b$, $b = \infty$ 也可以.)
\end{theorem}

\begin{proof}
这里不作证明. 请参考适当的文献.
\end{proof}

\subsection*{Lebesgue 可积函数}

作为调和分析的课题之一, 可举出函数各种可积性的测量. 这里暂且引入绝对可积的定义. 设 $E$ 为 Lebesgue 可测集合. Lebesgue 可测函数 $f: E \to [-\infty, \infty]$ 可积与 $|f|$ 可积等价.

$L^{1}(\mathbb{R}^{n})$ 表示满足
\[
\|f\|_{L^{1}} = \int_{\mathbb{R}^{n}}|f(x)|\dd x < \infty
\]
的可测函数 $f$ 全体构成的集合. 又, $L^{\infty}(\mathbb{R}^{n})$ 表示满足
\[
\|f\|_{L^{\infty}} = \inf\{\lambda > 0 : \text{在测度 } 0 \text{ 的集合 } N \text{ 之外, } |f(x)| < \lambda\} < \infty
\]
的全体构成的集合. 这里, $\|f\|_{L^{1}}$, $\|f\|_{L^{\infty}}$ 分别称为 $L^{1}$-范数, $L^{\infty}$-范数. 不过, 只在测度 $0$ 的集合上不同的函数视为同一. 这些范数在本书中经常使用. 记号 $\lesssim$ 是指 $\leq c$, 这里的 $c$ 在同一的定理、证明等中, 多个值也使用同一个字母. 例如 $3c \leq c$ 等就是其典型用例.

\begin{example}\label{exa:2.1}
$F(x) = (1 + |x|)^{-n-1}$, $x \in \mathbb{R}^{n}$ 是可积函数. 即 $F \in L^{1}(\mathbb{R}^{n})$. 实际上, 由单调收敛定理 (定理~\ref{thm:2.2}),
\[
\int_{\mathbb{R}^{n}}F(x)\dd x = \int_{|x| < 1}F(x)\dd x + \sum_{j=1}^{\infty}\int_{2^{j-1} < |x| < 2^{j}}F(x)\dd x
\]
成立, 但
\begin{align*}
\int_{2^{j-1} < |x| < 2^{j}}F(x)\dd x &\leq (1 + 2^{j-1})^{-n-1}\int_{2^{j-1} < |x| < 2^{j}}1\dd x \\
&\lesssim 2^{jn}(1 + 2^{j-1})^{-n-1} \\
&\leq 2^{jn-(j-1)(n+1)}
\end{align*}
成立, 所以把它关于 $j$ 相加即可.
\end{example}

\subsection*{Lebesgue 收敛定理及其变形}

\begin{theorem}[Lebesgue 收敛定理]\label{thm:2.5}
设 $E$ 为 Lebesgue 可测集合. 设 $f_{1}, f_{2}, \ldots: E \to [-\infty, \infty]$ 为逐点收敛到 $f$ 的 Lebesgue 可测函数列. 若存在可积函数 $g: E \to [-\infty, \infty]$ 满足对所有 $j \in \mathbb{N}$ 与 $x \in \mathbb{R}$ 有 $|f_{j}(x)| \leq g(x)$, 则
\begin{equation}\tag{2.4}
\lim_{j \to \infty}\int_{E}f_{j}(x)\dd x = \int_{E}f(x)\dd x
\label{eq:2.4}
\end{equation}
成立.
\end{theorem}

\begin{proof}
这里不作证明. 请参考适当的文献.
\end{proof}

\begin{theorem}\label{thm:2.6}
设 $E$ 为 Lebesgue 可测集合, $F: (a, b) \times E \to \mathbb{R}$ 是满足以下条件的函数.

(1) 对所有 $t \in (a, b)$ 与 $x \in \mathbb{R}$, $\partial_{t}F(t, x)$ 存在.

(2) $F(t, \cdot)$ 可积.

(3) 存在可积函数 $g: E \to \mathbb{R}$, 当 $x \in E$, $t \in (a, b)$ 时, $|\partial_{t}F(t, x)| \leq g(x)$ 成立.

此时,
\begin{equation}\tag{2.5}
\frac{\dd}{\dd t}\int_{E}F(t, x)\dd x = \int_{E}\partial_{t}F(t, x)\dd x
\label{eq:2.5}
\end{equation}
成立.
\end{theorem}

\begin{proof}
这里不作证明. 请参考适当的文献.
\end{proof}

\begin{theorem}\label{thm:2.7}
设 $E$ 为 Lebesgue 可测集合. 设 $F: E \times (a, b) \to \mathbb{R}$ 满足以下条件.

(A) 对所有 $t \in (a, b)$, $\int_{E}|F(x, t)|\dd x < \infty$ 成立.

(B) 对所有 $x \in E$, 函数 $t \in (a, b) \mapsto F(x, t) \in \mathbb{R}$ 连续.

(C) 存在某个可积函数 $G$, 对所有 $t \in (a, b)$ 与 $x \in E$, $|F(x, t)| \leq G(x)$ 成立.

此时,
\[
\lim_{t \to t_{0}}\int_{E}F(x, t)\dd x = \int_{E}F(x, t_{0})\dd x
\]
成立.
\end{theorem}

\begin{proof}
这里不作证明. 请参考适当的文献.
\end{proof}

\begin{theorem}\label{thm:2.8}
设 $E$ 为 Lebesgue 可测集合. 每固定 $t$, $F(\cdot, t)$ 可测的函数 $F: E \times (a, b) \to [0, \infty]$ 满足: 若 $a < t_{1} < t_{2} < b$, 则 $F(x, t_{1}) \leq F(x, t_{2})$. 此时,
\[
\lim_{t \uparrow b}\int_{E}F(x, t)\dd x = \int_{E}F(x, b)\dd x
\]
成立. (这里, $b = \infty$ 也可以.)
\end{theorem}

\begin{proof}
这里不作证明. 请参考适当的文献.
\end{proof}

\subsection*{Fubini 定理}

这里讨论等式
\begin{align*}
(\star): \quad \int_{\mathbb{R}^{m} \times \mathbb{R}^{n}}f(x, y)\dd x\dd y &= \int_{\mathbb{R}^{m}}\left(\int_{\mathbb{R}^{n}}f(x, y)\dd y\right)\dd x \\
&= \int_{\mathbb{R}^{n}}\left(\int_{\mathbb{R}^{m}}f(x, y)\dd x\right)\dd y
\end{align*}
的成立与否. 此后, 括弧写作``等式''时指这个等式. 这里, 关于等式, 也包含了
\[
x \mapsto \int_{\mathbb{R}^{n}}f(x, y)\dd y, \quad y \mapsto \int_{\mathbb{R}^{n}}f(x, y)\dd x
\]
的可测性.

用 $\mathcal{B}$ 表示 $\mathbb{R}^{n}$ 中包含开集的最小 $\sigma$-代数, 称为 Borel $\sigma$-代数. 也就是说, 把满足以下条件的集合族 $\mathcal{Z}$ 全体作为 $\mathcal{S}$:

\begin{itemize}
\item 直方体, $\mathbb{R}^{n}$ 中的空集合 $\emptyset$, $\mathbb{R}^{n}$ 本身属于 $\mathcal{Z}$. 即, 设 $R$ 为任意直方体时, $R, \emptyset, \mathbb{R}^{n} \in \mathcal{Z}$ 成立.

\item 若 $A \in \mathcal{Z}$, 则 $A^{\mathrm{c}} \in \mathcal{Z}$ 成立.

\item 若 $A_{1}, A_{2}, \ldots, A_{j}, \ldots \in \mathcal{Z}$, 则 $\bigcup_{j=1}^{\infty}A_{j} \in \mathcal{Z}$ 成立.
\end{itemize}

由此,
\[
\mathcal{B} = \{B \in 2^{\mathbb{R}^{n}} : \text{对所有 } \mathcal{Z} \in \mathcal{S} \text{ 有 } B \in \mathcal{Z} \text{ 成立}\}
\]
成立. 又, $\mathcal{B}$ 的集合称为 Borel 集合. 由于开集合可表为直方体的可数无限个并, 开集合和闭集合也都是 Borel 集合.

与 Lebesgue 可测函数一样, Borel 可测函数等也可用类似的方法定义. 即, $f: \mathbb{R}^{n} \to \overline{\mathbb{R}}$ 是 Borel 可测函数, 是指 $f^{-1}(\lambda, \infty]$ 是 Borel 可测集合. 例如, 连续函数是 Borel 可测函数. 设 $f$ 为 Borel 可测函数, 对 $f: \mathbb{R}^{m} \times \mathbb{R}^{n} \to [0, \infty]$ 或可积函数 $f: \mathbb{R}^{m} \times \mathbb{R}^{n} \to \mathbb{C}$, $(\star)$ 成立.

\subsection*{Lebesgue 积分与 Riemann 积分的比较}

已述在调和分析中自然地以 Lebesgue 积分的框架考虑积分, 但利用 Riemann 积分的知识来计算积分是重要的. 下面的定理成为这两种积分的桥梁.

\begin{theorem}\label{thm:2.9}
设 $f: [0, \infty) \to \mathbb{R}$ 为连续函数, 假设在广义 Riemann 积分的意义下, 或 Lebesgue 积分的意义下
\[
\int_{0}^{\infty}|f(x)|\dd x < \infty
\]
成立. 此时, 可以说以下事实.

(1) 在一方意义下成立的假设, 实际上在双方意义下都成立. 无论按哪种意义考虑, $\int_{0}^{\infty}|f(x)|\dd x$ 的值都相等.

(2) 在广义 Riemann 积分的意义下, 或 Lebesgue 积分的意义下
\[
\int_{0}^{\infty}f(x)\dd x
\]
存在, 且无论按哪种意义考虑值都相等.
\end{theorem}

\begin{proof}
这里不作证明. 请参考适当的文献.
\end{proof}

\subsection*{函数的支集}

为了不仅调查积分的大小, 还调查函数具有多大的扩展范围, 测量函数的大小是调和分析中的重要主题. 这里, 作为其基准, 说明函数的支集.

\begin{definition}[连续函数的支集]\label{def:2.8}
连续函数 $f: \mathbb{R}^{n} \to \mathbb{C}$ 的支集记作 $\operatorname{supp}(f)$, 定义为 $\{x \in \mathbb{R}^{n} : f(x) \neq 0\}$ 的闭包.
\end{definition}

在几何学教科书中, 把定义~\ref{def:2.8} 作为支集的定义给出, 但 Lebesgue 可测函数的支集也能定义. 此后, 对 $x \in \mathbb{R}^{n}$ 与 $r > 0$, 本书用 $B(x, r)$ 表示以 $x$ 为中心、半径为 $r$ 的开球. $x = 0$ 的情形省略``$x$, '', 记为 $B(r)$.

\begin{definition}[可测函数的支集]\label{def:2.9}
可测函数 $f: \mathbb{R}^{n} \to \mathbb{C}$ 的支集是指, 对任意 $r > 0$, 满足 $|\{y \in B(x, r) : f(y) \neq 0\}| > 0$ 的 $x \in \mathbb{R}^{n}$ 全体.
\end{definition}

对连续函数, 上述两种支集的定义一致. 虽然给出了两个定义, 但基于这一点, 两者都记作 $\operatorname{supp}(f)$. 用 $L_{\mathrm{c}}^{\infty}(\mathbb{R}^{n})$ 表示支集紧的 $L^{\infty}(\mathbb{R}^{n})$ 函数全体. 具有紧支集的函数一般多有可积性, 可以认为能以在有限测度空间上分析的感觉来处理函数. 恒等函数 $1$ 不可积, 而对球 $B(x, r)$, $\chi_{B(x, r)}$ 可积.

% ============ 第2章 §2.2 一般测度 (原书 pp.43-47) ============

\section{一般测度}\label{sec:2.2}
\index{0symbolsigma@$\sigma$-代数}

\index{H@荷重}

\index{C@测度}

\index{C@测度空间}










这里离开 Lebesgue 测度空间, 复习一般的测度空间. 本书中, 加权的情形与计数测度重要.

\subsection*{$\sigma$-代数}

调和分析中, 在配备 Lebesgue 测度的 Euclid 空间上分析是基本的, 但所考虑的集合不限于 Euclid 空间. 也有配备测度的距离空间的情形, 本书虽不涉及, 但也有 Riemann 流形的情形. 这里, 复习相当于为定义测度所需的定义域的 $\sigma$-代数, 以及广义化的测度的定义.

\begin{definition}[$\sigma$-代数]\label{def:2.10}
集合 $X$ 上的 $\sigma$-代数是指包含空集合、$X$ 本身, 并且对补集合、可数无限次的并及交的运算封闭的 $X$ 上的集合族.
\end{definition}

\begin{definition}[测度]\label{def:2.11}
设 $(X, \mathcal{M})$ 为可测空间. 即, $\mathcal{M}$ 是 $X$ 上的 $\sigma$-代数. $\mu$ 是测度是指

(1) $\mu: \mathcal{M} \to [0, \infty]$ 是集合函数.

(2) $\mu(\emptyset) = 0$.

(3) 当 $\{A_{j}\}_{j=1}^{\infty}$ 两两不交时,
\[
\mu\left(\sum_{j=1}^{\infty}A_{j}\right) = \sum_{j=1}^{\infty}\mu(A_{j})
\]
成立.
\end{definition}

\begin{definition}\label{def:2.12}
$(X, \mathcal{M}, \mu)$ 是测度空间是指

(1) $X$ 是集合.

(2) $\mathcal{M}$ 是 $X$ 上的 $\sigma$-代数.

(3) $\mu$ 是从 $\mathcal{M}$ 到 $[0, \infty]$ 的映射, 且是测度.
\end{definition}

作为 Lebesgue 测度以外的显著例子是下面这个. $2^{X}$ 是集合 $X$ 的全体子集构成的集合族.

\begin{definition}\label{def:2.13}
$\mu$ 是集合 $X$ 上的计数测度是指

(1) $\mu: 2^{X} \to [0, \infty]$ 是集合函数.

(2) $\mu(A) = \sharp A$.

成立.
\end{definition}

利用计数测度, 和作为积分被广义化. 作为从测度变换为别的测度的标准手法, 有加权的方法.

\begin{definition}\label{def:2.14}
设 $(X, \mathcal{M}, \mu)$ 为测度空间. $w: X \to (0, \infty)$ 为非负值可测函数时, 加权测度 $\nu = w\,\dd\mu$ 是指由
\[
\nu(E) = \int_{E}w(x)\dd\mu(x) \quad (E \in \mathcal{M})
\]
给出的测度.
\end{definition}

在这个定义中, $w$ 称为权.

\subsection*{几乎处处}

基于即使在测度 $0$ 的集合上函数不同也不影响积分值这一点, 给出下面的定义.

\begin{definition}\label{def:2.15}
某个性质 $P$ 在 $\mu$-a.e. (almost everywhere) 下成立, 是指性质 $P$ 不成立的点全体的 $\mu$-测度为 $0$. $\mu$ 为 Lebesgue 测度时简称为 a.e.. a.e. 称为``几乎处处''、``几乎所有'', 中文术语按副词、形容词分别使用.
\end{definition}

\subsection*{积分的定义等}

函数的可测性、单调收敛定理等也可以与 Lebesgue 测度的情形一样地定义. 这里给出其一例.

\begin{definition}\label{def:2.16}
设 $(X, \mathcal{M}, \mu)$ 为测度空间. $f: X \to [-\infty, \infty]$ 是 $\mathcal{M}$-可测 (或可测, $\mu$-可测) 是指
\[
\{f > \alpha\} = \{x \in X : f(x) > \alpha\} \in \mathcal{M} \quad (\alpha \in \mathbb{R})
\]
成立.
\end{definition}

\begin{theorem}[单调收敛定理]\label{thm:2.10}
设 $(X, \mathcal{M}, \mu)$ 为测度空间. 设 $E$ 为可测集合, $f_{1}, f_{2}, \ldots: E \to \overline{\mathbb{R}} = \mathbb{R} \cup \{\pm\infty\}$ 为可测函数. 若 $0 \leq f_{1} \leq f_{2} \leq \cdots$, 则
\[
\lim_{j \to \infty}\int_{E}f_{j}(x)\dd\mu(x) = \int_{E}\lim_{j \to \infty}f_{j}(x)\dd\mu(x)
\]
成立.
\end{theorem}

\begin{proof}
这里不作证明. 请参考适当的文献.
\end{proof}

\begin{theorem}[Fatou 引理]\label{thm:2.11}
设 $(X, \mathcal{M}, \mu)$ 为测度空间. 设 $E$ 为可测集合, $f_{1}, f_{2}, \ldots: E \to \overline{\mathbb{R}} = \mathbb{R} \cup \{\pm\infty\}$ 为可测函数. 若
\[
0 \leq f_{1}, f_{2}, \ldots,
\]
则
\[
\int_{E}\liminf_{j \to \infty}f_{j}(x)\dd\mu(x) \leq \liminf_{j \to \infty}\int_{E}f_{j}(x)\dd\mu(x)
\]
成立.
\end{theorem}

\begin{proof}
这里不作证明. 请参考适当的文献.
\end{proof}

\begin{theorem}[Lebesgue 收敛定理]\label{thm:2.12}
设 $(X, \mathcal{M}, \mu)$ 为测度空间. 设 $E$ 为可测集合, $f_{1}, f_{2}, \ldots: E \to \overline{\mathbb{R}} = \mathbb{R} \cup \{\pm\infty\}$ 为可测函数. 又设 (极限存在且) $\lim_{j \to \infty}f_{j}(x) = f(x)$ 在 $\mu$-a.e. 下成立. 又, 若可积函数 $g$ 对各自然数 $j$ 与 $x \in E$ 满足 $|f_{j}(x)| \leq g(x)$, 则
\[
\lim_{j \to \infty}\int_{E}f_{j}(x)\dd\mu(x) = \int_{E}\lim_{j \to \infty}f_{j}(x)\dd\mu(x)
\]
成立.
\end{theorem}

\begin{proof}
这里不作证明. 请参考适当的文献.
\end{proof}

\begin{theorem}\label{thm:2.13}
设 $(X, \mathcal{M}, \mu)$ 为测度空间. 若对 $a < t < b$, $f_{t}: E \to \overline{\mathbb{R}} = \mathbb{R} \cup \{\pm\infty\}$ 已被定义, 且满足以下条件, 则
\[
\frac{\dd}{\dd t}\int_{E}f_{t}(x)\dd\mu(x) = \int_{E}\partial_{t}f_{t}(x)\dd\mu(x)
\]
成立.

(1) 对 $t \in (a, b)$, $f_{t}$ 可积.

(2) $f_{t}(x)$ 关于 $\mu$-a.e. $x$ 可偏导.

(3) 存在可积函数 $g$, 当 $t \in (a, b)$ 时, 对 $\mu$-a.e. 有 $|\partial_{t}f_{t}(x)| \leq g(x)$ 成立.
\end{theorem}

\begin{proof}
这里不作证明. 请参考适当的文献.
\end{proof}

\subsection*{Fubini 定理}

作为特殊情形, 下面的概念重要. 特别是, $\sigma$-有限这一概念在考虑 Fubini 定理时重要. 在几乎所有情形下, 这个 $\sigma$-有限的条件都被满足.

\begin{definition}\label{def:2.17}
设 $(X, \mathcal{M}, \mu)$ 为测度空间.

(1) $\mu$ 是有限测度是指满足 $\mu(X) < \infty$.

(2) $\mu$ 是 $\sigma$-有限测度是指存在 $\mu$-测度有限的 $X_{1}, X_{2}, \ldots \in \mathcal{M}$ 使得
\[
X = \bigcup_{j=1}^{\infty}X_{j}.
\]
\end{definition}

单调收敛定理、Fatou 引理等特别地不需要对测度空间加限制, 但只有 Fubini 定理需要注意, 需要两个测度空间, 且这两个 $(X, \mathcal{M}, \mu)$ 与 $(Y, \mathcal{N}, \nu)$ 必须是 $\sigma$-有限的. 用 $\mathcal{M} \otimes \mathcal{N}$ 表示包含所有 $M \times N$ ($M \in \mathcal{M}$, $N \in \mathcal{N}$) 形状集合的 $X \times Y$ 上的最小 $\sigma$-代数. 又, 测度 $\mu \otimes \nu$ 满足 $\mu \otimes \nu(M \times N) = \mu(M)\nu(N)$, $M \in \mathcal{M}$, $N \in \mathcal{N}$.

\begin{theorem}[Fubini 定理]\label{thm:2.14}
设 $(X, \mathcal{M}, \mu)$ 与 $(Y, \mathcal{N}, \nu)$ 为 $\sigma$-有限的测度空间. $f$ 为非负值或复数值可积函数时,

(1) $f_{x} = f(x, \cdot)$ 关于 $\mu$-a.e. $x \in X$ 可积.

(2) 函数 $x \in X \mapsto \int_{Y}f(x, y)\dd\nu(y) \in [0, \infty]$ ($f$ 为非负值时), $x \in X \mapsto \int_{Y}f(x, y)\dd\nu(y) \in \mathbb{C}$ ($f$ 为复数值可积函数时) 是可积函数.

(3) 等式
\begin{align*}
\int_{X}\left(\int_{Y}f(x, y)\dd\nu(y)\right)\dd\mu(x) &= \int_{Y}\left(\int_{X}f(x, y)\dd\mu(x)\right)\dd\nu(y) \\
&= \int_{X \times Y}f(x, y)\dd\mu \otimes \nu(x, y)
\end{align*}
成立.
\end{theorem}

\begin{proof}
这里不作证明. 请参考适当的文献.
\end{proof}

\begin{example}\label{exa:2.2}
考虑配备计数测度的 $\mathbb{N}$.

(1) 函数 $f: \mathbb{N} \to \mathbb{C}$ 与复数列 $\{f(n)\}_{n=1}^{\infty}$ 具有完全相同的信息. 因此, $\mathbb{N}$ 上的函数与复数值数列被同一化.

(2) $L^{1}$-范数后述, 上述 $f$ 的 $L^{1}$-范数是 $\sum_{n=1}^{\infty}|f(n)| = \|\{f(n)\}_{n=1}^{\infty}\|_{\ell_{(\mathbb{N})}^{1}}$.

(3) 基于以上事实, 在这个测度空间上考虑 Fatou 引理时, 可如下表述. 对取值于 $[0, \infty]$ 的二重数列 $\{a_{m,n}\}_{m,n=1}^{\infty}$,
\[
\sum_{m=1}^{\infty}\liminf_{n \to \infty}a_{m,n} \leq \liminf_{n \to \infty}\sum_{m=1}^{\infty}a_{m,n}
\]
成立.
\end{example}

\begin{example}\label{exa:2.3}
对取值于 $[0, \infty]$ 的二重数列 $\{a_{m,n}\}_{m,n=1}^{\infty}$, 由 Fubini 定理, $\sum_{n=1}^{\infty}\sum_{m=1}^{\infty}a_{m,n} = \sum_{m=1}^{\infty}\sum_{n=1}^{\infty}a_{m,n}$ 成立.
\end{example}

% ============ 第2章 §2.3 Lebesgue 空间 (原书 pp.47-51) ============

\section{Lebesgue 空间}\label{sec:2.3}

不仅限于通常的绝对可积, 考虑种种可积性是有意义的. 例如, $L^{2}$ 空间具有内积这一运算, 由此得到在 $L^{1}$ 空间中不成立的反射性等. 这种性质得以成立的背景, 是由范数得到的开球的几何性质. 本书除 $L^{2}$ 之外不深入讨论这些. 又, 从调和分析学的立场, 仅考虑 $L^{1}$ 空间不够充分的迫切理由, 是 $L^{1}$ 中许多算子不连续, 而在 $L^{p}$ 空间 ($1 < p < \infty$) 中连续这一事实. 这里, 先列举考察 $L^{p}$ 空间所必需的积分不等式, 再定义 $L^{p}$ 空间并复习其性质.

\subsection*{积分不等式}

已述在调和分析学中, 一边大量使用积分一边进行种种分析, 但能精确求出积分值的情形几乎没有. 多数情形要看该积分值小于哪个数. 虽然无法求出值, 却要评价该积分的大小, 如本书所见, 这是项艰巨的作业. 作为基础的下述不等式是自然的, 但精确使用由它派生出的不等式是关键.

\begin{theorem}\label{thm:2.15}
设 $(X, \mathcal{B}, \mu)$ 为测度空间. 此时, 对可测函数 $f, g: X \to [0, \infty]$, 或可积函数 $f, g: X \to \mathbb{R}$, 若 $f(x) \leq g(x)$, 则积分的单调性
\[
\int_{X}f(x)\dd\mu(x) \leq \int_{X}g(x)\dd\mu(x)
\]
成立.
\end{theorem}

\begin{proof}
这里不作证明. 请参考适当的文献.
\end{proof}

由这个不等式得到的一连串基本不等式列举如下. 它们各自有名字, 本书中经常使用.

\begin{theorem}\label{thm:2.16}
设 $(X, \mathcal{B}, \mu)$ 为测度空间. 此时, 对可积函数 $f: X \to \mathbb{C}$, 积分的三角不等式
\[
\left|\int_{X}f(x)\dd\mu(x)\right| \leq \|f\|_{L^{1}(\mu)} = \int_{X}|f(x)|\dd\mu(x)
\]
成立.
\end{theorem}

\begin{proof}
这里不作证明. 请参考适当的文献.
\end{proof}

\begin{theorem}\label{thm:2.17}
设 $(X, \mathcal{B}, \mu)$ 为测度空间. 此时, 设 $f$ 为 $\mu$-可测函数且 $\lambda > 0$. 此时, Chebyshev 不等式
\begin{equation}\tag{2.6}
\mu\{|f| > \lambda\} \leq \frac{1}{\lambda}\int_{X}|f(x)|\dd\mu(x)
\label{eq:2.6}
\end{equation}
成立.
\end{theorem}

\begin{proof}
这里不作证明. 请参考适当的文献.
\end{proof}

\subsection*{$L^{p}$-范数}

来定义 $L^{p}$-范数. 设 $1 \leq p < \infty$. 对可测函数 $f: X \to \mathbb{C}$ 或 $\overline{\mathbb{R}}$,
\[
\|f\|_{L^{p}} = \|f\|_{L^{p}(\mu)} = \left(\int_{X}|f(x)|^{p}\dd\mu(x)\right)^{\frac{1}{p}}
\]
成立. $\|f\|_{L^{p}}$, $\|f\|_{L^{p}(\mu)}$ 称为 $f$ 的 $L^{p}$-范数. 用 $L^{p}(\mu)$ 表示这个量为有限的 Lebesgue 可测函数全体.

对 Lebesgue 可测函数 $f: X \to \mathbb{C}$ 或 $\overline{\mathbb{R}}$,
\[
\|f\|_{L^{\infty}(\mu)} = \inf\{\lambda > 0 : |f(x)| \leq \lambda \text{ a.e.}\}
\]
成立. $\|f\|_{L^{\infty}}$, $\|f\|_{L^{\infty}(\mu)}$ 称为 $f$ 的 $L^{\infty}(\mu)$-范数, 或略称为 $L^{\infty}$-范数. 用 $L^{\infty}(\mu)$ 表示这个量为有限的 Lebesgue 可测函数全体.

在 $L^{p}(\mu)$ 的定义中, 当 $(X, \mathcal{M}, \mu) = (\mathbb{R}^{n}, \mathcal{L}, \dd x)$ 时, 把 $L^{p}(\mu)$ 写作 $L^{p}(\mathbb{R}^{n})$, 用 $\|f\|_{L^{p}}$ 表示其范数.

\subsection*{Minkowski 不等式}

下面复习本书中经常使用的 Minkowski 不等式.

\begin{theorem}\label{thm:2.18}
设 $1 \leq p < \infty$. 设 $f$, $g$ 为可测集合 $X$ 上的可测函数. 此时, Minkowski 不等式 $\|f + g\|_{L^{p}(\mu)} \leq \|f\|_{L^{p}(\mu)} + \|g\|_{L^{p}(\mu)}$ 成立.
\end{theorem}

\begin{proof}
这里不作证明. 请参考适当的文献. 不过, 指出定理~\ref{thm:1.1} 的证明被变换到积分的文脉, 成为定理~\ref{thm:2.18} 的证明.
\end{proof}

\subsection*{Hölder 不等式}

下面复习本书中经常使用的 Hölder 不等式.

\begin{theorem}\label{thm:2.19}
设 $1 \leq p < \infty$, $q$ 为 $p$ 的调和共轭. 设 $f$, $g$ 为测度空间 $(X, \mathcal{B}, \mu)$ 上的可测函数.

(1) Hölder 不等式 $\|fg\|_{L^{1}(\mu)} \leq \|f\|_{L^{p}(\mu)}\|g\|_{L^{q}(\mu)}$ 成立. 不过, 在右边需要时, 请按 $0 \times \infty = \infty \times 0 = 0$ 解释.

(2) 【$L^{p}(\mu)$ 空间的对偶性】 $\|f\|_{L^{p}(\mu)} = \max\{\|f \cdot g\|_{L^{1}(\mu)} : \|g\|_{L^{q}(\mu)} \leq 1\}$ 成立.

(3) 【$L^{\infty}(\mu)$ 空间的对偶性】 $\|f\|_{L^{\infty}(\mu)} = \sup\{\|f \cdot g\|_{L^{1}(\mu)} : \|g\|_{L^{1}(\mu)} \leq 1\}$ 成立.
\end{theorem}

\begin{proof}
这里不作证明. 请参考适当的文献. 不过, 指出定理~\ref{thm:1.2} 及其直后相关命题的证明被变换到积分的文脉, 成为定理~\ref{thm:2.19} 的证明.
\end{proof}

\subsection*{卷积}

卷积在积分算子中也属于基本的, 出现在许多数学问题中. 偏微分方程的解算子多数由卷积得到, 这是显示卷积重要性的典型例子. 又, 分析数据时, 有通过取前后 $7$ 个数据的平均来去除数据间噪声的方法, 可以说这是在 $\mathbb{Z}$ 等离散空间上的卷积.

数字虽然没有精确的可信度, 数字也随时在变动, 但想说调和分析的 $45\%$ 左右的算子与卷积相关联的心情是有的.

设 $f$, $g$ 为 $\mathbb{R}^{n}$ 上的适当可测函数, $f$, $g$ 的卷积 $f*g$ 由
\[
f * g(x) = \int_{\mathbb{R}^{n}}f(x - y)g(y)\dd y
\]
定义. 一般, 当
\[
\int_{\mathbb{R}^{n}}|f(x - y)g(y)|\dd y = \infty
\]
时 $f * g(x)$ 没有意义, 但暂且只要有意义就考虑 $f * g(x)$. 本书特别是
\[
\int_{\mathbb{R}^{n}}|f(x - y)g(y)|\dd y = \infty
\]
的情形也多有考虑. 即使在这种情形, 在 $\pm$ 相互抵消的状况下也能考虑卷积, 许多情况下这极为有用, 在数学的其他领域带来重要结果. 设 $n \geq 3$, 偏微分方程之一的 Laplace 方程 $-\Delta u = f$ 的解具有
\[
u(x) = c_{n}\int_{\mathbb{R}^{n}}\frac{1}{|x - y|^{n-2}}f(y)\dd y \quad (x \in \mathbb{R}^{n})
\]
的形式, 所以可以看到卷积的形式. Schr\"odinger 方程 $\partial_{t}u - i\Delta u = 0$ 的解在初始值 $u_{0}$ 之下为
\[
u(x, t) = \frac{1}{\sqrt{(4\pi i t)^{n}}}\int_{\mathbb{R}^{n}}u_{0}(y)\exp\left(-\frac{|x - y|^{2}}{4it}\right)\dd y \quad (x \in \mathbb{R}^{n}, t > 0).
\]

卷积也出现在数学的其他领域. 设 $X$, $Y$ 为具有概率密度 $f$, $g$ 的独立随机变量. 因为 $X$ 具有概率密度 $f$,
\[
P(X > \lambda) = X \text{ 取大于 } \lambda \text{ 的值的概率} = \int_{\lambda}^{\infty}f(t)\dd t
\]
成立. 对 $Y$ 与 $g$ 也有对应的关系式成立. 此时, $X + Y$ 的概率密度由 $f * g$ 给出.

在上述情形中 $f * g(x)$ 没有意义时, 形式上令 $f * g(x) = 0$ 等, $f * g$ 暂且视为处处有定义.

\begin{theorem}[Young 不等式]\label{thm:2.20}
当 $1 \leq p, q, r \leq \infty$ 满足 $\frac{1}{r} = \frac{1}{p} + \frac{1}{q} - 1$ 时, 对 $f \in L^{p}(\mathbb{R}^{n})$, $g \in L^{q}(\mathbb{R}^{n})$, $f * g$ 有意义, 且 $\|f * g\|_{L^{r}} \leq \|f\|_{L^{p}}\|g\|_{L^{q}}$ 成立.
\end{theorem}

\begin{proof}
这里不作证明. 请参考适当的文献.
\end{proof}

\begin{exercise}\label{exer:2.1}
对 $A, B \geq 0$ 与 $0 < p < \infty$, 通过考虑 $A \leq B$ 与 $B \leq A$ 两种情形, 证明 $2^{-1}(A^{p} + B^{p}) \leq (A + B)^{p} \leq 2^{p}(A^{p} + B^{p})$. 【注意】由于具有对称性, 可假定前者来考虑.
\end{exercise}

\begin{exercise}\label{exer:2.2}
连续函数 $f: (0, \infty) \to \mathbb{R}$ 的 $L^{p}$-范数由
\[
\|f\|_{L^{p}(0, \infty)} = \left(\int_{0}^{\infty}|f(t)|^{p}\dd t\right)^{\frac{1}{p}}
\]
给出. 设 $a$ 为实数, 回答以下问题.

(1) 设 $f(x) = (1 + x)^{-1}$. 计算 $f$ 的 $L^{1}(0, \infty)$-范数.

(2) 设 $-\infty < a < 1$. 设 $f(x) = (1 + x)^{-a}$. 计算 $f$ 的 $L^{1}(0, \infty)$-范数.

(3) 设 $a > 1$. 设 $f(x) = (1 + x)^{-a}$. 计算 $f$ 的 $L^{1}(0, \infty)$-范数.

(4) 设 $f(x) = (1 + x)^{-a}$. $f$ 属于 $L^{1}(0, \infty)$ 的必要充分条件是什么?

(5) 设 $f(x) = (1 + x)^{-a}$. $f$ 属于 $L^{2}(0, \infty)$ 的必要充分条件是什么?
\end{exercise}

\begin{exercise}\label{exer:2.3}
令 $B(1) = \{|x| < 1\}$. 写出 $L^{p}(\mathbb{R}^{n})$ 的定义, 并证明当 $a < -\frac{n}{p}$ 时 $f(x) = \chi_{\mathbb{R}^{n}\setminus B(1)}(x)|x|^{a}$ 属于 $L^{p}(\mathbb{R}^{n})$.
\end{exercise}

% ============ 第2章 §2.4 赋范空间 (原书 pp.51-54) ============

\section{赋范空间}\label{sec:2.4}
\index{0symbolX*@$X^{*}$}

\index{0symbolAnormB@$\lVert A \rVert_{B(X,Y)}$}

\index{0symbolAnormXY@$\lVert A \rVert_{X \to Y}$}

\index{J@距离函数}

\index{J@距离空间}

\index{S@算子范数}

\index{X@线性泛函}

\index{D@对偶空间}

\index{F@赋范空间}

\index{B@Banach 空间}






















从这里开始转入函数解析的内容. 关于 Lebesgue 积分, 读本书时希望读者了解的事项已不作证明地给出, 但函数解析将给出证明. 函数解析是非常广范的数学领域, 书籍极多, 但本书不使用很多函数解析知识, 所以简洁地总结函数解析所需的知识.

\subsection*{距离空间, 赋范空间, Banach 空间}

本书因处理 Fourier 变换的关系, 对复值 Banach 空间感兴趣. 首先复习距离空间的定义. $(X, d)$ 是距离空间是指满足以下条件.

(1) $X$ 是集合.

(2) $d: X \times X \to [0, \infty)$ 是称为距离函数的映射. 即, 对所有 $x, y, z \in X$ 以下条件成立.

(A) $d(x, x) = 0$, 且 $x \neq y$ 时 $d(x, y) > 0$.

(B) $d(x, y) = d(y, x)$.

(C) $d(x, z) \leq d(x, y) + d(y, z)$. 这个不等式称为三角不等式.

作为距离空间的特殊例子, 赋范空间与具备完备性的赋范空间即 Banach 空间重要.

\begin{definition}[赋范空间]\label{def:2.18}
设 $X$ 为复线性空间. 映射 $\|\cdot\|: x \in X \mapsto \|x\| \in [0, \infty)$ 是范数是指满足以下 $2$ 个条件.

(1) 对 $x, y \in X$, 令 $d(x, y) = \|x - y\|$ 时, $d$ 是距离函数.

(2) 对所有 $x \in X$ 与 $a \in \mathbb{C}$, $\|ax\| = |a|\cdot\|x\|$ 成立.
\end{definition}

配备范数的复线性空间称为复赋范空间. 实赋范空间等也通过替换系数体来定义. 不特别注明系数体时简称为赋范空间.

考虑数列空间时没有深入思考完备性, 但无限维赋范空间未必完备 (参照练习问题~\ref{exer:1.30}), 所以明确说具有完备性很重要. 基于这一点, 定义 Banach 空间.

\begin{definition}[Banach 空间]\label{def:2.19}
在赋范空间中, 当对应的距离完备时, 该空间称为 Banach 空间. 即, 赋范空间 $X$ 是 Banach 空间是指满足以下条件: 对任意 $\varepsilon > 0$, 存在某个自然数 $N$, 使得若 $m, n > N$ 则 $\|x_{m} - x_{n}\| < \varepsilon$ 的 $X$ 的点列 $\{x_{n}\}_{n=1}^{\infty}$ 收敛.
\end{definition}

由于属于测度论的范畴, 下面的定理省略证明.

\begin{theorem}\label{thm:2.21}
当 $1 \leq p \leq \infty$, $(X, \mathcal{B}, \mu)$ 为测度空间时, $L^{p}(\mu)$ 是 Banach 空间.
\end{theorem}

\begin{proof}
这里不作证明. 请参考适当的文献.
\end{proof}

\subsection*{有界线性算子}

函数解析中重要的概念有有界线性算子. 关于这个概念, 已通过数列看过种种例子. 调和分析学中, 研究函数与算子的性质, 通过有界线性算子这一概念表达其结论的情形很多.

\begin{definition}\label{def:2.20}
设 $X$, $Y$ 为赋范空间. 线性算子 $A: X \to Y$ 是有界算子, 是指
\[
\|A\|_{X \to Y} = \|A\|_{B(X, Y)} = \sup\{\|Ax\|_{Y} : \|x\|_{X} = 1\}
\]
有限. 这个量称为算子范数. 把有界线性算子 $A: X \to Y$ 的全体构成的集合写作 $B(X, Y)$. $X = Y$ 时省略 $Y$, 写作 $B(X)$. $Y = \mathbb{C}$ 时, 把 $B(X, Y)$ 写作 $X^{*}$, 称为 $X$ 的对偶空间, $X^{*}$ 的元称为线性泛函.
\end{definition}

与赋范空间一样, 对实赋范空间 $X$, 也用相同记号 $X^{*}$ 表示 $B(X, \mathbb{R})$.

研究数列空间时没有意识到这一点, 但下面连接连续性与有界性的定理是基本的.

\begin{theorem}\label{thm:2.22}
设 $X$, $Y$ 为赋范空间. 线性算子 $A: X \to Y$ 是有界算子与 $A: X \to Y$ 连续等价.
\end{theorem}

\begin{proof}
假设 $A$ 有界, 由
\[
\|Ax - Ay\|_{Y} = \|A(x - y)\|_{Y} \leq \|A\|_{X \to Y}\|x - y\|_{X} \quad (x, y \in X)
\]
可知 $A$ 的连续性.

若 $A$ 连续, 则 $\{x \in X : \|Ax\|_{Y} < 1\} = A^{-1}(\{y \in Y : \|y\|_{Y} < 1\})$ 是包含 $0$ 的 $X$ 的开集合. 因此存在某个 $\delta > 0$, 若 $\|x\|_{X} < \delta$ 则 $\|Ax\|_{Y} < 1$. 由此得到 $\|A\|_{X \to Y} \leq \delta^{-1}$.
\end{proof}

用下面的方法构造有界线性算子的情形很多. $D$ 在 $X$ 中稠密是指 $\overline{D} = X$, 即 $D$ 的闭包与 $X$ 一致.

\begin{theorem}\label{thm:2.23}
设 $K > 0$ 为常数, $X$, $Y$ 为 Banach 空间, $D$ 是 $X$ 的稠密子集且为线性空间. $A$ 是从 $D$ 到 $Y$ 的线性算子, 若对所有 $x \in D$ 满足 $\|Ax\|_{Y} \leq K\|x\|_{X}$, 则存在唯一的有界线性算子 $B: X \to Y$, 对所有 $x \in D$ 有 $Bx = Ax$, 且对所有 $x \in X$ 有 $\|Bx\|_{Y} \leq K\|x\|_{X}$ 成立.
\end{theorem}

把上述两处 $\leq K$ 改写为 $= K$ 也可以, 但那种情况下的证明类似, 省略.

\begin{proof}
首先构造有界线性算子 $B: X \to Y$. 对 $x \in X$, 由 $D$ 的稠密性, 存在收敛于 $x$ 的 $D$ 的点列 $\{x_{k}\}_{k=1}^{\infty}$. 若 $k, k' \in \mathbb{N}$, 则 $\|Ax_{k} - Ax_{k'}\|_{Y} = \|A(x_{k} - x_{k'})\|_{Y} \leq K\|x_{k} - x_{k'}\|_{X}$, 所以 $\{Ax_{k}\}_{k=1}^{\infty}$ 在 $Y$ 中收敛. 又, 对收敛于 $x$ 的 $D$ 的其他点列 $\{v_{k}\}_{k=1}^{\infty}$, 若 $k \in \mathbb{N}$, 则 $\|Ax_{k} - Av_{k}\|_{Y} = \|A(x_{k} - v_{k})\|_{Y} \leq K\|x_{k} - v_{k}\|_{X}$, 所以 $\{Av_{k}\}_{k=1}^{\infty}$ 收敛到与 $\{Ax_{k}\}_{k=1}^{\infty}$ 相同的收敛目标. 因此, 与收敛于 $x$ 的 $D$ 的点列 $\{x_{k}\}_{k=1}^{\infty}$ 的取法无关, 可定义 $Bx = \lim_{k \to \infty}Ax_{k}$. 若 $x \in D$, 则对 $k = 1, 2, \ldots$ 令 $x_{k} = x$, $\{x_{k}\}_{k=1}^{\infty}$ 收敛于 $x$, 所以可知 $Bx = Ax$. 又, 对这种 $\{x_{k}\}_{k=1}^{\infty}$,
\[
\|Bx\|_{Y} = \lim_{k \to \infty}\|Ax_{k}\|_{Y} \leq K\lim_{k \to \infty}\|x_{k}\|_{X} = K\|x\|_{X}
\]
成立.

最后, 由 $D$ 的稠密性, 连续映射 $B$ 的定义由上述方法决定, 所以是唯一的.
\end{proof}

$X = Y = L^{p}(\mathbb{R}^{n})$ 时, 从 $X$ 到 $Y$ 的有界线性算子称为 $L^{p}$-有界 (线性) 算子.

关于算子的合成, 下面的公式经常使用.

\begin{theorem}\label{thm:2.24}
设 $X$, $Y$, $Z$ 为 Banach 空间. 对有界线性算子 $A: X \to Y$, $B: Y \to Z$, $\|BA\|_{X \to Z} \leq \|A\|_{X \to Y}\|B\|_{Y \to Z}$ 成立.
\end{theorem}

\begin{proof}
由于 $\|BAx\|_{Z} = \|B(Ax)\|_{Z} \leq \|B\|_{Y \to Z}\|Ax\|_{Y} \leq \|B\|_{Y \to Z}\|A\|_{X \to Y}\|x\|_{X}$, $x \in X$, 所以显然.
\end{proof}

举一个常用的例子.

\begin{example}\label{exa:2.4}
设 $(X, \mathcal{M}, \mu)$ 为 $\sigma$-有限测度空间, $g$ 为 $L^{\infty}(\mu)$ 函数. 此时, 对 $1 \leq p \leq \infty$, 乘法算子 $f \in L^{p}(\mu) \mapsto gf \in L^{p}(\mu)$ 满足 $\|gf\|_{L^{p}(\mu)} \leq \|g\|_{L^{\infty}(\mu)}\|f\|_{L^{p}(\mu)}$, 所以是有界线性算子, 但若 $N < \|g\|_{L^{\infty}(\mu)}$, 则 $E = \{x \in X : |g(x)| > N\}$ 具有正的测度, 特别由 $\sigma$-有限性, 存在可测集合 $F$ 使得 $0 < \mu(E \cap F) < \infty$ 成立. 令 $f = \chi_{F}$ 来考虑即可明白, 这个乘法算子的范数超过 $N$. 因此, 这个乘法算子的范数是 $\|g\|_{L^{\infty}(\mu)}$.
\end{example}

% ============ 第2章 §2.5 Hilbert 空间 (原书 pp.54-60) ============

\section{Hilbert 空间}\label{sec:2.5}
\index{0symbolAperp@$A^{\perp}$}

\index{0symbolAoplus@$A_{1} \oplus A_{2}$}

\index{0symbolAperp2@$A_{1} \perp A_{2}$}

\index{C@Cauchy--Schwarz 不等式}

\index{Z@准 Hilbert 空间}

\index{Z@正交直和}

\index{Y@有界自伴算子}

\index{Y@酉线性算子}


















调和分析学中, 算子不是 $L^{1}$-有界而是 $L^{2}$-有界的情形很多. 又, 利用 $L^{2}$-有界性导出 $L^{p}$-有界性的状况很多, 本书也会遭遇这种场面. 因此 $L^{2}$ 空间重要, 理解 $L^{2}$ 空间所具有的 Hilbert 空间结构也是重要的. 这里总结 Hilbert 空间. 附上证明列举 Hilbert 空间的基础知识.

\subsection*{准 Hilbert 空间}

一度忘掉完备性, 总结可代数处理的部分.

\begin{definition}\label{def:2.21}
设 $H$ 为复线性空间. $\langle \cdot, \cdot \rangle_{\mathcal{H}}: H \times H \to \mathbb{C}$ 是 Hermitian 内积, 是指对所有 $u, u_{1}, u_{2}, v, v_{1}, v_{2} \in H$ 与 $\alpha \in \mathbb{C}$ 满足以下条件.

(1) 若 $u \neq 0$, 则 $\langle u, u \rangle_{\mathcal{H}} > 0$. 【内积的正值性】

(2) $\langle \alpha u, v \rangle_{\mathcal{H}} = \langle u, \overline{\alpha}v \rangle_{\mathcal{H}} = \alpha\langle u, v \rangle_{\mathcal{H}}$.

(3) $\langle u, v \rangle_{\mathcal{H}} = \overline{\langle v, u \rangle_{\mathcal{H}}}$.

(4) $\langle u_{1} + u_{2}, v \rangle_{\mathcal{H}} = \langle u_{1}, v \rangle_{\mathcal{H}} + \langle u_{2}, v \rangle_{\mathcal{H}}$, $\langle u, v_{1} + v_{2} \rangle_{\mathcal{H}} = \langle u, v_{1} \rangle_{\mathcal{H}} + \langle u, v_{2} \rangle_{\mathcal{H}}$.

配备 Hermitian 内积的线性空间称为 (复) 准 Hilbert 空间. 又, 在上述定义中把 $\mathbb{C}$ 换成 $\mathbb{R}$ 而得到的运算 $\langle \cdot, \cdot \rangle_{\mathcal{H}}: \mathcal{H} \times \mathcal{H} \to \mathbb{R}$ 所配备的线性空间称为 (实) 准 Hilbert 空间. 又, 令 $\|x\| = \sqrt{\langle x, x \rangle_{\mathcal{H}}}$, $x \in \mathcal{H}$. (2), (3), (4) 称为双线性性.
\end{definition}

举一个用准 Hilbert 空间中的运算计算的例子.

\begin{example}\label{exa:2.5}
设 $H$ 为准 Hilbert 空间, $k \in \mathbb{C}$, $x, y \in H$. 此时,
\begin{align*}
\|x + ky\|^{2} &= \langle x + ky, x + ky \rangle_{\mathcal{H}} \\
&= \langle x, x \rangle_{\mathcal{H}} + \langle x, ky \rangle_{\mathcal{H}} + \langle ky, x \rangle_{\mathcal{H}} + \langle ky, ky \rangle_{\mathcal{H}} \\
&= \|x\|^{2} + \overline{k}\langle x, y \rangle_{\mathcal{H}} + k\overline{\langle x, y \rangle_{\mathcal{H}}} + |k|^{2}\|y\|^{2} \\
&= \|x\|^{2} + 2\mathrm{Re}(\overline{k}\langle x, y \rangle_{\mathcal{H}}) + |k|^{2}\|y\|^{2}
\end{align*}
成立.
\end{example}

准 Hilbert 空间的例子有下面这个. $L^{2}(\mathbb{R}^{n})$ 的内积取为
\[
\langle f, g \rangle_{L^{2}} = \int_{\mathbb{R}^{n}}f(x)\overline{g(x)}\dd x
\]
时, $L^{2}(\mathbb{R}^{n})$ 是准 Hilbert 空间. 不过, 如上所述, $L^{2}(\mathbb{R}^{n})$ 配备有完备距离, 所以由后述定义可知 $L^{2}(\mathbb{R}^{n})$ 是 Hilbert 空间. 本书不考虑是准 Hilbert 空间但不是 Hilbert 空间的例子, 所以只举这个例子.

\subsection*{Cauchy--Schwarz 不等式}

证明从 Hilbert 空间的性质可代数地看出的 Cauchy--Schwarz 不等式. 由此可知 $H$ 上可定义距离.

\begin{theorem}[Cauchy--Schwarz 不等式]\label{thm:2.25}
$H$ 为准 Hilbert 空间时, $|\langle x, y \rangle_{H}| \leq \|x\|\|y\|$, $x, y \in H$ 成立.
\end{theorem}

\begin{proof}
不妨设 $x, y \in H \setminus \{0\}$. 令 $\alpha = -\frac{\langle x, y \rangle_{H}}{\|y\|^{2}} \in \mathbb{C}$,
\[
\|x + \alpha y\|^{2} = \|x\|^{2} + 2\mathrm{Re}(\overline{\alpha}\langle x, y \rangle_{\mathcal{H}}) + |\alpha|^{2}\|y\|^{2} = \|x\|^{2} - \frac{|\langle x, y \rangle_{\mathcal{H}}|^{2}}{\|y\|^{2}}
\]
由于这个量非负, 所以显然.
\end{proof}

\subsection*{Hilbert 空间}

给定准 Hilbert 空间后, 由 Cauchy--Schwarz 不等式可知范数被确定, 但如前所述, 关于该范数完备时另有名称. 下面认真表述它.

\begin{definition}\label{def:2.22}
(复) Hilbert 空间是指配备 Hermitian 内积 $\langle \cdot, \cdot \rangle = \langle \cdot, \cdot \rangle_{H}$ 的内积空间, 且关于由此诱导的距离 $\|\cdot\| = \|\cdot\|_{H} = \sqrt{\langle \cdot, \cdot \rangle_{H}}$ 完备者.
\end{definition}

$H$ 的例子有 $L^{2}(\mathbb{R}^{n})$, $\ell^{2}(\mathbb{N})$. 此后设 $H$ 为 Hilbert 空间.

Hilbert 空间是 Banach 空间, 所以有界线性算子等概念自然可以设想, 但与 Hilbert 空间固有配备的内积相关联的算子变得重要. 重要的例子后面给出若干, 这里列举有界线性算子中的酉 (线性) 算子. $H$ 的线性算子 $U$ 是酉 (线性) 算子, 是指 $U$ 为双射且对所有 $x \in H$,
\begin{equation}\tag{2.7}
\|Ux\| = \|x\|
\label{eq:2.7}
\end{equation}
成立. 例如, 对 $\ell^{2}(\mathbb{N})$, 把 $\{a_{n}\}_{n=1}^{\infty} = (a_{1}, a_{2}, \ldots)$ 移到 $(a_{2}, a_{1}, a_{4}, a_{3}, a_{6}, a_{5}, \ldots)$ 的算子是酉算子.

此后, 设 $H$ 为 Hilbert 空间.

\subsection*{中线定理}

由代数计算可知下面的等式成立. Hilbert 空间的几何结构被概括在下面的中线定理中. 不作证明, 但已知满足中线定理的范数可用内积表示.

\begin{theorem}[中线定理]\label{thm:2.26}
对 $x, y \in H$, $\|x + y\|^{2} + \|x - y\|^{2} = 2\|x\|^{2} + 2\|y\|^{2}$ 成立.
\end{theorem}

\begin{proof}
由例~\ref{exa:2.5}, $\|x + y\|^{2} = \|x\|^{2} + \|y\|^{2} + 2\mathrm{Re}(\langle x, y \rangle_{\mathcal{H}})$ 与 $\|x - y\|^{2} = \|x\|^{2} + \|y\|^{2} - 2\mathrm{Re}(\langle x, y \rangle_{\mathcal{H}})$ 显然.
\end{proof}

\subsection*{Hilbert 空间上的有界线性算子}

设 $A$ 为从 $H$ 到 $H$ 的有界线性算子时, $\|A\|$ 已在定义~\ref{def:2.20} 中定义, 但给出由 $H$ 是 Hilbert 空间得到的另一个表示.

\begin{theorem}\label{thm:2.27}
对 $A \in B(\mathcal{H})$, $\|A\|_{\mathcal{H} \to \mathcal{H}} = \sup_{x, y \in \mathcal{H}, \|x\|, \|y\| \leq 1}|\langle Ax, y \rangle_{\mathcal{H}}|$ 成立.
\end{theorem}

\begin{proof}
若 $A = 0$, 则结论显然, 所以不妨设 $A \neq 0$. 由 Cauchy--Schwarz 不等式, 对 $x, y \in H$,
\[
|\langle Ax, y \rangle_{\mathcal{H}}| \leq \|Ax\|\|y\| \leq \|A\|_{\mathcal{H} \to \mathcal{H}}\|x\|\|y\|
\]
所以
\[
\|A\|_{\mathcal{H} \to \mathcal{H}} \geq \sup_{x, y \in \mathcal{H}, \|x\|, \|y\| \leq 1}|\langle Ax, y \rangle_{\mathcal{H}}|
\]
成立. 又,
\[
\sup_{x, y \in \mathcal{H}, \|x\|, \|y\| \leq 1}|\langle Ax, y \rangle_{\mathcal{H}}| = \sup_{x, y \in \mathcal{H} \setminus \{0\}}\frac{|\langle Ax, y \rangle_{\mathcal{H}}|}{\|x\|\|y\|}
\]
所以, 特别是对 $Ax \neq 0$ 的 $x \in H$, 通过令 $y = Ax$ 可部分到达上限,
\[
\frac{\|Ax\|}{\|x\|} = \frac{|\langle Ax, y \rangle_{\mathcal{H}}|}{\|x\|\|y\|}
\]
成立. 因此,
\[
\|A\|_{\mathcal{H} \to \mathcal{H}} \leq \sup_{x, y \in \mathcal{H}, \|x\|, \|y\| \leq 1}|\langle Ax, y \rangle_{\mathcal{H}}|
\]
成立.
\end{proof}

\subsection*{最短距离}

下面的定理保证图形的最短距离可达性. 有限维情形由于可用紧性, 结论显然. 但无限维情形不保证单位闭球的紧性 (参照练习问题~\ref{exer:1.29}), 即单位闭球不是全有界的, 所以证明需要技巧. 作为全有界性的代用物的是中线定理. 下面定理中出现的凸集合是指, 任意取该集合内的 $2$ 点, 以这 $2$ 点为端点的线段包含在该集合中.

\begin{theorem}\label{thm:2.28}
闭 (且) 凸集合 $C$ 与 $H$ 内的点 $x$ 的距离被 $C$ 内的点 $y$ 唯一地实现. 即, 存在唯一的 $y \in C$ 使得 $\|x - y\| = \inf_{z \in C}\|x - z\|$ 成立.
\end{theorem}

\begin{proof}
令 $d = \inf_{z \in C}\|x - z\|$. 对 $j \in \mathbb{N}$, 取 $y_{j} \in C$ 使得 $\|x - y_{j}\| \leq d + 2^{-j}$ 成立. 则由中线定理, 对各 $j, k \in \mathbb{N}$,
\[
\|2x - y_{j} - y_{k}\|^{2} + \|y_{j} - y_{k}\|^{2} = 2\|x - y_{j}\|^{2} + 2\|x - y_{k}\|^{2}
\]
成立. 这里, 由于 $C$ 是凸的, $\frac{y_{j} + y_{k}}{2} \in C$. 因此 $\|2x - y_{j} - y_{k}\| \geq 2d$, 得到
\[
\|y_{j} - y_{k}\|^{2} \leq 2(d + 2^{-j})^{2} + 2(d + 2^{-k})^{2} - 4d^{2}.
\]
令 $j, k \to \infty$, 则 $y_{j} - y_{k} \to 0$, 所以由 $H$ 的完备性, $\{y_{j}\}_{j=1}^{\infty}$ 收敛到某个 $y$. 因此 $y \in C$ 是所求的点.

假若 $y' \in C$ 也给出距离. 由于 $\frac{y + y'}{2} \in C$, $\|x - y'\| = d$,
\[
4d^{2} \leq \|2x - y - y'\|^{2} + \|y - y'\|^{2} = 2\|x - y\|^{2} + 2\|x - y'\|^{2} = 4d^{2}
\]
得到 $\|y - y'\| = 0$, 即 $y = y'$.
\end{proof}

\subsection*{正交补空间}

对 $H$ 的子集 $A \subset H$, 令 $A^{\perp} = \{x \in H : \langle x, z \rangle_{H} = 0, z \in A\}$. 对 $H$ 的子集 $A_{1}$, $A_{2}$, $A_{1} \perp A_{2}$ 是指对所有 $a_{1} \in A_{1}$, $a_{2} \in A_{2}$, $a_{1} \perp a_{2}$, 即 $\langle a_{1}, a_{2} \rangle_{H} = 0$ 成立. 进而, 当 $A_{1}$, $A_{2}$ 为线性空间时, 把和空间
\[
A_{1} + A_{2} = \{a_{1} + a_{2} \in \mathcal{H} : a_{1} \in A_{1}, a_{2} \in A_{2}\}
\]
表示为 $A_{1} \oplus A_{2}$. 这个 $A_{1} \oplus A_{2}$ 称为正交直和.

\begin{theorem}\label{thm:2.29}
$Z$ 是 $H$ 的闭线性子空间时, $H = Z \oplus Z^{\perp}$ 成立.
\end{theorem}

\begin{proof}
由内积的正值性, $Z \cap Z^{\perp} = \{0\}$, $H \supset Z + Z^{\perp}$, $Z \perp Z^{\perp}$, 即 $H \supset Z \oplus Z^{\perp}$ 显然. 因此必须证明 $H \subset Z \oplus Z^{\perp}$. 设 $x \in H$. 由定理~\ref{thm:2.28}, 存在 $z_{0} \in Z$ 使得 $z_{0}$ 给出 $x$ 与 $Z$ 之间的距离, 所以对所有 $z \in Z$, $\|x - z_{0}\| \leq \|x - z\|$ 成立. 通过把 $z$ 换成 $z_{0} + z$, 此条件与对 $z \in Z$ 有 $\|x - z_{0}\| \leq \|x - z_{0} - z\|$ 成立等价. 因此由例~\ref{exa:2.5},
\[
\|z\|^{2} \geq 2\mathrm{Re}(\langle x - z_{0}, z \rangle_{\mathcal{H}})
\]
成立. 通过把 $z$ 换成 $\kappa z$, $\kappa > 0$, 可知
\[
\kappa\|z\|^{2} \geq 2\mathrm{Re}(\langle x - z_{0}, z \rangle_{\mathcal{H}})
\]
成立. 令 $\kappa \downarrow 0$, 得 $0 \geq 2\mathrm{Re}(\langle x - z_{0}, z \rangle_{\mathcal{H}})$ 成立. 由于 $z$ 也可以换成 $-z$, $0 \leq 2\mathrm{Re}(\langle x - z_{0}, z \rangle_{\mathcal{H}})$ 也成立. 因此 $2\mathrm{Re}(\langle x - z_{0}, z \rangle_{\mathcal{H}}) = 0$. 由于 $z$ 也可以换成 $iz$, $2\mathrm{Im}(\langle x - z_{0}, z \rangle_{\mathcal{H}}) = 0$ 成立. 由这两个等式, 得 $\langle x - z_{0}, z \rangle_{\mathcal{H}} = 0$. 综上, 得 $x - z_{0} \perp z$. 因此 $x = z_{0} + (x - z_{0}) \in Z \oplus Z^{\perp}$, 结论得证.
\end{proof}

由代数计算可知, $\perp$ 取两次会回到原处.

\begin{corollary}\label{cor:2.1}
$Z$ 是 $H$ 的闭线性子空间时, $(Z^{\perp})^{\perp} = Z$ 成立.
\end{corollary}

\begin{proof}
$\mathcal{H} = Z \oplus Z^{\perp} = (Z^{\perp})^{\perp} \oplus Z^{\perp}$ 成立. 由 $Z^{\perp}$, $(Z^{\perp})^{\perp}$ 的定义, $(Z^{\perp})^{\perp} \supset Z$ 显然. 于是, 为证明 $(Z^{\perp})^{\perp} \subset Z$, 任取 $z \in (Z^{\perp})^{\perp}$. 则 $z = z_{0} + z_{1}$, $z_{0} \in Z$, $z_{1} \in Z^{\perp}$ 的分解是可能的. 因为 $z_{0}$, $z$ 均与 $z_{1}$ 正交, 在 $\langle z_{1}, z_{1} \rangle_{\mathcal{H}}$ 中将右侧的 $z_{1}$ 代入 $z_{1} = z - z_{0}$,
\[
\langle z_{1}, z_{1} \rangle_{\mathcal{H}} = \langle z_{1}, z \rangle_{\mathcal{H}} - \langle z_{1}, z_{0} \rangle_{\mathcal{H}} = 0
\]
所以得到 $z_{1} = 0$. 即 $z = z_{0} \in Z$.
\end{proof}

把系~\ref{cor:2.1} 稍微放宽, 以便使用.

\begin{corollary}\label{cor:2.2}
$Z$ 是不一定是闭子空间的 $H$ 的线性子空间时, $(Z^{\perp})^{\perp} = \overline{Z}$ 成立.
\end{corollary}

\begin{proof}
$\overline{Z}$ 是闭线性空间, 由系~\ref{cor:2.1}, $\overline{Z} = (\overline{Z}^{\perp})^{\perp} = (Z^{\perp})^{\perp}$, 即得.
\end{proof}

\subsection*{Riesz 表示定理}

与许多数学领域一样, 调和分析学中利用对偶性也是基本的. 已经说明的 $L^{p}$ 与 $L^{p'}$ 的对偶性即其一例. 下面证明相当于 Hilbert 空间中的自对偶性的 Riesz 表示定理. 回忆: 对线性空间 $U$, $V$ 与线性映射 $f: U \to V$, $f$ 的核 $\ker(f)$ 定义为 $\ker(f) = \{x \in U : f(x) = 0\}$, $f$ 的像 $\operatorname{im}(f)$ 定义为 $\operatorname{im}(f) = \{f(x) \in V : x \in U\}$. 证明中出现的同构定理是指断言 $\operatorname{im}(f) \simeq U/\ker(f)$ 成立的定理.

\begin{theorem}[Riesz 表示定理]\label{thm:2.30}
从 $H$ 到 $\mathbb{C}$ 的连续线性泛函可借助唯一的 $y \in H$ 由内积表示为 $\langle \cdot, y \rangle_{H}$ 的形式. 即, $H^{*}$ 与 $H$ 被同一化.
\end{theorem}

\begin{proof}
设 $\ell$ 为从 $\mathcal{H}$ 到 $\mathbb{C}$ 的连续线性泛函. $\ker(\ell) = \{x \in \mathcal{H} : \ell(x) = 0\} = \ell^{-1}(\{0\})$ 是 $\mathcal{H}$ 的闭子空间, 所以 $\mathcal{H} = \ker(\ell) \oplus \ker(\ell)^{\perp}$ 成立. 由同构定理, $\dim(\ker(\ell)^{\perp}) = \dim(\mathcal{H}/\ker(\ell)) = \dim(\operatorname{im}(\ell)) \leq 1$. 若 $\dim(\ker(\ell)^{\perp}) = 0$, 则 $\ker(\ell) = \mathcal{H}$, 所以 $\ell = 0$, $y = 0$.

若 $\dim(\ker(\ell)^{\perp}) = 1$, 取 $\ker(\ell)^{\perp}$ 的基底 $\{y_{0}\}$ 使 $\|y_{0}\| = 1$ 成立. 任意元 $x \in \mathcal{H}$ 满足 $x - \langle x, y_{0} \rangle_{\mathcal{H}}y_{0} \in (\ker(\ell)^{\perp})^{\perp} = \ker(\ell)$, 所以 $\ell(x) - \langle x, y_{0} \rangle_{\mathcal{H}}\ell(y_{0}) = 0$. 因此, 令 $y = \overline{\ell(y_{0})}y_{0}$, 得 $\ell(x) = \langle x, y \rangle_{\mathcal{H}}$.

最后, 与 $\dim(\ker(\ell)^{\perp})$ 的值无关, $y \in H$ 的唯一性由内积的正值性显然.
\end{proof}

\subsection*{伴随算子}

这里给出 Riesz 表示定理的应用. 伴随算子与对偶性相关联, 在调和分析学中重要.

\begin{theorem}\label{thm:2.31}
对从 $H$ 到 $H$ 的连续线性算子 $A$, 可以说以下事实.

(1) 称为 $A$ 的伴随算子的从 $H$ 到 $H$ 的连续线性算子 $A^{*}$ 唯一确定, 且 $\langle Ax, y \rangle_{H} = \langle x, A^{*}y \rangle_{H}$ 成立.

(2) $\|A^{*}A\|_{H \to H} = (\|A\|_{H \to H})^{2}$ 成立.

(3) 关于 $\ker(A^{*}) = \{x \in \mathcal{H} : A^{*}x = 0\}$ 与 $\operatorname{im}(A) = \{Ax \in \mathcal{H} : x \in \mathcal{H}\}$, 有关系 $\ker(A^{*}) = \operatorname{im}(A)^{\perp}$, $\ker(A^{*})^{\perp} = \overline{\operatorname{im}(A)}$ 成立.
\end{theorem}

\begin{proof}
(1) 固定 $x \in \mathcal{H}$. 线性映射 $\ell(y) = \langle Ay, x \rangle_{\mathcal{H}}$, $y \in \mathcal{H}$ 是从 $\mathcal{H}$ 到 $\mathbb{C}$ 的连续线性泛函, 所以由 Riesz 表示定理, 存在唯一的 $x' \in \mathcal{H}$ 满足 $\ell(y) = \langle y, x' \rangle_{\mathcal{H}}$. 由内积的性质 $\langle u, v \rangle_{\mathcal{H}} = \overline{\langle v, u \rangle_{\mathcal{H}}}$, $u, v \in \mathcal{H}$, 可知 $x \mapsto x' = A^{*}x$ 满足条件. 最后, 唯一性由内积的正值性显然.

(2) 由定理~\ref{thm:2.27} 与 Cauchy--Schwarz 不等式,
\begin{align*}
\|A^{*}A\|_{\mathcal{H} \to \mathcal{H}} &= \sup_{x, y \in \mathcal{H}, \|x\| \leq 1, \|y\| \leq 1}|\langle A^{*}Ax, y \rangle_{\mathcal{H}}| \\
&= \sup_{x, y \in \mathcal{H}, \|x\| \leq 1, \|y\| \leq 1}|\langle Ax, Ay \rangle_{\mathcal{H}}| \\
&= (\|A\|_{\mathcal{H} \to \mathcal{H}})^{2}.
\end{align*}

(3) 前者由伴随的定义显然. 后者用系~\ref{cor:2.2}, 对前者的关系取正交补空间即可.
\end{proof}

举调和分析中常出现的伴随算子的例子. 对函数 $g$, $\overline{g}$ 是复共轭. 另外, 这个算子的有界性由 Young 不等式得到.

\begin{example}\label{exa:2.6}
$g \in L^{1}(\mathbb{R})$ 时, 有界线性算子 $L^{2}(\mathbb{R}) \ni f \mapsto f * g \in L^{2}(\mathbb{R})$ 的伴随是 $L^{2}(\mathbb{R}) \ni f \mapsto f * \overline{g}(-\cdot) \in L^{2}(\mathbb{R})$. 参照练习问题~\ref{exer:4.4}.
\end{example}

\subsection*{Hilbert 空间上的有界自伴算子}

最后, 总结 Hilbert 空间上的有界自伴算子. 满足 $A = A^{*}$ 的算子 $A \in B(\mathcal{H})$ 称为有界自伴算子, 此类算子有称为谱分解的重要分解理论. 本书不需要谱分解, 但需要有界自伴算子的若干公式, 总结如下.

\begin{theorem}\label{thm:2.32}
设 $H$ 为 Hilbert 空间. 设 $A \in B(\mathcal{H})$ 为自伴算子. 此时, $\|A\|_{\mathcal{H} \to \mathcal{H}} = \sup_{x \in \mathcal{H}, \|x\| \leq 1}|\langle Ax, x \rangle_{\mathcal{H}}|$ 成立.
\end{theorem}

\begin{proof}
对任意范数不超过 $1$ 的 $x, y \in H$,
\[
\mathrm{Re}(\langle Ax, y \rangle_{\mathcal{H}}) = \frac{1}{4}\langle A(x + y), x + y \rangle_{\mathcal{H}} - \frac{1}{4}\langle A(x - y), x - y \rangle_{\mathcal{H}}
\]
成立. 因此由三角不等式与中线定理,
\begin{align*}
|\mathrm{Re}(\langle Ax, y \rangle_{\mathcal{H}})| &\leq \frac{1}{4}|\langle A(x + y), x + y \rangle_{\mathcal{H}}| + \frac{1}{4}|\langle A(x - y), x - y \rangle_{\mathcal{H}}| \\
&\leq \frac{1}{4}\sup_{u \in \mathcal{H}, \|u\| \leq 1}|\langle Au, u \rangle_{\mathcal{H}}|(\|x + y\|^{2} + \|x - y\|^{2}) \\
&\leq \sup_{u \in \mathcal{H}, \|u\| \leq 1}|\langle Au, u \rangle_{\mathcal{H}}|
\end{align*}
成立. 又, 另一方面, 由
\begin{align*}
&\sup_{x, y \in \mathcal{H}, \|x\| \leq 1, \|y\| \leq 1}|\mathrm{Re}(\langle Ax, y \rangle_{\mathcal{H}})| \\
&= \sup_{x, y \in \mathcal{H}, \theta \in \mathbb{R}, \|x\| \leq 1, \|y\| \leq 1}|\mathrm{Re}(\langle Ax, e^{i\theta}y \rangle_{\mathcal{H}})| \\
&= \sup_{x, y \in \mathcal{H}, \theta \in \mathbb{R}, \|x\| \leq 1, \|y\| \leq 1}|\mathrm{Re}(e^{-i\theta}\langle Ax, y \rangle_{\mathcal{H}})|
\end{align*}
得
\[
\sup_{x, y \in \mathcal{H}, \|x\| \leq 1, \|y\| \leq 1}|\mathrm{Re}(\langle Ax, y \rangle_{\mathcal{H}})| = \sup_{x, y \in \mathcal{H}, \|x\| \leq 1, \|y\| \leq 1}|\langle Ax, y \rangle_{\mathcal{H}}|
\]
成立. 因此由定理~\ref{thm:2.27}, 结论显然.
\end{proof}


% 第3章 Hardy-Littlewood 极大函数 (原书 pp.61-74, PDF pp.67-80)
\chapter{Hardy--Littlewood 极大函数}
% ============ 第3章 §3.1 Hardy--Littlewood 极大函数 (原书 pp.61-66) ============

可测函数的分析一般来说是极为困难的. 实际上, 作为无法把握可测函数这一概念的原因之一, 可以举出实在是有各种各样的东西都会成为可测的这一点. 因此, 寻找一种将可测函数在一定程度上放大, 使得变换后的函数易于处理的变换, 是很自然的. 作为这样的变换, 可以举出 Hardy--Littlewood 极大函数. 这里, 我们考察 Hardy--Littlewood 极大函数的性质, 并作为应用证明 Lebesgue 微分定理.

\section{Hardy--Littlewood 极大函数}\label{sec:3.1}
\index{W@5r 覆盖引理}

\index{H@Hardy--Littlewood 极大算子}

\index{H@Hardy--Littlewood 弱 (1,1) 极大不等式}








用积分定义的算子称为积分算子, 作为研究这些算子的基本工具, 有 Hardy--Littlewood 极大算子. 这里给出其定义, 并考察简单的例子. 特别是, 遗憾的是积分算子的像并不是可积函数, 因此考察取而代之能说些什么. 这时要使用称为覆盖引理的几何工具. 一般地, 经验上已知覆盖引理对研究函数的可微性是有用的, 而这一侧面到后面会变得清晰. 所谓立方体, 是指边与坐标轴平行的立方体.

\begin{definition}[Hardy--Littlewood 极大算子 I]\label{def:3.1}
对可测函数 $f: \mathbb{R}^{n} \to \mathbb{C}$, 函数 $Mf$ 由
\[
Mf(x) \equiv \sup\left\{\frac{1}{|Q|}\int_{Q}|f(y)|\dd y : Q \text{ 为包含 } x \text{ 的立方体}\right\}
\]
定义. 映射 $M: f \mapsto Mf$ 称为 Hardy--Littlewood 极大算子.
\end{definition}

通过将立方体稍微膨胀, 定义也可以改写为下面的形式.

\begin{definition}[Hardy--Littlewood 极大算子 II]\label{def:3.2}
对可测函数 $f: \mathbb{R}^{n} \to \mathbb{C}$, 函数 $Mf$ 由
\begin{equation}\tag{3.1}
Mf(x) \equiv \sup_{Q}\frac{\chi_{Q}(x)}{|Q|}\int_{Q}|f(y)|\dd y
\label{eq:3.1}
\end{equation}
定义. $Q$ 遍历其内部包含 $x$ 的全体立方体.
\end{definition}

\begin{example}\label{exa:3.1}
设 $n = 1$, $f = \chi_{[-1,1]}$, 则 $Mf(x) = \min\left(1, \frac{2}{1+|x|}\right)$. 实际上, 因为
\[
Mf(x) = \sup_{I \text{ 为区间}}\chi_{I}(x)\frac{|I \cap [-1,1]|}{|I|}.
\]
\end{example}

如这个例子所示(也参照练习题~\ref{exer:3.1}), 即使 $f \in L^{1}(\mathbb{R}^{n})$, $Mf \in L^{1}(\mathbb{R}^{n})$ 也未必成立. 另外, 对任意可测函数 $f: \mathbb{R}^{n} \to \mathbb{C}$, 有 $Mf(x) \leq \|f\|_{L^{\infty}}, x \in \mathbb{R}^{n}$.

那么, 当 $f \in L^{1}(\mathbb{R}^{n})$ 时, $Mf$ 满足什么呢? 回答这一问题的, 就是下面的覆盖引理.

\begin{theorem}[5r-覆盖引理]\label{thm:3.1}
在距离空间 $(X,d)$ 中给定由 $J$ 个球组成的族 $\{B_{\lambda}\}_{\lambda \in \Lambda}$ 时, 将下标重新编号为 $1, 2, \ldots, J$, 对重新编号后的球族 $B_{1}, B_{2}, \ldots, B_{J}$, 可以做到存在某个 $1 \leq K \leq J$ 满足下面的条件.

(1) 最初的 $K$ 个互不相交. 即 $B_{1}, B_{2}, \ldots, B_{K}$ 彼此没有公共部分.

(2) 存在某个映射 $\iota: \{1, 2, \ldots, J\} \to \{1, 2, \ldots, K\}$, 使得 $B_{j} \subset 3B_{\iota(j)}$ 成立.
\end{theorem}

\begin{note}\label{note:3.1}
所谓 5r, 是半径 $r$ 的 5 倍 (five times radius) 之意. 这里因为是有限族, 3 倍就足够了, 但若是无限族, 则需要 $3 + \varepsilon$ 倍. 典型地取 $\varepsilon = 2$, 因此总称为 5r.
\end{note}

粗略地说, 证明是反复进行「取最大半径的.」「舍弃最大半径的球以及包含在最大半径的球的 3 倍中的那些.」「在被舍弃的当中取最大半径的.」「舍弃最大半径的球以及包含在最大半径的球的 3 倍中的那些.」「在被舍弃的当中取最大半径的.」「舍弃最大半径的球以及包含在最大半径的球的 3 倍中的那些.」来推进的.

\begin{proof}
用关于 $J$ 的归纳法证明. 当 $J = 1$ 时, 实际上只要取 $K = 1$ 即可, 因此定理~\ref{thm:3.1} 成立.

归纳假设如下.

设 $J \leq J_{0}$. 假设当给定由 $J$ 个球组成的有限族 $\{B_{\lambda}^{*}\}_{\lambda \in \Lambda}$ 时, 将下标重新编号为 $1, 2, \ldots, J$ 后, 对重新编号的球族 $B_{1}^{*}, B_{2}^{*}, \ldots, B_{J}^{*}$, 可以做到满足下面的条件.

(1) 最初的 $K$ 个互不相交. 即 $B_{1}^{*}, B_{2}^{*}, \ldots, B_{K}^{*}$ 彼此没有公共部分.

(2) 存在某个映射 $\iota: \{1, 2, \ldots, J\} \to \{1, 2, \ldots, K\}$, 使得 $B_{j}^{*} \subset 3B_{\iota(j)}^{*}$ 成立.

于是设给定由 $J_{0}+1$ 个球组成的族. 取 $B_{1}$ 为半径最大的一个, 适当编号为 $B_{1}, B_{2}, \ldots, B_{J_{0}+1}$. 此后, 根据需要, 变更 $B_{2}, B_{3}, \ldots, B_{J_{0}+1}$ 的编号. 如果给出半径最大值的球有多个, 则从中任选一个. 对与 $B_{1}$ 相交的球 $B_{j}$ ($j = 2, 3, \ldots, J_{0}+1$), 由于 $r(B_{j}) \leq r(B_{1})$, 所以 $B_{j} \subset 3B_{1}$ 成立.

因此, $\{B_{j}\}_{j=1}^{J_{0}+1}$ 中不满足 $B_{j} \subset 3B_{1}$ 的($B_{1}$ 自动被排除)至多有 $J_{0}$ 个, 这些球与 $B_{1}$ 不相交. 所以, 将 $\{B_{j}\}_{j=1}^{J_{0}+1}$ 中满足 $B_{j} \cap B_{1} = \emptyset$ 的全体重新编号为 $\{B_{j}\}_{j=2}^{J+1}$, 将满足 $B_{j} \cap B_{1} \neq \emptyset$ 的全体重新编号为 $\{B_{j}\}_{j=J+2}^{J_{0}+1}$. 此时, 由 $J$ 的取法, 当 $j > J+1$ 时, $B_{j} \subset 3B_{1}$ 成立.

对 $\{B_{j}\}_{j=2}^{J+1}$ 使用归纳假设, 通过交换编号, 存在互不相交的球族 $B_{2}, \ldots, B_{K}$, 使得对 $2 \leq j \leq J+1$, 存在 $J(j) = 2, 3, \ldots, K$ 满足 $B_{j} \subset 3B_{J(j)}$. 此时, 令
\begin{equation}\tag{3.2}
\iota(j) \equiv \begin{cases} 1 & (j = 1 \text{ 时}) \\ J(j) & (2 \leq j \leq J+1 \text{ 时}) \\ 1 & (j > J+1 \text{ 时}) \end{cases}
\label{eq:3.2}
\end{equation}
即可知 $B_{1}, B_{2}, \ldots, B_{K}$ 满足定理的条件.

由此, 归纳法的证明完成.
\end{proof}

\begin{theorem}[Hardy--Littlewood 弱 (1,1) 极大不等式]\label{thm:3.2}
对 $f \in L^{1}(\mathbb{R}^{n})$, $\lambda > 0$, 有
\begin{equation}\tag{3.3}
|\{Mf > \lambda\}| \leq \frac{3^{n}}{\lambda}\int_{\mathbb{R}^{n}}|f(x)|\dd x
\label{eq:3.3}
\end{equation}
成立.
\end{theorem}

一般地, $\mathrm{w}L^{1}(\mathbb{R}^{n})$ 是指满足
\[
\|f\|_{\mathrm{w}L^{1}} = \sup_{\lambda > 0}\lambda\|\chi_{\{|f| > \lambda\}}\|_{L^{1}} = \sup_{\lambda > 0}\lambda|\{|f| > \lambda\}| < \infty
\]
的可测函数 $f: \mathbb{R}^{n} \to \mathbb{C}$ 所组成的线性空间. 对可测函数 $f: \mathbb{R}^{n} \to \mathbb{C}$, 有
\[
\|f\|_{\mathrm{w}L^{1}} \leq \|f\|_{L^{1}}.
\]
这个空间与 $L^{1}(\mathbb{R}^{n})$ 之间有包含关系 $\mathrm{w}L^{1}(\mathbb{R}^{n}) \supset L^{1}(\mathbb{R}^{n})$, 因此可知 $L^{1}(\mathbb{R}^{n})$ 是 $\mathrm{w}L^{1}(\mathbb{R}^{n})$ 的子集. 进而, 由于 $|\cdot|^{-n} \in \mathrm{w}L^{1}(\mathbb{R}^{n})$, $L^{1}(\mathbb{R}^{n})$ 是 $\mathrm{w}L^{1}(\mathbb{R}^{n})$ 的真子集. 基于这样的缘由, $\mathrm{w}L^{1}(\mathbb{R}^{n})$ 被称为弱 $L^{1}$ 空间, 从 $L^{1}(\mathbb{R}^{n})$ 到 $\mathrm{w}L^{1}(\mathbb{R}^{n})$ 的连续映射被称为弱 (1,1) 连续映射, 或具有弱 $L^{1}$ 有界性, 在 $L^{1}(\mathbb{R}^{n})$ 中弱有界, 具有弱 (1,1) 有界性的映射, 等等.

\begin{proof}
由 Lebesgue 测度的内正则性, 只需对满足 $K \subset \{Mf > \lambda\}$ 的任意紧集 $K$ 证明
\begin{equation}\tag{3.4}
|K| \leq \frac{3^{n}}{\lambda}\int_{\mathbb{R}^{n}}|f(x)|\dd x
\label{eq:3.4}
\end{equation}
即可. 设 $x \in K$. 由 $M$ 的定义, 存在(在内部)包含 $x$ 的立方体 $Q_{x}$, 使得
\begin{equation}\tag{3.5}
\frac{1}{|Q_{x}|}\int_{Q_{x}}|f(y)|\dd y > \lambda
\label{eq:3.5}
\end{equation}
成立. 由 $K$ 的紧性, 存在有限个点 $x_{1}, x_{2}, \ldots, x_{J}$, 使得 $K \subset \bigcup_{j=1}^{J}Q_{x_{j}}$ 成立. 为简化记号, 令 $Q_{j} \equiv Q_{x_{j}}$. 由定理~\ref{thm:3.1}, 重新编号后, 可以做到存在某个 $L \leq J$ 满足下面的条件.

(1) $\{Q_{j}\}_{j=1}^{L}$ 互不相交. 即
\begin{equation}\tag{3.6}
\sum_{j=1}^{L}\chi_{Q_{j}} \leq 1
\label{eq:3.6}
\end{equation}
成立.

(2) 关于 $K$ 的如下包含关系成立.
\begin{equation}\tag{3.7}
K \subset \bigcup_{j=1}^{J}Q_{j} \subset \bigcup_{j=1}^{L}3Q_{j}.
\label{eq:3.7}
\end{equation}

因此,
\[
\sum_{j=1}^{L}|3Q_{j}| = 3^{n}\sum_{j=1}^{L}|Q_{j}|
\]
所以依次使用式~\eqref{eq:3.7}, 式~\eqref{eq:3.5}, 式~\eqref{eq:3.6},
\begin{align*}
|K| &\leq \sum_{j=1}^{L}|3Q_{j}| \leq \sum_{j=1}^{L}\frac{3^{n}}{\lambda}\int_{Q_{j}}|f(y)|\dd y \\
&\leq \frac{3^{n}}{\lambda}\int_{\mathbb{R}^{n}}|f(y)|\dd y,
\end{align*}
定理得证.
\end{proof}

下面用例子来考虑 Hardy--Littlewood 极大函数是如何被使用的.

\begin{lemma}\label{lem:3.1}
设 $\Theta: \mathbb{R} \to [0, \infty)$ 是在 $(0, \infty)$ 上单调递减的偶函数, 且 $C^{1}$ 级并可积, 则对非负值局部可积函数 $F: \mathbb{R} \to [0, \infty)$, 有
\[
\int_{\mathbb{R}}\Theta(y)F(x-y)\dd y \leq MF(x)\|\Theta\|_{L^{1}} \quad (x \in \mathbb{R})
\]
成立. 特别地, 记在 $(0, \infty)$ 上单调递减、$C^{1}$ 级且 $L^{1}$ 范数不超过 $1$ 的偶函数的全体所组成的集合为 $\mathcal{Z}$. 此时, 对任意非负值可积函数 $F$, 有
\[
\left|\left\{x \in \mathbb{R}: \sup_{\Theta \in \mathcal{Z}}\int_{\mathbb{R}}\Theta(y)F(x-y)\dd y > \lambda\right\}\right| \leq \frac{3}{\lambda}\|F\|_{L^{1}}
\]
成立.
\end{lemma}

\begin{proof}
实际上, 由于 $\lim_{y \to \infty}\Theta(y) = 0$, 由分部积分得
\begin{align*}
\int_{\mathbb{R}}\Theta(y)F(x-y)\dd y &= \int_{\mathbb{R}}\Theta(y)\left(\frac{\dd}{\dd y}\int_{0}^{y}F(x-z)\dd z\right)\dd y \\
&= -\int_{\mathbb{R}}\Theta'(y)\left(\int_{0}^{y}F(x-z)\dd z\right)\dd y \\
&\leq -\int_{\mathbb{R}}\Theta'(y)y\,MF(x)\dd y \\
&= \int_{\mathbb{R}}\Theta(y)MF(x)\dd y = MF(x)\|\Theta\|_{L^{1}}.
\end{align*}
\end{proof}

\begin{exercise}\label{exer:3.1}
设 $n \in \mathbb{N}$. 设 $Q(x, r)$ 表示 $\mathbb{R}^{n}$ 中以 $x$ 为中心、体积为 $(2r)^{n}$ 的立方体.

(1) 设 $x \in \mathbb{R}^{n}$. 通过计算 $[-1,1]^{n} \cap Q(x, |x|+2)$ 及其体积, 证明 $M\chi_{[-1,1]^{n}}(x) \geq \frac{2^{n}}{(|x|+2)^{n}}$. 【提示】取 $n = 2$, $x = (x_{1}, x_{2}) \in (0, \infty)^{2}$, 尝试在坐标平面上画出 $Q(x, |x|+2)$ 与 $[-1,1]^{n}$.

(2) 利用 (1) 与 $(|x|+2)^{n} \geq \prod_{j=1}^{n}(|x_{j}|+2)$, 证明 $M\chi_{[-1,1]^{n}} \notin L^{1}(\mathbb{R}^{n})$.
\end{exercise}

\begin{exercise}\label{exer:3.2}
设立方体 $Q$, 可测函数 $f$ 与 $\lambda > 0$ 满足
\[
\frac{1}{|Q|}\int_{Q}|f(x)|\dd x > \lambda.
\]
此时, 证明存在某个 $\varepsilon \in (0, \infty)$, 使得
\[
\frac{1}{|(1+\varepsilon)Q|}\int_{Q}|f(x)|\dd x > \lambda
\]
成立. 【提示】令
\[
A = \frac{1}{|Q|}\int_{Q}|f(x)|\dd x,
\]
则假设与应当判定的条件分别可读作
\[
A > \lambda, \quad \frac{A}{(1+\varepsilon)^{n}} > \lambda.
\]
尝试关于 $\varepsilon$ 解这个不等式, 会得到什么? 不过, 要注意 $A = \infty$ 的可能性.
\end{exercise}

\begin{exercise}\label{exer:3.3}
对具有公共点的两个立方体 $Q$ 与 $Q'$, 证明 $Q' \subset 3Q$ 或 $Q \subset 3Q'$ 成立.
\end{exercise}

\begin{exercise}\label{exer:3.4}
给定 $\mathbb{R}$ 中由 6 个开区间 $A = (10, 14)$, $B = (10, 18)$, $C = (10, 16)$, $D = (3, 11)$, $E = (1, 9)$, $F = (9, 16)$ 组成的族, 选出其中一部分 $K_{1}, K_{2}, \ldots, K_{L}$, 使得下面的条件成立. 不过, $L$ 的值可以自由选取以适合条件. 另外, $3K_{l}$ 是保持 $K_{l}$ 的中心的 3 倍放大.

$\bullet$ $K_{1}, K_{2}, \ldots, K_{L}$ 互不相交.

$\bullet$ 存在某个映射 $\iota: \{A, B, C, D, E, F\} \to \{1, 2, \ldots, L\}$, 使得 $A \subset 3K_{\iota(A)}$, $B \subset 3K_{\iota(B)}$, $C \subset 3K_{\iota(C)}$, $D \subset 3K_{\iota(D)}$, $E \subset 3K_{\iota(E)}$, $F \subset 3K_{\iota(F)}$ 成立.

【注意】求出一组 $L, K_{1}, K_{2}, \ldots, K_{L}$ 以及 $\iota$.
\end{exercise}

\begin{exercise}\label{exer:3.5}
对 $n = 1, 2, \ldots$, $k = 0, 1, \ldots, n$, 定义 $b_{n,k}(x) = {}_{n}C_{k}x^{k}(1-x)^{n-k}$. 对可积函数 $f: [0,1] \to \mathbb{R}$, 将 $n$ 次 Kantorovich 多项式 $K_{n}(f, \cdot)$ 定义为
\[
K_{n}(f, x) = (n+1)\sum_{k=0}^{n}b_{n,k}(x)\int_{k/(n+1)}^{(k+1)/(n+1)}f(t)\dd t.
\]
回答下面的问题.

(1) 对 $x \in (0,1)$, $n \in \mathbb{N}$, 将 $k_{n,x} = 0, 1, \ldots, n$ 定义为满足 $k_{n,x} \leq (n+1)x < (k_{n,x}+1)$ 的唯一整数. 此时, 证明当 $k \leq k_{n,x}$ 时 $b_{n,k}(x) \leq b_{n,k+1}(x)$ 成立, 当 $k > k_{n,x}$ 时 $b_{n,k}(x) \geq b_{n,k+1}(x)$ 成立.

(2) 将函数 $K_{n}(x, t): [0,1] \times [0,1] \to [0, \infty)$ 定义为
\[
K_{n}(x, t) = (n+1)\sum_{k=0}^{n}b_{n,k}(x)\chi_{(k/(n+1), (k+1)/(n+1)]}(t).
\]
此时, 证明 $K_{n}(x, \cdot)$ 在 $t \leq x$ 上单调递增, 在 $t \geq x$ 上单调递减.

(3) 将符号测度 $\mu$ 定义为
\[
\mu_{x} = (n+1)\sum_{k=0}^{n}b_{n,k}(x)(\delta_{(k+1)/(n+1)} - \delta_{k/(n+1)}).
\]
这里 $\delta_{a}$ 是 $a$ 处的 Dirac delta. 此时, 对 $x \in [0,1]$ 与可积函数 $f: [0,1] \to \mathbb{R}$, 令 $F(t) = \int_{x}^{t}f(s)\dd s$, 证明
\[
K_{n}f(x, t) = \int_{[0,1]}F(t)\dd\mu_{x}(t)
\]
成立.

(4) 将作用于 $f \in L^{1}([0,1])$ 的极大算子 $M$ 定义为
\[
Mf(x) = \sup_{0 \leq a \leq x \leq b \leq 1, a < b}\frac{1}{b-a}\int_{a}^{b}f(s)\dd s,
\]
证明 $|K_{n}(f, t)| \leq Mf(t)$ 成立.
\end{exercise}

% ============ 第3章 §3.2 Hardy--Littlewood 极大函数的 $L^{p}$ 有界性 (原书 pp.67-70) ============

\section{\texorpdfstring{Hardy--Littlewood 极大函数的 $L^{p}$ 有界性}{Hardy-Littlewood 极大函数的 Lp 有界性}}\label{sec:3.2}
\index{H@Hardy--Littlewood 的 $L^{p}$ 不等式}


对 $1 \leq p < \infty$, (可测)函数 $f$ 的 $L^{p}$ 范数是由
\[
\|f\|_{L^{p}} = \left(\int_{\mathbb{R}^{n}}|f(x)|^{p}\dd x\right)^{1/p}
\]
给出的. 一般地, 对算子 $S$, 说 $S$ 是 $L^{p}$ 有界的, 是指存在常数 $C > 0$, 使得对一切 $f \in L^{p}(\mathbb{R}^{n})$, 不等式 $\|Sf\|_{L^{p}} \leq C\|f\|_{L^{p}}$ 成立. 我们已经知道 Hardy--Littlewood 极大算子不是 $L^{1}$ 有界的, 且是 $L^{\infty}$ 有界的. Hardy--Littlewood 极大算子不是 $L^{1}$ 有界的, 但取而代之, 有称为弱 (1,1) 不等式的不等式成立. 这里, 证明当 $1 < p < \infty$ 时, Hardy--Littlewood 极大算子是 $L^{p}$ 有界的.

从上面的定义式可以看出, $L^{p}$ 范数是难处理的. 这是因为会出现 $p$ 分之一次方与 $p$ 次方. 在 $p = 2$ 的情形, 由于 $|f(x)|^{2} = f(x)\overline{f(x)}$, 或许容易处理, 但除去这种特殊情形, 一般地要对这个量作代数的解释是困难的. 因此, 首先说明不借助函数的 $p$ 次方来表出 $L^{p}$ 范数的方法, 并以此为基础证明 Hardy--Littlewood 极大算子是 $L^{p}$ 有界的.

将定理~\ref{thm:1.3} 改写到函数的情形, 就得到下面的定理.

\begin{theorem}\label{thm:3.3}
设 $1 \leq p < \infty$. 此时, 对可测函数 $f: \mathbb{R}^{n} \to \mathbb{C}$, 有
\begin{equation}\tag{3.8}
\|f\|_{L^{p}} = \left(\int_{0}^{\infty}p\lambda^{p-1}|\{x \in \mathbb{R}^{n}: |f(x)| > \lambda\}|\dd\lambda\right)^{1/p}
\label{eq:3.8}
\end{equation}
成立.
\end{theorem}

\begin{proof}
对非负值函数使用 Fubini 定理, 直接计算.
\begin{align*}
\int_{\mathbb{R}^{n}}|f(x)|^{p}\dd x &= \int_{\mathbb{R}^{n}}\left(\int_{0}^{|f(x)|}p\lambda^{p-1}\dd\lambda\right)\dd x \\
&= \int_{\mathbb{R}^{n}}\left(\int_{0}^{\infty}p\lambda^{p-1}\chi_{\{(x,\lambda): |f(x)| > \lambda\}}(x,\lambda)\dd\lambda\right)\dd x \\
&= \int_{0}^{\infty}\left(\int_{\mathbb{R}^{n}}p\lambda^{p-1}\chi_{\{(x,\lambda): |f(x)| > \lambda\}}(x,\lambda)\dd x\right)\dd\lambda \\
&= \int_{0}^{\infty}p\lambda^{p-1}\left(\int_{\mathbb{R}^{n}}\chi_{\{(x,\lambda): |f(x)| > \lambda\}}(x,\lambda)\dd x\right)\dd\lambda \\
&= \int_{0}^{\infty}p\lambda^{p-1}|\{x \in \mathbb{R}^{n}: |f(x)| > \lambda\}|\dd\lambda.
\end{align*}
对此等式两边取 $p$ 次方根, 即得式~\eqref{eq:3.8}.
\end{proof}

以定理~\ref{thm:3.3} 为基础考察 $M$ 的 $L^{p}$ 有界性. 使用前一节所证明的
\[
|\{Mf > \lambda\}| \leq \frac{3^{n}}{\lambda}\int_{\mathbb{R}^{n}}|f(x)|\dd x, \quad Mf(x) \leq \|f\|_{L^{\infty}}.
\]

\begin{theorem}[Hardy--Littlewood $L^{p}$ 不等式]\label{thm:3.4}
设 $1 < p < \infty$. 此时,
\begin{equation}\tag{3.9}
\|Mf\|_{L^{p}} \leq \left(\frac{p2^{p}\cdot 3^{n}}{p-1}\right)^{1/p}\|f\|_{L^{p}}, \quad f \in L^{p}(\mathbb{R}^{n})
\label{eq:3.9}
\end{equation}
成立.
\end{theorem}

这里出现了
\[
\left(\frac{p2^{p}\cdot 3^{n}}{p-1}\right)^{1/p}
\]
这样的数, 但它无需记忆. 确实存在需要这种常数信息的局面, 但在本书以及相当一般的调和分析的讨论中, 很多情况下没有必要追踪常数. 为什么不必记忆常数呢, 是因为存在定理~\ref{thm:2.22}.

\begin{proof}
首先, 如叙述这个定理之前所述, 注意对一般的可测函数 $h$ 与 $\lambda > 0$, 当 $\|h\|_{\infty} \leq \lambda$ 时, 函数不等式 $Mh \leq \lambda$ 成立. 由这一点与次可加性 $M[g+h] \leq Mg + Mh$, 有 $\{Mf > 2\lambda\} \subset \{M[\chi_{\{|f| > \lambda\}}f] > \lambda\}$. 实际上, 令 $g = \chi_{\{|f| > \lambda\}}f$, $h = f - g = \chi_{\{|f| \leq \lambda\}}f$, 则 $Mh \leq \lambda$. 因此, 若 $x \in \mathbb{R}^{n}$ 满足 $Mf(x) > 2\lambda$, 则
\[
2\lambda < Mf(x) \leq M[\chi_{\{|f| > \lambda\}}f](x) + Mh(x) \leq M[\chi_{\{|f| > \lambda\}}f](x) + \lambda
\]
成立. 所以得到 $\lambda < M[\chi_{\{|f| > \lambda\}}f](x)$. 以上确证了包含关系
\[
\{Mf > 2\lambda\} \subset \{M[\chi_{\{|f| > \lambda\}}f] > \lambda\}.
\]

由这个包含关系与弱 (1,1) 有界性,
\begin{equation}\tag{3.10}
|\{Mf > 2\lambda\}| \leq |\{M[\chi_{\{|f| > \lambda\}}f] > \lambda\}| \leq \frac{3^{n}}{\lambda}\int_{\{|f| > \lambda\}}|f(x)|\dd x
\label{eq:3.10}
\end{equation}
成立. 由式~\eqref{eq:3.8} 作变量替换 $\sigma = 2\lambda$, 得
\begin{equation}\tag{3.11}
\|Mf\|_{L^{p}} = \left(p2^{p}\int_{0}^{\infty}\lambda^{p-1}|\{x \in \mathbb{R}^{n}: Mf(x) > 2\lambda\}|\dd\lambda\right)^{1/p}.
\label{eq:3.11}
\end{equation}
将式~\eqref{eq:3.10} 与式~\eqref{eq:3.11} 结合,
\begin{align*}
\int_{\mathbb{R}^{n}}Mf(x)^{p}\dd x &= p2^{p}\int_{0}^{\infty}\left(\int_{\mathbb{R}^{n}}\chi_{\{Mf(x) > 2\lambda\}}(x,\lambda)\dd x\right)\lambda^{p-1}\dd\lambda \\
&\leq p2^{p}\cdot 3^{n}\int_{0}^{\infty}\left(\int_{\{|f| > \lambda\}}|f(x)|\dd x\right)\lambda^{p-2}\dd\lambda.
\end{align*}
使用 Fubini 定理, 仔细计算一维定积分, 得
\begin{align*}
\int_{\mathbb{R}^{n}}Mf(x)^{p}\dd x
&\leq p2^{p}\cdot 3^{n}\int_{0}^{\infty}\left(\int_{\mathbb{R}^{n}}\chi_{\{(x,\lambda): |f(x)| > \lambda\}}(x,\lambda)|f(x)|\dd x\right)\lambda^{p-2}\dd\lambda \\
&= p2^{p}\cdot 3^{n}\int_{\mathbb{R}^{n}}\left(\int_{0}^{\infty}\chi_{\{(x,\lambda): |f(x)| > \lambda\}}(x,\lambda)\lambda^{p-2}\dd\lambda\right)|f(x)|\dd x \\
&= p2^{p}\cdot 3^{n}\int_{\mathbb{R}^{n}}\left(\int_{0}^{|f(x)|}\lambda^{p-2}\dd\lambda\right)|f(x)|\dd x \\
&= \frac{p2^{p}\cdot 3^{n}}{p-1}\int_{\mathbb{R}^{n}}|f(x)|^{p}\dd x.
\end{align*}
\end{proof}

前面说过, Hardy--Littlewood 极大算子对研究积分算子是有效的, 作为其一例, 来考虑分数次积分算子. 设 $0 < \alpha < n$. 对适当的可测函数 $f$, 定义
\begin{equation}\tag{3.12}
I_{\alpha}f(x) \equiv \int_{\mathbb{R}^{n}}\frac{f(y)}{|x-y|^{n-\alpha}}\dd y \quad (x \in \mathbb{R}^{n}).
\label{eq:3.12}
\end{equation}
由于积分必须被定义, 将 $f$ 取为非负值函数是一种想法, 但可知对 $L^{p}$ 函数这个积分算子也能定义. 下面证明下面的定理.

\begin{theorem}[Hardy--Littlewood--Sobolev]\label{thm:3.5}
当 $1 < p < q < \infty$, $0 < \alpha < n$ 满足 $\frac{1}{q} = \frac{1}{p} - \frac{\alpha}{n}$ 时, 对非负值 Lebesgue 可测函数 $f$, 有 $\|I_{\alpha}f\|_{L^{q}} \lesssim \|f\|_{L^{p}}$ 成立.
\end{theorem}

\begin{proof}
将积分核改写为
\[
\frac{1}{|x-y|^{n-\alpha}} = (n-\alpha)\int_{0}^{\infty}\chi_{B(x,r)}(y)\frac{\dd r}{r^{n-\alpha+1}}.
\]
于是, 由 Fubini 定理, 可以表为
\[
I_{\alpha}f(x) = (n-\alpha)\int_{0}^{\infty}\left(\frac{1}{r^{n-\alpha+1}}\int_{B(x,r)}f(y)\dd y\right)\dd r.
\]
这里, 由 Hardy--Littlewood 极大算子的定义与 Hölder 不等式, 得
\[
\int_{B(x,r)}f(y)\dd y \leq |B(x,r)|Mf(x), \quad \int_{B(x,r)}f(y)\dd y \leq |B(x,r)|^{1/p'}\|f\|_{L^{p}}.
\]
即
\[
\int_{B(x,r)}f(y)\dd y \leq \min(|B(x,r)|Mf(x), |B(x,r)|^{1/p'}\|f\|_{L^{p}}).
\]
这里, 关于正的未知数 $r$ 的方程 $r^{n}A = r^{n/p'}B$ 的解是 $r = \left(\frac{B}{A}\right)^{p/n}$. 所以,
\begin{align*}
\int_{0}^{\infty}\min(r^{n}A, r^{n/p'}B)r^{-n+\alpha-1}\dd r
&= A\int_{0}^{(B/A)^{p/n}}r^{\alpha-1}\dd r + B\int_{(B/A)^{p/n}}^{\infty}r^{-n/p+\alpha-1}\dd r \\
&= A\int_{0}^{(B/A)^{p/n}}r^{\alpha-1}\dd r + B\int_{(B/A)^{p/n}}^{\infty}r^{-n/q-1}\dd r \\
&= \frac{A}{\alpha}\left(\frac{B}{A}\right)^{\frac{p}{n}\alpha} + \frac{nB}{q}\left(\frac{A}{B}\right)^{p/q} \\
&= \left(\frac{1}{\alpha} + \frac{n}{q}\right)B\left(\frac{A}{B}\right)^{p/q} = \left(\frac{1}{\alpha} + \frac{n}{q}\right)A^{p/q}B^{1-p/q}.
\end{align*}
所以
\[
I_{\alpha}f(x) \lesssim Mf(x)^{p/q}\|f\|_{L^{p}}^{1-p/q}
\]
成立, 对此式两边取 $L^{q}$ 范数,
\begin{equation}
\begin{aligned}
\|I_{\alpha}f\|_{L^{q}} &\lesssim \|(Mf)^{p/q}\|_{L^{q}}\|f\|_{L^{p}}^{1-p/q} \\
&= (\|Mf\|_{L^{p}})^{p/q}\|f\|_{L^{p}}^{1-p/q}
\end{aligned}\tag{3.13}
\label{eq:3.13}
\end{equation}
成立. 由于 $1 < p < \infty$, $M$ 是 $L^{p}$ 有界的, $\|Mf\|_{L^{p}} \lesssim \|f\|_{L^{p}}$ 成立. 所以, 将此不等式代入式~\eqref{eq:3.13},
\[
\|I_{\alpha}f\|_{L^{q}} \lesssim (\|f\|_{L^{p}})^{p/q}\|f\|_{L^{p}}^{1-p/q} = \|f\|_{L^{p}}.
\]
通过分解 $f = \max(f, 0) - \max(-f, 0)$, 可得 $I_{\alpha}$ 的 $L^{p}$ 有界性.
\end{proof}

\begin{corollary}[Hardy--Littlewood--Sobolev]\label{cor:3.1}
当 $1 < p < q < \infty$, $0 < \alpha < n$ 满足 $\frac{1}{q} = \frac{1}{p} - \frac{\alpha}{n}$ 时, 对 $f \in L^{p}(\mathbb{R}^{n})$, $I_{\alpha}[|f|](x) < \infty$ 几乎处处成立, 且 $\|I_{\alpha}f\|_{L^{q}} \lesssim \|f\|_{L^{p}}$ 成立.
\end{corollary}

\begin{proof}
$I_{\alpha}[|f|](x) < \infty$ 几乎处处成立, 是由 $I_{\alpha}[|f|] \in L^{q}(\mathbb{R}^{n})$ 得到的. 不等式如上面所述, 通过分解 $f = \max(f, 0) - \max(-f, 0)$ 得到.
\end{proof}

\begin{exercise}\label{exer:3.6}
对 Hardy--Littlewood 极大算子 $M$ 与可测函数 $f$, $g$ 及 $a \in \mathbb{C}$, 证明下面的函数不等式或等式.

(1) $M[f+g] \leq Mf + Mg$. 不过, 可以承认积分的三角不等式.

(2) $M[af] = |a|Mf$.

【注意】应当证明的是, 对任意 $x \in \mathbb{R}^{n}$, (1) $M[f+g](x) \leq Mf(x) + Mg(x)$, (2) $M[af](x) = |a|Mf(x)$ 成立.
\end{exercise}

% ============ 第3章 §3.3 Lebesgue 微分定理 (原书 pp.71-74) ============

\section{Lebesgue 微分定理}\label{sec:3.3}
\index{0symbolBr@$B(r)$}

\index{0symbolLloc@$L_{\mathrm{loc}}^{1}$}

\index{L@Lebesgue 点}

\index{L@Lebesgue 微分定理}










前面说过覆盖引理对函数可微性的分析是有用的, 这里说明这一侧面. 这里考察当 $r \downarrow 0$ 时,
\[
\frac{1}{|Q(x,r)|}\int_{Q(x,r)}f(y)\dd y
\]
会呈现怎样的行为. 这里, $Q(x,r)$ 是以 $x$ 为中心、体积为 $(2r)^{n}$ 的立方体.

实际上, 因为想讨论的是与各点收敛相近的概念——几乎处处收敛, 虽然有些离题, 还是重新收录 Lebesgue 收敛定理.

\begin{theorem}[Lebesgue 收敛定理 I]\label{thm:3.6}
设 $\mathbb{R}^{n}$ 中的函数列 $\{f_{j}\}_{j=1}^{\infty}$ 几乎处处收敛于 $f$. 进而, 假设存在 $g \in L^{1}(\mathbb{R}^{n})$, 使得对一切自然数 $j$, 几乎处处有 $|f_{j}| \leq g$. 此时,
\[
\lim_{j \to \infty}\int_{\mathbb{R}^{n}}f_{j}(x)\dd x = \int_{\mathbb{R}^{n}}f(x)\dd x
\]
成立.
\end{theorem}

这一结论用 $L^{1}$ 范数也可以改写如下. 实际上, 这是因为下面的定理与最初的定理等价.

\begin{theorem}[Lebesgue 收敛定理 II]\label{thm:3.7}
设 $\mathbb{R}^{n}$ 中的函数列 $\{f_{j}\}_{j=1}^{\infty}$ 几乎处处收敛于 $f$. 进而, 假设存在 $g \in L^{1}(\mathbb{R}^{n})$, 使得对一切自然数 $j$, 几乎处处有 $|f_{j}| \leq g$. 此时, $f_{j}$ 在 $L^{1}$ 的拓扑下收敛于 $f$. 即
\[
\lim_{j \to \infty}\|f_{j} - f\|_{L^{1}} = 0
\]
成立.
\end{theorem}

进而, 这个定理的假设可以改写如下.

\begin{theorem}[Lebesgue 收敛定理 III]\label{thm:3.8}
设 $\mathbb{R}^{n}$ 中的函数列 $\{f_{j}\}_{j=1}^{\infty}$ 几乎处处收敛于 $f$. 此时, 若 $\sup_{j \in \mathbb{N}}|f_{j}| \in L^{1}(\mathbb{R}^{n})$, 则 $f_{j}$ 在 $L^{1}$ 的拓扑下收敛于 $f$. 即
\[
\lim_{j \to \infty}\|f_{j} - f\|_{L^{1}} = 0
\]
成立.
\end{theorem}

这样改写后, 就能理解极大算子承担着自然的角色. 也就是说, 作为 Lebesgue 收敛定理 I, II 中的函数 $g$, 它取的是自然的(极小的)对象.

$L_{\mathrm{loc}}^{1}(\mathbb{R}^{n})$ 是指乘以 $\chi_{K}$ 后成为可积函数的函数全体所组成的集合, 其中 $K$ 为任意紧集. 这样的函数称为局部可积函数.

\begin{theorem}[Lebesgue 微分定理]\label{thm:3.9}
设 $f \in L_{\mathrm{loc}}^{1}(\mathbb{R}^{n})$. 此时, 对几乎所有 $x \in \mathbb{R}^{n}$,
\begin{equation}\tag{3.14}
\lim_{r \downarrow 0}\frac{1}{|Q(x,r)|}\int_{Q(x,r)}|f(y) - f(x)|\dd y = 0
\label{eq:3.14}
\end{equation}
成立. 特别地, 对这样的点, $|f(x)| \leq Mf(x)$ 成立.
\end{theorem}

称使式~\eqref{eq:3.14} 成立的点为 Lebesgue 点.

\begin{proof}
后半部分的主张由极大算子的定义是显然的, 因此证明前半部分. 为此, 只需证明使式~\eqref{eq:3.14} 不成立的点的全体所组成的集合的测度为 $0$. 将使式~\eqref{eq:3.14} 不成立的点的全体按水平来考虑. 即, 对一切自然数 $k \in \mathbb{N}$, 证明
\[
E_{k} \equiv \left\{x \in \mathbb{R}^{n}: \limsup_{r \downarrow 0}\frac{1}{|Q(x,r)|}\int_{Q(x,r)}|f(y) - f(x)|\dd y > \frac{1}{k}\right\}
\]
的测度为 $0$ 就足够了. 取 $r > 0$, 令 $B(r) = \{|x| < r\}$. 由于式~\eqref{eq:3.14} 是否成立是局部的问题, 用 $f \cdot \chi_{B(r)}$ 代替 $f$ 来考虑, 可以假设 $f \in L^{1}(\mathbb{R}^{n})$. 进而, 注意对 $g \in L^{1}(\mathbb{R}^{n}) \cap \mathrm{BC}(\mathbb{R}^{n})$, 式~\eqref{eq:3.14} 是显然的\footnote{BC 是有界连续 (bounded continuous) 之意.}. 任意取 $\varepsilon > 0$. 通过证明 $|E_{k}| \leq 2(3^{n}+1)k\varepsilon$, 完成定理的证明.

取 $g \in L^{1}(\mathbb{R}^{n}) \cap \mathrm{BC}(\mathbb{R}^{n})$, 使得 $\|f - g\|_{L^{1}} \leq \varepsilon$. 这由稠密性是可能的. 略记 $H(x, y) = |f(y) - g(y) - f(x) + g(x)|$. 作为集合的包含关系
\begin{equation}
\begin{aligned}
E_{k} &= \left\{x \in \mathbb{R}^{n}: \limsup_{r \downarrow 0}\frac{1}{|Q(x,r)|}\int_{Q(x,r)}|H(x,y)|\dd y > \frac{1}{k}\right\} \\
&\subset \left\{x \in \mathbb{R}^{n}: M[f-g](x) + |f(x) - g(x)| > \frac{1}{k}\right\} \\
&\subset \left\{x \in \mathbb{R}^{n}: M[f-g](x) > \frac{1}{2k}\right\} \cup \left\{x \in \mathbb{R}^{n}: |f(x) - g(x)| > \frac{1}{2k}\right\}
\end{aligned}\tag{3.15}
\label{eq:3.15}
\end{equation}
成立, 因此由测度的次可加性、Chebyshev 不等式与弱 (1,1) 极大不等式,
\[
|E_{k}| \leq 2(3^{n}+1)k\|f - g\|_{L^{1}} \leq 2(3^{n}+1)k\varepsilon
\]
成立. 由于 $\varepsilon > 0$ 是任意的, 式~\eqref{eq:3.14} 随之成立.
\end{proof}

$C_{\mathrm{c}}^{\infty}(\mathbb{R}^{n})$ 是指 $C^{\infty}(\mathbb{R}^{n})$ 级函数中, 在有界集合之外取 $0$ 的函数全体所组成的线性空间.

\begin{note}\label{note:3.2}
与这里处理的 Lebesgue 微分定理相同, 对 $\varphi \in C_{\mathrm{c}}^{\infty}(\mathbb{R}^{n})$ 与 $f \in L_{\mathrm{loc}}^{1}(\mathbb{R}^{n})$, 将关于 $x \in \mathbb{R}^{n}$ 的函数定义为 $\varphi_{t}(x) = t^{-n}\varphi(t^{-1}x)$ 时, a.e. 有
\[
\varphi_{t} * f(x) \to \int_{\mathbb{R}^{n}}\varphi(y)\dd y \times f(x) \quad (t \downarrow 0)
\]
成立. 这也常常被称为 Lebesgue 微分定理. 为证明它, 注意下面两点. 一点是当 $f \in C_{\mathrm{c}}^{\infty}(\mathbb{R}^{n})$ 时结论对一切 $x \in \mathbb{R}^{n}$ 正确, 另一点是由于 $\varphi \in C_{\mathrm{c}}^{\infty}(\mathbb{R}^{n})$ 是具有紧支集的有界函数,
\[
\sup_{t > 0}|\varphi_{t} * f(x)| \lesssim Mf(x) \quad (x \in \mathbb{R}^{n})
\]
成立. 注意这两点, 就能用与定理~\ref{thm:3.9} 相同的方法证明.
\end{note}

\begin{exercise}\label{exer:3.7}
对下面的开区间 $I$ 与函数 $\varphi$, 令 $f(t) = \varphi(t)\chi_{I}(t)$. 另外, 取 $a \in \mathbb{R}$, 令 $F_{a}(t) = \int_{a}^{t}f(s)\dd s$. 在下面的 5 种情形下, 求出 $F_{a}$, 并判定 $F_{a}$ 在 $t = a$ 处是否可微.
\[
\begin{array}{l}
\bullet\ a < \inf I \\
\bullet\ a = \inf I \\
\bullet\ a \in I \\
\bullet\ a = \sup I \\
\bullet\ a > \sup I
\end{array}
\]
(1) $I = (5, 8)$, $\varphi(t) = t - 5$.

(2) $I = (6, 9)$, $\varphi(t) = t - 9$.
\end{exercise}

\begin{exercise}\label{exer:3.8}
回答下面的问题.

(1) 设 $f \in L_{\mathrm{loc}}^{1}(\mathbb{R})$. 此时, 证明对几乎所有 $x \in \mathbb{R}$,
\begin{equation}\tag{3.16}
\lim_{r \downarrow 0}\frac{1}{r}\int_{x}^{x+r}f(y)\dd y = f(x)
\label{eq:3.16}
\end{equation}
成立.

(2) 设 $f \in L_{\mathrm{loc}}^{1}(\mathbb{R}^{n})$. 此时, 证明对几乎所有 $x \in \mathbb{R}^{n}$,
\begin{equation}\tag{3.17}
\lim_{r \downarrow 0}\frac{1}{|B(x,r)|}\int_{B(x,r)}f(y)\dd y = f(x)
\label{eq:3.17}
\end{equation}
成立.
\end{exercise}

\begin{exercise}\label{exer:3.9}
设 $f \in L_{\mathrm{loc}}^{1}(\mathbb{R})$, 证明以下命题.

(1) 当 $x$ 是 $f$ 的 Lebesgue 点时,
\[
\lim_{r \downarrow 0}\frac{1}{|B(x,r)|}\int_{B(x,r)}|f(y) - f(x)|\dd y = 0
\]
成立.

(2) 反之, 使
\[
\lim_{r \downarrow 0}\frac{1}{|B(x,r)|}\int_{B(x,r)}|f(y) - f(x)|\dd y = 0
\]
成立的 $x \in \mathbb{R}$ 是 $f$ 的 Lebesgue 点.
\end{exercise}

\begin{exercise}\label{exer:3.10}
证明不存在 Lebesgue 可测集 $E$, 使得对一切开区间 $I$ 有 $|E \cap I| = \frac{1}{2}|I|$.
\end{exercise}


% 第4章 奇异积分算子 (原书 pp.75-97, PDF pp.81-103)
\chapter{奇异积分算子}
% ============ 第4章 §4.1 Cotlar 引理 (原书 pp.75-76) ============

\section{Cotlar 引理}\label{sec:4.1}

这一节在抽象的框架下展开讨论. 主要目的是处理奇异积分算子, 特别是 Hilbert 变换, 但在 Hilbert 空间的框架下作抽象的讨论.

这里证明后面要使用的正交性引理. 对 Hilbert 空间 $H$, $P(\mathcal{H})$ 表示全体正交射影组成的集合. 即, 有界线性算子 $P: H \to H$ 属于 $P(\mathcal{H})$, 是指 $P = P^{*}$ 与 $P^{2} = P$ 成立. 这里, 对不同的自然数 $j, k$, 对满足正交条件 $P_{j}P_{k} = 0$ 的正交射影列 $\{P_{j}\}_{j=1}^{\infty} \subset P(\mathcal{H})$,
\begin{equation}\tag{4.1}
\|P_{1} + P_{2} + \cdots + P_{N}\|_{\mathcal{H} \to \mathcal{H}} \leq 1
\label{eq:4.1}
\end{equation}
成立. 由于极限在适当的意义下存在,
\[
\|P_{1} + P_{2} + \cdots\|_{\mathcal{H} \to \mathcal{H}} \leq 1
\]
也成立. 但是, 像 $P_{1} = P_{2} = \cdots = A$ 的情形那样, 对不是正交射影的一般线性算子, 式~\eqref{eq:4.1} 不成立.

从应用上看, 这个正交条件经常出现, 但这个关系式只要稍加摄动就不再成立. 这里, 寻找正交条件的替代.

\begin{theorem}[Cotlar 引理]\label{thm:4.1}
设 $J \subset \mathbb{Z}$ 为有限集合. 对有界线性算子的族 $\{T_{j}\}_{j \in J} \subset B(\mathcal{H})$, 存在适当的正数列 $\{\gamma_{j}\}_{j \in \mathbb{Z}}$, 使得
\begin{equation}\tag{4.2}
A \equiv \max\left(\sup_{j' \in J}\sum_{j \in J}\gamma_{j-j'}, \sup_{j' \in J}\sum_{j \in J}\gamma_{j'-j}\right) < \infty
\label{eq:4.2}
\end{equation}
以及
\begin{equation}\tag{4.3}
\|T_{j'}^{*}T_{j}\|_{\mathcal{H} \to \mathcal{H}} + \|T_{j'}T_{j}^{*}\|_{\mathcal{H} \to \mathcal{H}} \leq \gamma_{j'-j}^{2} \quad (j, j' \in J)
\label{eq:4.3}
\end{equation}
成立. 此时, 算子 $T_{J} \equiv \sum_{j \in J}T_{j}$ 满足 $\|T_{J}\|_{\mathcal{H} \to \mathcal{H}} \leq A$.
\end{theorem}

\begin{proof}
设自然数 $m$ 满足 $\log_{2}m \in \mathbb{N}$. 认真地写下来的话,
\begin{equation}\tag{4.4}
(T_{J}^{*}T_{J})^{m} = \sum_{j_{1}, j_{2}, \ldots, j_{2m} \in J}T_{j_{1}}^{*}T_{j_{2}}T_{j_{3}}^{*} \cdots T_{j_{2m}}
\label{eq:4.4}
\end{equation}
成立.

评价各项. 首先, 将括号如下配对.
\[
T_{j_{1}}^{*}T_{j_{2}}T_{j_{3}}^{*} \cdots T_{j_{2m}} = \prod_{l=1}^{m}\left(T_{j_{2l-1}}^{*}T_{j_{2l}}\right).
\]
由不等式~\eqref{eq:4.3},
\begin{equation}\tag{4.5}
\|T_{j_{1}}^{*}T_{j_{2}}T_{j_{3}}^{*} \cdots T_{j_{2m}}\|_{\mathcal{H} \to \mathcal{H}} \leq \prod_{l=1}^{m}\gamma_{j_{2l-1}-j_{2l}}^{2}
\label{eq:4.5}
\end{equation}
成立.

接下来, 将括号如下配对.
\[
T_{j_{1}}^{*}T_{j_{2}}T_{j_{3}}^{*} \cdots T_{j_{2m}} = T_{j_{1}}^{*}\left(\prod_{l=1}^{m-1}\left(T_{j_{2l}}T_{j_{2l+1}}^{*}\right)\right)T_{j_{2m}}.
\]
这里, 关于两端, 由式~\eqref{eq:4.3},
\[
\|T_{j_{1}}\|_{\mathcal{H} \to \mathcal{H}} = \sqrt{\|T_{j_{1}}^{*}T_{j_{1}}\|_{\mathcal{H} \to \mathcal{H}}} \leq \gamma_{0} \leq A
\]
成立. 同样地, 由 $\|T_{j_{2m}}\|_{\mathcal{H} \to \mathcal{H}} \leq A$,
\begin{equation}\tag{4.6}
\|T_{j_{1}}^{*}T_{j_{2}}T_{j_{3}}^{*} \cdots T_{j_{2m}}\|_{\mathcal{H} \to \mathcal{H}} \leq A^{2}\prod_{l=1}^{m-1}\gamma_{j_{2l}-j_{2l+1}}^{2}
\label{eq:4.6}
\end{equation}
成立. 考虑式~\eqref{eq:4.5} 与式~\eqref{eq:4.6} 的几何平均, 代入式~\eqref{eq:4.4},
\[
\|(T_{J}^{*}T_{J})^{m}\|_{\mathcal{H} \to \mathcal{H}} \leq A\sum_{j_{1}, j_{2}, \ldots, j_{2m} \in J}\prod_{l=1}^{2m-1}\gamma_{j_{l}-j_{l+1}}
\]
成立.

这里, 先对 $j_{1}$ 求和, 利用 $\sum_{j_{1} \in J}\gamma_{j_{1}-j_{2}} \leq A$ 来控制. 接下来对 $j_{2}$ 求和, 利用 $\sum_{j_{2} \in J}\gamma_{j_{2}-j_{3}} \leq A$ 来控制. 如此反复, 只特别对待 $j_{2m}$, 得 $\|T_{J}\|_{\mathcal{H} \to \mathcal{H}}^{2m} = \|(T_{J}^{*}T_{J})^{m}\|_{\mathcal{H} \to \mathcal{H}} \leq A^{2m}\sharp J$. 这里, 由于 $J$ 是有限的, 得 $\|T_{J}\|_{\mathcal{H} \to \mathcal{H}} \leq A(\sharp J)^{\frac{1}{2m}}$. 最后令 $m \to \infty$, 即得结论.
\end{proof}

至此 $J$ 一直取有限集合, 但对 $J = \mathbb{Z}$ 定理也能适用. 此时 $T_{J} = T_{\mathbb{Z}}$ 是无穷和, 但可以以对一切 $x \in H$,
\[
T_{\mathbb{Z}}x = \sum_{j \in \mathbb{Z}}T_{j}x
\]
收敛的意义来理解无穷和. 实际上, 若存在 $j$ 使得 $x \in \mathrm{Im}(T_{j}^{*})$, 则定义无穷和的极限的存在性由式~\eqref{eq:4.2} 与式~\eqref{eq:4.3} 是显然的. 由于在有限和的情形已证明算子范数具有不依赖于 $J$ 的估计, 移到线性闭包,
\[
x \in \overline{\operatorname{Span}}\left(\mathrm{Im}(T_{j}^{*}), j \in \mathbb{Z}\right)
\]
的情形, 极限也同样存在. 最后, 考虑
\[
x \in \bigcap_{j \in \mathbb{Z}}\ker(T_{j}) = \overline{\operatorname{Span}}\left(\mathrm{Im}(T_{j}^{*}), j \in \mathbb{Z}\right)^{\perp}
\]
的情形. 这里 Span 是在代数的意义下使用的表示线性空间的生成的记号. 即, 若 $T_{j}x = 0$ 对一切 $j \in \mathbb{Z}$ 成立, 则由对一切有限集合 $J \subset \mathbb{Z}$ 有 $T_{J}x = 0$, 极限是存在的.

\begin{exercise}\label{exer:4.1}
用 $\|\cdot\|$ 也表示复 Hilbert 空间 $H$ 中全体有界线性算子组成的 Banach 空间 $B(H)$ 的范数. 此时, 证明 $(A^{*}A)^{2^{m}}$ 是自伴的, 且 $\|(A^{*}A)^{2^{m}}\| = \|A\|^{2^{m+1}}$ 成立.
\end{exercise}

% ============ 第4章 §4.2 Hilbert 变换及其 L^2 有界性 (原书 pp.77-81) ============

\section{\texorpdfstring{Hilbert 变换及其 $L^{2}$ 有界性}{Hilbert 变换及其 L2 有界性}}\label{sec:4.2}

应用 Cotlar 引理来考察 Hilbert 变换. 这里虽不涉及, 但 Hilbert 变换是分析 Fourier 级数的重要工具. Hilbert 变换属于称为奇异积分算子的数学变换的范畴, 因此积分核的处理很困难. 这里, 考察对具有紧支集的无穷次可微函数有什么样的性质, 并以此为立足点考察对 $L^{2}$ 函数的输入. Hilbert 变换是由
\[
Hf(x) = \lim_{\varepsilon \downarrow 0}\int_{|y| > \varepsilon}\frac{f(x-y)}{y}\dd y
\]
给出的变换.

令
\[
H_{\varepsilon}f(x) = \int_{|y| > \varepsilon}\frac{f(x-y)}{y}\dd y.
\]
一般地, 函数 $f, g$ 的卷积由
\[
f * g(x) = \int_{\mathbb{R}}f(x-y)g(y)\dd y
\]
给出, 但现在的情形是 $g(y) = y^{-1}$, 所以分析困难. 例如, 若 $f \in L^{p}(\mathbb{R})$, $g \in L^{1}(\mathbb{R})$, $1 \leq p \leq \infty$, 则 $\|f * g\|_{L^{p}} \leq \|f\|_{L^{p}}\|g\|_{L^{1}}$, 因而好处理, 但 $g(y) = y^{-1}$ 不可积, 所以需要（非常糟糕的）技巧. 如此在 $y = 0$ 具有奇异性, 所以 Hilbert 变换是奇异积分算子.

极限是否存在令人担心, 但至少 $f \in C_{\mathrm{c}}^{\infty}(\mathbb{R})$ 时极限存在. 实际上, 设 $\operatorname{supp}(f) \subset [-R, R]$, 令 $K = [|x|+1]$, 则 $x - y \in [-R, R]$ 时 $y \in [x-R, x+R] \subset [-K-R, K+R]$, 所以
\[
Hf(x) = \lim_{\varepsilon \downarrow 0}\int_{|y| > \varepsilon}\frac{f(x-y)}{y}\dd y = \lim_{\varepsilon \downarrow 0}\int_{\varepsilon < |y| < K+R}\frac{f(x-y)}{y}\dd y.
\]
注意到 $y^{-1}$ 是奇函数, 利用
\[
\int_{\varepsilon < |y| < R+K}\frac{\dd y}{y} = 0,
\]
得
\[
Hf(x) = \lim_{\varepsilon \downarrow 0}\int_{\varepsilon < |y| < K+R}\frac{f(x-y)-f(x)}{y}\dd y.
\]
这里, 由
\[
\lim_{y \to 0}\frac{f(x-y)-f(x)}{y} = -f'(x)
\]
可知, 插入 $-f(x)$ 后, $\frac{f(x-y)}{y}$ 在 $y = 0$ 的奇异性消失了. 因此,
\begin{equation}\tag{4.7}
Hf(x) = \int_{|y| < K+R}\frac{f(x-y)-f(x)}{y}\dd y
\label{eq:4.7}
\end{equation}
成立.

如此, 在可微的情形可知极限的存在, 但不可微的情形能说什么呢. 这是当前的目标. 这里先作为开端, 讨论在 $L^{2}(\mathbb{R})$ 中的极限.

设 $C^{1}$ 级偶函数 $\rho: \mathbb{R} \to \mathbb{R}$ 满足
\[
\operatorname{supp}(\rho) \subset [1, 4] \cup [-4, -1], \quad \sum_{j=-\infty}^{\infty}\rho(2^{j}\cdot) = \chi_{\mathbb{R}\setminus\{0\}}.
\]
对 $j \in \mathbb{Z}$, $x \neq 0$, 令 $\rho_{j}(x) = \rho(2^{-j}x)$, $\tau_{j}(x) = x^{-1}\rho_{j}(x)$. 进而, 令 $\rho_{j}(0) = \tau_{j}(0) = 0$, 则 $\tau_{j}$ 是 $C^{1}$ 级函数.

若忽略极限记号 $\lim_{\varepsilon \downarrow 0}$ 的存在性, 则 Hilbert 变换是与 $\frac{1}{x}$ 的卷积, 将此 $\frac{1}{x}$ 分解为
\[
\frac{1}{x} = \sum_{j=-\infty}^{\infty}\tau_{j}(x) \quad (x \neq 0).
\]
与此对应, 设想
\begin{equation}\tag{4.8}
Hf(x) = \sum_{j=-\infty}^{\infty}\tau_{j} * f(x)
\label{eq:4.8}
\end{equation}
来考虑. 这个设想是否正当, 对属于 $C_{\mathrm{c}}^{\infty}(\mathbb{R})$ 的函数 $f$ 没有问题, 对其他函数以后考察. 令 $H_{j}f(x) = \tau_{j} * f(x)$. 证明 $\{H_{j}\}_{j=-\infty}^{\infty}$ 满足 Cotlar 引理的假设.

由
\[
\int_{\mathbb{R}}H_{j}f(x)\overline{g(x)}\dd x = -\int_{\mathbb{R}}f(x)\overline{H_{j}g(x)}\dd x
\]
可知 $H_{j}^{*} = -H_{j}$, 所以考察 $H_{k}H_{j}$ 的算子范数. 由 Young 不等式 $\|f * g\|_{L^{2}} \leq \|f\|_{L^{2}}\|g\|_{L^{1}}$ 得 $\|H_{k}H_{j}\|_{L^{2} \to L^{2}} \leq \|\tau_{j} * \tau_{k}\|_{L^{1}}$. 卷积积的函数的支集可以由 Minkowski 和计算\footnote{设 $A, B$ 是线性空间 $V$ 的子集. Minkowski 和 $A+B$ 由 $A+B = \{a+b \in V : a \in A, b \in B\}$ 给出. 另外, $f$ 是具有紧支集的连续函数时, 对连续函数 $g$, $\operatorname{supp}(f * g) \subset \operatorname{supp}(f) + \operatorname{supp}(g)$ 成立.}, $\operatorname{supp}(\tau_{j} * \tau_{k}) \subset [-2^{\max(j,k)+3}, 2^{\max(j,k)+3}]$. 另一方面,
\[
\Theta_{j}(x) = \int_{-\infty}^{x}\tau_{j}(y)\dd y
\]
由 $\int_{\mathbb{R}}\tau_{j}(x)\dd x = 0$ 可知满足
\[
\operatorname{supp}(\Theta_{j}) \subset [-2^{j+2}, 2^{j+2}].
\]
因此, 由卷积公式 $(f * g)' = f' * g = f * g'$（练习题~\ref{exer:4.5} 参照）,
\[
\tau_{j} * \tau_{k} = (\Theta_{j})' * \tau_{k} = \Theta_{j} * (\tau_{k})'.
\]
这里, 仔细考察 $\tau_{k}$ 与 $\Theta_{j}$. 由 $\tau_{k}(x) = \frac{\rho(2^{-k}x)}{x}$,
\[
\tau_{k}'(x) = \frac{2^{-k}\rho'(2^{-k}x)}{x} - \frac{\rho(2^{-k}x)}{x^{2}},
\]
其结果是
\[
|\tau_{k}'| \leq C4^{-k}\chi_{[-2^{k+2}, 2^{k+2}]}.
\]
另一方面, 由变量替换
\[
\Theta_{j}(x) = \int_{-\infty}^{2^{-j}x}\frac{\rho(y)}{y}\dd y
\]
可知 $\|\Theta_{j}\|_{L^{\infty}} = -\Theta_{0}(1)$, 所以 $|\Theta_{j}| \leq C\chi_{[-2^{j+2}, 2^{j+2}]}$. 因此, 具体写出卷积可知
\[
\|\chi_{[-2^{j+2}, 2^{j+2}]} * \chi_{[-2^{k+2}, 2^{k+2}]}\|_{L^{\infty}} = \min(2^{j+3}, 2^{k+3}),
\]
所以由支集大小的关系,
\[
\|\tau_{j} * \tau_{k}\|_{L^{\infty}} \leq C2^{-2k}\|\chi_{[-2^{j+2}, 2^{j+2}]} * \chi_{[-2^{k+2}, 2^{k+2}]}\|_{L^{\infty}} \leq C2^{\min(j,k)-2k}.
\]
由于可以交换 $j, k$ 的角色,
\[
\|\tau_{j} * \tau_{k}\|_{L^{\infty}} \leq C2^{\min(j,k)-2\max(j,k)}.
\]
由刚才计算的支集大小的关系,
\[
\|\tau_{j} * \tau_{k}\|_{L^{1}} \leq C2^{\min(j,k)-\max(j,k)} = C2^{-|j-k|}
\]
可得. 因此, 由于满足 Cotlar 引理的假设,
\[
Hf(x) = \sum_{j=-\infty}^{\infty}H_{j} * f(x) = \sum_{j=-\infty}^{\infty}\tau_{j} * f(x)
\]
在 $L^{2}(\mathbb{R})$ 的拓扑下成立, 这一点得以证明.

从以上的式子变形可以看出, $\pm$ 的相消反映在 $\Theta_{j}$ 具有紧支集这一点上, 再将其与分部积分结合, 得到了有效的估计. 在估计 $\tau_{j} * \tau_{k}$ 时, 不写成 $|\tau_{j} * \tau_{k}| \leq |\tau_{j}| * |\tau_{k}|$, 而是变形为 $|\tau_{j} * \tau_{k}| \leq |\Theta_{j}| * |\tau_{k}'|$, 再对 $\Theta_{j}$ 与 $\tau_{k}$ 作估计, 这是要点.

作为这个事实的应用, 也能处理开头所述的 $\varepsilon \downarrow 0$ 的极限. 这里确认设想~\eqref{eq:4.8} 是正当的. 实际上, 对一般的 $f \in L^{2}(\mathbb{R})$, 存在不依赖于 $\varepsilon > 0$ 的一致常数 $C$, 使得
\[
\left|\int_{|y| > \varepsilon}\frac{f(x-y)}{y}\dd y - \sum_{j=[\log_{2}\varepsilon]}^{\infty}\tau_{j} * f(x)\right| \leq CMf(x)
\]
成立. 实际上, 将左边的函数表成 $K_{\varepsilon} * f$ 的形式时, 成立 $|K_{\varepsilon}(x)| \leq C\varepsilon^{-1}\chi_{[0,8\varepsilon]}(|x|)$ 这个估计, 原因就在于此. $f \in C_{\mathrm{c}}^{\infty}(\mathbb{R})$ 时, 由于
\[
\lim_{\varepsilon \downarrow 0}\sum_{j=[\log_{2}\varepsilon]}^{\infty}\tau_{j} * f(x) = \lim_{\varepsilon \downarrow 0}\int_{|y| > \varepsilon}\frac{f(x-y)}{y}\dd y,
\]
所以至少对 $f \in C_{\mathrm{c}}^{\infty}(\mathbb{R})$ 可以使用 Lebesgue 收敛定理, 在 $L^{2}(\mathbb{R})$ 中
\[
\lim_{\varepsilon \downarrow 0}\int_{|y| > \varepsilon}\frac{f(x-y)}{y}\dd y = \lim_{\varepsilon \downarrow 0}\sum_{j=[\log_{2}\varepsilon]}^{\infty}\tau_{j} * f(x)
\]
可得. 由 $C_{\mathrm{c}}^{\infty}(\mathbb{R})$ 在 $L^{2}(\mathbb{R})$ 中的稠密性, 这个关系式可以提升到 $L^{2}(\mathbb{R})$ 的一切元.

\begin{exercise}\label{exer:4.2}
必要时, 利用 $0 < \varepsilon < 1$ 时
\[
\int_{\varepsilon < |x| < 1}\frac{x_{j}}{|x|^{n+1}}\dd x = 0
\]
成立, 证明对 $f \in C_{\mathrm{c}}^{\infty}(\mathbb{R}^{n})$,
\[
\lim_{\varepsilon \downarrow 0}\int_{\mathbb{R}^{n}\setminus B(x, \varepsilon)}\frac{(x_{j}-y_{j})f(y)}{|x-y|^{n+1}}\dd y = \lim_{\varepsilon \downarrow 0}\int_{|x-y| > \varepsilon}\frac{(x_{j}-y_{j})f(y)}{|x-y|^{n+1}}\dd y
\]
的极限存在.
\end{exercise}

\begin{exercise}\label{exer:4.3}
设 $C^{1}$ 级偶函数 $\rho: \mathbb{R} \to \mathbb{R}$ 满足
\[
(\star)\quad \sum_{j=-\infty}^{\infty}\rho(2^{j}\cdot) = \chi_{\mathbb{R}\setminus\{0\}}, \quad \operatorname{supp}(\rho) = [1, 3].
\]
回答下面的问题.

(1) 证明 $t \in \mathbb{R}$ 的函数 $\varphi(t) = \max(0, t-1)^{3}\max(0, 3-t)^{3}$ 是 $C^{1}$ 级的. 并求 $\operatorname{supp}(\varphi)$.

(2) 必要时参考 (1) 中考虑的 $\varphi$, 举出实际满足这样的 $(\star)$ 的 $\rho$ 的例子.

(3) $t \neq 0$ 时, 令 $\tau(t) = t^{-1}\rho(t)$. 进而令 $\tau(0) = 0$, 证明 $\tau$ 是 $C^{1}$ 级函数.
\end{exercise}

\begin{exercise}\label{exer:4.4}
$g \in L^{1}(\mathbb{R})$ 时, 考虑有界线性算子 $L_{g}: L^{2}(\mathbb{R}) \ni f \mapsto f * g \in L^{2}(\mathbb{R})$.

(1) 用 Young 不等式证明 $L_{g}$ 是 $L^{2}(\mathbb{R})$ 有界的.

(2) 令 $\check{g}(x) = \overline{g(-\cdot)}$. 证明 $L_{g}$ 在 Hilbert 空间 $L^{2}(\mathbb{R})$ 中的伴随是 $L_{\check{g}}$.
\end{exercise}

\begin{exercise}\label{exer:4.5}
设 $f, g \in C_{\mathrm{c}}^{1}(\mathbb{R})$. 即, $f, g$ 是 1 次可微的, $f, g$ 的 1 次导函数与 $f, g$ 本身全都连续, 且 $|x| \gg 1$ 时 $f(x) = g(x) = 0$. 此时, 证明 $(f * g)' = f' * g = f * g'$ 成立.
\end{exercise}

% ============ 第4章 §4.3 Calderón--Zygmund 分解与弱 (1,1) 有界性 (原书 pp.82-88) ============

\section{Calderón--Zygmund 分解与弱 (1,1) 有界性}\label{sec:4.3}
\index{0symbolEjf@$E_{j}f$}

\index{0symbolMdyadic@$M_{\dyadic}f$}

\index{C@Calder\'on--Zygmund 分解}

\index{E@二进制极大算子}

\index{E@二进制区间}

\index{L@Lebesgue 微分定理}














可测函数 $f: \mathbb{R} \to \mathbb{C}$ 的 $L^{2}$ 范数由
\[
\|f\|_{L^{2}} = \left(\int_{\mathbb{R}}|f(x)|^{2}\dd x\right)^{\frac{1}{2}}
\]
给出, 同样地 $L^{1}$ 范数由
\[
\|f\|_{L^{1}} = \int_{\mathbb{R}}|f(x)|\dd x
\]
给出. 前一节证明了 Hilbert 变换 $H$ 的 $L^{2}$ 有界性, 这里想考虑 $H$ 的 $L^{1}$ 有界性. 遗憾的是（练习题~\ref{exer:4.10} 参照）,
\[
H\chi_{[-1,1]}(x) = -\log\left|\frac{1-x}{1+x}\right| = O(x^{-1}) \quad (|x| \to \infty),
\]
所以 $H$ 不是 $L^{1}$ 有界的. Hilbert 变换定义为
\[
Hf(x) = \lim_{\varepsilon \downarrow 0}\int_{|y| > \varepsilon}\frac{f(x-y)}{y}\dd y,
\]
但回顾: 若 $f \in L^{2}(\mathbb{R})$, 则在 $L^{2}$ 的拓扑下收敛, 且 $H$ 是（由 Cotlar 引理）$L^{2}$ 有界的. 关于 Hilbert 变换的显著性质仅此而已, 剩下的只是利用 Hilbert 变换大致上是由与 $x^{-1}$ 的卷积给出这一点来推进本节的讨论.

作为考察 Hilbert 变换性质的起点, 考察对 $L^{1}$ 函数的输入有什么样的性质. 此时, 使用 1950 年代构思的 Calderón--Zygmund 分解.

这里定义在调和分析中起重要作用的二进极大算子与二进区间.

\begin{definition}[二进区间]\label{def:4.1}
对 $j, k \in \mathbb{Z}$, 区间 $I_{jk}$ 由
\begin{equation}\tag{4.9}
I_{jk} \equiv \left[\frac{k}{2^{j}}, \frac{k+1}{2^{j}}\right)
\label{eq:4.9}
\end{equation}
定义. 如此用 $j, k \in \mathbb{Z}$ 表为 $I_{jk}$ 的集合称为二进区间. 全体二进区间组成的集合记为 $\mathcal{D} = \mathcal{D}(\mathbb{R}) \equiv \{I_{jk} : j \in \mathbb{Z}, k \in \mathbb{Z}\}$. $j$ 世代的区间是指这些区间中属于
\begin{equation}\tag{4.10}
\mathcal{D}_{j} = \mathcal{D}_{j}(\mathbb{R}) \equiv \{Q \in \mathcal{D}(\mathbb{R}) : \ell(Q) = 2^{-j}\}
\label{eq:4.10}
\end{equation}
的全体.
\end{definition}

对二进区间 Lebesgue 微分定理也成立. 即, 对每个 $x \in \mathbb{R}$, 用 $Q_{j,x}$ 表示满足 $x \in Q_{j,x} \in \mathcal{D}_{j}(\mathbb{R})$ 的唯一区间, 则对任意局部可积函数 $f$, 几乎处处
\[
f(x) = \lim_{j \to \infty}2^{j}\int_{Q_{j,x}}f(y)\dd y
\]
成立. 实际上, 对连续函数这件事是显然的. 这个关系也称为 Lebesgue 微分定理.

下面的引理几乎显然, 省略证明（定理~\ref{thm:1.7} 也参照）.

\begin{lemma}[二进区间的嵌套结构]\label{lem:4.1}
两个二进区间相交时, 一方包含另一方.
\end{lemma}

定义了二进区间, 下面定义二进极大算子. 对区间 $I$, 用 $m_{I}(f)$ 表示 $f$ 在 $I$ 上的平均.

\begin{definition}[二进极大算子]\label{def:4.2}
对 $\mathbb{R}$ 上的局部可积函数 $f$,
\begin{align*}
E_{j}f(x) &\equiv E_{j}[f](x) \equiv \sum_{Q \in \mathcal{D}_{j}(\mathbb{R})}m_{Q}(f) \cdot \chi_{Q}(x), \\
M_{\dyadic}f(x) &\equiv \sup_{j \in \mathbb{Z}}E_{j}[|f|](x),
\end{align*}
并称 $M_{\dyadic}$ 为二进极大算子. $E_{j}f(x)$ 是局部可积函数 $f$ 在包含 $x$ 的唯一 $j$ 世代区间 $Q_{j,x} \in \mathcal{D}_{j}(\mathbb{R})$ 上的平均.
\end{definition}

$E_{j}$ 是线性的, 但 $M_{\dyadic}$ 不是线性的.

对一切 $j \in \mathbb{Z}$ 与 $x \in \mathbb{R}$,
\[
|E_{j}f(x)| = \sum_{Q \in \mathcal{D}_{j}}|m_{Q}(f)\chi_{Q}(x)| \leq 2^{j}\|f\|_{L^{1}}\sum_{Q \in \mathcal{D}_{j}}\chi_{Q}(x) = 2^{j}\|f\|_{L^{1}}
\]
可知不等式 $|E_{j}f(x)| \leq 2^{j}\|f\|_{L^{1}}$ 成立, 所以下面的引理是显然的.

\begin{lemma}\label{lem:4.2}
对任意 $f \in L^{1}(\mathbb{R})$ 与 $x \in \mathbb{R}$, $\lim_{j \to -\infty}E_{j}f(x) = 0$ 成立.
\end{lemma}

下面的定理是 Hilbert 变换的弱 (1,1) 有界性的证明的关键的函数分解, 是调和分析学的重要工具之一. 用 $\coprod$ 表示不相交的并.

\begin{theorem}[Calderón--Zygmund 区间的构造]\label{thm:4.2}
对可积函数 $f$ 与 $\lambda > 0$, 存在互不相交的二进区间的集合 $\{I_{j}\}_{j \in J} \subset \mathcal{D}$, 使得
\begin{equation}\tag{4.11}
\{x \in \mathbb{R} : M_{\dyadic}f(x) > \lambda\} = \coprod_{j \in J}I_{j}, \quad \lambda < m_{I_{j}}(|f|) \leq 2\lambda
\label{eq:4.11}
\end{equation}
成立.
\end{theorem}

\begin{proof}
为简化记号, 令 $E_{\lambda} \equiv \{x \in \mathbb{R} : M_{\dyadic}f(x) > \lambda\}$. 由二进极大算子的定义式
\[
M_{\dyadic}f(x) = \sup_{Q \in \mathcal{D}}m_{Q}(|f|)\chi_{Q}(x),
\]
特别地由这里出现的 sup 的性质, $E_{\lambda} = \bigcup_{j \in \mathbb{Z}}\{E_{j}[|f|] > \lambda\}$ 成立. 由引理~\ref{lem:4.2}, 对每个 $x \in E_{\lambda}$, 使 $E_{j}[|f|](x) > \lambda$, 即
\[
x \in \{y \in \mathbb{R} : E_{j}[|f|](y) > \lambda\}
\]
成立的 $j$ 存在下界. 即, 存在 $j^{*}$, 使得对一切 $j < j^{*}$, $E_{j}[|f|](x) \leq \lambda$ 成立. 因此, $E_{\lambda}$ 可以分割为
\[
E_{\lambda} = \coprod_{j \in \mathbb{Z}}\{x \in \mathbb{R} : E_{j}[|f|](x) > \lambda, \text{ 且 } k < j \text{ 时 } E_{k}[|f|](x) \leq \lambda\}.
\]
这里, 设 $j_{M} \in \mathbb{N} \cup \{\infty\}$,
\begin{equation}\tag{4.12}
\{x \in \mathbb{R} : E_{j}[|f|](x) > \lambda, \text{ 且 } k < j \text{ 时 } E_{k}[|f|](x) \leq \lambda\}
\label{eq:4.12}
\end{equation}
可以分割为 $I_{m}^{j} \in \mathcal{D}_{j}, m = 1, 2, \ldots, j_{M}$ 的和. 让 $j, m$ 变动, 将 $I_{m}^{j}$ 重新排列, 得到 $E_{\lambda}$ 的分割
\begin{equation}\tag{4.13}
E_{\lambda} = \coprod_{j \in J}I_{j}.
\label{eq:4.13}
\end{equation}
设 $x \in I_{j} =: I_{m_{0}}^{j_{0}}$. 由 $I_{m_{0}}^{j_{0}}$ 的取法, $m_{I_{j}}(|f|) = E_{j_{0}}[|f|](x) > \lambda$ 成立. 用 $I_{j}^{*} \in \mathcal{D}_{j_{0}-1}$ 表示二进区间 $I_{j}$ 的父区间, 即包含 $I_{j}$ 的体积为 2 倍的唯一区间. 由 $I_{m_{0}}^{j_{0}}$ 的取法, 即由式~\eqref{eq:4.12}, 必须有 $m_{I_{j}^{*}}(|f|) = E_{j_{0}-1}[|f|](x) \leq \lambda$. 这里,
\begin{align*}
\frac{1}{2}m_{I_{j}}(|f|) &= \frac{1}{2|I_{j}|}\int_{I_{j}}|f(y)|\dd y \\
&= \frac{1}{|I_{j}^{*}|}\int_{I_{j}}|f(y)|\dd y \\
&= \frac{1}{|I_{j}^{*}|}\int_{I_{j}^{*}}|f(y)|\dd y \\
&= m_{I_{j}^{*}}(|f|),
\end{align*}
所以
\[
\frac{1}{2}m_{I_{j}}(|f|) \leq m_{I_{j}^{*}}(|f|) = E_{j_{0}-1}[|f|](x) \leq \lambda,
\]
证明结束.
\end{proof}

由 $\{I_{j}\}_{j \in J} \subset \mathcal{D}$ 互不相交与式~\eqref{eq:4.11} 可知
\begin{equation}\tag{4.14}
|\{x \in \mathbb{R} : M_{\dyadic}f(x) > \lambda\}| = \sum_{j \in J}|I_{j}| \leq \frac{1}{\lambda}\|f\|_{L^{1}}
\label{eq:4.14}
\end{equation}
成立.

下面的定理是奇异积分算子的有界性的关键. $g$ 是 good 的略写, $b$ 是 bad 的略写. 这里的讨论是把第 1.6 节的讨论集大成并应用于函数.

\begin{theorem}[Calderón--Zygmund 分解]\label{thm:4.3}
对 $L^{1}$ 函数 $f$ 与 $\lambda > 0$, 存在满足下列性质的至多可数个二进区间的集族 $\{I_{j}\}_{j \in J}$ 与至多可数个 $L^{1}$ 函数的集族 $\{g\} \cup \{b_{j}\}_{j \in J}$.

(1) $f$ 分解为下面的形式.
\begin{equation}\tag{4.15}
f = g + \sum_{j \in J}b_{j} = g + b.
\label{eq:4.15}
\end{equation}

(2)（$L^{1}$ 条件）$g$ 的 $L^{1}$ 范数小于 $f$ 的.
\begin{equation}\tag{4.16}
\int_{\mathbb{R}}|g(x)|\dd x \leq \int_{\mathbb{R}}|f(x)|\dd x.
\label{eq:4.16}
\end{equation}

(3)（$L^{\infty}$ 条件）$g$ 属于 $L^{\infty}(\mathbb{R})$.
\begin{equation}\tag{4.17}
\text{对几乎所有的 } x \in \mathbb{R}, \quad |g(x)| \leq 2\lambda.
\label{eq:4.17}
\end{equation}

(4)（支集条件）对每个 $j \in J$, $b_{j}$ 的支集包含于 $I_{j}$ 的闭包.
\begin{equation}\tag{4.18}
\operatorname{supp}(b_{j}) \subset \overline{I_{j}}.
\label{eq:4.18}
\end{equation}

(5)（矩条件）对每个 $j$, $b_{j}$ 是可积函数, 且满足
\begin{equation}\tag{4.19}
\int_{\mathbb{R}}b_{j}(x)\dd x = 0.
\label{eq:4.19}
\end{equation}

(6)（水平集 $\{M_{\dyadic}f > \lambda\}$ 的分割）$\{I_{j}\}_{j \in J}$ 互不相交, 且
\begin{equation}\tag{4.20}
\{M_{\dyadic}f > \lambda\} = \coprod_{j \in J}I_{j}
\label{eq:4.20}
\end{equation}
成立.

这个分解称为（高度为 $\lambda$ 的）Calderón--Zygmund 分解.
\end{theorem}

\begin{proof}
取互不相交的二进区间的可数序列 $\{I_{j}\}_{j \in J}$, 满足
\begin{equation}\tag{4.21}
\lambda \leq m_{I_{j}}(|f|) \leq 2\lambda, \quad \{x \in \mathbb{R} : M_{\dyadic}f(x) > \lambda\} = \bigcup_{j \in J}I_{j}.
\label{eq:4.21}
\end{equation}
这个区间的序列的存在性已在定理~\ref{thm:4.2} 中证明. 满足式~\eqref{eq:4.18} 与式~\eqref{eq:4.19} 的 $b_{j}$ 如下定义（练习题~\ref{exer:4.8} 参照）.
\begin{equation}\tag{4.22}
b_{j} \equiv \left(f(x) - m_{I_{j}}(f)\right)\chi_{I_{j}}.
\label{eq:4.22}
\end{equation}
与 $b_{j}$ 相应, 令 $g \equiv f - \sum_{j \in J}b_{j}$, 则至少可以得到 $f$ 的分解. 这里注意, 和至少在"几乎所有的"的意义与 $L^{1}$ 的意义两种意义下成立. 这是因为 $b_{j}$ 的支集不共有内部. 由 $b_{j}$ 的定义, $g$ 的 $L^{1}$ 条件式~\eqref{eq:4.16} 成立. 剩下的只是对 $g$ 证明式~\eqref{eq:4.17}.

设 $x \in I_{j}$. 由于 $\{I_{\tilde{j}}\}_{\tilde{j} \in J}$ 互不相交,
\[
g(x) = f(x) - \sum_{j \in J}b_{j}(x) = f(x) - f(x) + m_{I_{j}}(f) = m_{I_{j}}(f)
\]
成立. 因此, 由式~\eqref{eq:4.21}, $|g(x)| \leq 2\lambda$. 反之, 设 $x \notin \bigcup_{j \in J}I_{j}$. 由 Lebesgue 微分定理, 式~\eqref{eq:4.17} 由 $|g(x)| = |f(x)| \leq M_{\dyadic}f(x) \leq \lambda$ 得出.
\end{proof}

\begin{note}\label{note:4.1}
$f \in L^{1}(\mathbb{R}) \cap L^{2}(\mathbb{R})$ 时, 在 $L^{1}(\mathbb{R}) \cap L^{2}(\mathbb{R})$ 的拓扑下 $b = \sum_{j \in J}b_{j}$, 且 $g$ 属于 $L^{1}(\mathbb{R}) \cap L^{2}(\mathbb{R})$, 这由 $g$ 的构造可知.
\end{note}

\begin{note}\label{note:4.2}
引理~\ref{lem:4.2}, 定理~\ref{thm:4.2}, \ref{thm:4.3} 对 $f \in L^{2}(\mathbb{R})$ 也同样成立.
\end{note}

证明 Hilbert 变换的弱 (1,1) 有界性.

\begin{theorem}\label{thm:4.4}
对一切 $\lambda$ 与 $f \in L^{1}(\mathbb{R}) \cap L^{2}(\mathbb{R})$,
\begin{equation}\tag{4.23}
|\{|Hf| > \lambda\}| \lesssim \frac{1}{\lambda}\|f\|_{L^{1}}
\label{eq:4.23}
\end{equation}
成立.
\end{theorem}

首先, 取 $f \in L^{1}(\mathbb{R}) \cap L^{2}(\mathbb{R})$ 这样的形式, 不只要求 $f \in L^{1}(\mathbb{R})$ 还要求 $f \in L^{2}(\mathbb{R})$, 是因为要使 $Hf$ 这个记号有意义, 眼下需要如此.

\begin{proof}
将 $f$ 按高度 $\lambda$ 作 Calderón--Zygmund 分解, 表为 $f = g + \sum_{j \in J}b_{j} = g + b$. 由注意~\ref{note:4.1} 可知 $b, g \in L^{2}(\mathbb{R})$, 所以 $Hf = Hg + Hb$ 成立. 转到补集考虑显然地,
\begin{equation}\tag{4.24}
\{|Hf| > \lambda\} \subset \left\{|Hg| > \frac{\lambda}{2}\right\} \cup \left\{|Hb| > \frac{\lambda}{2}\right\}
\label{eq:4.24}
\end{equation}
成立. 因此, 只需证明
\begin{equation}\tag{4.25}
\left|\left\{|Hg| > \frac{\lambda}{2}\right\}\right| \lesssim \frac{1}{\lambda}\|f\|_{L^{1}},
\label{eq:4.25}
\end{equation}
\begin{equation}\tag{4.26}
\left|\left\{|Hb| > \frac{\lambda}{2}\right\}\right| \lesssim \frac{1}{\lambda}\|f\|_{L^{1}}
\label{eq:4.26}
\end{equation}
这两个. 对应好函数的式~\eqref{eq:4.25} 容易证明. 由 Chebyshev 不等式,
\[
\left|\left\{|Hg| > \lambda\right\}\right| \leq \frac{1}{\lambda^{2}}\int_{\mathbb{R}}|Hg(x)|^{2}\dd x
\]
成立. 因此, 依次使用 $H$ 的 $L^{2}$ 有界性, 式~\eqref{eq:4.17}, 式~\eqref{eq:4.16},
\[
|\{|Hg| > \lambda\}| \lesssim \frac{1}{\lambda^{2}}\int_{\mathbb{R}}|g(x)|^{2}\dd x \lesssim \frac{1}{\lambda}\|g\|_{L^{1}} \lesssim \frac{1}{\lambda}\|f\|_{L^{1}}
\]
成立. 式~\eqref{eq:4.26} 稍难. 如下证明. 将 $I_{j}$ 表为 $I_{j} = [a_{j}-r_{j}, a_{j}+r_{j})$, 令 $2I_{j} = [a_{j}-2r_{j}, a_{j}+2r_{j})$. 另外, 对这种 $I_{j}$, 令 $c(I_{j}) = a_{j}$. 首先, 由式~\eqref{eq:4.14} 与式~\eqref{eq:4.20},
\begin{equation}\tag{4.27}
\left|\bigcup_{j \in J}2I_{j}\right| \leq 2\sum_{j \in J}|I_{j}| = 2|\{M_{\dyadic}f > \lambda\}| \leq \frac{2}{\lambda}\|f\|_{L^{1}},
\label{eq:4.27}
\end{equation}
所以只需证明
\begin{equation}\tag{4.28}
\left|\{|Hb| > \lambda\} \setminus \bigcup_{j \in J}2I_{j}\right| \lesssim \frac{1}{\lambda}\|f\|_{L^{1}},
\label{eq:4.28}
\end{equation}
但式~\eqref{eq:4.28} 的左边由 Chebyshev 不等式（定理~\ref{thm:2.17}）$\leq \frac{1}{\lambda}\sum_{j \in J}\int_{\mathbb{R}\setminus 2I_{j}}|Hb_{j}(x)|\dd x$, 所以式~\eqref{eq:4.26} 归结为
\begin{equation}\tag{4.29}
\int_{\mathbb{R}\setminus 2I_{j}}|Hb_{j}(x)|\dd x \lesssim \int_{I_{j}}|b_{j}(x)|\dd x \quad (j \in J).
\label{eq:4.29}
\end{equation}
设 $x \in \mathbb{R}\setminus 2I_{j}$, $y \in I_{j}$, 注意
\[
2|x-y| \geq \ell(I_{j}) \geq |y-c_{j}|
\]
且
\[
\left|\frac{1}{x-y} - \frac{1}{x-c_{j}}\right| \leq \frac{|y-c_{j}|}{|(x-y)(x-c_{j})|} \leq \frac{2\ell(I_{j})}{|x-c_{j}|^{2}}.
\]
$x \in \mathbb{R}\setminus 2I_{j}$ 时, 由这个逐点估计与关于 $b_{j}$ 的支集条件式~\eqref{eq:4.18} 和矩条件式~\eqref{eq:4.19}, 制造正负相消,
\[
Hb_{j}(x) = \int_{I_{j}}\frac{1}{x-y}b_{j}(y)\dd y = \int_{I_{j}}\left(\frac{1}{x-y} - \frac{1}{x-c(I_{j})}\right)b_{j}(y)\dd y,
\]
对其使用三角不等式,
\[
|Hb_{j}(x)| \leq \int_{I_{j}}\left|\frac{1}{x-y} - \frac{1}{x-c(I_{j})}\right| \cdot |b_{j}(y)|\dd y \lesssim \frac{\ell(I_{j})^{2}m_{I_{j}}(|b_{j}|)}{|x-c(I_{j})|^{2}},
\]
所以将其代入式~\eqref{eq:4.29}, 在 $\mathbb{R}\setminus 2I_{j}$ 上积分即可. 即, 将正负相消巧妙地累加即可.
\end{proof}

\begin{exercise}\label{exer:4.6}
将 $\mathcal{D}_{0}(\mathbb{R})$, $\mathcal{D}_{1}(\mathbb{R})$, $\mathcal{D}_{2}(\mathbb{R})$ 在数直线 $\mathbb{R}$ 上图示.
\end{exercise}

\begin{exercise}\label{exer:4.7}
对 $\mathbb{R}$ 上的可测函数 $f$,
\[
M_{\dyadic}f = \chi_{(-\infty, 0)}M_{\dyadic}[\chi_{(-\infty, 0)}f] + \chi_{[0, \infty)}M_{\dyadic}[\chi_{[0, \infty)}f]
\]
成立, 证明之.
\end{exercise}

\begin{exercise}\label{exer:4.8}
设 $f$ 是局部可积函数, $Q$ 是立方体. 令
\[
B(x) = \chi_{Q}(x)f(x) - \frac{\chi_{Q}(x)}{|Q|}\int_{Q}f(z)\dd z,
\]
证明 $B$ 是可积函数, 且满足 $\int_{Q}B(x)\dd x = 0$.
\end{exercise}

\begin{exercise}\label{exer:4.9}
设 $E$ 是具有有限测度的 Lebesgue 可测集合, 且能表为 $\mathcal{D}(\mathbb{R})$ 的元的并. 即, 存在 $\mathcal{A} \subset \mathcal{D}(\mathbb{R})$, 使得
\[
E = \bigcup_{Q \in \mathcal{A}}Q.
\]
(1) 证明 $\sup_{Q \in \mathcal{A}}|Q| < \infty$ 成立.

(2) 证明存在 $\mathcal{A}_{0} \subset \mathcal{A}$, 使得 $\sum_{Q \in \mathcal{A}_{0}}\chi_{Q} = \chi_{E}$ 成立.
\end{exercise}

\begin{exercise}\label{exer:4.10}
计算 $t > 2$ 时的 $H\chi_{[-1,1]}(t)$, 证明 $H\chi_{[-1,1]} \notin L^{1}([2, \infty))$.
\end{exercise}

% ============ 第4章 §4.4 Hilbert 变换的 L^p 有界性 (原书 pp.89-90) ============

\section{\texorpdfstring{Hilbert 变换的 $L^{p}$ 有界性}{Hilbert 变换的 Lp 有界性}}\label{sec:4.4}

这里, 利用 $L^{p}$ 空间的对偶性（定理~\ref{thm:2.19}）
\[
\|f\|_{L^{p}} = \sup\left\{\|f \cdot g\|_{L^{1}} : g \in L^{p'}(\mathbb{R}), \|g\|_{L^{p'}} = 1\right\},
\]
与 $H$ 在 $L^{2}(\mathbb{R})$ 有界, 以及对一切 $\lambda > 0$, $f \in L^{1}(\mathbb{R}) \cap L^{2}(\mathbb{R})$, $\lambda|\{|Hf| > \lambda\}| \lesssim \|f\|_{L^{1}}$, 以及 $f, g \in L^{2}(\mathbb{R})$ 时,
\[
\int_{\mathbb{R}}f(x)Hg(x)\dd x = -\int_{\mathbb{R}}g(x)Hf(x)\dd x,
\]
来考察 Hilbert 变换的 $L^{p}$ 有界性. 虽然是用极限记号定义的, 但实质上 $H$ 是卷积, 所以这个等式成立. 即, 目标是: 对 $f \in L^{p}(\mathbb{R}) \cap L^{2}(\mathbb{R})$ 与 $1 < p < \infty$,
\[
\|Hf\|_{L^{p}} \lesssim \|f\|_{L^{p}}
\]
成立. 这里, $\lesssim$ 中包含的常数不依赖于 $f$, 只依赖于 $p$. 如果能做到这一点, 就能由稠密性将 $H$ 定义为 $L^{p}(\mathbb{R})$ 上的有界线性算子.

考虑 $\chi_{[-1,1]}$ 可知 Hilbert 变换不是 $L^{1}$ 有界的, 但证明对其他的 $p \in (1, \infty)$, Hilbert 变换是 $L^{p}$ 有界的. 重要的是, 从这以后, 除上述对偶公式外, 并不特别需要 $H$ 的结构. 下面的命题称为弱 (1,1) 与 $L^{p}$ 有界性的插值, 即 Marcinkiewicz 插值定理, 它不依赖于维数, 所以作为一般维数的命题来证明.

\begin{theorem}[Marcinkiewicz 插值定理]\label{thm:4.5}
设 $1 < q < p < \infty$. 算子 $S: L^{p}(\mathbb{R}^{n}) \to L^{p}(\mathbb{R}^{n})$ 满足下列条件.

(1) $f \in L^{p}(\mathbb{R}^{n})$ 时, $Sf$ 是非负值函数.

(2)（齐次性）$a \in \mathbb{C}$, $f \in L^{p}(\mathbb{R}^{n})$ 时, 对几乎所有的 $x \in \mathbb{R}^{n}$, $S[af](x) = |a|Sf(x)$ 成立.

(3)（次可加性）$f, g \in L^{p}(\mathbb{R}^{n})$ 时, 对几乎所有的 $x \in \mathbb{R}^{n}$, $S[f+g](x) \leq Sf(x) + Sg(x)$ 成立.

(4)（弱 (1,1) 有界性）对一切 $f \in L^{1}(\mathbb{R}^{n}) \cap L^{p}(\mathbb{R}^{n})$ 与 $\lambda > 0$, $|\{Sf > \lambda\}| \lesssim \frac{1}{\lambda}\|f\|_{L^{1}}$ 成立.

(5)（$L^{p}$ 有界性）对一切 $f \in L^{p}(\mathbb{R}^{n})$, $\|Sf\|_{L^{p}} \lesssim \|f\|_{L^{p}}$ 成立.

此时, 对一切 $f \in L^{p}(\mathbb{R}^{n}) \cap L^{q}(\mathbb{R}^{n})$, $\|Sf\|_{L^{q}} \lesssim \|f\|_{L^{q}}$ 成立. 不过, $\lesssim$ 中包含的常数只依赖于弱 (1,1) 有界性的常数, $L^{p}$ 有界性的常数, $p$, $q$.
\end{theorem}

作为这样的 $S$ 的例子, 除极大算子 $M$ 之外, 还有 Hilbert 变换 $H$ 与绝对值复合而成的 $f \mapsto |Hf|$. 把定理的叙述和证明与定理~\ref{thm:1.3} 比较也是一件乐事.

\begin{proof}
用定理~\ref{thm:3.3} 改写 $L^{q}$ 范数.
\begin{equation}\tag{4.30}
\|Sf\|_{L^{q}}^{q} = q2^{q}\int_{0}^{\infty}\lambda^{q-1}|\{Sf > 2\lambda\}|\dd\lambda.
\label{eq:4.30}
\end{equation}
这里, 对 $\lambda > 0$, 由次可加性, 弱 (1,1) 有界性, $L^{p}$ 有界性以及 Chebyshev 不等式
\[
|\{Sf > \lambda\}| \lesssim \frac{1}{\lambda^{p}}\int_{\mathbb{R}^{n}}Sf(x)^{p}\dd x,
\]
得
\begin{align*}
|\{Sf > 2\lambda\}| &\leq |\{S[f\chi_{\{|f| > \lambda\}}] > \lambda\}| + |\{S[f\chi_{\{|f| \leq \lambda\}}] > \lambda\}| \\
&\lesssim \frac{1}{\lambda}\int_{\{|f| > \lambda\}}|f(x)|\dd x + \frac{1}{\lambda^{p}}\int_{\{|f| \leq \lambda\}}|f(x)|^{p}\dd x \\
&\lesssim \int_{\mathbb{R}^{n}}\left(\frac{\chi_{\{|f| > \lambda\}}(x)|f(x)|}{\lambda} + \frac{\chi_{\{|f| \leq \lambda\}}(x)|f(x)|^{p}}{\lambda^{p}}\right)\dd x.
\end{align*}
因此, 对此不等式乘以 $\lambda^{q-1}$ 再积分,
\begin{align*}
\|Sf\|_{L^{q}} &\lesssim \int_{0}^{\infty}\left\{\int_{\mathbb{R}^{n}}\left(\frac{\chi_{\{|f| > \lambda\}}(x)|f(x)|}{\lambda^{2-q}} + \frac{\chi_{\{|f| \leq \lambda\}}(x)|f(x)|^{p}}{\lambda^{p-q+1}}\right)\dd x\right\}\dd\lambda
\end{align*}
成立. 用 Fubini 定理交换积分顺序后, 关于 $\lambda$ 积分, 得 $\|Sf\|_{L^{q}} \lesssim \|f\|_{L^{q}}$.
\end{proof}

\begin{theorem}[Hilbert 变换的 $L^{p}$ 有界性]\label{thm:4.6}
对一切 $1 < p < \infty$, 存在满足关系式 $c_{p} = c_{p'}$ 的常数族 $\{c_{p}\}_{1 < p < \infty}$, 使得对一切 $f \in L^{2}(\mathbb{R}) \cap L^{p}(\mathbb{R})$, $\|Hf\|_{L^{p}} \leq c_{p}\|f\|_{L^{p}}$ 成立.
\end{theorem}

注意, $p'$ 的共轭指数是 $p$, 且 $p \in (1, 2)$ 与 $p' \in (2, \infty)$ 等价, 所以对一切 $1 < p < \infty$ 关系式 $c_{p} = c_{p'}$ 成立与对一切 $1 < p < 2$ 关系式 $c_{p} = c_{p'}$ 成立是等价的.

\begin{proof}
分别证明 $1 < p < 2$ 的情形与 $2 < p < \infty$ 的情形即可. $1 < p < 2$ 的情形用定理~\ref{thm:4.5} 插值定理~\ref{thm:4.4} 与 $L^{2}$ 有界性. 如此, 对 $1 < p \leq 2$, 对一切 $f \in L^{2}(\mathbb{R}) \cap L^{p}(\mathbb{R})$, $\|Hf\|_{L^{p}} \leq c_{p}\|f\|_{L^{p}}$ 成立. 对 $2 < p < \infty$, 令 $c_{p'} = c_{p}$, 用对偶性证明对 $2 < p < \infty$, 对一切 $f \in L^{2}(\mathbb{R}) \cap L^{p}(\mathbb{R})$, $\|Hf\|_{L^{p}} \leq c_{p}\|f\|_{L^{p}}$ 也成立. 具体地如下计算. 任意固定 $\varepsilon > 0$, 取满足
\begin{equation}\tag{4.31}
\|Hf\|_{L^{p}} \leq (1+\varepsilon)\left|\int_{\mathbb{R}}Hf(x) \cdot \overline{g(x)}\dd x\right|, \quad \|g\|_{L^{p'}} = 1
\label{eq:4.31}
\end{equation}
的 $g \in L^{2}(\mathbb{R}) \cap L^{p'}(\mathbb{R})$. 由于 $f, g \in L^{2}(\mathbb{R})$, 满足 $\int_{\mathbb{R}}Hf(x) \cdot \overline{g(x)}\dd x = -\int_{\mathbb{R}}f(x) \cdot \overline{Hg(x)}\dd x$. 用 Hölder 不等式, 得 $\|Hf\|_{L^{p}} \leq (1+\varepsilon)\|f\|_{L^{p}} \cdot \|Hg\|_{L^{p'}}$. 由 $1 < p' < 2$, 得 $\|Hf\|_{L^{p}} \leq c_{p'}(1+\varepsilon)\|f\|_{L^{p}}$. 由于 $\varepsilon > 0$ 任意, 得 $\|Hf\|_{L^{p}} \leq c_{p'}\|f\|_{L^{p}}$. 因此, 定理~\ref{thm:4.6} 证明完毕.
\end{proof}

\begin{exercise}\label{exer:4.11}
设 $1 \leq p < \infty$, 证明 $L^{p}(\mathbb{R}) \cap L^{2}(\mathbb{R})$ 在 $L^{p}(\mathbb{R})$ 中稠密.
\end{exercise}

% ============ 第4章 §4.5 Hilbert 变换的几乎处处收敛 (原书 pp.91-94) ============

\section{Hilbert 变换的几乎处处收敛}\label{sec:4.5}

$1 \leq p < \infty$, $f \in L^{p}(\mathbb{R})$ 时, Hilbert 变换 $H$ 由稠密性连续地扩张, 对这样扩张的 $H$, 考察 $Hf(x)$ 与
\[
\lim_{\varepsilon \downarrow 0}\int_{|y| > \varepsilon}\frac{f(x-y)}{y}\dd y
\]
的关系. $L^{2}$ 收敛的函数列取子列后也在逐点收敛的意义下收敛, 所以 $f \in C_{\mathrm{c}}^{\infty}(\mathbb{R})$ 的情形, 对几乎所有的 $x \in \mathbb{R}$ 这些一致. 这里证明 $f \in L^{p}(\mathbb{R})$, $1 \leq p < \infty$ 时, 对几乎所有的 $x \in \mathbb{R}$ 它们一致.

这里, 至今得到的成果被总动员.

(1) $1 < p < \infty$ 时, Hardy--Littlewood 极大算子 $M$ 与 Hilbert 变换 $H$ 在 $L^{p}(\mathbb{R})$ 中有界.

(2) $M$ 与 $H$ 在 $L^{1}(\mathbb{R})$ 中弱有界, 因此, 例如对一切 $\lambda > 0$ 与 $f \in L^{1}(\mathbb{R})$, $|\{Mf > \lambda\}| \lesssim \lambda^{-1}\|f\|_{L^{1}}$ 成立.

(3) 证明 $H$ 的弱 (1,1) 有界性时使用的 Calderón--Zygmund 分解这一技法.

(4) 证明 Lebesgue 微分定理时使用的稠密性讨论.

为了证明几乎处处极限存在, 考察
\[
H^{*}f(x) = \sup_{\varepsilon > 0}\left|\int_{|y| > \varepsilon}\frac{f(x-y)}{y}\dd y\right|
\]
具有什么样的性质是有效的. 作为关于 $H^{*}$ 的性质的第一条, 证明下面的不等式.

\begin{theorem}\label{thm:4.7}
对 $f \in C_{\mathrm{c}}^{\infty}(\mathbb{R})$, 对 $x \in \mathbb{R}$ 的逐点估计
\begin{equation}\tag{4.32}
H^{*}f(x) \leq M[Hf](x) + CMf(x)
\label{eq:4.32}
\end{equation}
成立, 因此, 对一切 $f \in L^{p}(\mathbb{R})$, $1 < p < \infty$,
\begin{equation}\tag{4.33}
\|H^{*}f\|_{L^{p}} \leq C\|f\|_{L^{p}}
\label{eq:4.33}
\end{equation}
成立.
\end{theorem}

若承认式~\eqref{eq:4.32}, 则式~\eqref{eq:4.33} 由 $M$, $H$ 都 $L^{p}$ 有界是显然的.

\begin{proof}
需要对每个 $\varepsilon > 0$ 证明
\[
\left|\int_{|y| > \varepsilon}\frac{f(x-y)}{y}\dd y\right| \leq M[Hf](x) + CMf(x).
\]
由于 $M[f(\varepsilon\cdot)] = Mf(\varepsilon\cdot)$, $H[f(\varepsilon\cdot)] = Hf(\varepsilon\cdot)$, 作变量替换只需考虑 $\varepsilon = 2$ 的情形.

令
\[
\tau(x) = \frac{15}{16}\chi_{(-1,1)}(x)(x^{2}-1)^{2} \quad (x \in \mathbb{R}).
\]
由积分公式
\[
\int_{\alpha}^{\beta}(x-\alpha)^{m}(\beta-x)^{n}\dd x = \frac{m!n!(\beta-\alpha)^{m+n+1}}{(m+n+1)!} \quad (\alpha, \beta \in \mathbb{R}, m, n \in \mathbb{N})
\]
可知 $\|\tau\|_{L^{1}} = 1$, 注意这一点. 于是,
\[
\int_{|y| > 2}\frac{f(x-y)}{y}\dd y = \int_{\mathbb{R}}\frac{f(x-y)}{y}\chi_{\mathbb{R}\setminus(-2,2)}(y)\dd y,
\]
所以
\begin{align*}
\int_{|y| > 2}\frac{f(x-y)}{y}\dd y - f * H\tau(x) &= \int_{\mathbb{R}}f(x-y)\left(\frac{\chi_{\mathbb{R}\setminus(-2,2)}(y)}{y} - H\tau(y)\right)\dd y.
\end{align*}
这里, 令
\[
\Psi(y) = \frac{\chi_{\mathbb{R}\setminus(-2,2)}(y)}{y} - H\tau(y), \quad \Phi(y) = \frac{1}{1+y^{2}}.
\]
$|x| > 3$ 时, 由支集的关系积分范围只需考虑 $[-1, 1]$, 所以
\[
|H\tau(x)| \leq \int_{\mathbb{R}}\frac{\tau(y)}{|y-x|}\dd y \leq \frac{1}{2}\|\tau\|_{L^{1}} = \frac{1}{2}.
\]
因此, 由于 $\tau$ 是具有紧支集的可微函数, 由式~\eqref{eq:4.7}, $\Psi \in L^{\infty}(\mathbb{R})$. 另外, $|y| > 2$ 时, $\chi_{\mathbb{R}\setminus(-2,2)}(y) = 1$, 所以
\begin{align*}
\Psi(y) &= \frac{1}{y} - \int_{\mathbb{R}}\frac{\tau(z)}{y-z}\dd z \\
&= \int_{\mathbb{R}}\left(\frac{\tau(z)}{y} - \frac{\tau(z)}{y-z}\right)\dd z \\
&= -\int_{\mathbb{R}}\frac{z\tau(z)}{y(y-z)}\dd z \\
&= O(y^{-2})
\end{align*}
成立. 因此, 存在常数 $C$ 使得 $|\Psi(y)| \leq \frac{C}{1+y^{2}} = C\Phi(y)$. 由这个不等式,
\begin{align*}
\left|\int_{|y| > \varepsilon}\frac{f(x-y)}{y}\dd y\right| &\leq |f * H\tau(x)| + \left|\int_{|y| > \varepsilon}\frac{f(x-y)}{y}\dd y - f * H\tau(x)\right| \\
&\leq |f * H\tau(x)| + C\Phi * |f|(x)
\end{align*}
可知成立. 由引理~\ref{lem:3.1} 与 $\|\Phi\|_{1} = \pi$,
\begin{align*}
|f * H\tau(x)| &= |Hf * \tau(x)| \leq M[Hf](x), \\
\Phi * |f|(x) &\leq \pi Mf(x)
\end{align*}
成立. 代入以上即得结论.
\end{proof}

这个命题是 $p > 1$ 时对 $L^{p}$ 函数的性质, $p = 1$ 时下面的结果成立.

\begin{theorem}\label{thm:4.8}
对一切 $\lambda > 0$ 与 $f \in L^{1}(\mathbb{R})$,
\begin{equation}\tag{4.34}
|\{H^{*}f > \lambda\}| \lesssim \frac{1}{\lambda}\|f\|_{L^{1}}
\label{eq:4.34}
\end{equation}
成立.
\end{theorem}

\begin{proof}
将 $f$ 按高度 $\lambda$ 作 Calderón--Zygmund 分解, 表为 $f = g + \sum_{j \in J}b_{j} = g + b$. 此时, 如练习题~\ref{exer:4.12} 所述, $H^{*}f \leq H^{*}g + H^{*}b$ 成立. 转到补集考虑显然地,
\begin{equation}\tag{4.35}
\{H^{*}f > \lambda\} \subset \left\{H^{*}g > \frac{\lambda}{2}\right\} \cup \left\{H^{*}b > \frac{\lambda}{2}\right\}
\label{eq:4.35}
\end{equation}
成立. 因此, 只需证明
\begin{equation}\tag{4.36}
\left|\left\{H^{*}g > \frac{\lambda}{2}\right\}\right| \lesssim \frac{1}{\lambda}\|f\|_{L^{1}},
\label{eq:4.36}
\end{equation}
\begin{equation}\tag{4.37}
\left|\left\{H^{*}b > \frac{\lambda}{2}\right\}\right| \lesssim \frac{1}{\lambda}\|f\|_{L^{1}}
\label{eq:4.37}
\end{equation}
这两个. 式~\eqref{eq:4.36} 容易证明. 与 $H$ 的弱 (1,1) 有界性的证明一样, 用 Chebyshev 不等式,
\[
|\{H^{*}g > \lambda\}| \leq \frac{1}{\lambda^{2}}\int_{\mathbb{R}}H^{*}g(x)^{2}\dd x
\]
开始. 依次使用定理~\ref{thm:4.7}, 式~\eqref{eq:4.17}, 式~\eqref{eq:4.16},
\[
|\{H^{*}g > \lambda\}| \leq \frac{1}{\lambda^{2}}\int_{\mathbb{R}}|g(x)|^{2}\dd x \lesssim \frac{1}{\lambda}\|g\|_{L^{1}} \lesssim \frac{1}{\lambda}\|f\|_{L^{1}}
\]
成立. 但确证式~\eqref{eq:4.37} 稍难. 实际上式~\eqref{eq:4.37} 如下证明. 首先, 与 $2I_{j}$ 同样地定义 $3I_{j}$, 由式~\eqref{eq:4.14},
\begin{equation}\tag{4.38}
\left|\bigcup_{j \in J}3I_{j}\right| \leq 3\sum_{j \in J}|I_{j}| = 3|\{M_{\dyadic}f > \lambda\}| \leq \frac{3}{\lambda}\|f\|_{L^{1}},
\label{eq:4.38}
\end{equation}
所以只需证明
\begin{equation}\tag{4.39}
\left|\{H^{*}b > \lambda\} \setminus \bigcup_{j \in J}3I_{j}\right| \lesssim \frac{1}{\lambda}\|f\|_{L^{1}}.
\label{eq:4.39}
\end{equation}
固定 $x \in \mathbb{R}\setminus\bigcup_{j \in J}3I_{j}$ 与 $\varepsilon > 0$.

(1) 对满足 $[x-\varepsilon, x+\varepsilon] \supset I_{j}$ 的 $j \in J$, $H_{\varepsilon}b_{j}(x) = 0$ 成立.

(2) 对满足 $[x-\varepsilon, x+\varepsilon] \cap I_{j} = \emptyset$ 的 $j \in J$, $H_{\varepsilon}b_{j}(x) = Hb_{j}(x)$ 成立.

(3) 设不满足这两个条件的 $j$ 的全体为 $Z$. 则对 $j \in Z$, $x - \varepsilon \notin I_{j}$ 且 $x + \varepsilon \in \overline{I_{j}}$, 或 $x - \varepsilon \in \overline{I_{j}}$ 且 $x + \varepsilon \notin I_{j}$ 成立. 因此, $\{I_{j}\}_{j=1}^{\infty}$ 不相交,
\[
I_{j} \subset (x-10\varepsilon, x+10\varepsilon) \quad (j \in Z),
\]
且对 $j \in Z$, $y \in I_{j}$, 若 $x - \varepsilon \in \overline{I_{j}}$, 则 $|c(I_{j}) - y| \lesssim \varepsilon$, $|x - y| \gtrsim \varepsilon$, $|x - c(I_{j})| \gtrsim \varepsilon$, 所以
\begin{align*}
\sum_{j \in Z}|H_{\varepsilon}b_{j}(x)| &\leq \frac{C}{\varepsilon}\sum_{j \in Z}\int_{I_{j}}|b_{j}(y)|\dd y \\
&\leq \frac{C}{\varepsilon}\int_{(x-10\varepsilon, x+10\varepsilon)}|f(y)|\dd y \\
&\leq CMf(x)
\end{align*}
成立.

因此, 式~\eqref{eq:4.39} 归结为
\[
\left|\left\{\sum_{j \in J}|Hb_{j}| + CMf > \lambda\right\} \setminus \bigcup_{j \in J}3I_{j}\right| \lesssim \frac{1}{\lambda}\|f\|_{L^{1}}.
\]
到了这一步, 与通常的 $H$ 的弱 (1,1) 有界性的证明相同.
\end{proof}

\begin{theorem}\label{thm:4.9}
$1 \leq p < \infty$, $f \in L^{p}(\mathbb{R})$ 时, 对几乎所有的 $x \in \mathbb{R}$, $H_{\varepsilon}f(x)$ 在 $\varepsilon \downarrow 0$ 时的极限存在.
\end{theorem}

\begin{proof}
$p > 1$ 的情形与 $p = 1$ 的情形类似, 所以证明 $p = 1$ 的情形. 只需证明极限不存在的点全体组成的集合 $E$ 的测度为 $0$. 于是, 用上极限 $\limsup$ 给 $E$ 的各点赋予级别, 将 $E$ 按级别分开考虑. 即, 对一切自然数 $k \in \mathbb{N}$,
\[
E_{k} \equiv \left\{x \in \mathbb{R} : \limsup_{\varepsilon, \varepsilon' \downarrow 0}|H_{\varepsilon}f(x) - H_{\varepsilon'}f(x)| > \frac{1}{k}\right\},
\]
则
\[
E = \bigcup_{k=1}^{\infty}E_{k},
\]
所以只需证明各 $E_{k}$ 的测度为 $0$. 注意, 极限的存在性对 $g \in C_{\mathrm{c}}^{\infty}(\mathbb{R})$ 是显然的. 任意取 $\varepsilon > 0$. 通过证明 $|E_{k}| \leq Ck\varepsilon$ 来完成定理的证明.

取 $g \in C_{\mathrm{c}}^{\infty}(\mathbb{R})$ 使得 $\|f - g\|_{L^{1}} \leq \varepsilon$. 这由稠密性是可能的. 作为集合的包含关系
\begin{equation}\tag{4.40}
E_{k} = \left\{x \in \mathbb{R} : \limsup_{\varepsilon, \varepsilon' \downarrow 0}|H_{\varepsilon}[f-g](x) - H_{\varepsilon'}[f-g](x)| > \frac{1}{k}\right\} \subset \left\{x \in \mathbb{R} : 2H^{*}[f-g](x) > \frac{1}{k}\right\}
\label{eq:4.40}
\end{equation}
成立, 所以由弱 (1,1) 不等式,
\[
|E_{k}| \leq Ck\|f - g\|_{L^{1}} \leq Ck\varepsilon.
\]
由于 $\varepsilon > 0$ 任意, 得 $|E_{k}| = 0$. 由于 $k \in \mathbb{N}$ 也任意, 得 $|E| = 0$. 即, 对几乎所有的 $x \in \mathbb{R}$, $H_{\varepsilon}f(x)$ 在 $\varepsilon \downarrow 0$ 时的极限存在.
\end{proof}

\begin{exercise}\label{exer:4.12}
对 $f, g \in L^{1}(\mathbb{R}^{n})$, 证明 $H^{*}[f+g](x) \leq H^{*}f(x) + H^{*}g(x)$ 成立.
\end{exercise}

% ============ 第4章 §4.6 高维情形的推广 (原书 pp.95-97) ============

\section{高维情形的推广}\label{sec:4.6}
\index{0symbolCZ@CZ 核 (Calder\'on--Zygmund)}

\index{H@H\"ormander 条件}

\index{R@Riesz 变换}








以上的理论推广到高维. 直接的推广是 Riesz 变换, 由下式给出.
\begin{equation}\tag{4.41}
R_{j}f(x) \equiv \lim_{\varepsilon \downarrow 0}\int_{\mathbb{R}^{n}\setminus B(x, \varepsilon)}\frac{x_{j}-y_{j}}{|x-y|^{n+1}}f(y)\dd y.
\label{eq:4.41}
\end{equation}
再将其一般化也是可能的, 并且有用.

\begin{definition}[CZ (Calderón--Zygmund) 核]\label{def:4.3}
设 $0 < \varepsilon \leq 1$. 令 $\mathrm{diag} = \{(x, x) \in \mathbb{R}^{2n} : x \in \mathbb{R}^{n}\}$. $L^{2}$ 有界线性算子 $T$ 称为广义 Calderón--Zygmund 算子, 是指存在满足以下条件的可测函数 $K: \mathbb{R}^{2n}\setminus\mathrm{diag} \to \mathbb{C}$. 另外, 这个 $K$ 称为 CZ 核.

(1)（尺寸条件）
\begin{equation}\tag{4.42}
|K(x, y)| \lesssim |x-y|^{-n}, \quad (x, y) \in \mathbb{R}^{2n}\setminus\mathrm{diag}
\label{eq:4.42}
\end{equation}
成立.

(2) 当满足 $|x-z| > 2|y-z| > 0$ 时, Hörmander 条件
\begin{equation}\tag{4.43}
|K(x, y) - K(x, z)| + |K(y, x) - K(z, x)| \lesssim |y-z|^{\varepsilon}|x-z|^{-n-\varepsilon}
\label{eq:4.43}
\end{equation}
成立.

(3) $f \in L^{2}(\mathbb{R}^{n})$ 的支集紧时, 在 $\mathrm{supp}(f)$ 之外几乎处处
\[
Tf(x) = \int_{\mathbb{R}^{n}}K(x, y)f(y)\dd y
\]
成立.
\end{definition}

满足这个定义的积分算子唯一地扩张为弱 (1,1) 有界, $L^{p}$ 有界（$1 < p < \infty$）的算子, 这也与 Hilbert 变换一样可以证明. 另外, 已知几乎处处收敛的充分必要条件, 是极限
\[
\lim_{\varepsilon \downarrow 0}\int_{\varepsilon < |x-y| < 1}K(x, y)\dd y
\]
对几乎所有的 $x$ 存在.

\begin{exercise}\label{exer:4.13}
设 $1 < p < \infty$, $\varphi \in C^{\infty}(\mathbb{R}^{n}\setminus\{0\})$. 令
\[
A = \sup_{x \in \mathbb{R}^{n}\setminus\{0\}}|x|^{n}|\varphi(x)|, \quad B = \sup_{x \in \mathbb{R}^{n}\setminus\{0\}}|x|^{n+1}|\nabla\varphi(x)|.
\]
对任意 $f \in L^{2}(\mathbb{R}^{n})$, 关于在 $\mathrm{supp}(f)$ 之外几乎处处可以表为
\[
Tf(x) = \lim_{\varepsilon \downarrow 0}\int_{\mathbb{R}^{n}\setminus B(\varepsilon)}\varphi(x-y)f(y)\dd y
\]
的 $L^{2}$ 有界线性算子 $T$, 回答下面的问题.

(1) 对每个固定的 $P, Q \geq 0$, 存在使得 $\zeta(\cdot, P, Q)$, $\zeta(P, \cdot, Q)$, $\zeta(P, Q, \cdot)$ 成为单调递增函数的函数 $\zeta = \zeta_{p}: [0, \infty)^{3} \to [0, \infty)$, 对 $f \in L^{p}(\mathbb{R}^{n})$,
\[
\|Tf\|_{L^{p}} \leq \zeta(\|T\|_{L^{2} \to L^{2}}, A, B)\|f\|_{L^{p}}
\]
形式的估计成立, 证明之.

(2) 实际上, 证明存在常数 $C_{p}$, 使得
\[
\|\varphi * f\|_{L^{p}} \leq C_{p}(\|T\|_{L^{2} \to L^{2}} + A + B)\|f\|_{L^{p}}
\]
成立.
\end{exercise}


% 第5章 Schwartz 分布及其性质 (原书 pp.98-122, PDF pp.104-128)
\chapter{Schwartz 分布及其性质}
% ============ 第5章 §5.1 Schwartz 函数 (原书 pp.98-103) ============

这里, 作为深入考察至今为止出现的数学变换, 特别是 Hilbert 变换等的工具, 引入 Schwartz 函数空间与 Schwartz 分布空间.

\section{Schwartz 函数}\label{sec:5.1}
\index{0symboldx@$\mathrm{d}x'$}

\index{G@Gaussian}

\index{S@Schwartz 函数空间 $S(\mathbb{R}^{n})$}








在处理函数的性质时, 以 $L^{p}(\mathbb{R}^{n})$ 为代表的积分值等量的方面, 与以有界连续函数全体构成的 Banach 空间 $\mathrm{BC}(\mathbb{R}^{n})$ 为代表的连续性等质的方面, 都是重要的. 巧妙地组合这些性质, 就能进一步分析函数的性质. 兼备两者的函数是最理想的, 但作为这样的典型函数例子, 有具有紧支集的无限次可微函数全体构成的线性空间 $C_{\mathrm{c}}^{\infty}(\mathbb{R}^{n})$. 然而, $C_{\mathrm{c}}^{\infty}(\mathbb{R}^{n})$ 与数学中出现的变换, 例如积分变换等的相性不好. 这是因为 $C_{\mathrm{c}}^{\infty}(\mathbb{R}^{n})$ 的函数经过变换后, 往往不再是 $C_{\mathrm{c}}^{\infty}(\mathbb{R}^{n})$ 的函数. 因此, 考察比 $C_{\mathrm{c}}^{\infty}(\mathbb{R}^{n})$ 更大, 但被认为保持良好性质的函数空间. 令 $\mathbb{N}_{0} = \{0, 1, \ldots\}$, 用 $\mathbb{N}_{0}^{n}$ 表示多重指标全体构成的集合. 即,
\[
\mathbb{N}_{0}^{n} = \left\{(\alpha_{1}, \alpha_{2}, \ldots, \alpha_{n}) : \alpha_{j} \in \mathbb{N}_{0}, j = 1, 2, \ldots, n\right\}.
\]
对多重指标 $\beta = (\beta_{1}, \beta_{2}, \ldots, \beta_{n}) \in \mathbb{N}_{0}^{n}$ 与 $f \in C^{\infty}(\mathbb{R}^{n})$, 定义
\[
|\beta| = \sum_{j=1}^{n}\beta_{j} \quad (\beta \text{ 的长度}),
\]
\[
x^{\beta} \equiv x_{1}^{\beta_{1}}x_{2}^{\beta_{2}}\cdots x_{n}^{\beta_{n}} \quad (x \text{ 的 } \beta \text{ 次乘方}),
\]
\[
\partial^{\beta}f(x) \equiv \left(\frac{\partial}{\partial x_{1}}\right)^{\beta_{1}}\left(\frac{\partial}{\partial x_{2}}\right)^{\beta_{2}}\cdots\left(\frac{\partial}{\partial x_{n}}\right)^{\beta_{n}}f(x) \quad (f \text{ 的 } \beta \text{ 次偏导数}).
\]
又, 设 $\alpha = (\alpha_{1}, \alpha_{2}, \ldots, \alpha_{n})$, $\beta = (\beta_{1}, \beta_{2}, \ldots, \beta_{n})$ 为多重指标时, $\alpha \geq \beta$ 是指各分量满足相同的大小关系 $\alpha_{j} \geq \beta_{j}$. 并且, 此时用通常的二项式系数 ${}_{n}\mathrm{C}_{r}$ 定义一般的二项式系数 ${}_{\alpha}\mathrm{C}_{\beta}$:
\[
{}_{\alpha}\mathrm{C}_{\beta} = \prod_{j=1}^{n}{}_{\alpha_{j}}\mathrm{C}_{\beta_{j}}.
\]

\begin{definition}[Schwartz 函数空间 $\mathcal{S}(\mathbb{R}^{n})$]\label{def:5.1}
$\mathcal{S}(\mathbb{R}^{n})$ 是如下定义的函数集合. 对 $\alpha, \beta \in \mathbb{N}_{0}^{n}$, 暂且令
\begin{equation}\tag{5.1}
\varphi_{(\alpha,\beta)}(x) \equiv x^{\alpha}\partial^{\beta}\varphi(x) \quad (x \in \mathbb{R}^{n}),
\label{eq:5.1}
\end{equation}
将 $\mathcal{S}(\mathbb{R}^{n})$ 定义为对所有 $\alpha, \beta \in \mathbb{N}_{0}^{n}$ 满足 $\varphi_{(\alpha,\beta)} \in L^{\infty}(\mathbb{R}^{n})$ 的 $\varphi \in C^{\infty}(\mathbb{R}^{n})$ 全体构成的集合, $\mathcal{S}(\mathbb{R}^{n})$ 的元称为 Schwartz 函数.
\end{definition}

对 $\varphi \in C^{\infty}(\mathbb{R}^{n})$, 令
\[
p_{\alpha,\beta}(\varphi) = \sup_{x \in \mathbb{R}^{n}}|\varphi_{(\alpha,\beta)}(x)|.
\]

\begin{example}\label{exa:5.1}
对 $x \in \mathbb{R}^{n}$, 由 $E(x) \equiv e^{-\frac{1}{2}|x|^{2}}$ 给出的函数 $E$ 称为 Gaussian, 属于 $\mathcal{S}(\mathbb{R}^{n})$. 又 $C_{\mathrm{c}}^{\infty}(\mathbb{R}^{n}) \subset \mathcal{S}(\mathbb{R}^{n})$. 因此, $C_{\mathrm{c}}^{\infty}(\mathbb{R}^{n})$ 是 $\mathcal{S}(\mathbb{R}^{n})$ 的真子集.
\end{example}

在下面的例子中, 使用多项式定理
\[
(A_{1} + A_{2} + \cdots + A_{n})^{k} = \sum_{k_{1}, k_{2}, \ldots, k_{n} \in \mathbb{N}_{0}, k_{1} + k_{2} + \cdots + k_{n} = k}\frac{k!}{k_{1}!k_{2}!\cdots k_{n}!}A_{1}^{k_{1}}A_{2}^{k_{2}}\cdots A_{n}^{k_{n}}.
\]

\begin{example}\label{exa:5.2}
设 $\varphi \in \mathcal{S}(\mathbb{R}^{n})$. 由上述注意,
\begin{align*}
(1+|x|)^{n+1} &\leq 2^{n+1}(1+|x|^{2})^{n+1} \\
&= 2^{n+1}\sum_{\substack{k_{0}, k_{1}, k_{2}, \ldots, k_{n} \in \mathbb{N}_{0} \\ k_{0}+k_{1}+k_{2}+\cdots+k_{n}=n+1}}\frac{(n+1)!|x^{2k}|}{k_{1}!k_{2}!\cdots k_{n}!k_{0}!}
\end{align*}
所以 $\Phi(x) \equiv (1+|x|)^{n+1}\varphi(x)$ 是有界的. 因此, $|\varphi(x)| \leq (1+|x|)^{-n-1}\|\Phi\|_{L^{\infty}}$, 由例~\ref{exa:2.1}, $\varphi \in L^{1}(\mathbb{R}^{n})$. 由于对任意多重指标 $\alpha, \beta$, $\varphi_{(\alpha,\beta)} \in \mathcal{S}(\mathbb{R}^{n})$, 所以 $\varphi_{(\alpha,\beta)} \in L^{1}(\mathbb{R}^{n})$.
\end{example}

\begin{example}\label{exa:5.3}
设 $f, g \in C^{\infty}(\mathbb{R}^{n})$. 由广义 Leibniz 法则,
\[
\partial^{\alpha}[f(x)g(x)] = \sum_{\beta \leq \alpha}{}_{\alpha}\mathrm{C}_{\beta}\partial^{\beta}f(x) \cdot \partial^{\alpha-\beta}g(x)
\]
成立. 因此, 设 $f, g \in \mathcal{S}(\mathbb{R}^{n})$, 则 $f \cdot g \in \mathcal{S}(\mathbb{R}^{n})$.
\end{example}

\begin{example}\label{exa:5.4}
设 $M > 0$. 对满足 $|h(x)| + |\nabla h(x)| \lesssim (1+|x|)^{M}$, $x \in \mathbb{R}^{n}$ 的 $C^{1}$ 级函数 $h$ 与 $f \in \mathcal{S}(\mathbb{R}^{n})$, 分部积分公式
\[
\int_{\mathbb{R}^{n}}\partial_{j}f(x) \cdot h(x)\dd x = -\int_{\mathbb{R}^{n}}f(x) \cdot \partial_{j}h(x)\dd x \quad (j = 1, 2, \ldots, n)
\]
成立. 由例~\ref{exa:2.1}, 注意两边都是有限的.

实际上, 考虑乘积的微分, 只需证明
\begin{equation}\tag{5.2}
\int_{\mathbb{R}^{n}}\partial_{j}[f(x) \cdot h(x)]\dd x = 0
\label{eq:5.2}
\end{equation}
即可.

由于具有对称性, 以 $j = n$ 的情形证明. 记 $x' = (x_{1}, x_{2}, \ldots, x_{n-1})$, $\dd x' = \dd x_{1}\dd x_{2}\cdots\dd x_{n-1}$. 由例~\ref{exa:2.1} 与 Fubini 定理,
\[
\int_{\mathbb{R}^{n}}\partial_{n}[f(x) \cdot h(x)]\dd x = \int_{\mathbb{R}^{n-1}}\left(\int_{\mathbb{R}}\partial_{n}[f(x) \cdot h(x)]\dd x_{n}\right)\dd x'
\]
可以这样分解. 固定 $x' \in \mathbb{R}^{n-1}$. 将微积分基本定理
\[
\int_{-R}^{R}\varphi'(t)\dd t = \varphi(R) - \varphi(-R)
\]
应用于 $\varphi(t) = f(x', t)h(x', t)$,
\begin{align*}
\int_{\mathbb{R}}\partial_{n}[f(x) \cdot h(x)]\dd x_{n} &= \lim_{R \to \infty}\int_{-R}^{R}\partial_{n}[f(x) \cdot h(x)]\dd x_{n} \\
&= \lim_{R \to \infty}(f(x', R)h(x', R) - f(x', -R)h(x', -R)) \\
&= 0
\end{align*}
成立. 由以上, 可知式~\eqref{eq:5.2} 成立.
\end{example}

作为集合的 $\mathcal{S}(\mathbb{R}^{n})$ 虽然如此定义, 但可以给 $\mathcal{S}(\mathbb{R}^{n})$ 赋予距离. 即, 当 $f, g \in \mathcal{S}(\mathbb{R}^{n})$ 时, 通过定义 f, g 的距离 $d(f, g)$ 为
\[
d(f, g) = \sum_{\alpha, \beta \in \mathbb{N}_{0}^{n}}\frac{p_{\alpha,\beta}(f-g)}{2^{|\alpha|+|\beta|}p_{\alpha,\beta}(f-g)+1},
\]
$(\mathcal{S}(\mathbb{R}^{n}), d)$ 是距离空间. 由
\[
\frac{p_{\alpha,\beta}(f-g)}{2^{|\alpha|+|\beta|}p_{\alpha,\beta}(f-g)+1} \leq \frac{1}{2^{|\alpha|+|\beta|}}, \quad \sum_{\alpha, \beta \in \mathbb{N}_{0}^{n}}\frac{1}{2^{|\alpha|+|\beta|}} < \infty
\]
可以使用 Lebesgue 收敛定理, 从而可以说出下面的事.

\begin{theorem}\label{thm:5.1}
$\{f_{j}\}_{j=1}^{\infty} \subset \mathcal{S}(\mathbb{R}^{n})$ 收敛于 $f \in \mathcal{S}(\mathbb{R}^{n})$ 的必要充分条件是, 对所有 $\alpha, \beta \in \mathbb{N}_{0}^{n}$,
\[
\lim_{j \to \infty}p_{\alpha,\beta}(f_{j}-f) = 0
\]
成立.
\end{theorem}

对 $\varphi \in \mathcal{S}(\mathbb{R}^{n})$, 令
\[
p_{N}(\varphi) \equiv \sum_{|\beta| \leq N}\sup_{x \in \mathbb{R}^{n}}(1+|x|)^{N}|\partial^{\beta}\varphi(x)|.
\]
也可以令
\[
p_{N}(\varphi) \equiv \sum_{|\alpha|, |\beta| \leq N}\sup_{x \in \mathbb{R}^{n}}|x^{\alpha}\partial^{\beta}\varphi(x)|,
\]
但暂且为统一起见采用前者. 虽然无论采用哪一个都是相同的, 但作为作者个人的感觉, 有着想在可能的定义中采用给出最大值的定义作为定义的心理. 即, 基于
\[
\sum_{|\alpha|, |\beta| \leq N}\sup_{x \in \mathbb{R}^{n}}|x^{\alpha}\partial^{\beta}\varphi(x)| \leq (N+1)^{n}\sum_{|\beta| \leq N}\sup_{x \in \mathbb{R}^{n}}(1+|x|)^{N}|\partial^{\beta}\varphi(x)|
\]
想采用前者作为定义.

\begin{example}\label{exa:5.5}
$p_{N}, p_{\alpha,\beta}$ 等被称为泛函, 但引入 $p_{N}, p_{\alpha,\beta}$ 的目的, 并不是具体计算各 $\varphi$ 的值 $p_{N}(\varphi), p_{\alpha,\beta}(\varphi)$, 而是为了讨论这些值的有限性, 或者用 $p_{N}, p_{\alpha,\beta}$ 讨论某种函数列的收敛性. 因此, 将这些自由替换来考虑也可以. 立足于这一观点, 可知除这些之外还可能有各种各样的泛函. 举例来看. 对 $\varphi \in \mathcal{S}(\mathbb{R}^{n})$, 定义
\[
\begin{aligned}
p_{N}(\varphi) = p_{N}^{1}(\varphi) &= \sup_{\beta \in \mathbb{N}_{0}^{n}, |\beta| \leq N}\sup_{x \in \mathbb{R}^{n}}(1+|x|)^{N}|\partial^{\beta}\varphi(x)|, \\
p_{N}^{2}(\varphi) &= \sup_{\beta \in \mathbb{N}_{0}^{n}, |\beta| \leq N}\sup_{x \in \mathbb{R}^{n}}(1+|x|^{2})^{N}|\partial^{\beta}\varphi(x)|
\end{aligned}
\]
由于 $1+|x| \leq 2(1+|x|^{2})$, $1+|x|^{2} \leq (1+|x|)^{2}$,
\[
p_{N}^{1}(\varphi) \leq 2^{N}p_{N}^{2}(\varphi) \leq 2^{N}p_{2N}^{1}(\varphi)
\]
成立. 因此, $\{p_{N}^{1}\}_{N=1}^{\infty}$, $\{p_{N}^{2}\}_{N=1}^{\infty}$ 无论将哪一个作为 $\{p_{N}\}_{N=1}^{\infty}$ 的定义都可以.
\end{example}

\begin{example}\label{exa:5.6}
设 $O \subset \mathcal{S}(\mathbb{R}^{n})$ 是包含 0 函数的开集. 此时存在某个 $L > 0$, 使得
\[
\{\varphi \in \mathcal{S}(\mathbb{R}^{n}) : Lp_{L}(\varphi) < 1\} \subset O
\]
成立. 实际上, 由于开集由 d 生成, 存在某个 $r > 0$, 使得
\[
\{\varphi \in \mathcal{S}(\mathbb{R}^{n}) : d(\varphi, 0) < r\} \subset O
\]
成立. 取自然数 N 使得
\[
\sum_{\alpha, \beta \in \mathbb{N}_{0}^{n}}\frac{\chi_{(N,\infty)}(\max(|\alpha|, |\beta|))}{2^{|\alpha|+|\beta|}} < \frac{r}{2}.
\]
对这样的 N, 选取 $L \geq N$ 使得
\[
L^{-1}\sharp\{\alpha, \beta \in \mathbb{N}_{0}^{n} : \max(|\alpha|, |\beta|) \leq N\} < \frac{r}{2}.
\]
于是, 若 $Lp_{L}(\varphi) < 1$, 则
\begin{align*}
d(\varphi, 0) &= \sum_{\alpha, \beta \in \mathbb{N}_{0}^{n}}\frac{p_{\alpha,\beta}(\varphi)\chi_{(N,\infty)}(\max(|\alpha|, |\beta|))}{2^{|\alpha|+|\beta|}p_{\alpha,\beta}(\varphi)+1} \\
&\quad + \sum_{\alpha, \beta \in \mathbb{N}_{0}^{n}}\frac{p_{\alpha,\beta}(\varphi)\chi_{[0,N]}(\max(|\alpha|, |\beta|))}{2^{|\alpha|+|\beta|}p_{\alpha,\beta}(\varphi)+1} \\
&\leq \sum_{\alpha, \beta \in \mathbb{N}_{0}^{n}}\frac{\chi_{(N,\infty)}(\max(|\alpha|, |\beta|))}{2^{|\alpha|+|\beta|}} \\
&\quad + L^{-1}\sharp\{\alpha, \beta \in \mathbb{N}_{0}^{n} : \max(|\alpha|, |\beta|) \leq N\} \\
&< \frac{r}{2} + \frac{r}{2} = r
\end{align*}
所以
\[
\{\varphi \in \mathcal{S}(\mathbb{R}^{n}) : Lp_{L}(\varphi) < 1\} \subset \{\varphi \in \mathcal{S}(\mathbb{R}^{n}) : d(\varphi, 0) < r\} \subset O
\]
成立.
\end{example}

由下面的例子可知, $C_{\mathrm{c}}^{\infty}(\mathbb{R}^{n})$ 在 $\mathcal{S}(\mathbb{R}^{n})$ 中是稠密的.

\begin{example}\label{exa:5.7}
设 $f \in \mathcal{S}(\mathbb{R}^{n})$, $\psi \in \mathcal{S}(\mathbb{R}^{n})$ 满足 $\chi_{B(1)} \leq \psi \leq \chi_{B(2)}$. 此时, 在 $\mathcal{S}(\mathbb{R}^{n})$ 的拓扑中, 当 $j \to \infty$ 时, $\psi(2^{-j}\cdot)f \to f$ 成立. 实际上, 取 $\alpha, \beta \in \mathbb{N}_{0}^{n}$, 基于定理~\ref{thm:5.1} 证明
\[
\lim_{j \to \infty}p_{\alpha,\beta}(\psi(2^{-j}\cdot)f - f) = 0.
\]
由 Leibniz 法则,
\begin{align*}
&x^{\alpha}(\partial^{\beta}(\psi(2^{-j}x)f(x) - f(x))) \\
&= \sum_{0 \neq \gamma \leq \beta}2^{-j|\gamma|}{}_{\beta}\mathrm{C}_{\gamma}x^{\alpha}\partial^{\gamma}\psi(2^{-j}x) \cdot \partial^{\beta-\gamma}f(x) \\
&\quad + x^{\alpha}(\psi(2^{-j}x)\partial^{\beta}f(x) - \partial^{\beta}f(x))
\end{align*}
这里, 使用
\begin{align*}
&\sup_{x \in \mathbb{R}^{n}}|x^{\alpha}\partial^{\gamma}\psi(2^{-j}x) \cdot \partial^{\beta-\gamma}f(x)| \leq \|\partial^{\gamma}\psi\|_{L^{\infty}}\sup_{x \in \mathbb{R}^{n}}|x^{\alpha}\partial^{\beta-\gamma}f(x)| < \infty, \\
&\sup_{x \in \mathbb{R}^{n}}|x^{\alpha}[\psi(2^{-j}x)\partial^{\beta}f(x) - \partial^{\beta}f(x)]| \\
&= \sup_{x \in \mathbb{R}^{n}\setminus B(2^{j})}|x^{\alpha}[\psi(2^{-j}x)\partial^{\beta}f(x) - \partial^{\beta}f(x)]| \\
&\leq \sup_{x \in \mathbb{R}^{n}\setminus B(2^{j})}|x^{\alpha}\partial^{\beta}f(x)| \\
&\leq 2^{-j}\sup_{x \in \mathbb{R}^{n}\setminus B(2^{j})}|x| \cdot |x^{\alpha}\partial^{\beta}f(x)|
\end{align*}
得到
\[
p_{\alpha,\beta}(\psi(2^{-j}\cdot)f - f) = \sup_{x \in \mathbb{R}^{n}}|x^{\alpha}(\partial^{\beta}(\psi(2^{-j}x)f(x) - f(x)))| = O(2^{-j}).
\]
因此 $\lim_{j \to \infty}p_{\alpha,\beta}(\psi(2^{-j}\cdot)f - f) = 0$.
\end{example}

\begin{exercise}\label{exer:5.1}
设 $A, B, C > 0$. 证明
\[
\frac{B}{AB+1} + \frac{C}{AC+1} \geq \frac{B+C}{A(B+C)+1}.
\]
【提示】当 $t \geq 0$ 时, 令 $\varphi(t) = \frac{t}{t+1}$, 则问题中的不等式可以改写为
\[
\varphi(AB + AC) \leq \varphi(AB) + \varphi(AC)
\]
($\iff \varphi(AB + AC) + \varphi(0) \leq \varphi(AB) + \varphi(AC)$).
利用
\[
\varphi(AB) - \varphi(AB + AC) = -\int_{AB}^{AB+AC}\varphi'(t)\dd t,
\]
\[
\varphi(AC) - \varphi(0) = \int_{0}^{AC}\varphi'(t)\dd t,
\]
证明起来很容易.
\end{exercise}

\begin{exercise}\label{exer:5.2}
证明 $(\mathcal{S}(\mathbb{R}^{n}), d)$ 是完备的距离空间.
\end{exercise}

\begin{exercise}\label{exer:5.3}
设 $1 \leq p \leq \infty$. 利用 $\mathcal{S}(\mathbb{R}^{n})$ 的任意元 $\varphi$ 满足
\[
q(\varphi) = \sup_{x \in \mathbb{R}^{n}}|\varphi(x)|, \quad r_{j}(\varphi) = \sup_{x \in \mathbb{R}^{n}}|x_{j}|^{n+1}|\varphi(x)| < \infty
\]
这一点, 回答下面的问题.

(1) 必要时使用
\[
(1+|x|) \geq \prod_{j=1}^{n}\sqrt[n]{1+|x_{j}|},
\]
证明 $f(x) = \frac{1}{(1+|x|)^{n+1}} \in L^{p}(\mathbb{R}^{n})$.

(2) 设 $g \in \mathcal{S}(\mathbb{R}^{n})$. 必要时使用
\[
\begin{aligned}
(1+|x|)^{n+1} &\leq \left(1+\sum_{j=1}^{n}|x_{j}|\right)^{n+1} \\
&\leq (n+1)^{n+1} + (n+1)^{n+1}\sum_{j=1}^{n}|x_{j}|^{n+1},
\end{aligned}
\]
证明 $M = \sup_{x \in \mathbb{R}^{n}}(1+|x|)^{n+1}|g(x)| < \infty$.

(3) 证明 $\mathcal{S}(\mathbb{R}^{n}) \subset L^{p}(\mathbb{R}^{n})$.
\end{exercise}

\begin{exercise}\label{exer:5.4}
当 $f, g \in \mathcal{S}(\mathbb{R}^{n})$ 时, 证明 $f \cdot g \in \mathcal{S}(\mathbb{R}^{n})$. (补充上述例5.3 的细节.)
\end{exercise}

% ============ 第5章 §5.2 Schwartz 函数的卷积 (原书 pp.104-107) ============

\section{Schwartz 函数的卷积}\label{sec:5.2}

一般地, 函数的卷积由
\[
f * g(x) = \int_{\mathbb{R}^{n}}f(x-y)g(y)\dd y
\]
给出. 对 Schwartz 函数这个定义也有效. 并且, Schwartz 函数的卷积仍是 Schwartz 函数. 积分在某种意义下可以用``Riemann 和''写出, 但重要的是, 用``Riemann 和''逼近卷积在 Schwartz 函数空间的拓扑中也是有效的. 为了避免计算过程变长, 稍加巧思, 决定按``Riemann 和风格''来分割. 首先, 证明刚才所述``Schwartz 函数的卷积仍是 Schwartz 函数''这一事实.

卷积是用积分号表示的运算, 必须证明卷积后的函数可微, 所以想先确认微分号与积分号的交换 (定理~\ref{thm:2.13}). 那么, 这里确认 $f * g$ 是 $\mathcal{S}(\mathbb{R}^{n})$ 的元.

\begin{theorem}[* 在 $\mathcal{S}(\mathbb{R}^{n})$ 中的乘法性]\label{thm:5.2}
卷积运算 * 在 $\mathcal{S}(\mathbb{R}^{n})$ 中是封闭的. 即, 当 $f, g \in \mathcal{S}(\mathbb{R}^{n})$ 时, $f * g \in \mathcal{S}(\mathbb{R}^{n})$.
\end{theorem}

\begin{proof}
首先必须证明 $f * g$ 属于 $C^{\infty}(\mathbb{R}^{n})$. 为此, 考察 $f * g(x)$ 的被积函数 $f(x-y)g(y)$, 记其为 $F(x, y)$. $\partial_{x_{j}}F(x, y) = \partial_{j}f(x-y)g(y)$ 成立, 对各 $x, y \in \mathbb{R}^{n}$, $F(x, y)$ 满足下面的条件.
\[
|F(x, y)|(1+|y|)^{n+1} = |f(x-y)g(y)|(1+|y|)^{n+1}
\]
\begin{align*}
&\leq \sup_{z \in \mathbb{R}^{n}}|f(z)| \cdot \sup_{w \in \mathbb{R}^{n}}(1+|w|)^{n+1}|g(w)|, \\
&|\partial_{x_{j}}F(x, y)|(1+|y|)^{n+1} = |\partial_{j}f(x-y)g(y)|(1+|y|)^{n+1} \\
&\leq \sup_{z \in \mathbb{R}^{n}}|\partial_{j}f(z)| \cdot \sup_{w \in \mathbb{R}^{n}}(1+|w|)^{n+1}|g(w)|.
\end{align*}
因此可以交换微分号与积分号, $\partial_{j}[f*g](x) = (\partial_{j}f)*g(x)$ 成立. 特别地, $f * g$ 至少 1 次可微, 且该等式成立. 由这个计算也可知, $f * g$, $\partial_{j}[f*g]$ 是有界函数. 由于 $\partial_{j}f$ 与 $f$ 一样是 $\mathcal{S}(\mathbb{R}^{n})$ 的元, 这个考察可以反复使用, 可知 $f * g$ 任意次可微, 且包括 $f * g$ 在内连续, $f * g \in C^{\infty}(\mathbb{R}^{n})$ 成立.

取 $\alpha \in \mathbb{N}_{0}^{n}$ 与 $N \in \mathbb{N}$. 此时, 由积分的三角不等式, $1+|x| \leq (1+|x-y|)(1+|y|)$, $x, y \in \mathbb{R}^{n}$, 以及
\[
\int_{\mathbb{R}^{n}}\frac{\dd y}{(1+|y|)^{n+1}} < \infty
\]
得到
\begin{align}
&\sup_{x \in \mathbb{R}^{n}}(1+|x|)^{N}|\partial^{\alpha}[f*g](x)| \nonumber \\
&\leq \sup_{x \in \mathbb{R}^{n}}(1+|x|)^{N}\int_{\mathbb{R}^{n}}|\partial^{\alpha}f(x-y)| \cdot |g(y)|\dd y \nonumber \\
&\leq \sup_{x \in \mathbb{R}^{n}}\int_{\mathbb{R}^{n}}(1+|x-y|)^{N}|\partial^{\alpha}f(x-y)| \cdot (1+|y|)^{N}|g(y)|\dd y \nonumber \\
&\lesssim \sup_{w \in \mathbb{R}^{n}}(1+|w|)^{N}|\partial^{\alpha}f(w)| \cdot \sup_{y \in \mathbb{R}^{n}}(1+|y|)^{N+n+1}|g(y)|.
\tag{5.3}
\label{eq:5.3}
\end{align}
因此, 可以结论 $f * g \in \mathcal{S}(\mathbb{R}^{n})$.
\end{proof}

卷积是由
\[
f * g(x) = \int_{\mathbb{R}^{n}}f(x-y)g(y)\dd y
\]
给出的函数, 但通过分割积分区域, 可以将其分割为
\[
f * g(x) = \sum_{k \in \mathbb{Z}^{n}}\int_{k+[0,1]^{n}}f(x-y)g(y)\dd y.
\]
重要的是, 这并非各点收敛的意义, 而是在 $\mathcal{S}(\mathbb{R}^{n})$ 的拓扑中进行的.

\begin{lemma}\label{lem:5.1}
设 $f, g \in \mathcal{S}(\mathbb{R}^{n})$. 对 $k \in \mathbb{Z}^{n}$, 令
\[
\Phi_{k}(x) = \int_{k+[0,1]^{n}}f(x-y)g(y)\dd y.
\]
(1) 对各 $k \in \mathbb{Z}^{n}$, $\Phi_{k}$ 属于 $\mathcal{S}(\mathbb{R}^{n})$.
(2) $f * g = \sum_{k \in \mathbb{Z}^{n}}\Phi_{k}$ 成立. 准确地说,
\[
\lim_{R \to \infty}\sum_{k \in \mathbb{Z}^{n} \cap B(R)}\Phi_{k} = f * g
\]
在 $\mathcal{S}(\mathbb{R}^{n})$ 的拓扑中成立.
\end{lemma}

\begin{proof}
(1) 与前面一样, $\Phi_{k}$ 可微, 且
\begin{align*}
&\sup_{x \in \mathbb{R}^{n}}(1+|x|)^{N}|\partial^{\alpha}\Phi_{k}(x)| \\
&= \sup_{x \in \mathbb{R}^{n}}(1+|x|)^{N}\int_{k+[0,1]^{n}}|\partial^{\alpha}f(x-y)| \cdot |g(y)|\dd y \\
&\leq \sup_{x \in \mathbb{R}^{n}}\int_{k+[0,1]^{n}}(1+|x-y|)^{N}|\partial^{\alpha}f(x-y)| \cdot (1+|y|)^{N}|g(y)|\dd y \\
&\lesssim (1+|k|)^{-n-1}\sup_{y \in \mathbb{R}^{n}}(1+|y|)^{N}|\partial^{\alpha}f(y)|\sup_{y \in k+[0,1]^{n}}(1+|y|)^{N}|g(y)|
\end{align*}
所以得到结论.

(2) 基于刚才的计算, 设 $R > 0$, 由
\begin{align*}
\sum_{k \in \mathbb{Z}^{n}\setminus B(R)}(1+|k|)^{-n-1} &\lesssim \int_{\mathbb{R}^{n}\setminus B(R)}\frac{\dd x}{(1+|x|)^{n+1}} \\
&\lesssim \frac{1}{\sqrt{R}}\int_{\mathbb{R}^{n}\setminus B(R)}\frac{\dd x}{(1+|x|)^{n+1/2}} \\
&\lesssim \frac{1}{\sqrt{R}}
\end{align*}
得到
\begin{align*}
&\sup_{x \in \mathbb{R}^{n}}(1+|x|)^{N}\left|\partial^{\alpha}[f*g](x) - \sum_{k \in \mathbb{Z}^{n} \cap B(R)}\partial^{\alpha}\Phi_{k}(x)\right| \\
&\leq \sum_{k \in \mathbb{Z}^{n}\setminus B(R)}(1+|k|)^{-n-1}\sup_{y \in \mathbb{R}^{n}}(1+|y|)^{N}|\partial^{\alpha}f(y)| \\
&\quad \times \sup_{y \in k+[0,1]^{n}}(1+|y|)^{N+n+1}|g(y)| \\
&\lesssim R^{-1/2}
\end{align*}
所以得到所希望的收敛.
\end{proof}

这里, 进一步分析 $\Phi_{k}$. 记 $m = (m_{1}, m_{2}, \ldots, m_{n})$. 将 $\Phi_{k}$ 按 Riemann 和风格分割,
\[
\Phi_{k}(x) = \lim_{N \to \infty}\sum_{m \in \{1, 2, \ldots, N\}^{n}}f\left(x - k - \frac{m}{N}\right)\int_{k+\frac{m}{N}+[-N^{-1},0]^{n}}g(y)\dd y
\]
对所有 x 成立, 但可知该等式也是在 $\mathcal{S}(\mathbb{R}^{n})$ 的拓扑中发生的.

\begin{lemma}\label{lem:5.2}
在 $\mathcal{S}(\mathbb{R}^{n})$ 的拓扑中,
\[
\Phi_{k} = \lim_{N \to \infty}\sum_{m \in \{1, 2, \ldots, N\}^{n}}f\left(\cdot - k - \frac{m}{N}\right)\int_{k+\frac{m}{N}+[-N^{-1},0]^{n}}g(y)\dd y
\]
成立.
\end{lemma}

\begin{proof}
由练习题~\ref{exer:5.5}, 在一致收敛的意义下, 只要能证明
\[
\lim_{N \to \infty}\sum_{m \in \{1, 2, \ldots, N\}^{n}}(1+|\cdot|)^{N}\partial^{\beta}f\left(\cdot - k - \frac{m}{N}\right)\int_{k+\frac{m}{N}+[-N^{-1},0]^{n}}g(y)\dd y
\]
收敛于 $(1+|\cdot|)^{N}\partial^{\beta}\Phi_{k}$ 即可, 但将 $\partial^{\beta}f$ 替换为 f, 设 $\beta = 0$ 来考虑也可以.

令
\[
h(x, y, m, N, k) = \left(f(x-y) - f\left(x - k - \frac{m}{N}\right)\right)g(y).
\]
于是,
\begin{align*}
&\Phi_{k}(x) - \sum_{m \in \{1, 2, \ldots, N\}^{n}}f\left(x - k - \frac{m}{N}\right)\int_{k+\frac{m}{N}+[-N^{-1},0]^{n}}g(y)\dd y \\
&= \sum_{m \in \{1, 2, \ldots, N\}^{n}}\int_{k+\frac{m}{N}+[-N^{-1},0]^{n}}h(x, y, m, N, k)\dd y.
\end{align*}
与前面一样计算, 由于 $|y - k - N^{-1}m| \lesssim N^{-1}$, 存在依赖于 k, x, y, m, N 的 $t \in [0, 1]$, 使得
\begin{align*}
&\int_{k+\frac{m}{N}+[-N^{-1},0]^{n}}|h(x, y, m, N, k)|\dd y \\
&\lesssim N^{-1}\int_{k+\frac{m}{N}+[-N^{-1},0]^{n}}\left|\nabla f\left(x - (1-t)y - tk - t\frac{m}{N}\right)\right| \cdot |g(y)|\dd y
\end{align*}
成立. 因此, 注意
\[
1 + |x| \leq \left(1 + \left|x - y + ty - tk - t\frac{m}{N}\right|\right)\left(1 + \left|(1-t)y + tk + t\frac{m}{N}\right|\right),
\]
得到
\begin{align*}
&\sup_{x \in \mathbb{R}^{n}}(1+|x|)^{N}\int_{k+\frac{m}{N}+[-N^{-1},0]^{n}}|h(x, y, m, N, k)|\dd y \\
&\lesssim \frac{1}{N}\sup_{x \in \mathbb{R}^{n}}\int_{k+\frac{m}{N}+[-N^{-1},0]^{n}}\left(1 + \left|x - (1-t)y - tk - t\frac{m}{N}\right|\right)^{N} \\
&\quad \times (1+|y|)^{N}\left|\nabla f\left(x - (1-t)y - tk - t\frac{m}{N}\right)\right| \cdot |g(y)|\dd y \\
&\lesssim \frac{1}{N(1+|k|)^{n+1}}\sup_{y \in \mathbb{R}^{n}}(1+|y|)^{N}|\nabla f(y)|\sup_{y \in k+[0,1]^{n}}(1+|y|)^{N}|g(y)|.
\end{align*}
之后只需将其代入再取极限.
\end{proof}

\begin{exercise}\label{exer:5.5}
设 $\{\varphi_{j}\}_{j=1}^{\infty} \subset \mathcal{S}(\mathbb{R}^{n})$ 收敛于 $\varphi \in \mathcal{S}(\mathbb{R}^{n})$ 的必要充分条件是, 对所有自然数 N 与 $\alpha \in \mathbb{N}_{0}^{n}$,
\begin{equation}\tag{5.4}
\lim_{j \to \infty}\left(\sup_{x \in \mathbb{R}^{n}}(1+|x|)^{N}|\partial^{\alpha}\varphi(x) - \partial^{\alpha}\varphi_{j}(x)|\right) = 0
\label{eq:5.4}
\end{equation}
成立, 证明这一点.
\end{exercise}

% ============ 第5章 §5.3 Schwartz 分布 (原书 pp.108-110) ============

\section{Schwartz 分布}\label{sec:5.3}
\index{S@Schwartz 分布(空间)}




Schwartz 函数由性质良好的函数构成, 所以对象狭窄. 为了扩大考虑的对象, 考虑对偶空间是有用的想法. 这里, 将对偶空间作为拓扑对偶线性空间来考虑, 考察其性质.

$\mathcal{S}'(\mathbb{R}^{n})$ 作为拓扑对偶, 即作为 $\mathcal{S}(\mathbb{R}^{n})$ 的连续线性泛函全体构成的集合来定义.

\begin{definition}[Schwartz 分布 (空间)]\label{def:5.2}
$\mathcal{S}'(\mathbb{R}^{n})$ 是从 $\mathcal{S}(\mathbb{R}^{n})$ 到 $\mathbb{C}$ 的连续复线性算子全体构成的集合. 即,
\[
\begin{aligned}
\mathcal{S}'(\mathbb{R}^{n}) \equiv \{f &: \mathcal{S}(\mathbb{R}^{n}) \to \mathbb{C} : \langle f, \varphi + \psi \rangle = \langle f, \varphi \rangle + \langle f, \psi \rangle, \\
&\quad \langle f, a \cdot \varphi \rangle = a \cdot \langle f, \varphi \rangle, f \text{ 连续}\}.
\end{aligned}
\]
这里 a 遍历 $\mathbb{C}$ 全体, $\varphi, \psi$ 遍历 $\mathcal{S}(\mathbb{R}^{n})$ 全体. f 在 $\varphi$ 处的值 $f(\varphi)$ 按惯例记为 $\langle f, \varphi \rangle$.
\end{definition}

``小''集合的对偶一般地``大'', 如下面的例子所示, $\mathcal{S}'(\mathbb{R}^{n})$ 包含通常函数以外的东西, 实际上非常大.

\begin{example}[Dirac 分布 (Dirac delta)]\label{exa:5.8}
对 $a \in \mathbb{R}^{n}$, 代入映射 $\delta_{a}: \varphi \in \mathcal{S}(\mathbb{R}^{n}) \mapsto \varphi(a) \in \mathbb{C}$ 属于 $\mathcal{S}'(\mathbb{R}^{n})$. 实际上, 由于 $|\langle\delta_{a}, \varphi\rangle| = |\varphi(a)| \leq p_{0}(\varphi) = \sup_{x \in \mathbb{R}^{n}}|\varphi(x)|$, 从 $\mathcal{S}(\mathbb{R}^{n})$ 到 $\mathbb{C}$ 的连续性得到保证. $\delta_{a}$ 在物理学等中也是重要的. 当 $a = 0$ 时, 记为 $\delta_{a} = \delta$.
\end{example}

\begin{example}\label{exa:5.9}
对 $\varphi \in \mathcal{S}(\mathbb{R}^{n})$, $f \in L^{p}(\mathbb{R}^{n})$, $1 \leq p \leq \infty$, 由 Hölder 不等式,
\[
\int_{\mathbb{R}^{n}}|f(x)\varphi(x)|\dd x \leq \|f\|_{L^{p}}\|\varphi\|_{L^{p'}} \lesssim \|f\|_{L^{p}}\sup_{x \in \mathbb{R}^{n}}(1+|x|)^{n+1}|\varphi(x)|
\]
所以映射
\[
F_{f}: f \in L^{p}(\mathbb{R}^{n}) \mapsto \int_{\mathbb{R}^{n}}f(x)\varphi(x)\dd x \in \mathbb{C}
\]
是 $\mathcal{S}'(\mathbb{R}^{n})$ 的元. 如此, 通过由 f 给出 $F_{f}$ 生成 $\mathcal{S}'(\mathbb{R}^{n})$ 的元, 但重要的是, 由 Lebesgue 微分定理 (参照注意~\ref{note:3.2}), 这个对应是单射. 实际上, 若 $F_{f} = 0$, 则
\[
\int_{\mathbb{R}^{n}}f(y)\varphi(y)\dd y = F_{f}(\varphi) = 0.
\]
由于 $\varphi \in \mathcal{S}(\mathbb{R}^{n})$ 是任意的, 取 $\varphi_{t}(x - \cdot)$ 作为 $\varphi$,
\[
f * \varphi_{t}(x) = \int_{\mathbb{R}^{n}}f(y)\varphi_{t}(x-y)\dd y = 0.
\]
因此, 由注意~\ref{note:3.2}, 对几乎所有 $x \in \mathbb{R}^{n}$, 得到 $f(x) = \lim_{t \downarrow 0}f * \varphi_{t}(x) = 0$. 由以上的经过, $F_{f}$ 可以记为 f.
\end{example}

虽然不是所有的元都具有 $F_{f}$ 的形式, 但相当多的元可以用 $F_{f}$ 的形式表示, 所以常用与函数 f 相同的字母来表示 $\mathcal{S}'(\mathbb{R}^{n})$ 的元. $\mathcal{S}'(\mathbb{R}^{n})$ 的下面的特征刻画具体且容易理解.

\begin{theorem}[$\mathcal{S}'(\mathbb{R}^{n})$ 的特征刻画]\label{thm:5.3}
设 F 是从 $\mathcal{S}(\mathbb{R}^{n})$ 到 $\mathbb{C}$ 的 $\mathbb{C}$-线性映射. 此时, $F \in \mathcal{S}'(\mathbb{R}^{n})$ 的必要充分条件是, 存在某个 $N \in \mathbb{N}$, 使得对所有 $\varphi \in \mathcal{S}(\mathbb{R}^{n})$,
\begin{equation}\tag{5.5}
|F(\varphi)| \leq Np_{N}(\varphi)
\label{eq:5.5}
\end{equation}
成立.
\end{theorem}

\begin{proof}
式~\eqref{eq:5.5} 是充分条件. 实际上, 当 $\mathcal{S}(\mathbb{R}^{n})$ 的列 $\{\varphi_{j}\}_{j=1}^{\infty}$ 收敛于 $\varphi \in \mathcal{S}(\mathbb{R}^{n})$ 时使用式~\eqref{eq:5.5},
\[
|F(\varphi_{j}) - F(\varphi)| = |F(\varphi_{j} - \varphi)| \leq Np_{N}(\varphi_{j} - \varphi)
\]
成立, 所以当 $j \to \infty$ 时得到 $F(\varphi_{j}) \to F(\varphi)$. 下面证明必要性. 假设 $F: \mathcal{S}(\mathbb{R}^{n}) \to \mathbb{C}$ 连续, 证明满足式~\eqref{eq:5.5} 的 $N \in \mathbb{N}$ 的存在性. 由 $F(0) = 0$ 与 F 的连续性, 注意
\begin{equation}\tag{5.6}
F^{-1}(\{z \in \mathbb{C} : |z| < 1\}) = \{\varphi \in \mathcal{S}(\mathbb{R}^{n}) : |F(\varphi)| < 1\}
\label{eq:5.6}
\end{equation}
是包含 0 的 $\mathcal{S}(\mathbb{R}^{n})$ 的开集. 因此, 由例~\ref{exa:5.6}, 取 $L \in \mathbb{N}$ 足够大, 使得
\[
\{\varphi \in \mathcal{S}(\mathbb{R}^{n}) : Lp_{L}(\varphi) < 1\} \subset \{\varphi \in \mathcal{S}(\mathbb{R}^{n}) : |F(\varphi)| < 1\}.
\]
因此, 若 $\varphi \in \mathcal{S}(\mathbb{R}^{n})$ 满足 $Lp_{2L}(\varphi) = \frac{1}{2}$, 则 $|F(\varphi)| \leq 1$.

那么, 现在任意给定 $\varphi \in \mathcal{S}(\mathbb{R}^{n}) \setminus \{0\}$, $\psi \equiv \frac{1}{2Lp_{2L}(\varphi)}\varphi$ 满足 $Lp_{2L}(\psi) = \frac{1}{2}$, 所以得到 $|F(\psi)| \leq 1$. 由 $\psi$ 的定义式, 对 $\varphi \in \mathcal{S}(\mathbb{R}^{n}) \setminus \{0\}$, $|F(\varphi)| \leq 2Lp_{2L}(\varphi)$ 成立. 该式对 $\varphi = 0$ 也当然成立. 因此, 取 $N = 2L$, 就取到了满足式~\eqref{eq:5.5} 的 N, 所以式~\eqref{eq:5.5} 成立.
\end{proof}

$\mathcal{S}(\mathbb{R}^{n})$ 中引入了距离拓扑, 也给 $\mathcal{S}'(\mathbb{R}^{n})$ 引入拓扑. 一般地, 给定集合 X 的子集族 Z 时, 考虑包含 Z 的开集族全体 V, 其交集
\[
\bigcap_{\mathcal{W} \in \mathcal{V}}\{U \in 2^{X} : U \in \mathcal{W}\}
\]
构成开集族, 回忆这被称为由 Z 生成的拓扑.

\begin{definition}\label{def:5.3}
设 $\mathcal{S}'(\mathbb{R}^{n})$ 中的拓扑是由下面的集合族 Z 生成的拓扑:
\[
\mathcal{Z} = \{g \in \mathcal{S}'(\mathbb{R}^{n}) : |\langle f - g, \varphi\rangle| < \varepsilon\}.
\]
这里考虑满足 $f \in \mathcal{S}'(\mathbb{R}^{n})$, $\varphi \in \mathcal{S}(\mathbb{R}^{n})$, $\varepsilon > 0$ 的所有 f, $\varphi$, $\varepsilon$.
\end{definition}

换句话说, $\mathcal{S}'(\mathbb{R}^{n})$ 的拓扑是使映射 $F \in \mathcal{S}'(\mathbb{R}^{n}) \mapsto F(\varphi) \in \mathbb{C}$ 对所有 $\varphi \in \mathcal{S}(\mathbb{R}^{n})$ 连续的最弱拓扑.

\begin{example}\label{exa:5.10}
对 $f \in \mathcal{S}(\mathbb{R}^{n})$ 的 $\partial_{x_{j}}f$ 也可以在 $\mathcal{S}(\mathbb{R}^{n})$ 的拓扑中考虑. 即, 设 $e_{j}$ 为第 j 个基本向量,
\[
\partial_{x_{j}}f = \lim_{h \to 0}\frac{1}{h}(f(\cdot + h\mathbf{e}_{j}) - f)
\]
在 $\mathcal{S}(\mathbb{R}^{n})$ 的拓扑中成立. 实际上, 需要证明的是对所有 $\alpha, \beta \in \mathbb{N}_{0}^{n}$,
\[
\lim_{h \to 0}p_{\alpha,\beta}\left(\frac{1}{h}(f(\cdot + h\mathbf{e}_{j}) - f) - \partial_{x_{j}}f\right) = 0,
\]
但由于
\begin{align*}
&\lim_{h \to 0}p_{\alpha,\beta}\left(\frac{1}{h}(f(\cdot + h\mathbf{e}_{j}) - f) - \partial_{x_{j}}f\right) \\
&= \lim_{h \to 0}p_{\alpha,0}\left(\frac{1}{h}(\partial^{\beta}f(\cdot + h\mathbf{e}_{j}) - \partial^{\beta}f) - \partial^{\beta}\partial_{x_{j}}f\right)
\end{align*}
所以可以设 $\beta = 0$. 由于 $\beta = 0$, 可以用 $p_{N}$ 代替 $p_{\alpha,0}$ 来考虑. 由于不需要考虑偏导数, 考虑
\[
q_{N}(\varphi) = \sup_{x \in \mathbb{R}^{n}}(1+|x|)^{N}|\varphi(x)|.
\]
此时, 由中值定理
\[
F(t+s) = F(t) + sF'(t) + \frac{s^{2}}{2}F''(t+\theta s), \theta \in (0, 1)
\]
可以取到依赖于 x, h, j 等的 $t \in (0, 1)$, 使得
\[
\frac{1}{h}(f(x + h\mathbf{e}_{j}) - f(x)) - \partial_{x_{j}}f(x) = \frac{h}{2}\partial_{x_{j}}^{2}f(x + ht\mathbf{e}_{j}).
\]
一般地, 对 $y_{1}, y_{2} \in \mathbb{R}^{n}$,
\[
1 + |y_{1} + y_{2}| \leq 1 + |y_{1}| + |y_{2}| \leq (1+|y_{1}|)(1+|y_{2}|)
\]
所以设 $N \in \mathbb{N}$, $|h| < 1$,
\[
(1+|x|)^{N} \leq (1+|x + ht\mathbf{e}_{j}|)^{N}(1+|ht\mathbf{e}_{j}|)^{N} \leq 2^{N}(1+|x + ht\mathbf{e}_{j}|)^{N}.
\]
这里, 右边的不等式中使用了
\[
|h| < 1, t \in (0, 1) \Longrightarrow 1 + |ht\mathbf{e}_{j}| \leq 1 + |h| < 2.
\]
由
\[
q_{N}(\varphi) = \sup_{x \in \mathbb{R}^{n}}(1+|x|)^{N}|\varphi(x)|
\]
可知, 若 $|h| < 1$, 则
\begin{align*}
&q_{N}\left(\frac{1}{h}(f(\cdot + h\mathbf{e}_{j}) - f) - \partial_{x_{j}}f\right) \\
&= \sup_{x \in \mathbb{R}^{n}}(1+|x|)^{N}\left|\frac{1}{h}(f(x + h\mathbf{e}_{j}) - f(x)) - \partial_{x_{j}}f(x)\right| \\
&= \sup_{x \in \mathbb{R}^{n}}(1+|x|)^{N}\left|\frac{h}{2}\partial_{x_{j}}^{2}f(x + ht\mathbf{e}_{j})\right| \\
&\leq 2^{N}\sup_{x \in \mathbb{R}^{n}}(1+|x + ht\mathbf{e}_{j}|)^{N}\left|\frac{h}{2}\partial_{x_{j}}^{2}f(x + ht\mathbf{e}_{j})\right| \\
&= 2^{N-1}|h|\sup_{x \in \mathbb{R}^{n}}(1+|x + ht\mathbf{e}_{j}|)^{N}|\partial_{x_{j}}^{2}f(x + ht\mathbf{e}_{j})| \\
&= 2^{N-1}|h|\sup_{z \in \mathbb{R}^{n}}(1+|z|)^{N}|\partial_{x_{j}}^{2}f(z)|
\end{align*}
即,
\[
q_{N}\left(\frac{1}{h}(\partial^{\beta}f(\cdot + h\mathbf{e}_{j}) - \partial^{\beta}f) - \partial^{\beta}\partial_{x_{j}}f\right) = O(|h|)
\]
证明结束.
\end{example}

\begin{exercise}\label{exer:5.6}
设 $\mathcal{O} = 2^{\mathcal{S}'(\mathbb{R}^{n})}$. 证明: O 是使对所有 $\varphi \in \mathcal{S}(\mathbb{R}^{n})$ 的映射 $F \in \mathcal{S}'(\mathbb{R}^{n}) \mapsto F(\varphi) \in \mathbb{C}$ 连续的 $\mathcal{S}'(\mathbb{R}^{n})$ 中的最强拓扑.
\end{exercise}

% ============ 第5章 §5.4 Fourier 变换 (原书 pp.111-117) ============

\section{Fourier 变换}\label{sec:5.4}
\index{G@Gaussian}

\index{S@Schwartz 分布的 Fourier 变换}

\index{F@Fourier 逆变换}

\index{F@Fourier 变换}

\index{P@Plancherel 定理}

\index{R@Riemann--Lebesgue 定理}














引入 Schwartz 函数 $\mathcal{S}(\mathbb{R}^{n})$ 与 Schwartz 分布 $\mathcal{S}'(\mathbb{R}^{n})$ 的一个动机, 是能够很好地控制 Fourier 变换. 即, Fourier 变换在这些函数空间中是封闭的 (自足). 这里定义 Fourier 变换, 并实际证明 Fourier 变换在 $\mathcal{S}(\mathbb{R}^{n})$, $\mathcal{S}'(\mathbb{R}^{n})$ 中成为同构.

\begin{definition}[Fourier 变换, Fourier 逆变换]\label{def:5.4}
对可积函数 f, 连续函数 $\mathcal{F}f$, $\mathcal{F}^{-1}f$ 定义为
\[
\begin{aligned}
\mathcal{F}f(\xi) &\equiv (2\pi)^{-\frac{n}{2}}\int_{\mathbb{R}^{n}}f(x)e^{-ix \cdot \xi}\dd x \quad (\xi \in \mathbb{R}^{n}), \\
\mathcal{F}^{-1}f(x) &\equiv (2\pi)^{-\frac{n}{2}}\int_{\mathbb{R}^{n}}f(\xi)e^{ix \cdot \xi}\dd \xi \quad (x \in \mathbb{R}^{n}).
\end{aligned}
\]
将函数 f 对应到 $\mathcal{F}f$ 的算子 $\mathcal{F}$ 称为 Fourier 变换, 将函数 f 对应到 $\mathcal{F}^{-1}f$ 的算子 $\mathcal{F}^{-1}$ 称为 Fourier 逆变换.
\end{definition}

$\mathcal{F}f$, $\mathcal{F}^{-1}f$ 实际上成为连续函数这一点将在定理~\ref{thm:5.6} 中证明, 先看几个例子.

\begin{example}\label{exa:5.11}
设 n = 1.

(1) 计算 $\chi_{[-1,1]}$ 的 Fourier 变换. 当 $\xi \neq 0$ 时,
\begin{align*}
\mathcal{F}\chi_{[-1,1]}(\xi) &= \frac{1}{\sqrt{2\pi}}\int_{\mathbb{R}}\chi_{[-1,1]}(x)e^{-ix\xi}\dd x \\
&= \frac{1}{\sqrt{2\pi}}\int_{-1}^{1}e^{-ix\xi}\dd x \\
&= \frac{1}{\sqrt{2\pi}}\left[\frac{e^{-ix\xi}}{-i\xi}\right]_{-1}^{1} \\
&= \frac{2\sin\xi}{\sqrt{2\pi}\xi}.
\end{align*}
当 $\xi = 0$ 时,
\[
\mathcal{F}\chi_{[-1,1]}(\xi) = \frac{1}{\sqrt{2\pi}}\int_{\mathbb{R}}\chi_{[-1,1]}(x)\dd x = \frac{1}{\sqrt{2\pi}}\int_{-1}^{1}\dd x = \sqrt{\frac{2}{\pi}}.
\]
或者, $\xi = 0$ 的情形也可以考虑为 $\xi \neq 0$ 的情形在 $\xi \to 0$ 时的极限.

(2) 计算 Gaussian $E(x) = e^{-\frac{1}{2}x^{2}}$ 的 Fourier 变换. 使用复分析, 但省略与复分析相关的细节. 设 $\xi \in \mathbb{R}$, 由定义,
\[
\mathcal{F}E(\xi) = \frac{1}{\sqrt{2\pi}}\int_{\mathbb{R}}e^{-\frac{1}{2}x^{2}-ix\xi}\dd x = \frac{1}{\sqrt{2\pi}}\int_{\mathbb{R}}e^{-\frac{1}{2}(x-i\xi)^{2}-\frac{1}{2}\xi^{2}}\dd x.
\]
因此, 使用复积分 (由 Cauchy 积分定理),
\[
\mathcal{F}E(\xi) = \frac{1}{\sqrt{2\pi}}\int_{\mathbb{R}}e^{-\frac{1}{2}x^{2}-\frac{1}{2}\xi^{2}}\dd x = E(\xi).
\]
\end{example}

利用可积函数 F 的积分的平移不变性
\[
\int_{\mathbb{R}^{n}}F(x+a)\dd x = \int_{\mathbb{R}^{n}}F(x)\dd x
\]
计算卷积的 Fourier 变换. 通过 Fourier 变换, 卷积被变换为乘积.

\begin{theorem}\label{thm:5.4}
设 $f, g \in L^{1}(\mathbb{R}^{n})$. 此时, 对 $\xi \in \mathbb{R}^{n}$, $\mathcal{F}[f*g](\xi) = (2\pi)^{\frac{n}{2}}\mathcal{F}f(\xi)\mathcal{F}g(\xi)$ 成立.
\end{theorem}

\begin{proof}
左边是
\[
(2\pi)^{-\frac{n}{2}}\int_{\mathbb{R}^{n}}e^{-ix\xi}\left(\int_{\mathbb{R}^{n}}f(y)g(x-y)\dd y\right)\dd x.
\]
确认可以交换积分顺序, 再进一步变形.
\begin{align*}
&\int\int_{\mathbb{R}^{n} \times \mathbb{R}^{n}}|f(x-y)g(y)e^{-i\xi x}|\dd x\dd y \\
&= \int\int_{\mathbb{R}^{n} \times \mathbb{R}^{n}}|f(x-y)g(y)|\dd x\dd y \\
&= \int_{\mathbb{R}^{n}}\left(\int_{\mathbb{R}^{n}}|f(x-y)g(y)|\dd x\right)\dd y \\
&= \int_{\mathbb{R}^{n}}|g(y)|\left(\int_{\mathbb{R}^{n}}|f(x-y)|\dd x\right)\dd y.
\end{align*}
这里, 为了得到第2个等号, 使用了被积函数是正的这一点. 由积分的平移不变性,
\begin{align*}
&\int\int_{\mathbb{R}^{n} \times \mathbb{R}^{n}}|f(x-y)g(y)e^{-i\xi x}|\dd x\dd y \\
&= \int_{\mathbb{R}^{n}}|g(y)|\left(\int_{\mathbb{R}^{n}}|f(x)|\dd x\right)\dd y \\
&= \|f\|_{L^{1}}\|g\|_{L^{1}} < \infty.
\end{align*}
因此, 可以改写为
\begin{align*}
\mathcal{F}[f*g](\xi) &= \frac{1}{(2\pi)^{\frac{n}{2}}}\int_{\mathbb{R}^{n}}\left(\int_{\mathbb{R}^{n}}f(x-y)g(y)e^{-i\xi x}\dd y\right)\dd x \\
&= \frac{1}{(2\pi)^{\frac{n}{2}}}\int_{\mathbb{R}^{n}}\left(\int_{\mathbb{R}^{n}}f(x-y)g(y)e^{-i\xi x}\dd x\right)\dd y.
\end{align*}
作变量替换 x - y = z,
\[
\mathcal{F}[f*g](\xi) = \frac{1}{(2\pi)^{\frac{n}{2}}}\int_{\mathbb{R}^{n}}\left(\int_{\mathbb{R}^{n}}f(z)g(y)e^{-i\xi \cdot (z+y)}\dd z\right)\dd y.
\]
分解累次积分, 就得到 $\mathcal{F}[f*g](\xi) = (2\pi)^{\frac{n}{2}}\mathcal{F}f(\xi)\mathcal{F}g(\xi)$.
\end{proof}

由积分号下求导的公式, 下面的事是显然的. 预定在定理~\ref{thm:5.6} 中说明 $\mathcal{F}$, $\mathcal{F}^{-1}$ 将 $L^{1}(\mathbb{R}^{n})$ 移到 $L^{\infty}(\mathbb{R}^{n})$, 这里确认 $\mathcal{F}$, $\mathcal{F}^{-1}$ 将作为 $L^{1}(\mathbb{R}^{n})$ 的子集的 $\mathcal{S}(\mathbb{R}^{n})$ 移到 $\mathcal{S}(\mathbb{R}^{n})$.

\begin{theorem}[Fourier 变换对运算的变换]\label{thm:5.5}
对 $\varphi \in \mathcal{S}(\mathbb{R}^{n})$, 微分与单项式乘法之间的关系式
\begin{equation}\tag{5.7}
\mathcal{F}\left[\frac{\partial\varphi}{\partial x_{j}}\right](\xi) = i\xi_{j}\mathcal{F}\varphi(\xi),
\label{eq:5.7}
\end{equation}
\begin{equation}\tag{5.8}
\mathcal{F}[x_{j}\varphi](\xi) = i\frac{\partial[\mathcal{F}\varphi]}{\partial\xi_{j}}(\xi)
\label{eq:5.8}
\end{equation}
成立. 并且, 对任意 $\alpha, \beta \in \mathbb{N}_{0}^{n}$, 存在 $C_{\alpha,\beta}, N_{\alpha,\beta} > 0$, 各点估计
\begin{equation}\tag{5.9}
|\xi^{\alpha}\partial^{\beta}\mathcal{F}\varphi(\xi)| \leq C_{\alpha,\beta}p_{N_{\alpha,\beta}}(\varphi)
\label{eq:5.9}
\end{equation}
成立. 特别地, Fourier 变换与 Fourier 逆变换是从 $\mathcal{S}(\mathbb{R}^{n})$ 到自身的连续映射.
\end{theorem}

\begin{proof}
使用定理~\ref{thm:2.13}, 例~\ref{exa:5.4} 与 $\partial_{j}e^{-ix \cdot \xi} = -i\xi_{j}e^{-ix \cdot \xi}$. 省略细节, 作为练习题~\ref{exer:5.9} 处理, 省略证明.
\end{proof}

作为定理~\ref{thm:5.5} 的应用, 证明 Riemann--Lebesgue 定理.

\begin{theorem}[Riemann--Lebesgue 定理]\label{thm:5.6}
设 $f \in L^{1}(\mathbb{R}^{n})$.
(1) $\mathcal{F}f \in L^{\infty}(\mathbb{R}^{n})$, 且 $\|\mathcal{F}f\|_{L^{\infty}} \leq (2\pi)^{-\frac{n}{2}}\|f\|_{L^{1}}$ 成立.
(2) $\mathcal{F}f$ 连续.
(3) $\lim_{\xi \to \infty}\mathcal{F}f(\xi) = 0$ 成立.
\end{theorem}

\begin{proof}
(1) 由积分的三角不等式显然. 实际上, 设 $\xi \in \mathbb{R}^{n}$,
\begin{align*}
|\mathcal{F}f(\xi)| &= \left|(2\pi)^{-\frac{n}{2}}\int_{\mathbb{R}^{n}}f(x)e^{-ix \cdot \xi}\dd x\right| \\
&\leq (2\pi)^{-\frac{n}{2}}\int_{\mathbb{R}^{n}}|f(x)e^{-ix \cdot \xi}|\dd x \\
&= (2\pi)^{-\frac{n}{2}}\|f\|_{L^{1}}
\end{align*}
所以.

(2) 取 $f_{j} \in \mathcal{S}(\mathbb{R}^{n})$ 使得在 $L^{1}(\mathbb{R}^{n})$ 的拓扑中 $f_{j} \to f$. 于是, 当 $j \to \infty$ 时, $\|\mathcal{F}f - \mathcal{F}f_{j}\|_{L^{\infty}} \leq (2\pi)^{-\frac{n}{2}}\|f - f_{j}\|_{L^{1}} \to 0$, 所以 $\mathcal{F}f$ 可以表示为连续函数列 $\{\mathcal{F}f_{j}\}_{j=1}^{\infty}$ 的一致收敛极限. 因此, $\mathcal{F}f$ 连续.

(3) 在上述中, $\mathcal{F}f_{j}(\xi)$ 在 $\xi \to \infty$ 时收敛于 0. 因此, 由刚才证明的一致收敛性, 该性质遗传给函数 $\mathcal{F}f$.
\end{proof}

下面转入 Fourier 逆变换的考察. 虽然简单, 但下面的公式是 Fourier 逆变换的关键.

\begin{theorem}[Fourier 变换的基本公式]\label{thm:5.7}
对 $\varphi, \psi \in \mathcal{S}(\mathbb{R}^{n})$,
\begin{equation}\tag{5.10}
\int_{\mathbb{R}^{n}}\varphi(x)\mathcal{F}\psi(x)\dd x = \int_{\mathbb{R}^{n}}\mathcal{F}\varphi(\xi)\psi(\xi)\dd \xi
\label{eq:5.10}
\end{equation}
成立.
\end{theorem}

\begin{proof}
写下两边, 由 Fubini 定理可知两边一致于 $(2\pi)^{-\frac{n}{2}}\int_{\mathbb{R}^{n}}\int_{\mathbb{R}^{n}}\varphi(x)\psi(\xi)e^{-ix \cdot \xi}\dd x\dd \xi$.
\end{proof}

用定理~\ref{thm:5.7} 证明 $\mathcal{F}$ 与 $\mathcal{F}^{-1}$ 互为逆变换.

\begin{theorem}[Fourier 变换的逆公式]\label{thm:5.8}
对所有 $\varphi \in \mathcal{S}(\mathbb{R}^{n})$, $\mathcal{F}^{-1}[\mathcal{F}\varphi] = \mathcal{F}[\mathcal{F}^{-1}\varphi] = \varphi$ 成立.
\end{theorem}

\begin{proof}
由于 $\mathcal{F}[\mathcal{F}^{-1}\varphi] = \varphi$ 的证明也相同, 只证明 $\mathcal{F}^{-1}[\mathcal{F}\varphi] = \varphi$. 使用例~\ref{exa:5.11} 中考虑的 $E(x) \equiv \exp\left(-\frac{1}{2}|x|^{2}\right)$. 对 $x \in \mathbb{R}^{n}$, $t > 0$, 令 $E_{t}(x) \equiv E(tx)$, $x \in \mathbb{R}^{n}$. 关于一元函数,
\[
\int_{\mathbb{R}}E(\sqrt{2}s)\dd s = \sqrt{\pi}
\]
成立这一点是已知的. 通过变量替换,
\[
\int_{\mathbb{R}}E(x)\dd x = \sqrt{2\pi}
\]
被导出. 多用 Fubini 定理,
\[
\int_{\mathbb{R}^{n}}e^{-\frac{1}{2}|x|^{2}}\dd x = \sqrt{(2\pi)^{n}}.
\]
再次通过变量替换,
\begin{align}
\int_{\mathbb{R}^{n}}t^{-n}E\left(\frac{\xi}{t}\right)\dd \xi &= \int_{\mathbb{R}^{n}}t^{-n}\exp\left(-\frac{|\xi|^{2}}{2t^{2}}\right)\dd \xi \nonumber \\
&= \int_{\mathbb{R}^{n}}\exp\left(-2^{-1}|\xi|^{2}\right)\dd \xi \nonumber \\
&= \sqrt{(2\pi)^{n}}. \tag{5.11}
\label{eq:5.11}
\end{align}
成立.

此时, 由定理~\ref{thm:5.7},
\begin{equation}\tag{5.12}
\int_{\mathbb{R}^{n}}\mathcal{F}E_{t}(\xi)\varphi(\xi)\dd \xi = \int_{\mathbb{R}^{n}}\mathcal{F}\varphi(x)E_{t}(x)\dd x.
\label{eq:5.12}
\end{equation}
Gaussian $E(x) \equiv \exp\left(-\frac{1}{2}|x|^{2}\right)$ 的 Fourier 变换仍是 Gaussian E, 所以作缩放变换得到 $\mathcal{F}E_{t}(\xi) = t^{-n}E\left(\frac{\xi}{t}\right)$. 将其代入式~\eqref{eq:5.12},
\begin{equation}\tag{5.13}
\int_{\mathbb{R}^{n}}t^{-n}E\left(\frac{\xi}{t}\right)\varphi(\xi)\dd \xi = \int_{\mathbb{R}^{n}}\mathcal{F}\varphi(x)E(tx)\dd x.
\label{eq:5.13}
\end{equation}
于是, 减去 $\varphi(0)$, 考察
\begin{align*}
&(2\pi)^{-\frac{n}{2}}\int_{\mathbb{R}^{n}}t^{-n}E\left(\frac{\xi}{t}\right)\varphi(\xi)\dd \xi - \varphi(0) \\
&= (2\pi)^{-\frac{n}{2}}\int_{\mathbb{R}^{n}}t^{-n}E\left(\frac{\xi}{t}\right)(\varphi(\xi) - \varphi(0))\dd \xi \\
&= (2\pi)^{-\frac{n}{2}}\int_{\mathbb{R}^{n}}E(\xi)(\varphi(t\xi) - \varphi(0))\dd \xi.
\end{align*}
由于对 $\xi \in \mathbb{R}^{n}$ 与 $t > 0$,
\[
|E(\xi)(\varphi(t\xi) - \varphi(0))| \leq 2E(\xi)\|\varphi\|_{L^{\infty}}
\]
成立, 由 Lebesgue 收敛定理, 当 $t \downarrow 0$ 时它收敛于 0. 因此, 由 $|E(t\cdot)\mathcal{F}\varphi| \leq |\mathcal{F}\varphi|$,
\begin{equation}\tag{5.14}
\varphi(0) = \lim_{t \downarrow 0}\int_{\mathbb{R}^{n}}\mathcal{F}\varphi(x)E(tx)\dd x = (2\pi)^{-\frac{n}{2}}\int_{\mathbb{R}^{n}}\mathcal{F}\varphi(\xi)\dd \xi
\label{eq:5.14}
\end{equation}
得到. 式~\eqref{eq:5.14} 是对 x = 0 的逆变换公式.

用这个式~\eqref{eq:5.14} 完成定理~\ref{thm:5.8} 的证明. 设 $x \in \mathbb{R}^{n}$. 首先, 用平移, 由式~\eqref{eq:5.14} 得到 $\varphi(x) = (2\pi)^{-\frac{n}{2}}\int_{\mathbb{R}^{n}}\mathcal{F}[\varphi(\cdot + x)](\xi)\dd \xi$. 使用 $\mathcal{F}[\varphi(\cdot + x)](\xi) = \mathcal{F}\varphi(\xi)e^{ix \cdot \xi}$,
\[
\varphi(x) = (2\pi)^{-\frac{n}{2}}\int_{\mathbb{R}^{n}}\mathcal{F}\varphi(\xi)e^{ix \cdot \xi}\dd \xi = \mathcal{F}^{-1}[\mathcal{F}\varphi](x)
\]
得到.
\end{proof}

以定理~\ref{thm:5.7} 为基础定义 Schwartz 分布的 Fourier 变换. 想要定义 $f \in \mathcal{S}'(\mathbb{R}^{n})$ 的 Fourier 变换 $\mathcal{F}f \in \mathcal{S}'(\mathbb{R}^{n})$, 所以重要的是 $\mathcal{F}f \in \mathcal{S}'(\mathbb{R}^{n})$ 将 $\varphi \in \mathcal{S}(\mathbb{R}^{n})$ 映射到什么. 在式~\eqref{eq:5.15} 中将其具体化.

\begin{definition}[Schwartz 分布的 Fourier 变换]\label{def:5.5}
将 $f \in \mathcal{S}'(\mathbb{R}^{n})$ 的 Fourier 变换 $\mathcal{F}f \in \mathcal{S}'(\mathbb{R}^{n})$ 定义为
\begin{equation}\tag{5.15}
\langle\mathcal{F}f, \varphi\rangle \equiv \langle f, \mathcal{F}\varphi\rangle, \varphi \in \mathcal{S}(\mathbb{R}^{n}).
\label{eq:5.15}
\end{equation}
\end{definition}

Fourier 变换是从 $\mathcal{S}'(\mathbb{R}^{n})$ 到自身的连续线性同构映射, 这一点由定义显然. (作为练习题~\ref{exer:5.7}, 省略验证.)

现在, 证明主张 Fourier 变换是 $L^{2}$ 的酉变换的 Plancherel 定理. 该公式是定理~\ref{thm:5.7}, \ref{thm:5.8} 的系.

\begin{theorem}[Plancherel 定理]\label{thm:5.9}
设 $\varphi, \psi \in \mathcal{S}(\mathbb{R}^{n})$, 则
\begin{equation}\tag{5.16}
\int_{\mathbb{R}^{n}}\mathcal{F}\varphi(x)\overline{\mathcal{F}\psi(x)}\dd x = \int_{\mathbb{R}^{n}}\varphi(x)\overline{\psi(x)}\dd x
\label{eq:5.16}
\end{equation}
成立.
\end{theorem}

\begin{proof}
只需证明将式~\eqref{eq:5.16} 右边的 $\psi$ 替换为 $\overline{\mathcal{F}\psi}$ 后得到的等式. Fourier 变换与取共轭的运算通过
\begin{equation}\tag{5.17}
\mathcal{F}[\overline{\mathcal{F}\psi}] = \mathcal{F}\mathcal{F}^{-1}[\overline{\psi}] = \overline{\psi}
\label{eq:5.17}
\end{equation}
相联系. 由定理~\ref{thm:5.7} 与式~\eqref{eq:5.17}, 式~\eqref{eq:5.16} 显然.
\end{proof}

由于 $C_{\mathrm{c}}^{\infty}(\mathbb{R}^{n})$, 从而 $\mathcal{S}(\mathbb{R}^{n})$ 在 $L^{2}(\mathbb{R}^{n})$ 中稠密, 由定理~\ref{thm:2.23} 与定理~\ref{thm:5.9}, $\mathcal{F}$ 被扩张为 $L^{2}(\mathbb{R}^{n})$ 的酉算子. 此时, 定理~\ref{thm:5.4} 对 $f \in L^{1}(\mathbb{R}^{n})$, $g \in L^{2}(\mathbb{R}^{n})$ 成立.

已经说过 Fourier 变换是分析数学变换的便利工具, 这里重新考察 Hilbert 变换.

\begin{example}\label{exa:5.12}
Hilbert 变换是由
\[
Hf(x) = \lim_{\varepsilon \downarrow 0}\int_{|y| > \varepsilon}\frac{f(x-y)}{y}\dd y
\]
给出的变换. 令
\[
H_{\varepsilon}f(x) = \int_{|y| > \varepsilon}\frac{f(x-y)}{y}\dd y,
\]
通过分析这个极限得到了种种结论. 特别是 $L^{2}(\mathbb{R})$ 中的 Hilbert 变换是
\[
L^{2}(\mathbb{R}) \ni f \mapsto H_{\varepsilon}f = \chi_{\mathbb{R}\setminus(-\varepsilon, \varepsilon)}(\cdot)\cdot^{-1} * f \in L^{2}(\mathbb{R})
\]
在 $\varepsilon \to 0$ 时 $L^{2}(\mathbb{R})$ 中的极限, 但设 $f \in C_{\mathrm{c}}^{\infty}(\mathbb{R})$, 则
\begin{align*}
\mathcal{F}[H_{\varepsilon}f] &= \lim_{R \to \infty}\mathcal{F}[\chi_{(-R,R)\setminus(-\varepsilon, \varepsilon)}(\cdot)\cdot^{-1} * f] \\
&= \sqrt{2\pi}\lim_{R \to \infty}\mathcal{F}[\chi_{(-R,R)\setminus(-\varepsilon, \varepsilon)}(\cdot)\cdot^{-1}] \cdot \mathcal{F}f
\end{align*}
可以表示. 这里,
\[
\mathcal{F}[\chi_{(-R,R)\setminus(-\varepsilon, \varepsilon)}(\cdot)\cdot^{-1}](t) = -2i\int_{\varepsilon}^{R}\frac{\sin(ts)}{s}\dd s
\]
但这里出现的右边的积分关于 t, $\varepsilon$, R 有界, 所以 (作为练习题~\ref{exer:5.8}, 省略验证.) 处于可以对
\[
-2i\mathcal{F}f \cdot \int_{\varepsilon}^{R}\frac{\sin(\cdot s)}{s}\dd s
\]
的 $L^{2}$-范数使用 Lebesgue 收敛定理的状况, 其结果是, 可以证明 $H_{\varepsilon}f$ 在 $L^{2}(\mathbb{R})$ 的拓扑中收敛于某个函数 Hf.
\end{example}

\begin{exercise}\label{exer:5.7}
证明 Fourier 变换是从 $\mathcal{S}'(\mathbb{R}^{n})$ 到自身的连续线性同构映射.
\end{exercise}

\begin{exercise}\label{exer:5.8}
求给出 $f(R) = \int_{0}^{R}\frac{\sin s}{s}\dd s$ 的最大值的 R. 并且, 用这一点证明
\[
\sup_{t, \varepsilon, R > 0}\left|\int_{\varepsilon}^{R}\frac{\sin(ts)}{s}\dd s\right| < \infty.
\]
\end{exercise}

\begin{exercise}\label{exer:5.9}
证明定理~\ref{thm:5.5}.
\end{exercise}

\begin{exercise}\label{exer:5.10}
设 O(n) 为 n 次正交矩阵全体构成的集合. 即, 设 E 为单位矩阵, 各 $A \in \mathrm{O}(n)$ 满足 $AA^{t} = A^{t}A = E$. 回答下面的问题.

(1) 对 $f \in L^{1}(\mathbb{R}^{n})$, $A \in \mathrm{O}(n)$, 证明 $\mathcal{F}[f(A\cdot)](\xi) = \mathcal{F}f(A^{-1}\xi)$.

(2) 设函数 u 是在 $\mathbb{R}^{n} \setminus \{0\}$ 上定义的 $C^{\infty}$ 级函数. 假设对 $A \in \mathrm{O}(n)$, $x \in \mathbb{R}^{n} \setminus \{0\}$ 满足 $u(x) = u(Ax)$. 对 $k \in \mathbb{N}$, 函数 $\|\nabla^{k}u\|: \mathbb{R}^{n} \setminus \{0\} \to [0, \infty)$ 定义为
\begin{align*}
\|\nabla^{k}u\| &= \left(\sum_{\alpha \in \{1, 2, \ldots, n\}^{k}}|\partial^{\alpha}u|^{2}\right)^{\frac{1}{2}} \\
&= \left(\sum_{j_{1}=1}^{n}\sum_{j_{2}=1}^{n}\cdots\sum_{j_{k}=1}^{n}\left(\frac{\partial}{\partial x_{j_{1}}}\frac{\partial}{\partial x_{j_{2}}}\cdots\frac{\partial}{\partial x_{j_{k}}}u\right)^{2}\right)^{1/2}.
\end{align*}
此时, 证明对 $A \in \mathrm{O}(n)$ 与 $x \in \mathbb{R}^{n} \setminus \{0\}$ 满足 $\|\nabla^{k}u\|(x) = \|\nabla^{k}u\|(Ax)$.
\end{exercise}

% ============ 第5章 §5.5 Schwartz 分布的微分 (原书 pp.118-119) ============

\section{Schwartz 分布的微分}\label{sec:5.5}
\index{0symbolSdiff@$S'(\mathbb{R}^{n})$ 中的微分}




对函数可以考虑种种微分, 但对 Schwartz 分布基本上考虑偏导数. 交换偏导数的顺序需要某种条件, 但重要的是在 Schwartz 分布的情形可以无条件地进行. 与函数的情形一样, 通过平移取差分再取极限可以得到偏导数. 但是, 这样必须讨论极限的存在性, 而且以差分定义的分布本身在很多情形下找不到那么大的用处, 所以基于 $\mathcal{S}(\mathbb{R}^{n})$ 通过积分
\[
F_{f}: g \in \mathcal{S}(\mathbb{R}^{n}) \mapsto \int_{\mathbb{R}^{n}}f(x)g(x)\dd x \quad (f \in \mathcal{S}(\mathbb{R}^{n}))
\]
可以嵌入 $\mathcal{S}'(\mathbb{R}^{n})$ 这一点, 采用用分部积分定义的方法. 如例~\ref{exa:5.4} 所示, 若 $f, g \in \mathcal{S}(\mathbb{R}^{n})$, 则由分部积分,
\[
\int_{\mathbb{R}^{n}}\partial^{\alpha}f(x) \cdot g(x)\dd x = (-1)^{|\alpha|}\int_{\mathbb{R}^{n}}f(x) \cdot \partial^{\alpha}g(x)\dd x
\]
成立, 这一点是关键.

想要对 $f \in \mathcal{S}'(\mathbb{R}^{n})$ 定义 $\partial^{\alpha}f \in \mathcal{S}'(\mathbb{R}^{n})$, 但由于 $\mathcal{S}'(\mathbb{R}^{n})$ 是从 $\mathcal{S}(\mathbb{R}^{n})$ 到 $\mathbb{C}$ 的映射, 必须考虑 $\partial^{\alpha}f$ 在 $\varphi$ 处的值. 基于这一点, 做下面的定义.

\begin{definition}[$\mathcal{S}'(\mathbb{R}^{n})$ 中的微分]\label{def:5.6}
对 $f \in \mathcal{S}'(\mathbb{R}^{n})$ 与多重指标 $\alpha \in \mathbb{N}_{0}^{n}$, 定义
\[
\langle\partial^{\alpha}f, \varphi\rangle \equiv (-1)^{|\alpha|}\langle f, \partial^{\alpha}\varphi\rangle.
\]
\end{definition}

Schwartz 分布的平移也一样, 由于要对 $f \in \mathcal{S}'(\mathbb{R}^{n})$ 定义 $f(\cdot + h) \in \mathcal{S}'(\mathbb{R}^{n})$, 必须考虑 $f(\cdot + h)$ 在 $\varphi$ 处的值.
\[
\langle f(\cdot + h), \varphi\rangle = \langle f, \varphi(\cdot - h)\rangle \quad (\varphi \in \mathcal{S}(\mathbb{R}^{n}))
\]
这样定义.

分布的微分与通常的微分有很大不同. 实际上, 如下面的定理所示, 在 $\mathcal{S}'(\mathbb{R}^{n})$ 中可以自由交换极限与微分的顺序.

\begin{theorem}[$\mathcal{S}'(\mathbb{R}^{n})$ 中极限与微分顺序的交换]\label{thm:5.10}
设 $\alpha, \beta \in \mathbb{N}_{0}^{n}$, $f \in \mathcal{S}'(\mathbb{R}^{n})$.
(1) 若 $\mathcal{S}'(\mathbb{R}^{n})$ 的点列 $\{f_{j}\}_{j=1}^{\infty}$ 收敛于 f, 则 $\lim_{j \to \infty}\partial^{\alpha}f_{j} = \partial^{\alpha}f$ 成立.
(2) $\partial^{\alpha}(\partial^{\beta}f) = \partial^{\beta}(\partial^{\alpha}f) = \partial^{\alpha+\beta}f$ 成立.
\end{theorem}

\begin{proof}
(1) 当 $j \to \infty$ 时 $f_{j} \to f$ 在 $\mathcal{S}'(\mathbb{R}^{n})$ 的意义下成立, 与对任意 $\varphi \in \mathcal{S}(\mathbb{R}^{n})$, $\langle f_{j}, \varphi\rangle \to \langle f, \varphi\rangle$ 成立是等价的. 因此, 只需证明 $\langle\partial^{\alpha}f_{j}, \varphi\rangle \to \langle\partial^{\alpha}f, \varphi\rangle$, 而这与 $(-1)^{|\alpha|}\langle f_{j}, \partial^{\alpha}\varphi\rangle \to (-1)^{|\alpha|}\langle f, \partial^{\alpha}\varphi\rangle$ 同义.

(2) 看 $\varphi \in \mathcal{S}(\mathbb{R}^{n})$ 处的值. 最左边是
\begin{align*}
\langle\partial^{\alpha}(\partial^{\beta}f), \varphi\rangle &= (-1)^{|\alpha|}\langle\partial^{\beta}f, \partial^{\alpha}\varphi\rangle \\
&= (-1)^{|\alpha+\beta|}\langle f, \partial^{\beta}\partial^{\alpha}\varphi\rangle \\
&= (-1)^{|\alpha+\beta|}\langle f, \partial^{\alpha+\beta}\varphi\rangle.
\end{align*}
$\langle\partial^{\beta}(\partial^{\alpha}f), \varphi\rangle$ 也变形为同样的式子, 所以得到结论.
\end{proof}

分布的微分的定义可以如下改写.

\begin{theorem}\label{thm:5.11}
对 $f \in \mathcal{S}'(\mathbb{R}^{n})$ 的 $\partial_{x_{j}}f$ 是古典定义的自然扩张. 即, 设 $e_{j}$ 为第 j 个基本向量,
\[
\partial_{x_{j}}f = \lim_{h \to 0}\frac{1}{h}(f(\cdot + h\mathbf{e}_{j}) - f)
\]
在 $\mathcal{S}'(\mathbb{R}^{n})$ 的拓扑中成立.
\end{theorem}

\begin{proof}
对 $\varphi \in \mathcal{S}(\mathbb{R}^{n})$, 由例~\ref{exa:5.10},
\begin{align*}
\left\langle\frac{1}{h}(f(\cdot + h\mathbf{e}_{j}) - f), \varphi\right\rangle &= \left\langle f, \frac{1}{h}(\varphi(\cdot - h\mathbf{e}_{j}) - \varphi)\right\rangle \\
&\to -\langle f, \partial_{x_{j}}\varphi\rangle = \langle\partial_{x_{j}}f, \varphi\rangle
\end{align*}
成立, 所以得到结论.
\end{proof}

\begin{exercise}\label{exer:5.11}
证明 $f(t) = e^{t}\sin(e^{t})$ 可以看作 $\mathcal{S}'(\mathbb{R})$ 的元. 【提示】找出满足
\[
\left|\int_{\mathbb{R}}f(t)\varphi(t)\dd t\right| \leq Cp_{2,1}(\varphi) \quad (\varphi \in C_{\mathrm{c}}^{\infty}(\mathbb{R}))
\]
的 $C > 0$.
\end{exercise}

% ============ 第5章 §5.6 Schwartz 分布的卷积 (原书 pp.120-122) ============

\section{Schwartz 分布的卷积}\label{sec:5.6}
\index{0symbolSconv@$S(\mathbb{R}^{n})$ 与 $S'(\mathbb{R}^{n})$ 的卷积}




作为属于 $\mathcal{S}(\mathbb{R}^{n})$ 的函数之间卷积的定义的动机, 如下定义 $f \in \mathcal{S}(\mathbb{R}^{n})$ 与 $g \in \mathcal{S}'(\mathbb{R}^{n})$ 的卷积.

\begin{definition}[$\mathcal{S}(\mathbb{R}^{n})$ 与 $\mathcal{S}'(\mathbb{R}^{n})$ 的卷积]\label{def:5.7}
对 $f \in \mathcal{S}(\mathbb{R}^{n})$ 与 $g \in \mathcal{S}'(\mathbb{R}^{n})$, 函数 $f * g$ 定义为
\[
f * g(x) \equiv \langle g, f(x - \cdot)\rangle \quad (x \in \mathbb{R}^{n}).
\]
\end{definition}

在定义~\ref{def:5.7} 中, 即使 $g \in \mathcal{S}'(\mathbb{R}^{n})$ 也能取到值 $f * g(x)$, 这是重要的点. 对 $f \in L^{1}(\mathbb{R}^{n})$ 的 $F_{f}$ 在例~\ref{exa:5.9} 中给出. 汇总一下这个运算的基本性质.

\begin{theorem}[卷积的基本性质]\label{thm:5.12}
设 $f, g \in \mathcal{S}(\mathbb{R}^{n})$, $h \in \mathcal{S}'(\mathbb{R}^{n})$.
(1) 对 $x \in \mathbb{R}^{n}$, $L^{1}$ 函数 f 与 $L^{1}$ 函数 g 的卷积 $f * g(x) = \int_{\mathbb{R}^{n}}f(x-y)g(y)\dd y$ 与定义~\ref{def:5.7} 的 $F_{f} * g(x)$ 一致.
(2) $f * h \in C^{\infty}(\mathbb{R}^{n})$. 并且, 对 $\alpha \in \mathbb{N}_{0}^{n}$,
\begin{equation}\tag{5.18}
\partial^{\alpha}[f*h] = (\partial^{\alpha}f)*h = f*\partial^{\alpha}h
\label{eq:5.18}
\end{equation}
成立. 又, 对所有 $\alpha \in \mathbb{N}_{0}^{n}$, 存在 $N = N_{\alpha} \in \mathbb{N}_{0}$, 使得
\begin{equation}\tag{5.19}
\sup_{x \in \mathbb{R}^{n}}|\partial^{\alpha}[f*h](x)| \lesssim (1+|x|)^{N}
\label{eq:5.19}
\end{equation}
成立. 因此, 特别地, $f*h = F_{f*h} \in \mathcal{S}'(\mathbb{R}^{n})$ 可以这样看待.
(3) 结合律 $(f*g)*h = f*(g*h)$ 成立.
\end{theorem}

\begin{proof}
(1) 由 $f \in L^{1}(\mathbb{R}^{n}) \mapsto F_{f} \in \mathcal{S}'(\mathbb{R}^{n})$ 与卷积的定义显然.

(2) 只要证明
\begin{equation}\tag{5.20}
\partial^{\alpha}[f*h] = f*\partial^{\alpha}h, |\alpha| = 1
\label{eq:5.20}
\end{equation}
就可以反复使用定义~\ref{def:5.6} 与式~\eqref{eq:5.20} 来证明, 所以证明式~\eqref{eq:5.20} 就足够了. $f*h \in C^{\infty}(\mathbb{R}^{n})$ 的证明使用例~\ref{exa:5.10}. 若 $|\alpha| = 1$, 即存在 $j = 1, 2, \ldots$ 使得 $\partial^{\alpha} = \partial_{x_{j}}$, 则式~\eqref{eq:5.18} 由分布的微分表示为差分这一点是显然的. 如上所述, $|\alpha| \geq 2$ 的情形可以用数学归纳法证明式~\eqref{eq:5.18}.

为了证明不等式~\eqref{eq:5.19}, 由定理~\ref{thm:5.3} 首先注意存在自然数 N, 使得对所有 $\varphi \in \mathcal{S}(\mathbb{R}^{n})$,
\begin{equation}\tag{5.21}
|\langle h, \varphi\rangle| \leq Np_{N}(\varphi)
\label{eq:5.21}
\end{equation}
成立. 在式~\eqref{eq:5.21} 中取 $f(x - \cdot)$ 作为 $\varphi$. 由 $p_{N}$ 的定义, 通过简单计算可知 $p_{N}(\partial^{\alpha}f(x - \cdot)) \lesssim (1+|x|)^{N}$, 所以 $|\partial^{\alpha}[f*h](x)| \lesssim (1+|x|)^{N}$ 随之成立.

(3) 由于
\[
f*g = \sum_{k \in \mathbb{Z}^{n}}\int_{k+[0,1]^{n}}g(\cdot - y)f(y)\dd y
\]
在 $\mathcal{S}(\mathbb{R}^{n})$ 的拓扑中成立, 固定 $x \in \mathbb{R}^{n}$ 后,
\[
(f*g)*h(x) = \sum_{k \in \mathbb{Z}^{n}}\left\langle h, \int_{k+[0,1]^{n}}g(x - \cdot - y)f(y)\dd y\right\rangle.
\]
由引理~\ref{lem:5.2}, 将积分改写为 Riemann 和风格的操作在 $\mathcal{S}(\mathbb{R}^{n})$ 的拓扑中也有效, 所以
\[
(f*g)*h(x) = \sum_{k \in \mathbb{Z}^{n}}\left(\int_{k+[0,1]^{n}}f(y)g*h(x-y)\dd y\right).
\]
由于被积函数 $f(\cdot)g*h(x - \cdot)$ 在 $\mathbb{R}^{n}$ 上可积,
\[
(f*g)*h(x) = \int_{\mathbb{R}^{n}}f(y)g*h(x-y)\dd y.
\]
由以上, 证明了 $f*(g*h) = (f*g)*h$.
\end{proof}

\begin{exercise}\label{exer:5.12}
当 $f \in \mathcal{S}(\mathbb{R}^{n})$, $g \in \mathcal{S}'(\mathbb{R}^{n})$ 时, 证明 $\mathcal{F}[f*g] = (2\pi)^{\frac{n}{2}}\mathcal{F}g \cdot \mathcal{F}f$. 但是, 对 $F \in \mathcal{S}'(\mathbb{R}^{n})$, $\Phi \in \mathcal{S}(\mathbb{R}^{n})$, 通过 $\langle F \cdot \Phi, \varphi\rangle = \langle F, \Phi\varphi\rangle$ 定义 $F \cdot \Phi \in \mathcal{S}'(\mathbb{R}^{n})$.
\end{exercise}


% 第6章 Littlewood-Paley 理论 (原书 pp.123-132, PDF pp.129-138) + 结束语
\chapter{Littlewood--Paley 理论}
% ============ 第6章 §6.1 Littlewood--Paley 理论 (原书 pp.123-130) ============

这里, 作为奇异积分算子有界性理论的应用, 说明 $L^{p}$ 范数的有用改写是可能的. 这一应用称为 Littlewood--Paley 理论, 由此可以考察种种算子的性质.

\section{Littlewood--Paley 理论}\label{sec:6.1}

对 $\varphi \in \mathcal{S}(\mathbb{R}^{n})$ 与 $f \in \mathcal{S}'(\mathbb{R}^{n})$, 令 $\varphi(D)f = \mathcal{F}^{-1}[\varphi\mathcal{F}f] \in C^{\infty}(\mathbb{R}^{n})$. 这与定义为 $\varphi(D)f = (2\pi)^{-\frac{n}{2}}\mathcal{F}^{-1}\varphi * f$ 是同一回事. 另外, 对 $\varphi \in \mathcal{S}(\mathbb{R}^{n})$ 与 $j \in \mathbb{Z}$, 令 $\varphi_{j} = \varphi(2^{-j}\,\cdot\,)$.

\begin{theorem}\label{thm:6.1}
设 $1 < p < \infty$. 设只依赖于 $|x|$ 的实值函数 $\varphi \in \mathcal{S}(\mathbb{R}^{n})$ 满足 $\supp(\varphi) \subset B(4) \setminus B(1)$ 与 $\sum_{j=-\infty}^{\infty}\varphi_{j}^{2} = \chi_{\mathbb{R}^{n}\setminus\{0\}}$. 此时, 若 $f \in L^{p}(\mathbb{R}^{n})$, 则
\[
\|f\|_{L^{p}} \sim \left\|\left(\sum_{j=-\infty}^{\infty}|\varphi_{j}(D)f|^{2}\right)^{1/2}\right\|_{L^{p}}
\]
成立.
\end{theorem}

\begin{note}\label{note:6.1}
考察 $\sum_{j=-\infty}^{\infty}\varphi_{j}^{2} = \chi_{\mathbb{R}^{n}\setminus\{0\}}$ 这个条件.

(1) 由集合特征函数的定义可知, 这个条件在 $x \neq 0$ 时, 可读作
\[
\sum_{j=-\infty}^{\infty}\varphi_{j}(x)^{2} = 1,
\]
在 $x = 0$ 时, 可读作
\[
\sum_{j=-\infty}^{\infty}\varphi_{j}(x)^{2} = 0.
\]

(2) 因为 $\varphi(0) = 0$, 所以 $\varphi_{j}(0) = 0$. 因此, (1) 的 $x = 0$ 情形的条件是平凡的.

(3) 想考察是否真的存在这样的函数. 如从在 $t > 0$ 时, 用次数低于 $k$ 次的多项式 $p_{k}$ 可以表示为
\begin{equation}\tag{6.1}
\varphi^{(k)}(t) = \frac{p_{k}(t)}{t^{2k}}e^{-1/t}
\label{eq:6.1}
\end{equation}
的形式可以看出, 由 $\lim_{t \to \infty}t^{k}e^{-t} = 0$ ($k = 0, 1, 2, \ldots$), 对 $t \leq 0$ 满足 $\varphi(t) = 0$, 对 $t > 0$ 满足 $\varphi(t) = e^{-1/t}$ 的函数 $\varphi$ 是 $C^{\infty}$ 级函数. 也就是说, 无论微分多少次都可以. 因此,
\[
\Phi(x) = \frac{\varphi(1 - |x|^{2})}{\varphi(1 - |x|^{2}) + \varphi(|x|^{2})}
\]
是在 $|x| < 1$ 取正值, 在 $|x| \geq 1$ 取 $0$ 的, 定义在全空间 $\mathbb{R}^{n}$ 上的 $C_{c}^{\infty}$ 级函数. 另外,
\[
\Psi(x) = \frac{\varphi(9|x|^{2} - 1)}{\varphi(1 - |x|^{2}) + \varphi(9|x|^{2} - 1)}
\]
是在 $|x| > \frac{1}{3}$ 取正值, 在 $|x| \leq \frac{1}{3}$ 取 $0$ 的, 定义在全空间 $\mathbb{R}^{n}$ 上的 $C^{\infty}$ 级函数. 因此,
\[
\Theta(x) = \Phi(x)\Psi(x) \quad (x \in \mathbb{R}^{n})
\]
是仅在 $\frac{1}{3} < |x| < 1$ 取正值的, 定义在全空间 $\mathbb{R}^{n}$ 上的 $C^{\infty}$ 级函数. 于是, 对 $x \in \mathbb{R}^{n}$, 令
\[
\zeta(x) = \sqrt{\sum_{j=-\infty}^{\infty}\Theta(2^{-j}x)^{2}},
\]
若 $x \neq 0$, 则
\[
\varphi(x) = \frac{1}{\zeta(x)}\Theta\left(\frac{2}{7}x\right)
\]
通过补充定义 $\varphi(0) = 0$, 成为 $C^{\infty}$ 级函数, 支集包含在 $B(4) \setminus B(1)$ 中, 且满足 $\sum_{j=-\infty}^{\infty}\varphi_{j}^{2} = \chi_{\mathbb{R}^{n}\setminus\{0\}}$.
\end{note}

$p = 2$ 的情形可以由 Fourier 变换的性质得到结论.

\begin{lemma}\label{lem:6.1}
若 $p = 2$, 则定理~\ref{thm:6.1} 成立. 更详细地, 对一切 $f \in L^{2}(\mathbb{R}^{n})$, 有
\begin{equation}\tag{6.2}
\left\|\left(\sum_{j=-\infty}^{\infty}|\varphi_{j}(D)f|^{2}\right)^{1/2}\right\|_{L^{2}} = \|f\|_{L^{2}}
\label{eq:6.2}
\end{equation}
成立.
\end{lemma}

\begin{proof}
由单调收敛定理,
\[
\left\|\left(\sum_{j=-\infty}^{\infty}|\varphi_{j}(D)f|^{2}\right)^{1/2}\right\|_{L^{2}} = \left(\sum_{j=-\infty}^{\infty}(\|\varphi_{j}(D)f\|_{L^{2}})^{2}\right)^{1/2}.
\]
由于 $\mathcal{F}^{-1}$ 由 Plancherel 定理保持 $L^{2}$ 范数,
\[
\left\|\left(\sum_{j=-\infty}^{\infty}|\varphi_{j}(D)f|^{2}\right)^{1/2}\right\|_{L^{2}} = \left(\sum_{j=-\infty}^{\infty}(\|\varphi_{j}\mathcal{F}f\|_{L^{2}})^{2}\right)^{1/2}.
\]
再次由单调收敛定理,
\[
\left\|\left(\sum_{j=-\infty}^{\infty}|\varphi_{j}(D)f|^{2}\right)^{1/2}\right\|_{L^{2}} = \left\|\left(\sum_{j=-\infty}^{\infty}|\varphi_{j}\mathcal{F}f|^{2}\right)^{1/2}\right\|_{L^{2}}.
\]
这里使用 $\sum_{j=-\infty}^{\infty}\varphi_{j}^{2} = \chi_{\mathbb{R}^{n}\setminus\{0\}}$, 得
\[
\left\|\left(\sum_{j=-\infty}^{\infty}|\varphi_{j}(D)f|^{2}\right)^{1/2}\right\|_{L^{2}} = \|\mathcal{F}f\|_{L^{2}}.
\]
因此, 再次由 Plancherel 定理, 由于 $\mathcal{F}$ 保持 $L^{2}$ 范数, 得到式~\eqref{eq:6.2}.
\end{proof}

证明 $p \neq 2$ 的情形时, 下面的引理很管用.

\begin{lemma}\label{lem:6.2}
对 $a > 0$, 由
\[
\Xi(a) \equiv \sum_{j=-\infty}^{\infty}(2^{j}a)^{n}(1 + 2^{j}a)^{-n-1}
\]
给出的函数 $\Xi: (0,\infty) \to \mathbb{R}$ 是有界的.
\end{lemma}

\begin{proof}
假设 $1 < a \leq 2$, 则
\begin{align*}
\Xi(a) &\leq 2^{n}\sum_{j=-\infty}^{\infty}2^{jn}(1 + 2^{j})^{-n-1} \\
&\leq 2^{n}\sum_{j=0}^{\infty}2^{jn}\cdot(2^{j})^{-n-1} + 2^{n}\sum_{j=-\infty}^{-1}2^{jn} < \infty.
\end{align*}
由于 $\Xi(2a) = \Xi(a)$, $\Xi(0) = 0$, 所以对一切 $a \geq 0$, $\Xi(a) \lesssim 1$ 成立.
\end{proof}

\begin{proof}
来证明 $p \neq 2$ 的情形下的定理~\ref{thm:6.1}. 首先证明式~\eqref{eq:6.3}:
\begin{equation}\tag{6.3}
\left\|\left(\sum_{j=-\infty}^{\infty}|\varphi_{j}(D)f|^{2}\right)^{1/2}\right\|_{L^{p}} \lesssim \|f\|_{L^{p}}
\label{eq:6.3}
\end{equation}
关于 $j$ 适当地截断, 之后用 Fatou 引理令 $J \to \infty$, 从而只要证明与自然数 $J$ 无关的一致估计
\[
\left\|\left(\sum_{j=-J}^{J}|\varphi_{j}(D)f|^{2}\right)^{1/2}\right\|_{L^{p}} \lesssim \|f\|_{L^{p}}
\]
即可.

设 $\{r_{j}\}_{j=1}^{\infty}$ 为先前考虑过的 Rademacher 序列. 将 $\{r_{j}\}_{j=1}^{\infty}$ 不重复地重新排列, 使得 $\{\tilde{r}_{j}\}_{j=-\infty}^{\infty} = \{r_{j}\}_{j=1}^{\infty}$. 由此,
\[
\left\|\left(\sum_{j=-J}^{J}|\varphi_{j}(D)f|^{2}\right)^{1/2}\right\|_{L^{p}} \lesssim \left\|\sum_{j=-J}^{J}\tilde{r}_{j} \otimes \varphi_{j}(D)f\right\|_{L^{p}([0,1]_{t} \times \mathbb{R}^{n}_{x})}.
\]
所以, 只要证明固定 $t \in [0,1]$ 时, 算子
\[
T_{t}: f \mapsto \sum_{j=-J}^{J}\tilde{r}_{j}(t)\mathcal{F}^{-1}\varphi_{j} * f
\]
是 $L^{p}$ 有界的, 且其算子范数可用与 $t$ 无关的数来估计即可. 对应的积分核 $K_{t}$ 具体呈如下形状.
\[
K_{t}(z) = \sum_{j=-J}^{J}\tilde{r}_{j}(t)\mathcal{F}^{-1}\varphi_{j}(z).
\]

(1) 关于 $T_{t}$ 的 $B(L^{2}(\mathbb{R}^{n}))$ 范数 $\|T_{t}\|_{L^{2} \to L^{2}}$, 由
\[
A \equiv \sup_{\xi \in \mathbb{R}^{n}}\sum_{j=-\infty}^{\infty}|\varphi_{j}(\xi)| \lesssim 1
\]
与 Plancherel 定理, 可得 $\|T_{t}\|_{L^{2} \to L^{2}} \leq A$.

(2) 因为 $\mathcal{F}^{-1}\varphi_{j} = 2^{jn}\mathcal{F}\varphi(2^{j}\,\cdot\,)$, 所以使用引理~\ref{lem:6.2} 的函数 $\Xi$,
\[
\begin{aligned}
|x-y|^{n}|K_{t}(z)| &\leq \sum_{j=-\infty}^{\infty}|x-y|^{n}|\mathcal{F}^{-1}\varphi_{j}(x-y)| \\
&\lesssim \sum_{j=-\infty}^{\infty}(2^{j}|x-y|)^{n}(1 + 2^{j}|x-y|)^{-n-1} \\
&= \Xi(|x-y|).
\end{aligned}
\]
因此, 由引理~\ref{lem:6.2}, 可得
\[
|x-y|^{n}|K_{t}(z)| \lesssim 1.
\]

(3) 因为 $\nabla K_{t}(z) = 2^{j(n+1)}(\nabla\mathcal{F}^{-1}\varphi)(2^{j}\,\cdot\,)$, 所以使用 $|\nabla\mathcal{F}^{-1}\varphi(z)| \lesssim (1+|z|)^{-n-2}$, 与刚才一样, 可得
\[
\begin{aligned}
|x-y|^{n+1}|\nabla K_{t}(z)| &\leq \sum_{j=-\infty}^{\infty}|x-y|^{n+1}|\nabla(\mathcal{F}^{-1}\varphi_{j})(x-y)| \\
&\lesssim \sum_{j=-\infty}^{\infty}(2^{j}|x-y|)^{n+1}(1 + 2^{j}|x-y|)^{-n-2} \\
&\lesssim 1.
\end{aligned}
\]

由以上可知, $T_{t}$ 是一致地 $L^{p}$ 有界的.

来证明反向的不等式
\[
\left\|\left(\sum_{j=-\infty}^{\infty}|\varphi_{j}(D)f|^{2}\right)^{1/2}\right\|_{L^{p}} \gtrsim \|f\|_{L^{p}}.
\]
暂时假设 $f \in L^{2}(\mathbb{R}^{n}) \cap L^{p}(\mathbb{R}^{n})$. 由对偶性,
\[
\|f\|_{L^{p}} = \sup\left\{\left|\int_{\mathbb{R}^{n}}f(x)g(x)\dd x\right|: g \in L^{p'}(\mathbb{R}^{n}), \|g\|_{L^{p'}} \leq 1\right\}.
\]
因为 $\mathcal{S}(\mathbb{R}^{n})$ 在 $L^{p}(\mathbb{R}^{n})$ 中稠密,
\[
\|f\|_{L^{p}} = \sup\left\{\left|\int_{\mathbb{R}^{n}}f(x)g(x)\dd x\right|: g \in \mathcal{S}(\mathbb{R}^{n}), \|g\|_{L^{p'}} \leq 1\right\}.
\]
基于这一点, 固定 $g \in \mathcal{S}(\mathbb{R}^{n})$. 因为 $f, g \in L^{2}(\mathbb{R}^{n})$, 由引理~\ref{lem:6.1},
\[
\begin{aligned}
\int_{\mathbb{R}^{n}}\sum_{j=-\infty}^{\infty}|\varphi_{j}(D)f(x)\varphi_{j}(D)g(x)|\dd x &\leq \sqrt{\int_{\mathbb{R}^{n}}\sum_{j=-\infty}^{\infty}|\varphi_{j}(D)f(x)|^{2}\dd x\int_{\mathbb{R}^{n}}\sum_{j=-\infty}^{\infty}|\varphi_{j}(D)g(x)|^{2}\dd x} \\
&\leq \|f\|_{L^{2}}\|g\|_{L^{2}} < \infty.
\end{aligned}
\]
因此, 由
\[
\sum_{j=-\infty}^{\infty}\varphi_{j}^{2} = \chi_{\mathbb{R}^{n}\setminus\{0\}},
\]
可变形为
\[
\begin{aligned}
\int_{\mathbb{R}^{n}}f(x)\overline{g(x)}\dd x &= \int_{\mathbb{R}^{n}}\mathcal{F}f(\xi)\overline{\mathcal{F}g(\xi)}\dd\xi \\
&= \sum_{j=-\infty}^{\infty}\int_{\mathbb{R}^{n}}\varphi_{j}(\xi)\mathcal{F}f(\xi)\overline{\varphi_{j}(\xi)\mathcal{F}g(\xi)}\dd\xi \\
&= \sum_{j=-\infty}^{\infty}\int_{\mathbb{R}^{n}}\varphi_{j}(D)f(x)\overline{\varphi_{j}(D)g(x)}\dd x \\
&= \int_{\mathbb{R}^{n}}\left(\sum_{j=-\infty}^{\infty}\varphi_{j}(D)f(x)\overline{\varphi_{j}(D)g(x)}\right)\dd x.
\end{aligned}
\]
这里, 由 Cauchy--Schwarz 不等式,
\[
\begin{aligned}
\left|\int_{\mathbb{R}^{n}}f(x)\overline{g(x)}\dd x\right| &\leq \int_{\mathbb{R}^{n}}\sum_{j=-\infty}^{\infty}|\varphi_{j}(D)f(x)\varphi_{j}(D)g(x)|\dd x \\
&\leq \int_{\mathbb{R}^{n}}\left(\sum_{j=-\infty}^{\infty}|\varphi_{j}(D)f(x)|^{2}\right)^{1/2}\left(\sum_{j=-\infty}^{\infty}|\varphi_{j}(D)g(x)|^{2}\right)^{1/2}\dd x.
\end{aligned}
\]
进而, 通过使用 Hölder 不等式, 可得
\[
\begin{aligned}
\left|\int_{\mathbb{R}^{n}}f(x)\overline{g(x)}\dd x\right| &\leq \left\|\left(\sum_{j=-\infty}^{\infty}|\varphi_{j}(D)f|^{2}\right)^{1/2}\right\|_{L^{p}}\left\|\left(\sum_{j=-\infty}^{\infty}|\varphi_{j}(D)g|^{2}\right)^{1/2}\right\|_{L^{p'}}.
\end{aligned}
\]
这里, 因为 $g \in L^{p'}(\mathbb{R}^{n})$,
\[
\left\|\left(\sum_{j=-\infty}^{\infty}|\varphi_{j}(D)g|^{2}\right)^{1/2}\right\|_{L^{p'}} \lesssim \|g\|_{L^{p'}} \leq 1
\]
可得. 因此, 若 $f \in L^{2}(\mathbb{R}^{n}) \cap L^{p}(\mathbb{R}^{n})$, 则
\begin{equation}\tag{6.4}
\left\|\left(\sum_{j=-\infty}^{\infty}|\varphi_{j}(D)f|^{2}\right)^{1/2}\right\|_{L^{p}} \gtrsim \|f\|_{L^{p}}
\label{eq:6.4}
\end{equation}
就得到证明.

最后, 考虑把假设 $f \in L^{2}(\mathbb{R}^{n}) \cap L^{p}(\mathbb{R}^{n})$ 放宽为 $f \in L^{p}(\mathbb{R}^{n})$. 设 $f \in L^{p}(\mathbb{R}^{n})$. 因为 $L^{2}(\mathbb{R}^{n}) \cap L^{p}(\mathbb{R}^{n})$ 在 $L^{p}(\mathbb{R}^{n})$ 中稠密, 可以取列 $f_{k} \in L^{2}(\mathbb{R}^{n}) \cap L^{p}(\mathbb{R}^{n})$, $k = 1, 2, \ldots$, 使得 $\lim_{k \to \infty}f_{k} = f$ 在 $L^{p}(\mathbb{R}^{n})$ 中成立. 由三角不等式与式~\eqref{eq:6.4},
\[
\begin{aligned}
\|f\|_{L^{p}} &\leq \|f - f_{k}\|_{L^{p}} + \|f_{k}\|_{L^{p}} \\
&\lesssim \|f - f_{k}\|_{L^{p}} + \left\|\left(\sum_{j=-\infty}^{\infty}|\varphi_{j}(D)f_{k}|^{2}\right)^{1/2}\right\|_{L^{p}}.
\end{aligned}
\]
进而, 由三角不等式与式~\eqref{eq:6.3},
\[
\begin{aligned}
\|f\|_{L^{p}} &\leq \|f - f_{k}\|_{L^{p}} + \left\|\left(\sum_{j=-\infty}^{\infty}|\varphi_{j}(D)[f_{k} - f]|^{2}\right)^{1/2}\right\|_{L^{p}} \\
&\quad + \left\|\left(\sum_{j=-\infty}^{\infty}|\varphi_{j}(D)f|^{2}\right)^{1/2}\right\|_{L^{p}} \\
&\lesssim \|f - f_{k}\|_{L^{p}} + \left\|\left(\sum_{j=-\infty}^{\infty}|\varphi_{j}(D)f|^{2}\right)^{1/2}\right\|_{L^{p}},
\end{aligned}
\]
所以令 $k \to \infty$, 就得到关于 $f$ 的结论式
\[
\left\|\left(\sum_{j=-\infty}^{\infty}|\varphi_{j}(D)f|^{2}\right)^{1/2}\right\|_{L^{p}} \gtrsim \|f\|_{L^{p}}.
\]
\end{proof}

作为 Littlewood--Paley 理论的简单应用, 给出下面这一结果. 用
\[
\mathcal{S}_{\infty}(\mathbb{R}^{n}) = \bigcap_{\alpha \in \mathbb{N}_{0}^{n}}\left\{f \in \mathcal{S}(\mathbb{R}^{n}): \int_{\mathbb{R}^{n}}x^{\alpha}f(x)\dd x = 0\right\}
\]
定义 $\mathcal{S}_{\infty}(\mathbb{R}^{n})$. $\mathcal{S}_{\infty}(\mathbb{R}^{n})$ 的元素有时也称为 Lizorkin 分布.

\begin{corollary}\label{cor:6.1}
若 $1 < p < \infty$, 则 $\mathcal{S}_{\infty}(\mathbb{R}^{n})$ 在 $L^{p}(\mathbb{R}^{n})$ 中稠密.
\end{corollary}

\begin{proof}
已知 $C_{c}^{\infty}(\mathbb{R}^{n})$ 在 $L^{p}(\mathbb{R}^{n})$ 中稠密, 所以只要证明任意 $f \in C_{c}^{\infty}(\mathbb{R}^{n})$ 可以用 $\mathcal{S}_{\infty}(\mathbb{R}^{n})$ 的拓扑来近似即可. 设 $\psi \in C^{\infty}(\mathbb{R}^{n})$ 满足 $\chi_{B(4)} \leq \psi \leq \chi_{B(8)}$.

由注意~\ref{note:3.2}, Hardy--Littlewood 极大算子的 $L^{p}$ 有界性以及 Lebesgue 收敛定理, 在 $L^{p}(\mathbb{R}^{n})$ 的拓扑中,
\[
\lim_{j \to \infty}\psi_{j}(D)f = (2\pi)^{-\frac{n}{2}}\lim_{j \to \infty}2^{jn}\mathcal{F}^{-1}(2^{j}\,\cdot\,) * f = f.
\]
另外, 由 Littlewood--Paley 理论,
\[
\limsup_{k \to -\infty}\|\psi_{k}(D)f\|_{L^{p}} \lesssim \limsup_{k \to -\infty}\left\|\left(\sum_{j=-\infty}^{\infty}|\varphi_{j}(D)\psi_{k}(D)f|^{2}\right)^{1/2}\right\|_{L^{p}}.
\]
这里, $\varphi_{j}(D)\psi_{k}(D)f$ 在 $j > k+2$ 时为 $0$, 在 $j < k-2$ 时与 $\varphi_{j}(D)f$ 一致, 所以可以分解为
\[
\begin{aligned}
\limsup_{k \to -\infty}\|\psi_{k}(D)f\|_{L^{p}} &\lesssim \limsup_{k \to -\infty}\left\{\left\|\left(\sum_{j=-\infty}^{k-2}|\varphi_{j}(D)f|^{2}\right)^{1/2}\right\|_{L^{p}}\right. \\
&\qquad \left. + \sum_{j=k-2}^{k+2}\|\varphi_{j}(D)\psi_{k}(D)f\|_{L^{p}}\right\}.
\end{aligned}
\]
这里, 由 Young 不等式,
\[
\begin{aligned}
\limsup_{k \to -\infty}\|\psi_{k}(D)f\|_{L^{p}} &\lesssim \limsup_{k \to -\infty}\left\{\left\|\left(\sum_{j=-\infty}^{k-2}|\varphi_{j}(D)f|^{2}\right)^{1/2}\right\|_{L^{p}}\right. \\
&\qquad \left. + \sum_{j=k-2}^{k+2}\|\varphi_{j}(D)f\|_{L^{p}}\right\},
\end{aligned}
\]
所以由积分的单调性, 可归纳为
\[
\limsup_{k \to -\infty}\|\psi_{k}(D)f\|_{L^{p}} \lesssim \limsup_{k \to -\infty}\left\|\left(\sum_{j=-\infty}^{k+2}|\varphi_{j}(D)f|^{2}\right)^{1/2}\right\|_{L^{p}}.
\]
因此,
\[
\left\|\left(\sum_{j=-\infty}^{k+2}|\varphi_{j}(D)f|^{2}\right)^{1/2}\right\|_{L^{p}} \lesssim \|f\|_{L^{p}} < \infty,
\]
所以由 Lebesgue 收敛定理, 当 $k \to -\infty$ 时, 得到 $\psi_{k}(D)f \to 0$. 这意味着
\[
f = \lim_{k \to \infty}(\psi_{k}(D)f - \psi_{-k}(D)f),
\]
且因为 $(\psi_{k}(D)f - \psi_{-k}(D)f) \in \mathcal{S}_{\infty}(\mathbb{R}^{n})$, 所以 $f$ 可以由 $\mathcal{S}_{\infty}(\mathbb{R}^{n})$ 的元素在 $L^{p}(\mathbb{R}^{n})$ 中近似.
\end{proof}

\begin{exercise}\label{exer:6.1}
对 $a > 0$, 证明由
\[
\Xi(a) \equiv \sum_{j=-\infty}^{\infty}\frac{1}{2^{-j}a^{-1} + 2^{j}a}
\]
给出的函数 $\Xi: (0,\infty) \to \mathbb{R}$ 是有界的.
\end{exercise}

% ============ 第6章 §6.2 Littlewood--Paley 理论的应用 (原书 pp.130-132) ============

\section{Littlewood--Paley 理论的应用}\label{sec:6.2}
\index{R@Riesz 变换}




这里, 设 $m$ 是在除去原点的 $\mathbb{R}^{n}$ 上定义的 $C^{\infty}$ 级函数, 对所有多重指数 $\alpha$ 满足 $|\partial^{\alpha}m(\xi)| \lesssim |\xi|^{-|\alpha|}$. 不过, $\lesssim$ 中所含的常数可以依赖于 $\alpha$. $T_{m}f = \mathcal{F}^{-1}[m\mathcal{F}f]$ 由 Plancherel 定理是 $L^{2}$ 有界的, 但要对 $1 < p < \infty$ 证明 $T_{m}$ 可以扩张为 $L^{p}$ 有界线性算子.

\begin{theorem}\label{thm:6.2}
对 $f \in \mathcal{S}_{\infty}(\mathbb{R}^{n})$, $T_{m}f \in \mathcal{S}_{\infty}(\mathbb{R}^{n})$ 成立, 且若 $1 < p < \infty$, 则成立不依赖于 $f \in \mathcal{S}_{\infty}(\mathbb{R}^{n})$ 的常数估计 $\|T_{m}f\|_{L^{p}} \lesssim \|f\|_{L^{p}}$.
\end{theorem}

作为这样的 $m$ 的例子, 对 $j = 1, 2, \ldots, n$ 有
\[
m(\xi) = \frac{\xi_{j}}{|\xi|},
\]
与这个 $m$ 对应的 $T_{m}$ 称为第 $j$ Riesz 变换.

\begin{proof}
设 $g \in \mathcal{S}(\mathbb{R}^{n})$. $g \in \mathcal{S}_{\infty}(\mathbb{R}^{n})$ 与 $\partial^{\alpha}\mathcal{F}g(0) = 0$, $\alpha \in \mathbb{N}_{0}^{n}$ 是等价的. 因此, 使用这个判定, 可得 $T_{m}f \in \mathcal{S}_{\infty}(\mathbb{R}^{n})$.

来证明范数不等式 $\|T_{m}f\|_{L^{p}} \lesssim \|f\|_{L^{p}}$. 为此, 只要证明
\[
\left\|\left(\sum_{j=-\infty}^{\infty}|\varphi_{j}(D)T_{m}f|^{2}\right)^{1/2}\right\|_{L^{p}} \lesssim \|f\|_{L^{p}}
\]
即可. 通过使用 Rademacher 序列来考虑, 只要证明
\[
L_{t}(z) = \sum_{j=-J}^{J}r_{j}(t)\mathcal{F}^{-1}[m\cdot\varphi_{j}](z)
\]
满足以下条件即可.

(1) $\|\mathcal{F}L_{t}\|_{L^{\infty}} \lesssim 1$.

(2) $\sup_{z \in \mathbb{R}^{n}, t \in (0,1)}|z|^{n}|L_{t}(z)| \lesssim 1$.

(3) $\sup_{z \in \mathbb{R}^{n}, t \in (0,1)}|z|^{n+1}|\nabla L_{t}(z)| \lesssim 1$.

这与证明 Littlewood--Paley 理论时所考虑的 $K_{t}$ 的估计相同, 因此省略细节.
\end{proof}

\begin{exercise}\label{exer:6.2}
若 $m \in L^{\infty}(\mathbb{R}^{n})$, 则证明 $T_{m}$ 是 $L^{2}(\mathbb{R}^{n})$ 有界的.
\end{exercise}

\begin{exercise}\label{exer:6.3}
设 $m$ 是在除去原点的 $\mathbb{R}^{n}$ 上定义的 $C^{\infty}$ 级函数, 对所有多重指数 $\alpha$ 满足 $|\partial^{\alpha}m(\xi)| \lesssim |\xi|^{-|\alpha|}$. 对 $g \in \mathcal{S}_{\infty}(\mathbb{R}^{n})$, 证明 $T_{m}g \in \mathcal{S}_{\infty}$.
\end{exercise}


% 结束语 (原书 最後に, pp.132-133)
\section*{结束语}
\addcontentsline{toc}{section}{结束语}

本书以说明调和分析的基本事项为主旨, 但作为调和分析的研究对象, 极大算子, 奇异积分算子, Fourier 变换等是典型的. 在这一周边, 存在着函数空间, 积分不等式等理论. 算子不成为 $L^{1}$ 有界的场合屡见不鲜, 这一点已通过极大算子, 奇异积分算子作了说明, 但如何克服这一点, 是很有趣的问题. 由于非有界性已经被证明, 考虑替代的, 尽可能接近有界性的结论是什么, 是有意义的, 这也是这个领域的目的之一.

调和分析的讨论乍一看不透明的居多, 但说是为实现这一目的而形成的, 也不为过. 调和分析的技法日新月异. 本书在其中说明了下面这些事情.

(1) \textbf{【覆盖引理】}实际说明了 5r 覆盖引理, 并说明了它给极大算子带来的有界性.

(2) \textbf{【分解理论】}Calderón--Zygmund 理论是处理奇异积分算子有界性的有效工具.

(3) \textbf{【二进制分析】}为了作 Calderón--Zygmund 分解, 使用了二进制区间, 像这样使用 $\{0,1\}$ 的分析是有效的.

(4) \textbf{【Littlewood--Paley 理论】}从奇异积分算子的有界性可以看出 $L^{p}$ 范数的改写是可能的, 由此能够考察种种算子的性质.


% 练习题解答
% ============ 练习问题解答 (原书 pp.133-155) ============
\chapter*{练习题解答}

\section*{第1章}

\paragraph*{练习题 1.1}
$\|a\|_{\ell^{2}} = \sqrt{\frac{n(n+1)}{2}}$, $\|a\|_{\ell^{4}} = \sqrt[4]{\frac{n(n+1)(2n+1)}{6}}$.

\paragraph*{练习题 1.2}
由 $\vec{a}+\vec{b} = (5,7,3)$, 有 $\|\vec{a}\| = \sqrt{14}$, $\|\vec{b}\| = \sqrt{29}$, $\|\vec{a}+\vec{b}\| = \sqrt{83}$ 成立. 因此, 由积分的三角不等式是显然的. $\blacksquare$

\paragraph*{练习题 1.3}
有若干种解法.

(解法 1) 利用恒等式
\[
(a^{3}+b^{3})(c^{3}+d^{3})^{2}-(ac^{2}+bd^{2})^{3} = (-bc+ad)^{2}(bc^{4}+2ac^{3}d+2bcd^{3}+ad^{4}).
\]

(解法 2) 由 Hölder 不等式, $(a^{3}+b^{3})^{\frac{1}{3}}(C^{\frac{3}{2}}+D^{\frac{3}{2}})^{\frac{2}{3}} \geq aC+bD$, $C, D > 0$ 成立, 代入 $C = c^{2}$, $D = d^{2}$, 两边三次方. $\blacksquare$

\paragraph*{练习题 1.4}
有若干种解法.

(解法 1) 画出 $s = e^{t}$ 的图形, 考察 $e^{(1-\theta)x+\theta y}$ 与 $(1-\theta)e^{x} + \theta e^{y}$ 出现在何处. 【凸性的应用】

(解法 2) 在引理~\ref{lem:1.1} 中, 取 $a = e^{(1-\theta)x}$, $b = e^{\theta y}$, $p = \frac{1}{1-\theta}$, $q = \frac{1}{\theta}$.

(解法 3) 由微积分基本定理,
\begin{align*}
(1-\theta)e^{x} + \theta e^{y} - e^{(1-\theta)x+\theta y} &= \int_{y}^{x}(1-\theta)(e^{z} - e^{(1-\theta)z+\theta y})\dd z \\
&= \int_{y}^{x}(1-\theta)(1 - e^{\theta(y-z)})e^{z}\dd z \\
&\geq 0.
\end{align*}
$\blacksquare$

\paragraph*{练习题 1.5}
对 $1 \leq p < \infty$,
\[
\|ka\|_{\ell^{p}} = \left(\sum_{n=1}^{N}|ka_{n}|^{p}\right)^{\frac{1}{p}}
\]
\[
= \left(\sum_{n=1}^{N}|k|^{p}|a_{n}|^{p}\right)^{\frac{1}{p}} = |k|\left(\sum_{n=1}^{N}|a_{n}|^{p}\right)^{\frac{1}{p}} = |k|\|a\|_{\ell^{p}}
\]
由此得到 $|k|$. $\blacksquare$

\paragraph*{练习题 1.6}
当 $a > b$ 时, $\min(a, b) = b$, $\frac{a+b-|a-b|}{2} = \frac{a+b-a+b}{2} = b$, $\max(a, b) = a$, $\frac{a+b+|a-b|}{2} = \frac{a+b+a-b}{2} = a$, 故得结论. 另一方面, 当 $a \leq b$ 时, $\min(a, b) = a$, $\frac{a+b-|a-b|}{2} = \frac{a+b-b+a}{2} = a$, $\max(a, b) = b$, $\frac{a+b+|a-b|}{2} = \frac{a+b+b-a}{2} = b$, 故得结论. $\blacksquare$

\paragraph*{练习题 1.7}
(1) $\infty$, (2) 1, (3) $\infty$, (4) $\infty$, (5) $\infty$, (6) 2.

\paragraph*{练习题 1.8}
图形从略, 积分值为 $-3$. $\blacksquare$

\paragraph*{练习题 1.9}
(1) 由定理~\ref{thm:1.3},
\[
\sum_{n=1}^{N}|a_{n}| = \int_{0}^{\infty}\sharp\{1 \leq n \leq N: |a_{n}| > t\}\dd t
\]
成立. 这里, 当 $t > 0$ 时, $|a_{n}| > t$ 与 $|a_{n}| > [t]$ 等价, 因此
\[
\sum_{n=1}^{N}|a_{n}| = \int_{0}^{\infty}\sharp\{1 \leq n \leq N: |a_{n}| > [t]\}\dd t.
\]
将积分分割来考虑,
\begin{align*}
\int_{0}^{\infty}\sharp\{1 \leq n \leq N: |a_{n}| > [t]\}\dd t &= \sum_{k=0}^{\infty}\int_{k}^{k+1}\sharp\{1 \leq n \leq N: |a_{n}| > k\}\dd t \\
&= \sum_{k=0}^{\infty}\sharp\{1 \leq n \leq N: |a_{n}| > k\}.
\end{align*}

(2) 设 $a_{n}$ 为将 $n$ 素因数分解所得因子 $5$ 的个数, 则 $n!$ 用十进制书写时末尾出现的 $0$ 的个数表为
\[
\sum_{k=1}^{n}a_{k}.
\]
由 (1),
\[
\sum_{k=1}^{n}a_{k} = \sum_{k=0}^{\infty}\sharp\{1 \leq m \leq n: a_{m} > k\}
\]
成立. 这里, 由于 $a_{n} > k$ 与 $a_{n} \geq k+1$ 等价, 且与 $n$ 含有因子 $5^{k+1}$ 等价, 所以
\[
\sharp\{1 \leq m \leq n: a_{m} > k\} = \left[\frac{n}{5^{k+1}}\right].
\]
将其代入, 即得结论. $\blacksquare$

\paragraph*{练习题 1.10}
(1) 1.

(2) 2.

(3) 算子范数由 $\sup\{\sqrt{3x^{2}+2xy+y^{2}} : x^{2}+y^{2} = 1\}$ 给出. 用三角函数计算, 该值为 $\sup\{\sqrt{2+\sqrt{2}\cos(\theta+\pi/4)} : \theta \in \mathbb{R}\} = \sqrt{2+\sqrt{2}}$. $\blacksquare$

\paragraph*{练习题 1.11}
(1) 当 $x^{2}+y^{2} = 1$ 时, $(4x-2y)^{2}+(8x-4y)^{2} = 20(2x-y)^{2} \leq 100$, 故算子范数为 $10$.

(2) 因 $f^{2} = 0$, 算子范数为 $0$.

(3) $\|f^{2}\| = 0 < 100 = \|f\|^{2}$.

\paragraph*{练习题 1.12}
给出相应的条件最大值问题.

(1) 求当 $\sqrt{x^{2}+y^{2}} = 1$ 时 $\sqrt{(4x-3y)^{2}+(x-4y)^{2}}$ 的最大值. 其最大值为 $\sqrt{21+4\sqrt{17}} = 2+\sqrt{17}$.

(2) 求当 $|x|+|y| = 1$ 时 $|4x-3y|+|x-4y|$ 的最大值. 当 $x = 0$, $y = 1$ 时为 7. 当 $|x|+|y| = 1$ 时, 由 $|4x-3y|+|x-4y| \leq 5|x|+7|y| \leq 7(|x|+|y|) = 7$, 为 7.

(3) 求当 $\sqrt{x^{2}+y^{2}} = 1$ 时 $\max(|4x-3y|, |x-4y|)$ 的最大值. 当 $x^{2}+y^{2} = 1$ 时, $|4x-3y|$ 的最大值为 $5$. 同样, $|x-4y|$ 的最大值为 $\sqrt{17}$. 因此, 算子范数为 $5$.

(4) 求当 $\max(|x|, |y|) = 1$ 时 $\sqrt{(4x-3y)^{2}+(x-4y)^{2}}$ 的最大值. 实际上, 因 $(4x-3y)^{2} \leq 49$, $(x-4y)^{2} \leq 25$, 所求值为 $\sqrt{74}$. 范数值在 $(x,y) = (1,-1)$ 处取得. $\blacksquare$

\paragraph*{练习题 1.13}
在条件 $x^{2}+y^{2} = 1$ 下求 $(2x-3y)^{2}+(3x+4y)^{2} = 13x^{2}+12xy+25y^{2}$ 的最大值. 若 $x = \cos\theta$, $y = \sin\theta$, 则 $13x^{2}+12xy+25y^{2} = 19-6\cos 2\theta+6\sin 2\theta$, 故第一个算子范数为 $\sqrt{19+6\sqrt{2}} = 1+3\sqrt{2}$. 又, 在条件 $\max(|x|, |y|) = 1$ 下 $(2x-3y)^{2}+(3x+4y)^{2} = 13x^{2}+12xy+25y^{2}$ 的最大值为 50, 故第二个算子范数为 $5\sqrt{2}$. $\blacksquare$

\paragraph*{练习题 1.14}
在条件 $x^{2}+y^{2} = 1$ 下, $(2x+y)^{2}+(x-2y)^{2}+(x+y)^{2} = 6x^{2}+6y^{2}+2xy$ 的最大值为 7, 故所求算子范数为 $\sqrt{7}$. $\blacksquare$

\paragraph*{练习题 1.15}
由 Hölder 不等式, $\|(2,1)\|_{\ell^{3/2}} = (2\sqrt{2}+1)^{\frac{2}{3}}$. $\blacksquare$

\paragraph*{练习题 1.16}
设 $\vec{p} = (x,y)$, $\vec{a} = (1,-4)$, $\vec{b} = (2,-3)$, 由 Minkowski 不等式,
\[
\|\vec{p}+\vec{a}\|_{\ell^{3}} \leq \|\vec{p}+\vec{b}\|_{\ell^{3}} + \|\vec{a}-\vec{b}\|_{\ell^{3}}
\]
成立. 因此 $f(x,y) \leq \|\vec{a}-\vec{b}\|_{\ell^{3}} = \sqrt[3]{2}$. 这里, 若 $\vec{p} = -\vec{b}$, 则等号成立, 故最大值为 $\sqrt[3]{2}$. $\blacksquare$

\paragraph*{练习题 1.17}
(1) 对实数列 $a = \{a_{n}\}_{n=1}^{N}$,
\[
\|U_{N}a\|_{\ell^{1}} = \sum_{n=1}^{N}|b_{n}| = \sum_{n=1}^{N}\frac{|a_{1}+a_{2}+\cdots+a_{n}|}{n} \leq \sum_{n=1}^{N}\frac{\|a\|_{\ell^{1}}}{n}
\]
由 $(1,0,0,\ldots,0)$ 达成等号, 故 $\|U_{N}\|_{\ell^{1}\to\ell^{1}} = \sum_{n=1}^{N}\frac{1}{n}$.

(2) 取平均的操作不增大最大值, 故显然.

(3) 当 $\lambda \geq 1$ 或 $\lambda \leq \frac{1}{N}$ 时, 由所考虑的集合分别为空集与 $\{1,2,\ldots,N\}$ 的子集可知是显然的. 设 $\frac{1}{N} < \lambda < 1$. 取满足 $\frac{1}{l} < \lambda \leq \frac{1}{l-1}$ 的自然数 $l = 2, 3, \ldots, N$. 于是,
\[
\{k \in \{1, 2, \ldots, N\}: |b_{k}| > \lambda\|a\|_{\ell^{1}}\} \subset \left\{k \in \{1, 2, \ldots, N\}: |b_{k}| > \frac{1}{l}\|a\|_{\ell^{1}}\right\} \subset \{1, 2, \ldots, l-1\}
\]
所以
\[
\sharp\{k \in \{1, 2, \ldots, N\}: |b_{k}| > \lambda\|a\|_{\ell^{1}}\} \leq l-1 \leq \frac{1}{\lambda}.
\]

(4) 用 Marcinkiewicz 插值定理对 (2) 与 (3) 进行插值.

(5) 设字母变量 $a$, $b$, $c$ 为实数. 设 $S = |a|^{2}+|b|^{2}+|c|^{2}$. 由
\begin{align*}
\|U_{1}\|_{\ell^{2}\to\ell^{2}} &= \max\{|a|: |a| = 1\} = 1, \\
\|U_{2}\|_{\ell^{2}\to\ell^{2}} &= \max\{\sqrt{|a|^{2}+4^{-1}|a+b|^{2}}: |a|^{2}+|b|^{2} = 1\} \\
&= \frac{1+\sqrt{5}}{2\sqrt{2}} < 1.15,
\end{align*}
\begin{align*}
\|U_{3}\|_{\ell^{2}\to\ell^{2}} &= \max\{\sqrt{|a|^{2}+4^{-1}|a+b|^{2}+9^{-1}|a+b+c|^{2}}: S = 1\} \\
&\geq \|U_{3}\mathbf{e}_{1}\|_{\ell^{2}} = \frac{7}{6} > 1.15.
\end{align*}

\paragraph*{练习题 1.18}
由
\[
r_{1}(t)r_{2}(t)r_{3}(t) = \chi_{[0,1/8)\cup[3/8,1/2)\cup[5/8,7/8)}(t) - \chi_{[1/8,3/8)\cup[1/2,5/8)\cup[7/8,1)}(t),
\]
所求积分值为 $-\frac{1}{8}$. $\blacksquare$

\paragraph*{练习题 1.19}
由
\[
2e^{t^{2}} = \sum_{n=0}^{\infty}\frac{2}{n!}t^{2n}, \quad e^{t}+e^{-t} = \sum_{n=0}^{\infty}\frac{2}{(2n)!}t^{2n}
\]
是显然的. $\blacksquare$

\paragraph*{练习题 1.20}
只需在 $0 \leq t < 1$ 的范围内考虑, 所求为 $\sum_{j=1}^{N}|a_{j}|$. $\blacksquare$

\paragraph*{练习题 1.21}
\begin{align*}
\mathcal{U}_{0} &= \{\{1\}, \{3\}, \{6\}, \{8\}, \{12\}, \{13\}, \{8,9,10,11,12,13,14,15\}, \\
&\qquad \{14\}, \{6,7\}, \{8,9\}, \{12,13\}, \{14,15\}, \{8,9,10,11\}, \{12,13,14,15\}, \\
&\qquad \{0,1,2,3,4,5,6,7,8,9,10,11,12,13,14,15\}\},
\end{align*}
\[
\mathcal{U} = \{\{0,1,2,3,4,5,6,7,8,9,10,11,12,13,14,15\}\}, \quad E = I(4,0),
\]
\[
g_{0} = g_{1} = \cdots = g_{15} = 3,
\]
\[
b: -3, 1, -3, 1, -1, -3, 3, -3, 9, -3, -1, -3, 3, 3, 3, -3,
\]
\[
b_{\{0,1,2,3,4,5,6,7,8,9,10,11,12,13,14,15\}} = b.
\]
$\blacksquare$

\paragraph*{练习题 1.22}
\begin{itemize}
\item $E_{0}(a)$: $a$ 本身
\item $E_{2}(a)$: $2, 2, 2, 2, 2, 2, 2, 2, 3.5, 3.5, 3.5, 3.5, 4.5, 4.5, 4.5, 4.5$
\item $E_{3}(a)$: $2, 2, 2, 2, 2, 2, 2, 2, 4, 4, 4, 4, 4, 4, 4, 4$
\item $E_{4}(a)$: $3, 3, 3, 3, 3, 3, 3, 3, 3, 3, 3, 3, 3, 3, 3, 3$
\item $D_{0}(a)$: $-2, 2, -2, 2, 1, -1, 3, -3, 6, -6, 1, -1, 0, 0, 3, -3$
\item $D_{1}(a)$: $0, 0, 0, 0, -1, -1, 1, 1, 2.5, 2.5, -2.5, -2.5, 1.5, 1.5, -1.5, -1.5$
\item $D_{2}(a)$: $0, 0, 0, 0, 0, 0, 0, 0, -0.5, -0.5, -0.5, -0.5, 0.5, 0.5, 0.5, 0.5$
\item $D_{3}(a)$: $-1, -1, -1, -1, -1, -1, -1, -1, 1, 1, 1, 1, 1, 1, 1, 1$ $\blacksquare$
\end{itemize}

\paragraph*{练习题 1.23}
如直接计算所见, 为 $\sum_{k=0}^{7}a_{k}^{2}$, 但这可视为定理~\ref{thm:1.9}(3) 中 $N = 3$ 情形的结果. 实际上,
\[
E_{3}(a) = \{A, A, A, A, A, A, A, A\},
\]
\[
D_{0}(a) = \{B, -B, C, -C, D, -D, E, -E\},
\]
\[
D_{1}(a) = \{F, F, -F, -F, G, G, -G, -G\},
\]
\[
D_{2}(a) = \{H, H, H, H, -H, -H, -H, -H\}
\]
故 $\|a\|_{\ell^{2}}^{2} = \|E_{N}(a)\|_{\ell^{2}}^{2}+\sum_{k=0}^{N-1}\|D_{k}(a)\|_{\ell^{2}}^{2}$ 的右边为 $8A^{2}+2B^{2}+2C^{2}+2D^{2}+2E^{2}+4F^{2}+4G^{2}+8H^{2}$, 左边为 $\sum_{k=0}^{7}a_{k}^{2}$. $\blacksquare$

\paragraph*{练习题 1.24}
对常数数列 $m^{k} = (-1)^{k}$ 使用定理~\ref{thm:1.12}. $\blacksquare$

\paragraph*{练习题 1.25}
\begin{align*}
\|g\|_{\ell^{2}} &= \sqrt{\sum_{j=1}^{N}|g_{j}|^{2}} \leq \sqrt{\sup_{j=1,2,\ldots,N}|g_{j}|\sum_{j=1}^{N}|g_{j}|} \\
&= \sqrt{\|g\|_{\ell^{1}}\|g\|_{\ell^{\infty}}}.
\end{align*}

\paragraph*{练习题 1.26}
关于前半部分, 由 $E_{k}(e)_{j} = 2^{-k}$, $E_{k+1}(e)_{j} = 2^{-k-1}$ 是显然的. 后半部分如下证明.
\begin{align*}
\|g(e)\|_{\ell^{1}} &= \sum_{j=0}^{2^{N}-1}|g(e)_{j}| \\
&\geq |E_{N}(e)| + \sum_{k=0}^{N-1}\sum_{j=2^{k}}^{2^{k+1}-1}|D_{k}(e)_{j}| \\
&\geq \sum_{k=0}^{N-1}\sum_{j=2^{k}}^{2^{k+1}-1}2^{-k-2} \\
&= \frac{N}{4}.
\end{align*}

\paragraph*{练习题 1.27}
设 $a = \{a_{j_{1}j_{2}}\}_{j_{1},j_{2}=0,1,\ldots,2^{N}-1}$, $E_{k}^{2}(a) = \{b_{j_{1}j_{2}}\}_{j_{1},j_{2}=0,1,\ldots,2^{N}-1}$, 则
\[
b_{j_{1}j_{2}} = \frac{1}{2^{k}}\sum_{\mu=2^{k}[2^{-k}l]}^{2^{k}[2^{-k}l]+2^{k}-1}a_{j_{1}\mu}.
\]
又, 设 $E_{k'}^{1}(E_{k}^{2}(a)) = \{c_{j_{1}j_{2}}\}_{j_{1},j_{2}=0,1,\ldots,2^{N}-1}$, 则
\begin{align*}
c_{j_{1}j_{2}} &= \frac{1}{2^{k'}}\sum_{\nu=2^{k'}[2^{-k'}l]}^{2^{k'}[2^{-k'}l]+2^{k'}-1}a_{\nu j_{2}} \\
&= \frac{1}{2^{k'+k}}\sum_{\mu=2^{k}[2^{-k}l]}^{2^{k}[2^{-k}l]+2^{k}-1}\sum_{\nu=2^{k'}[2^{-k'}l]}^{2^{k'}[2^{-k'}l]+2^{k'}-1}a_{\nu\mu}.
\end{align*}
因此, 可知关于 $\mu$, $\nu$ 的求和顺序可以交换, 故得 $E_{k}^{2} \circ E_{k'}^{1}(a) = E_{k'}^{1} \circ E_{k}^{2}(a)$. $\blacksquare$

\paragraph*{练习题 1.28}
固定 $j$, 任意取 $J < j$. 由
\[
ja_{j} \leq Ja_{j} + \sum_{k=J+1}^{j}a_{k}
\]
令 $j \to \infty$, 得
\[
\limsup_{j \to \infty}ja_{j} \leq \sum_{k=J+1}^{\infty}a_{k}.
\]
令 $J \to \infty$ 即得结论. $\blacksquare$

\paragraph*{练习题 1.29}
(1) $2^{\frac{1}{p}}$.

(2) 若其为紧集, 则本应能从 $\{\mathbf{e}_{j}\}_{j=0}^{\infty}$ 中取出收敛子列, 但由 (1), 任何子列都不收敛, 故矛盾. $\blacksquare$

\paragraph*{练习题 1.30}
(1) 由 $x_{k} - x_{k'}$ 的具体形状
\[
x_{k} - x_{k'} = (0, 0, \ldots, 0, 2^{-k-1}, 2^{-k-2}, \ldots, 2^{-k'}, 0, 0, \ldots)
\]
是显然的.

(2) 对任意 $x = (x^{0}, x^{1}, \ldots) \in \operatorname{Span}(\{\mathbf{e}_{j}\}_{j=0}^{\infty})$, 设 $k = \max\{j \in \mathbb{N}_{0} : x^{j} \neq 0\}$. 当 $j > k$ 时, $\|x - x_{k}\|_{\ell^{\infty}} \geq 2^{-k}$, 故 $\{x_{k}\}_{k=1}^{\infty}$ 不收敛. $\blacksquare$

\section*{第2章}

\paragraph*{练习题 2.1}
设 $A \leq B$. 此时,
\[
2^{-1}(A^{p}+B^{p}) \leq B^{p} \leq (A+B)^{p} \leq (2B)^{p} \leq 2^{p}(A^{p}+B^{p})
\]
故得结论. $\blacksquare$

\paragraph*{练习题 2.2}
(1) $\|f\|_{L^{1}(0,\infty)} = \int_{0}^{\infty}\frac{\dd x}{1+x} = [\log(1+x)]_{0}^{\infty} = \infty$.

(2) $\|f\|_{L^{1}(0,\infty)} = \int_{0}^{\infty}\frac{\dd x}{(1+x)^{a}} = \left[\frac{(1+x)^{1-a}}{1-a}\right]_{0}^{\infty} = \infty$.

(3) $\|f\|_{L^{1}(0,\infty)} = \int_{0}^{\infty}\frac{\dd x}{(1+x)^{a}} = \left[\frac{(1+x)^{1-a}}{1-a}\right]_{0}^{\infty} = \frac{1}{a-1}$.

(4) 由 (1), (2), (3), $\|f\|_{L^{1}(0,\infty)} < \infty$ 成立的必要充分条件是 $a > 1$.

(5) 由 (4), $2a > 1$, 即 $a > \frac{1}{2}$. $\blacksquare$

\paragraph*{练习题 2.3}
省略 $L^{p}(\mathbb{R}^{n})$ 的定义. 为证明 $f \in L^{p}(\mathbb{R}^{n})$, 需要证明
\[
\int_{\mathbb{R}^{n}}\chi_{\mathbb{R}^{n}\setminus B(1)}(x)|x|^{ap}\dd x < \infty,
\]
即
\[
\int_{\mathbb{R}^{n}\setminus B(1)}|x|^{ap}\dd x < \infty.
\]
由 $\mathbb{R}^{n}\setminus B(1) = \bigcup_{j=1}^{\infty}B(2^{j})\setminus B(2^{j-1})$ 与变量替换 $x \mapsto 2^{-j}x$,
\begin{align*}
\int_{\mathbb{R}^{n}\setminus B(1)}|x|^{ap}\dd x &= \sum_{j=1}^{\infty}\int_{B(2^{j})\setminus B(2^{j-1})}|x|^{ap}\dd x \\
&= \sum_{j=1}^{\infty}2^{j(ap+n)}\int_{B(1)\setminus B(2^{-1})}|x|^{ap}\dd x \\
&= \frac{2^{ap+n}}{1-2^{ap+n}}\int_{B(1)\setminus B(2^{-1})}|x|^{ap}\dd x < \infty
\end{align*}
即得. 从该级数的计算可以看出, 若换为 $a = -n/p$, 则最后一个式子的分母变为 $0$, 论证便行不通了. $\blacksquare$

\section*{第3章}

\paragraph*{练习题 3.1}
(1) 画图便知, $[-1,1]^{n} \cap Q(x, |x|+2) = [-1,1]^{n}$. 因此,
\begin{align*}
M\chi_{[-1,1]^{n}}(x) &= \sup_{Q: Q \text{ 是立方体}}\frac{|Q \cap [-1,1]^{n}|}{|Q|} \\
&\geq \frac{|[-1,1]^{n} \cap Q(x, |x|+2)|}{|Q(x, |x|+2)|} \\
&= \frac{2^{n}}{(|x|+2)^{n}}.
\end{align*}

(2) 由提示中的不等式,
\begin{align*}
\int_{\mathbb{R}^{n}}Mf(x)\dd x &\geq \int_{\mathbb{R}^{n}}\frac{2^{n}}{(|x|+2)^{n}}\dd x \\
&\geq \int_{\mathbb{R}^{n}}\prod_{j=1}^{n}\frac{2}{(|x_{j}|+2)}\dd x \\
&= \prod_{j=1}^{n}\int_{\mathbb{R}}\frac{2}{(|x_{j}|+2)}\dd x_{j} \\
&= \infty
\end{align*}
即得. $\blacksquare$

\paragraph*{练习题 3.2}
如提示所述用 $A$ 表示条件等, 则要判定的条件与
\[
\varepsilon < \sqrt[n]{\frac{A}{\lambda}} - 1
\]
相同. 因此, 可知可取
\[
\varepsilon = \min\left(1, \frac{1}{2}\left(\sqrt[n]{\frac{A}{\lambda}} - 1\right)\right).
\]
$\blacksquare$

\paragraph*{练习题 3.3}
从略.

\paragraph*{练习题 3.4}
因 $B$ 最长, 取 $K_{1} = B$. 在与 $B$ 不相交者中最长的是 $E$, 故取 $K_{2} = E$. 因 $A \subset 3B$, $C \subset 3B$, $D \subset 3B$, $F \subset 3B$, 取 $L = 2$, $K_{1} = B$, $K_{2} = E$, $\iota(A) = \iota(B) = \iota(C) = \iota(D) = \iota(F) = 1$, $\iota(E) = 2$ 即可. 【注意】也有取 $K_{1} = E$ 的方法. $\blacksquare$

\paragraph*{练习题 3.5}
(1) 由
\[
b_{n,k}(x) - b_{n,k+1}(x) = \frac{n!}{(k+1)!(n-k)!}x^{k}(1-x)^{n-k-1}(k+1-nx-x)
\]
是显然的.

(2) 由 (1), 在 $t \leq \frac{k_{n,x}}{n+1}$ 上 $K(x,\cdot)$ 单调递增, 在 $\frac{k_{n,x}}{n+1} \leq t \leq \frac{k_{n,x}+1}{n+1}$ 上为常数函数, 在 $t \geq \frac{k_{n,x}+1}{n+1}$ 上单调递减.

(3) 计算右边,
\begin{align*}
\int_{[0,1]}F(t)\dd\mu_{x}(t) &= (n+1)\sum_{k=0}^{n}b_{n,k}(x)\left\{F\left(\frac{k+1}{n+1}\right) - F\left(\frac{k}{n+1}\right)\right\} \\
&= (n+1)\sum_{k=0}^{n}b_{n,k}(x)\int_{\frac{k}{n+1}}^{\frac{k+1}{n+1}}f(t)\dd t \\
&= K_{n}f(x, t).
\end{align*}

(4) 因 $|K_{n}f(x,t)| \leq K_{n}[|f|](x,t)$, 若必要, 可考虑以 $|f|$ 代替 $f$, 设 $f \geq 0$. 对 (3) 中的 $F$, 若 $t \geq x$, 则 $0 \leq F(t) \leq (t-x)Mf(x)$; 若 $t \leq x$, 则 $0 \geq F(t) \geq (t-x)Mf(x)$. 又由 (1), 当 $t \geq x$ 时 $\mu_{x} \geq 0$, 当 $t \leq x$ 时 $\mu_{x} \leq 0$. 因此由 (3), $0 \leq K_{n}f(x,t) \leq \int_{[0,1]}(t-x)Mf(x)\dd\mu_{x}(t)$. 将 (3) 应用于常数函数 $Mf(x)$ 以代替 $f$, 得 $0 \leq K_{n}f(x,t) \leq K_{n}[Mf(x)](x,t) = Mf(x)$. $\blacksquare$

\paragraph*{练习题 3.6}
(1) 任意取包含 $x$ 的立方体 $Q$. 由积分的三角不等式,
\[
\frac{1}{|Q|}\int_{Q}|f(y)+g(y)|\dd y \leq \frac{1}{|Q|}\int_{Q}|f(y)|\dd y + \frac{1}{|Q|}\int_{Q}|g(y)|\dd y.
\]
由 $Mf(x)$, $Mg(x)$ 的定义,
\[
\frac{1}{|Q|}\int_{Q}|f(y)|\dd y + \frac{1}{|Q|}\int_{Q}|g(y)|\dd y \leq Mf(x) + Mg(x).
\]
将这些组合起来, 得
\[
\frac{1}{|Q|}\int_{Q}|f(y)+g(y)|\dd y \leq Mf(x) + Mg(x).
\]
对包含 $x$ 的立方体 $Q$ 取上确界, 即得结论.

(2) 由绝对值的性质,
\[
\frac{1}{|Q|}\int_{Q}|af(y)|\dd y = |a| \times \frac{1}{|Q|}\int_{Q}|f(y)|\dd y
\]
故结论显然. $\blacksquare$

\paragraph*{练习题 3.7}
(1) \textbf{【当 $a < 5$ 时】}
\begin{align*}
F_{a}(t) &= \int_{5}^{\min(8, \max(5, t))}(s-5)\dd s \\
&= \frac{1}{2}(\min(8, \max(5, t))-5)^{2} \\
&= \frac{1}{2}\min(3, \max(0, t-5))^{2}.
\end{align*}
若 $t < 5$, 则 $F_{a}(t) = 0$, 故在 $t = a$ 处可微. 以后, 设
\[
\Psi(t) = \frac{1}{2}\min(3, \max(0, t-5))^{2}.
\]
此时, 对一切实数 $t$, $F_{a}(t) = \Psi(t) - \Psi(a)$.

\textbf{【当 $a = 5$ 时】} $F_{5}(t) = \Psi(t)$, 由 $\Psi'(t-0) = 0$, $\Psi'(t+0) = 0$ 可微.

\textbf{【当 $5 < a < 8$ 时】} 当 $t \in (5, 8)$ 时, $\Psi(t) = \frac{1}{2}(t-5)^{2}$, 故 $\Psi$ 在 $t = a$ 处可微. 因此 $F_{a}$ 也在 $t = a$ 处可微.

\textbf{【当 $a = 8$ 时】} 若 $t < 8$, 则 $F_{8}(t) = \frac{1}{2}(t-5)^{2} - \frac{9}{2}$. 因此, 若 $t < 8$, 则 $F_{8}'(8-0) = 3$. 另一方面, 若 $t > 8$, 则 $F_{8}(t) = -\frac{9}{2}$, 故 $F_{8}'(8+0) = 0$. 右导数与左导数不一致, 故 $F_{8}$ 在 $t = 8$ 处不可微.

\textbf{【当 $a > 8$ 时】} 当 $t > 8$ 时, $F_{a}(t) = \Psi(t) - \Psi(a) = -\frac{9}{2} + \frac{9}{2} = 0$, 故 $F_{a}'(a) = 0$.

(2) \textbf{【当 $a < 6$ 时】}
\begin{align*}
F_{a}(t) &= \int_{6}^{\min(9, \max(6, t))}(s-9)\dd s \\
&= \frac{1}{2}(\min(9, \max(6, t))-9)^{2} - \frac{9}{2} \\
&= \frac{1}{2}\min(0, \max(-3, t-9))^{2} - \frac{9}{2}.
\end{align*}
若 $t < 6$, 则 $F_{a}(t) = 0$, 故在 $t = a$ 处可微. 以后, 设
\[
\Psi(t) = \frac{1}{2}\min(0, \max(-3, t-9))^{2} - \frac{9}{2}.
\]
此时, 对一切实数 $t$, $F_{a}(t) = \Psi(t) - \Psi(a)$.

\textbf{【当 $a = 6$ 时】} 若 $t < 6$, 则 $F_{6}(t) = 0$. 因此, 若 $t < 6$, 则 $F_{6}'(6-0) = 0$. 另一方面, 若 $t > 6$, 则 $F_{6}(t) = \frac{1}{2}(t-9)^{2} - \frac{9}{2}$, 故 $F_{6}'(6+0) = -3$. 右导数与左导数不一致, 故 $F_{6}$ 在 $t = 6$ 处不可微.

\textbf{【当 $6 < a < 9$ 时】} 当 $t \in (6, 9)$ 时, $\Psi(t) = \frac{1}{2}(t-9)^{2}$, 故 $\Psi$ 在 $t = a$ 处可微. 因此 $F_{a}$ 也在 $t = a$ 处可微.

\textbf{【当 $a = 9$ 时】} $F_{9}(t) = \Psi(t) - \Psi(9)$, 由 $F_{9}'(t-0) = 0$, $F_{9}'(t+0) = 0$ 可微.

\textbf{【当 $a > 9$ 时】} 当 $t > 9$ 时, $F_{a}(t) = \Psi(t) - \Psi(a) = -\frac{9}{2} + \frac{9}{2} = 0$, 故 $F_{a}'(t) = 0$. $\blacksquare$

\paragraph*{练习题 3.8}
(1) 由 $\frac{1}{r}\int_{x}^{x+r}|f(y)-f(x)|\dd y \leq \frac{1}{r}\int_{x-r}^{x+r}|f(y)-f(x)|\dd y$ 与定理~\ref{thm:3.9}, 结论显然.

(2) 用 $v_{n}$ 表示 $n$ 维单位球 $B(1)$ 的体积. 由 $Q(x,r) \supset B(x,r)$, $|Q(x,r)| : |B(x,r)| = 2^{n} : v_{n}$,
\[
\frac{1}{|B(x,r)|}\int_{B(x,r)}|f(y)-f(x)|\dd y \leq \frac{2^{n}}{v_{n}|Q(x,r)|}\int_{Q(x,r)}|f(y)-f(x)|\dd y
\]
再次由定理~\ref{thm:3.9}, 结论显然. $\blacksquare$

\paragraph*{练习题 3.9}
(1) 注意 $Q(x,r) \supset B(x,r)$. 利用不等式
\[
\frac{1}{|B(x,r)|}\int_{B(x,r)}|f(y)-f(x)|\dd y \leq \frac{|Q(0,1)|}{|B(0,1)|} \times \frac{1}{|Q(x,r)|}\int_{Q(x,r)}|f(y)-f(x)|\dd y
\]
得
\begin{align*}
\lim_{r \downarrow 0}\frac{1}{|B(x,r)|}\int_{B(x,r)}|f(y)-f(x)|\dd y &\leq \frac{|Q(0,1)|}{|B(0,1)|}\lim_{r \downarrow 0}\frac{1}{|Q(x,r)|}\int_{Q(x,r)}|f(y)-f(x)|\dd y \\
&= 0.
\end{align*}

(2) 因 $Q(x,r) \subset B(x, \sqrt{n}r)$,
\begin{align*}
\frac{1}{|Q(x,r)|}\int_{Q(x,r)}|f(y)-f(x)|\dd y &\leq \frac{1}{|Q(x,r)|}\int_{B(x, \sqrt{n}r)}|f(y)-f(x)|\dd y \\
&= \frac{|B(0, \sqrt{n})|}{|Q(0,1)|} \times \frac{1}{|B(x, \sqrt{n}r)|}\int_{B(x, \sqrt{n}r)}|f(y)-f(x)|\dd y
\end{align*}
即得. 因此,
\[
\lim_{r \downarrow 0}\frac{1}{|Q(x,r)|}\int_{Q(x,r)}|f(y)-f(x)|\dd y \leq \frac{|B(0, \sqrt{n})|}{|Q(0,1)|}\lim_{r \downarrow 0}\frac{1}{|B(x, \sqrt{n}r)|}\int_{B(x, \sqrt{n}r)}|f(y)-f(x)|\dd y = 0.
\]
$\blacksquare$

\paragraph*{练习题 3.10}
若这样的 $E$ 存在, 则由 $E$ 所满足的条件, $\lim_{r\downarrow 0}\frac{|E\cap(x-r,x+r)|}{2r} = \frac{1}{2}$. 另一方面, 由 Lebesgue 微分定理, 在 $\chi_{E}$ 的 Lebesgue 点处应有 $\lim_{r\downarrow 0}\frac{|E\cap(x-r,x+r)|}{2r} = 0, 1$, 故矛盾. $\blacksquare$

\section*{第4章}

\paragraph*{练习题 4.1}
对 $m$ 用数学归纳法证明. 等式在 $m = 1$ 的情形就是定理~\ref{thm:2.31}(2) 本身. 由定理~\ref{thm:2.31}(1), $A^{*}A$ 是自伴的. 设 $M \in \mathbb{N}$, 若 $(A^{*}A)^{2^{M}}$ 自伴且 $\|(A^{*}A)^{2^{M}}\| = \|A\|^{2^{M+1}}$ 成立, 则 $(A^{*}A)^{2^{M}}$ 自伴, 由定理~\ref{thm:2.31}(2),
\begin{align*}
\|(A^{*}A)^{2^{M+1}}\| &= \|(A^{*}A)^{2^{M}}(A^{*}A)^{2^{M}}\| \\
&= \|(A^{*}A)^{2^{M}}\| \cdot \|(A^{*}A)^{2^{M}}\| \\
&= (\|A\|^{2^{M}})^{2} \\
&= \|A\|^{2^{M+1}}
\end{align*}
成立. 因此, 由数学归纳法, 证明完成. $\blacksquare$

\paragraph*{练习题 4.2}
由 $f \in C_{\mathrm{c}}^{\infty}(\mathbb{R}^{n})$, 积分
\[
\int_{\mathbb{R}^{n}\setminus B(x,1)}\frac{(x_{j}-y_{j})f(y)}{|x-y|^{n+1}}\dd y
\]
绝对收敛. 因此, 只需证明
\[
\lim_{\varepsilon \downarrow 0}\int_{B(x,1)\setminus B(x,\varepsilon)}\frac{(x_{j}-y_{j})f(y)}{|x-y|^{n+1}}\dd y
\]
存在. 如提示所述,
\[
\int_{B(x,1)\setminus B(x,\varepsilon)}\frac{(x_{j}-y_{j})f(x)}{|x-y|^{n+1}}\dd y = 0
\]
所以只需证明
\[
\lim_{\varepsilon \downarrow 0}\int_{B(x,1)\setminus B(x,\varepsilon)}\frac{(x_{j}-y_{j})(f(y)-f(x))}{|x-y|^{n+1}}\dd y
\]
存在. 这里由中值定理,
\begin{align*}
\int_{B(x,1)\setminus B(x,\varepsilon)}\frac{|(x_{j}-y_{j})(f(y)-f(x))|}{|x-y|^{n+1}}\dd y &\leq \sup_{z \in \mathbb{R}^{n}}|\nabla f(z)|\int_{B(x,1)\setminus B(x,\varepsilon)}\frac{1}{|x-y|^{n-1}}\dd y \\
&\leq \sup_{z \in \mathbb{R}^{n}}|\nabla f(z)|\int_{B(x,1)}\frac{1}{|x-y|^{n-1}}\dd y < \infty
\end{align*}
故由 Lebesgue 收敛定理极限存在. $\blacksquare$

\paragraph*{练习题 4.3}
(1) 因 $\varphi'(t) = 3(2-t)\max(0,2t-t^{2})^{2}$, $\varphi'$ 是 $C^{1}$ 级的.

(2) 与 (3) 实际上, 令 $\rho(0) = 0$, 设
\[
\rho(t) = \varphi(t)\left(\sum_{j=-\infty}^{\infty}\varphi(2^{-j}t)\right)^{-1} \quad (t \in \mathbb{R}\setminus\{0\})
\]
即可. 由这个具体形式, 得 $\tau \in C^{1}(\mathbb{R})$. $\blacksquare$

\paragraph*{练习题 4.4}
(1) 设 $f \in L^{2}(\mathbb{R})$. 此时, $\|g*f\|_{L^{2}} \leq \|g\|_{L^{1}}\|f\|_{L^{2}}$ 成立. 因此 $\|L_{g}\|_{L^{2}\to L^{2}} \leq \|g\|_{L^{1}}$ 成立. 故 $L_{g}$ 有界.

(2) 对 $f, h \in L^{2}(\mathbb{R})$, 只需证明 $\langle L_{g}f, h\rangle = \langle f, L_{\check{g}}h\rangle$. 因 $L_{g}$, $L_{\check{g}}$ 有界, 由 $C_{\mathrm{c}}^{\infty}(\mathbb{R})$ 在 $L^{2}$ 中的稠密性, 可假定 $f, h \in C_{\mathrm{c}}^{\infty}(\mathbb{R})$. 此时, 由
\[
\int_{\mathbb{R}^{2}}|g(x-y)f(y)h(x)|\dd x\dd y \leq \|f\|_{L^{\infty}}\|h\|_{L^{1}}\|g\|_{L^{1}} < \infty
\]
可交换积分顺序, 对 $f, h \in C_{\mathrm{c}}^{\infty}(\mathbb{R})$, 得 $\langle L_{g}f, h\rangle = \langle f, L_{\check{g}}h\rangle$. 因此, $L_{g}$ 的伴随为 $L_{\check{g}}$. $\blacksquare$

\paragraph*{练习题 4.5}
由 $f, g \in C_{\mathrm{c}}^{1}(\mathbb{R})$, 存在某个 $R > 1$, 使得 $\operatorname{supp}(f) \subset [-R, R]$, $\operatorname{supp}(g) \subset [-R, R]$ 成立. 设 $F(t, s) = f(t-s)g(s)$, 由 $F_{t}(t, s) = f'(t-s)g(s)$, 得 $|F_{t}(t, s)| \leq \chi_{[-2R,2R]}(t)\max|f'|\max|g|$. 又 $|F(t, s)| \leq \chi_{[-2R,2R]}(t)\max|f|\max|g|$. 这里,
\begin{align*}
&\int_{\mathbb{R}}\chi_{[-2R,2R]}(t)\max|f'|\max|g|\dd t + \int_{\mathbb{R}}\chi_{[-2R,2R]}(t)\max|f|\max|g|\dd t \\
&= 2R(\max|f| + \max|f'|)\max|g| < \infty
\end{align*}
故由定理~\ref{thm:2.6}, 微分记号与积分记号可以交换, 得 $(f*g)' = f'*g$. 同理, 因 $f*g = g*f$, 也证明 $(f*g)' = f*g'$.

\paragraph*{练习题 4.6}
从略.

\paragraph*{练习题 4.7}
若 $I \in \mathcal{D}(\mathbb{R})$, 则 $I \subset [0, \infty)$ 或 $I \subset (-\infty, 0)$. 因此, 设 $x \in [0, \infty)$,
\begin{align*}
M_{\dyadic}f(x) &= \sup_{I \in \mathcal{D}(\mathbb{R})}\frac{\chi_{I}(x)}{|I|}\int_{I}|f(y)|\dd y \\
&= \sup_{I \in \mathcal{D}(\mathbb{R}),\, I \subset [0, \infty)}\frac{\chi_{I}(x)}{|I|}\int_{I}|f(y)|\dd y \\
&= \sup_{I \in \mathcal{D}(\mathbb{R})}\frac{\chi_{I}(x)}{|I|}\int_{I}\chi_{[0, \infty)}(y)|f(y)|\dd y \\
&= M_{\dyadic}[\chi_{[0, \infty)}f](x)
\end{align*}
即得. 同样, 设 $x \in (-\infty, 0)$,
\[
M_{\dyadic}f(x) = M_{\dyadic}[\chi_{(-\infty, 0)}f](x)
\]
即得. $\blacksquare$

\paragraph*{练习题 4.8}
因 $f$ 局部可积且 $Q$ 紧, 有 $\chi_{Q}f \in L^{1}(\mathbb{R}^{n})$, $\chi_{Q} \in L^{1}(\mathbb{R}^{n})$. 因此 $B$ 可积. 又
\[
\int_{Q}B(x)\dd x = \int_{Q}\chi_{Q}(x)f(x)\dd x - \frac{1}{|Q|}\int_{Q}f(z)\dd z \times \int_{Q}\chi_{Q}(x)\dd x = \int_{Q}\chi_{Q}(x)f(x)\dd x - \int_{Q}f(z)\dd z = 0.
\]
$\blacksquare$

\paragraph*{练习题 4.9}
(1) 证明逆否命题. 设 $\sup_{Q\in\mathcal{A}}|Q| = \infty$, 则对任意 $Q \in \mathcal{A}$, 因 $E \supset Q$, 有 $|E| \geq |Q|$, 由假定得 $|E| = \infty$.

(2) 设 $\mathcal{A}_{1} = \{Q \in \mathcal{A} : \text{存在 } R \in \mathcal{A} \text{ 使得 } R \ni Q\}$, 并定义 $\mathcal{A}_{0} = \mathcal{A}\setminus\mathcal{A}_{1}$. 于是, 一般地, 当 $Q_{1}, Q_{2} \in \mathcal{D}$ 不同时, $Q_{1}\cap Q_{2} = \emptyset$, $Q_{2} \subsetneq Q_{1}$, $Q_{1} \subsetneq Q_{2}$ 三者之一成立, 因此特别地, 当 $Q_{1}, Q_{2} \in \mathcal{A}_{0}$ 时, $Q_{1}\cap Q_{2} = \emptyset$. 又若 $Q \in \mathcal{A}_{1}$, 则存在 $R \in \mathcal{A}$ 使得 $Q \subset R$, 而由 (1), 这样的对象中存在体积最大者. 将其重新记为 $R$, 则 $R \in \mathcal{A}_{0}$. 因此,
\[
\bigcup_{Q \in \mathcal{A}}Q = \bigcup_{Q \in \mathcal{A}_{0}}Q \cup \bigcup_{Q \in \mathcal{A}_{1}}Q \subset \bigcup_{Q \in \mathcal{A}_{0}}Q \subset \bigcup_{Q \in \mathcal{A}}Q.
\]
即 $\bigcup_{Q\in\mathcal{A}}Q = \bigsqcup_{Q\in\mathcal{A}_{0}}Q$ 成立. 由此得结论.

\paragraph*{练习题 4.10}
当 $t > 2$ 时,
\[
H\chi_{[-1,1]}(t) = \int_{-1}^{1}\frac{1}{t-s}\dd s = [-\log(t-s)]_{-1}^{1} = \log\frac{1+t}{t-1}.
\]
由
\[
\lim_{t \to \infty}tH\chi_{[-1,1]}(t) = 2
\]
可知, 对充分大的 $t > 2$ (设 $t > t_{0} > 2$),
\[
|tH\chi_{[-1,1]}(t)| > 1.
\]
因此,
\[
\int_{t_{0}}^{\infty}|H\chi_{[-1,1]}(t)|\dd t \geq \int_{t_{0}}^{\infty}\frac{\dd t}{t} = \infty.
\]
故 $H\chi_{[-1,1]} \notin L^{1}([2,\infty))$. $\blacksquare$

\paragraph*{练习题 4.11}
当 $f \in L^{2}(\mathbb{R})$ 时, 由 $f_{j} = \chi_{B(j)}\chi_{[0,j]}(|f|)f$ 给出的函数满足 $|f - f_{j}|^{2} \leq |f|^{2}$. 因此, 由于 $\{f_{j}\}_{j=1}^{\infty}$ 逐点收敛到 $f$, 由 Lebesgue 收敛定理, $\{f_{j}\}_{j=1}^{\infty}$ 在 $L^{2}$ 中收敛到 $f$. 因 $f_{j} \in L^{p}(\mathbb{R}) \cap L^{2}(\mathbb{R})$, 故在 $L^{p}(\mathbb{R})$ 中稠密. $\blacksquare$

\paragraph*{练习题 4.12}
任意固定 $\varepsilon > 0$.
\[
\left|\int_{|y|>\varepsilon}\frac{f(x-y)+g(x-y)}{y}\dd y\right| \leq \left|\int_{|y|>\varepsilon}\frac{f(x-y)}{y}\dd y\right| + \left|\int_{|y|>\varepsilon}\frac{g(x-y)}{y}\dd y\right| \leq H^{*}f(x) + H^{*}g(x).
\]
因此, 对 $\varepsilon > 0$ 取上确界, 得 $H^{*}[f+g](x) \leq H^{*}f(x) + H^{*}g(x)$. $\blacksquare$

\paragraph*{练习题 4.13}
(1) 仔细考察定理~\ref{thm:4.6} 的证明.

(2) 因两边均为 0, 当 $T = 0$ 时结论显然. 假定并非如此. 此时, 因 $A+B+\|T\|_{L^{2}\to L^{2}} > 0$, 设
\[
S = \frac{T}{A+B+\|T\|_{L^{2}\to L^{2}}}
\]
则由 (1), 对每个 $f \in L^{p}(\mathbb{R}^{n})$, 得 $\|Sf\|_{L^{p}} \leq \zeta(1,1,1)\|f\|_{L^{p}}$. 因此, 对每个 $f \in L^{p}(\mathbb{R}^{n})$, 得 $\|Tf\|_{L^{p}} \leq \zeta(1,1,1)(A+B+\|T\|_{L^{2}\to L^{2}})\|f\|_{L^{p}}$. $\blacksquare$

\section*{第5章}

\paragraph*{练习题 5.1}
利用提示中的等式与 $\varphi'(t+k) \leq \varphi'(t)$, $t, k \geq 0$ 即可证明. $\blacksquare$

\paragraph*{练习题 5.2}
设 $\{f_{j}\}_{j=1}^{\infty}$ 是 $(\mathcal{S}(\mathbb{R}^{n}), d)$ 中的 Cauchy 列. 因 $\max_{x\in\mathbb{R}^{n}}|f_{j}(x)-f_{k}(x)| = p_{0,0}(f_{j}-f_{k})$, $\{f_{j}\}_{j=1}^{\infty}$ 一致收敛到某个连续函数 $f$. 同样, 对任意多重指数 $\alpha \in \mathbb{N}_{0}^{n}$, 由于 $\{\partial^{\alpha}f_{j}\}_{j=1}^{\infty}$ 是 $(\mathcal{S}(\mathbb{R}^{n}), d)$ 中的 Cauchy 列, $\{\partial^{\alpha}f_{j}\}_{j=1}^{\infty}$ 一致收敛. 因此 $\{\partial^{\alpha}f_{j}\}_{j=1}^{\infty}$ 一致收敛到 $\partial^{\alpha}f$, 从而 $\partial^{\alpha}f$ 连续. 特别地, $f \in C^{\infty}(\mathbb{R}^{n})$. 这里, 对任意多重指数 $\alpha, \beta \in \mathbb{N}_{0}^{n}$,
\[
p_{\alpha,\beta}(f-f_{j}) \leq \limsup_{k \to \infty}p_{\alpha,\beta}(f_{k}-f_{j})
\]
令 $j \to \infty$, 由定理~\ref{thm:5.1} 可知 $\{f_{j}\}_{j=1}^{\infty}$ 在 $(\mathcal{S}(\mathbb{R}^{n}), d)$ 中收敛. $\blacksquare$

\paragraph*{练习题 5.3}
(1) 由提示中的不等式,
\begin{align*}
\|f\|_{L^{p}}^{p} &= \int_{\mathbb{R}^{n}}\frac{1}{(1+|x|)^{p(n+1)}}\dd x \\
&\leq \int_{\mathbb{R}^{n}}\frac{1}{(1+|x|)^{n+1}}\dd x \\
&\leq \int_{\mathbb{R}^{n}}\prod_{j=1}^{n}\frac{1}{(1+|x_{j}|)^{\frac{n+1}{n}}}\dd x \\
&= \left(\int_{\mathbb{R}}\frac{1}{(1+|t|)^{\frac{n+1}{n}}}\dd t\right)^{n} \\
&= (2n)^{n} < \infty.
\end{align*}

(2) 由提示中的不等式, $M \leq q(g) + \sum_{j=1}^{n}r_{j}(g) < \infty$.

(3) 若 $g \in \mathcal{S}(\mathbb{R}^{n})$, 由 (2), $|g(x)| \leq \frac{M}{(1+|x|)^{n+1}}$. 由 (1), 右边属于 $L^{p}(\mathbb{R}^{n})$, 故 $g$ 也属于同一空间. $\blacksquare$

\paragraph*{练习题 5.4}
对函数 $\varphi$, 使用记号~\eqref{eq:5.1}. 取 $\alpha$, $\gamma$ 为多重指数, 由
\[
(f \cdot g)_{(\gamma, \alpha)} = \sum C_{\beta}f_{(\gamma, \beta)} \cdot g_{(0, \alpha-\beta)}
\]
得
\[
p_{\gamma,\beta}(f \cdot g) \leq \sum C_{\beta}p_{\gamma,\beta}(f)p_{0,\alpha-\beta}(g) < \infty.
\]
$\blacksquare$

\paragraph*{练习题 5.5}
由定理~\ref{thm:5.1}, 只需证明其与式~\eqref{eq:5.4} 对一切 $\alpha, \beta \in \mathbb{N}_{0}^{n}$ 满足
\[
\lim_{j \to \infty}p_{\alpha,\beta}(\varphi_{j}-\varphi) = 0
\]
是等价的. 若式~\eqref{eq:5.4} 成立, 则对一切 $\alpha, \beta \in \mathbb{N}_{0}^{n}$ 满足上式可由
\[
p_{\alpha,\beta}(\varphi_{j}-\varphi) \leq \sup_{x \in \mathbb{R}^{n}}(1+|x|)^{|\alpha|}|\partial^{\beta}\varphi(x) - \partial^{\beta}\varphi_{j}(x)|
\]
看出, 这是显然的. 来证明其逆. 对 $k_{1}, k_{2}, \ldots, k_{n} \in \mathbb{N}_{0}$, 记 $|k| = k_{1}+k_{2}+\cdots+k_{n}$, $x^{k} = x_{1}^{k_{1}}x_{2}^{k_{2}}\cdots x_{n}^{k_{n}}$, $k = (k_{1}, k_{2}, \ldots, k_{n})$, 由
\begin{align*}
(1+|x|)^{N} &\leq (1+|x_{1}|+|x_{2}|+\cdots+|x_{n}|)^{N} \\
&\leq \sum_{k \in \mathbb{N}_{0}^{n},\, |k| \leq N}\frac{(N-|k|)!k_{1}!k_{2}!\cdots k_{n}!}{N!}|x^{k}| \\
&\leq (N!)^{n}\sum_{k \in \mathbb{N}_{0}^{n},\, |k| \leq N}|x^{k}|
\end{align*}
得
\[
\sup_{x \in \mathbb{R}^{n}}(1+|x|)^{N}|\partial^{\alpha}\varphi(x) - \partial^{\beta}\varphi_{j}(x)| \leq (N!)^{n}\sum_{k \in \mathbb{N}_{0}^{n},\, |k| \leq N}\frac{(N-|k|)!k_{1}!k_{2}!\cdots k_{n}!}{N!}p_{k,\alpha}(\varphi_{j}-\varphi)
\]
故若 $\{\varphi_{j}\}_{j=1}^{\infty} \subset \mathcal{S}(\mathbb{R}^{n})$ 收敛到 $\varphi \in \mathcal{S}(\mathbb{R}^{n})$, 即对一切 $\alpha, \beta \in \mathbb{N}_{0}^{n}$ 满足上式, 则式~\eqref{eq:5.4} 成立. $\blacksquare$

\paragraph*{练习题 5.6}
$O$ 本身是开集族, 且 $\mathcal{S}'(\mathbb{R}^{n})$ 的任意开集族都包含在 $O$ 中, 故显然.

\paragraph*{练习题 5.7}
\begin{itemize}
\item $\mathcal{F}: \mathcal{S}'(\mathbb{R}^{n}) \to \mathcal{S}'(\mathbb{R}^{n})$ 是单射. 若 $\mathcal{F}f = 0$, 则 $\langle f, \varphi\rangle = \langle \mathcal{F}^{-1}\mathcal{F}f, \varphi\rangle = \langle \mathcal{F}f, \mathcal{F}^{-1}\varphi\rangle = \langle 0, \mathcal{F}^{-1}\varphi\rangle = 0$ 即得.

\item $\mathcal{F}: \mathcal{S}'(\mathbb{R}^{n}) \to \mathcal{S}'(\mathbb{R}^{n})$ 是满射. 任意给定 $F \in \mathcal{S}'(\mathbb{R}^{n})$, 定义 $G \in \mathcal{S}'(\mathbb{R}^{n})$ 为 $\langle G, \varphi\rangle = \langle F, \mathcal{F}^{-1}\varphi\rangle$, 则 $\langle \mathcal{F}G, \varphi\rangle = \langle G, \mathcal{F}\varphi\rangle = \langle F, \mathcal{F}^{-1}\mathcal{F}\varphi\rangle = \langle F, \varphi\rangle$, 故得 $F = \mathcal{F}G$.

\item $\mathcal{F}: \mathcal{S}'(\mathbb{R}^{n}) \to \mathcal{S}'(\mathbb{R}^{n})$ 连续. 只需证明在 $0 \in \mathcal{S}'(\mathbb{R}^{n})$ 处的连续性. 设 $\mathcal{Z} = \{g \in \mathcal{S}'(\mathbb{R}^{n}) : |\langle g, \varphi\rangle| < \varepsilon\}$. 这里, $\varepsilon > 0$, $\varphi \in \mathcal{S}(\mathbb{R}^{n})$. 于是, 因 $\mathcal{F}^{-1}[\mathcal{Z}] = \{g \in \mathcal{S}'(\mathbb{R}^{n}) : |\langle g, \mathcal{F}\varphi\rangle| < \varepsilon\}$, 可知 $\mathcal{F}^{-1}$ 连续.

\item $\mathcal{F}: \mathcal{S}'(\mathbb{R}^{n}) \to \mathcal{S}'(\mathbb{R}^{n})$ 的逆映射连续. 这可以与 $\mathcal{F}: \mathcal{S}'(\mathbb{R}^{n}) \to \mathcal{S}'(\mathbb{R}^{n})$ 连续同样地证明, 故省略. $\blacksquare$
\end{itemize}

\paragraph*{练习题 5.8}
在 $R = \pi$ 处取最大值. 因 $|f(t)| \leq f(\pi)$,
\[
\sup_{t, \varepsilon, R > 0}\left|\int_{\varepsilon}^{R}\frac{\sin(ts)}{s}\dd s\right| \leq 2f(\pi) < \infty.
\]
$\blacksquare$

\paragraph*{练习题 5.9}
设 $\varphi(x) = \frac{1}{(1+|x|)^{n+1}}$.

\begin{itemize}
\item 将例~\ref{exa:5.4} 应用于 $h(x) = \exp(-i\xi \cdot x)$, 结合 $\partial_{j}e^{-ix \cdot \xi} = -i\xi_{j}e^{-ix \cdot \xi}$, 可知式~\eqref{eq:5.7} 成立.

\item 设 $F(x, \xi) = f(x)e^{-ix \cdot \xi}$. 与练习题~\ref{exer:5.3}(2) 相同, $M = \sup_{x \in \mathbb{R}^{n}}(1+|x|)^{n+2}|f(x)| < \infty$. $|F(x, \xi)| \leq M(1+|x|)^{-n-2}$ 成立, 且由练习题~\ref{exer:5.3}(1), $\varphi \in L^{1}(\mathbb{R}^{n})$, 故
\[
\int_{\mathbb{R}^{n}}|F(x, \xi)|\dd x < \infty.
\]
又因 $F_{\xi_{j}}(x, \xi) = -ix_{j}f(x)e^{-ix \cdot \xi}$, $|F_{\xi_{j}}(x, \xi)| \leq M|x_{j}|(1+|x|)^{-n-2} \leq M\varphi(x)$ 成立, 再次由练习题~\ref{exer:5.3}(1), $\varphi \in L^{1}(\mathbb{R}^{n})$, 故可用定理~\ref{thm:2.13}, 式~\eqref{eq:5.8} 成立.

\item 由 $f \in \mathcal{S}(\mathbb{R}^{n})$, 式~\eqref{eq:5.7}, \eqref{eq:5.8} 可反复使用, 故利用绝对值为 1 的适当复常数 $c_{\alpha,\beta}$, 可作
\[
\xi^{\alpha}\partial^{\beta}\mathcal{F}\varphi(\xi) = c_{\alpha,\beta}\mathcal{F}(\partial^{\alpha}[x^{\beta}\varphi])(x)
\]
形式的变形. 因此只需考虑 $\alpha = \beta = 0$ 的情形. 当 $\alpha = \beta = 0$ 时, 由练习题~\ref{exer:5.3}(1), (2), 可取 $N = n+1$. $\blacksquare$
\end{itemize}

\paragraph*{练习题 5.10}
(1) 使用 $Ax \cdot \xi = x \cdot A^{\mathrm{t}}\xi = x \cdot A^{-1}\xi$ 与 $\det A = \pm 1$.

(2) 由于等式局部成立, 特别地, 取 $0 < \varepsilon < R < \infty$, 只需对 $\varepsilon < |x| < R$ 证明等式. 取满足 $0 \notin \operatorname{supp}(g)$ 且在 $\varepsilon < |x| < R$ 上取值为 1 的 $g \in C_{\mathrm{c}}^{\infty}(\mathbb{R}^{n})$, 考虑以乘积 $g \cdot u$ 代替 $u$. 因 $\|\nabla^{k}u\| = \|\nabla^{k}[g \cdot u]\|$, 若必要, 用 $g \cdot u$ 替换 $u$, 可假定 $u \in C_{\mathrm{c}}^{\infty}(\mathbb{R}^{n})$.
\[
\partial^{\alpha}u(x) = \partial_{\alpha_{1}}\partial_{\alpha_{2}}\cdots\partial_{\alpha_{k}}[\mathcal{F}^{-1}[\mathcal{F}u]](x) = \frac{(-i)^{k}}{(2\pi)^{\frac{n}{2}}}\int_{\mathbb{R}^{n}}e^{ix \cdot \xi}\xi^{\alpha}\mathcal{F}u(\xi)\dd\xi \quad (x \in \mathbb{R}^{n})
\]
以及
\[
\partial^{\alpha}u(x) = \partial_{\alpha_{1}}\partial_{\alpha_{2}}\cdots\partial_{\alpha_{k}}[\mathcal{F}[\mathcal{F}^{-1}u]](x) = \frac{i^{k}}{(2\pi)^{\frac{n}{2}}}\int_{\mathbb{R}^{n}}e^{-ix \cdot \eta}\eta^{\alpha}\mathcal{F}^{-1}u(\eta)\dd\eta \quad (x \in \mathbb{R}^{n}).
\]
因此,
\[
(2\pi)^{n}(\partial^{\alpha}u(x))^{2} = \int\int_{\mathbb{R}^{n} \times \mathbb{R}^{n}}e^{ix \cdot (\xi-\eta)}\xi^{\alpha} \cdot \eta^{\alpha}\mathcal{F}u(\xi)\mathcal{F}^{-1}u(\eta)\dd\xi\dd\eta.
\]
所以,
\begin{align*}
(2\pi)^{n}\|\nabla^{k}u(x)\|^{2} &= (2\pi)^{n}\sum_{\alpha \in \{1,2,\ldots,n\}^{k}}(\partial^{\alpha}u(x))^{2} \\
&= \sum_{\alpha \in \{1,2,\ldots,n\}^{k}}\int\int_{\mathbb{R}^{n} \times \mathbb{R}^{n}}e^{ix \cdot (\xi-\eta)}\xi^{\alpha} \cdot \eta^{\alpha}\mathcal{F}u(\xi)\mathcal{F}^{-1}u(\eta)\dd\xi\dd\eta \\
&= \int\int_{\mathbb{R}^{n} \times \mathbb{R}^{n}}e^{ix \cdot (\xi-\eta)}(\xi \cdot \eta)^{k}\mathcal{F}u(\xi)\mathcal{F}^{-1}u(\eta)\dd\xi\dd\eta \quad (x \in \mathbb{R}^{n}).
\end{align*}
这里, 设 $A \in O(n)$,
\[
Ax \cdot y = x \cdot {}^{\mathrm{t}}Ay, \quad Ax \cdot Ay = x \cdot y \quad (x, y \in \mathbb{R}^{n})
\]
成立. 由 (1),
\[
\mathcal{F}u(A\xi) = \mathcal{F}u(\xi), \quad \mathcal{F}^{-1}u(A\xi) = \mathcal{F}^{-1}u(\xi) \quad (\xi \in \mathbb{R}^{n}).
\]
由 $A$ 的正交性与变量替换,
\begin{align*}
\|\nabla^{k}u(Ax)\|^{2} &= \frac{1}{(2\pi)^{n}}\int\int_{\mathbb{R}^{n} \times \mathbb{R}^{n}}e^{iAx \cdot (\xi-\eta)}(\xi \cdot \eta)^{k}\mathcal{F}u(\xi)\mathcal{F}^{-1}u(\eta)\dd\xi\dd\eta \\
&= \frac{1}{(2\pi)^{n}}\int\int_{\mathbb{R}^{n} \times \mathbb{R}^{n}}e^{ix \cdot A(\xi-\eta)}(\xi \cdot \eta)^{k}\mathcal{F}u(\xi)\mathcal{F}^{-1}u(\eta)\dd\xi\dd\eta \\
&= \frac{1}{(2\pi)^{n}}\int\int_{\mathbb{R}^{n} \times \mathbb{R}^{n}}e^{ix \cdot (\tilde{\xi}-\tilde{\eta})}(A\tilde{\xi} \cdot A\tilde{\eta})^{k}\mathcal{F}u(A\tilde{\xi})\mathcal{F}^{-1}u(A\tilde{\eta})\dd\tilde{\xi}\dd\tilde{\eta} \\
&= \frac{1}{(2\pi)^{n}}\int\int_{\mathbb{R}^{n} \times \mathbb{R}^{n}}e^{ix \cdot (\tilde{\xi}-\tilde{\eta})}(\tilde{\xi} \cdot \tilde{\eta})^{k}\mathcal{F}u(\tilde{\xi})\mathcal{F}^{-1}u(\tilde{\eta})\dd\tilde{\xi}\dd\tilde{\eta} \\
&= \|\nabla^{k}u(x)\|^{2}.
\end{align*}
因此, $\|\nabla^{k}u\|^{2}$ 确实满足条件. $\blacksquare$

\paragraph*{练习题 5.11}
由分部积分
\begin{align*}
\left|\int_{\mathbb{R}}f(t)\varphi(t)\dd t\right| &= \left|\int_{\mathbb{R}}\varphi'(t)\cos(e^{t})\dd t\right| \\
&= \int_{\mathbb{R}}|\varphi'(t)\cos(e^{t})|\dd t \\
&\leq \pi p_{2,1}(\varphi)
\end{align*}
即得. 因此, 由 $C_{\mathrm{c}}^{\infty}(\mathbb{R}^{n})$ 在 $\mathcal{S}(\mathbb{R}^{n})$ 中的稠密性 (参照例~\ref{exa:5.7}) 得结论. $\blacksquare$

\paragraph*{练习题 5.12}
对任意 $\Phi \in \mathcal{S}(\mathbb{R}^{n})$, 只需证明
\[
\langle\mathcal{F}[f*g], \Phi\rangle = (2\pi)^{\frac{n}{2}}\langle\mathcal{F}g \cdot \mathcal{F}f, \Phi\rangle.
\]
由 $\mathcal{S}'(\mathbb{R}^{n})$ 中 Fourier 变换的定义,
\[
\langle\mathcal{F}[f*g], \Phi\rangle = \langle f*g, \mathcal{F}\Phi\rangle.
\]
这里, 设 $\Psi = \mathcal{F}\Phi(-\cdot)$, 由卷积的定义,
\[
\langle\mathcal{F}[f*g], \Phi\rangle = \langle f*g, \Psi(-\cdot)\rangle = (f*g)*\Psi(0).
\]
由定理~\ref{thm:5.12}(3),
\[
\langle\mathcal{F}[f*g], \Phi\rangle = g*(f*\Psi)(0) = \langle g, f*[\Psi(-\cdot)]\rangle.
\]
再次由 $\mathcal{S}(\mathbb{R}^{n})$ 中 Fourier 变换的定义,
\[
\langle\mathcal{F}[f*g], \Phi\rangle = \langle\mathcal{F}g, \mathcal{F}^{-1}[f*[\Psi(-\cdot)]]\rangle.
\]
这里,
\begin{align*}
\mathcal{F}^{-1}[f*[\Psi(-\cdot)]](\xi) &= \frac{1}{(2\pi)^{\frac{n}{2}}}\int_{\mathbb{R}^{n}}f*[\Psi(-\cdot)](x)e^{ix \cdot \xi}\dd x \\
&= \frac{1}{(2\pi)^{\frac{n}{2}}}\int_{\mathbb{R}^{n}}\left(\int_{\mathbb{R}^{n}}f(y)\Psi(-x+y)e^{ix \cdot \xi}\dd y\right)\dd x
\end{align*}
但由于 $f, \Psi \in \mathcal{S}(\mathbb{R}^{n}) \subset L^{1}(\mathbb{R}^{n})$, 可以交换积分顺序. 因此,
\[
\mathcal{F}[f*\Psi(-\cdot)](\xi) = \frac{1}{(2\pi)^{\frac{n}{2}}}\int_{\mathbb{R}^{n}}\left(\int_{\mathbb{R}^{n}}f(y)\Psi(-x+y)e^{ix \cdot \xi}\dd x\right)\dd y
\]
即得. 作变量替换,
\[
\mathcal{F}[f*\Psi(-\cdot)](\xi) = (2\pi)^{\frac{n}{2}}\mathcal{F}f(\xi)\mathcal{F}\Psi(\xi) = (2\pi)^{\frac{n}{2}}\mathcal{F}f(\xi)\Phi(\xi)
\]
将其代入, 运用 $\mathcal{F}$ 的定义, 得
\[
\langle\mathcal{F}[f*g], \Phi\rangle = (2\pi)^{\frac{n}{2}}\langle\mathcal{F}g \cdot \mathcal{F}f, \Phi\rangle.
\]
$\blacksquare$

\section*{第6章}

\paragraph*{练习题 6.1}
仿照引理~\ref{lem:6.2} 的证明即可. $\blacksquare$

\paragraph*{练习题 6.2}
因三个线性算子 $\mathcal{F}$, $f \mapsto mf$, $\mathcal{F}^{-1}$ 都是 $L^{2}$ 有界的, 故 $T_{m}$ 是 $L^{2}$ 有界的. $\blacksquare$

\paragraph*{练习题 6.3}
$\mathcal{F}[T_{m}g] = m\mathcal{F}g$. 对任意 $\alpha, \beta \in \mathbb{N}_{0}^{n}$ 与 $N \in \mathbb{N}$, $\partial^{\alpha}m(\xi) = O(|\xi|^{-|\alpha|})$ 且 $\partial^{\beta}\mathcal{F}g(\xi) = O(|\xi|^{N})$, 故由 Leibniz 法则, $\partial^{\alpha}\mathcal{F}[T_{m}g] = O(|\xi|^{N})$ 成立. 因此 $T_{m}g \in \mathcal{S}_{\infty}$. $\blacksquare$


% 参考文献
% ============ 参考文献 (原书 pp.156) + 作者简介 (原书卷末) ============

\chapter*{参考文献}

就日语教材而言, [2], [3] 是新刊. 这两本书是长达 600 页的巨著, 在读完本书内容之后再读为好. [4]-[6] 也可以说是与 [2], [3] 不同的, 从本书延伸的文献. 就英语教材而言, 有 [7]-[9], [12] 等. [7] 是由约 200 页构成的相对较短的书籍, 但在读同样由两部构成的 [8], [9] 以及巨著 [12] 之前, 先读这本书的话, 虽然无法概观这些巨著, 却能体会其工作的大致氛围. 数列空间中的调和分析请参考 [1]. 另外, 练习问题 3.5 与 5.10 分别采录自 [10], [11].

\begin{itemize}
\item[{[1]}] 新井仁之, フーリエ解析学 (講座 数学の考え方 17), 朝倉書店, 2003.
\item[{[2]}] 宮地晶彦, ユークリッド空間上のフーリエ解析 I (朝倉数学大系 13), 朝倉書店, 2021.
\item[{[3]}] 宮地晶彦, ユークリッド空間上のフーリエ解析 II (朝倉数学大系 14), 朝倉書店, 2021.
\item[{[4]}] 藪田公三, 佐藤圓治, 田中仁, 宮地晶彦, 中路貴彦, 古典調和解析 (解析学百科 I), 朝倉書店, 2008.
\item[{[5]}] 藪田公三, 特異積分 (岩波数学叢書), 岩波書店, 2010.
\item[{[6]}] 澤野嘉宏, ベゾフ空間論, 日本評論社, 2011.
\item[{[7]}] J. Duoandikoetxea, \emph{Fourier Analysis}, translated and revised from the 1995 Spanish original by D. Cruz-Uribe, Graduate Studies in Mathematics, 29, American Mathematical Society, Providence, RI, 2001.
\item[{[8]}] L. Grafakos, \emph{Classical Fourier analysis}, Graduate Texts in Mathematics, 249, Springer, New York, 2014.
\item[{[9]}] L. Grafakos, \emph{Modern Fourier analysis}, Graduate Texts in Mathematics, 250, Springer, New York, 2014.
\item[{[10]}] V. Burenkov, A. Ghorbanalizadeh and Y. Sawano, Uniform boundedness of Kantorovich operators in Morrey spaces, \emph{Positivity}, 22, 1097--1107 (2018).
\item[{[11]}] K. Morii, T. Sato and Y. Sawano, Certain identities on derivatives of radial homogeneous and logarithmic functions, \emph{Commun.\ Math.\ Anal.}, 9, no.\ 2, 51--66 (2010).
\item[{[12]}] E.M. Stein, \emph{Harmonic Analysis, Real-variable Methods, Orthogonality, and Oscillatory Integrals}, Princeton University Press, Princeton, NJ, 1993.
\end{itemize}

\clearpage
\chapter*{作者简介}

\begin{center}
\large 泽野嘉宏 \quad (さわのよしひろ)
\end{center}

\begin{itemize}
\item 1979 年 生于静冈县
\item 2006 年 东京大学大学院数理科学研究科博士后期课程修了, 博士(数理科学)
\item 2008 年 学习院大学理学部助教
\item 2009 年 京都大学理学部助教
\item 2012 年 首都大学东京理学部准教授
\item 2020 年 中央大学理工学部教授, 至今
\end{itemize}

2009 年获日本数学会奖建部贤弘特别奖.

专业, 研究方向: 调和分析, 函数空间论.

主要著作: \emph{Theory of Besov spaces}, (Springer, 2018); \emph{Morrey Spaces: Introduction and Applications to Integral Operators and PDE's}, Volume I, II, (Chapman and Hall/CRC, 2020); 『ベゾフ空間論』(日本评论社, 2011), 等.


% 索引
\printindex

\end{document}
